\documentclass[12pt]{book}

% Preámbulo oficial del repositorio
\input{../../../_assets/preambulo}

% Idioma italiano y paquetes adicionales
\usepackage{yhmath}
\usepackage{centernot}
\newcommand{\widearc}[1]{\wideparen{#1}}

% Limpieza completa de contadores y entornos de preambulo.tex para ponerlos en italiano
\makeatletter
\let\c@teo\relax \let\teo\relax \let\endteo\relax
\let\c@definicion\relax \let\definicion\relax \let\enddefinicion\relax
\let\c@ejercicio\relax \let\ejercicio\relax \let\endejercicio\relax
\expandafter\let\csname ejercicio*\endcsname\relax
\expandafter\let\csname endejercicio*\endcsname\relax
\let\c@coro\relax \let\coro\relax \let\endcoro\relax
\let\c@prop\relax \let\prop\relax \let\endprop\relax
\let\c@lema\relax \let\lema\relax \let\endlema\relax
\let\c@observacion\relax \let\observacion\relax \let\endobservacion\relax
\let\c@notacion\relax \let\notacion\relax \let\endnotacion\relax
\let\c@ejemplo\relax \let\ejemplo\relax \let\endejemplo\relax
\makeatother

\theoremstyle{plain}
\newtheorem{teo}{Teorema}[section]
\newtheorem{prop}[teo]{Proposizione}
\newtheorem{lema}[teo]{Lemma}
\newtheorem{coro}[teo]{Corollario}

\theoremstyle{definition}
\newtheorem{definicion}[teo]{Definizione}
\newtheorem{ejemplo}[teo]{Esempio}
\newtheorem{ejercicio}[teo]{Esercizio}
\newtheorem*{ejercicio*}{Esercizio}

\theoremstyle{remark}
\newtheorem{observacion}[teo]{Osservazione}
\newtheorem{notacion}[teo]{Nota}
\newtheorem*{intuizione}{Intuizione}

\newenvironment{soluzione}{\begin{proof}[Soluzione]}{\end{proof}}

\begin{document}

    \selectlanguage{italian}
    \input{../../../_assets/portada}

    % Versión en italiano de \portadaFotoDif
    \renewcommand{\portadaFotoDif}[8]{%
        \thispagestyle{empty}
        \newgeometry{margin=0cm}
        \fancyhead[L]{\helv #3}
        \fancyhead[R]{\helv \nouppercase{\leftmark}}
        \pagecolor{#4}
        \begin{tikzpicture}[remember picture, overlay]
            % Imagen centrada cubriendo todo el alto del A4 sin deformarse
            \node[anchor=center, opacity=0.3] at (current page.center) {
                \includegraphics[height=\paperheight]{#1}
            };
            % Título grande arriba a la izquierda (independiente de la imagen)
            \node[anchor=north west, text=white, opacity=1, font=\fontsize{52}{68}\selectfont\bfseries\sffamily, align=left] 
                at ([xshift=2.2cm, yshift=-3cm]current page.north west) {#2};
            % Textos abajo a la derecha
            \node[anchor=south east, text=white, opacity=1, font=\fontsize{12}{18}\selectfont\sffamily, align=right] 
                at ([xshift=-1.8cm, yshift=3.2cm]current page.south east) {\href{https://losdeldgiim.github.io/}{\textbf{Los Del DGIIM}, \texttt{\footnotesize losdeldgiim.github.io}}};
            \node[anchor=south east, text=white, opacity=1, font=\fontsize{12}{15}\selectfont\sffamily, align=right] 
                at ([xshift=-1.8cm, yshift=1.8cm]current page.south east) {Doble Grado en Ingeniería Informática y Matemáticas\\Universidad de Granada};
        \end{tikzpicture}
        \null\newpage
        \restoregeometry
        \nopagecolor

        \newpage
        \thispagestyle{empty}
        \begin{tikzpicture}[remember picture, overlay]
            \node[anchor=south west, inner sep=3cm] at (current page.south west) {
                \begin{minipage}{0.5\paperwidth}
                    Quest'opera è distribuita con
                    \href{https://creativecommons.org/licenses/by-nc-nd/4.0/}{
                        Licenza Creative Commons Attribuzione - Non commerciale - Non opere derivate 4.0 Internazionale (CC BY-NC-ND 4.0).
                    }\\

                    Sei libero di condividere e ridistribuire il contenuto di quest'opera con qualsiasi mezzo o formato, a condizione di riconoscere la paternità agli autori originali e di non utilizzarla per fini commerciali.
                \end{minipage}
            };
        \end{tikzpicture}

        \pretitle{\begin{center}\bfseries\fontsize{42}{56}\selectfont}
        \posttitle{\par\end{center}\vspace{2em}}
        \preauthor{\begin{center}\Large}
        \postauthor{\par\end{center}\vfill}
        \predate{\begin{center}\huge}
        \postdate{\par\end{center}\vspace{2em}}

        \title{#2}
        \author{\href{https://losdeldgiim.github.io/}{Los Del DGIIM, \texttt{\large losdeldgiim.github.io}}
        \\ \vspace{0.5cm}#8}
        \date{Lecce, #7}
        \maketitle

        \tableofcontents
    }

    \portadaFotoDif{../../unisalento.jpg}{Istituzioni di\\Geometria\\Superiore}{Istituzioni di Geometria Superiore}{MidnightBlue}{-8}{28}{2026/27}{Ángel Caro Montalbán}

    \chapter{Il Gruppo Fondamentale di Poincar\'e}

\section{Cammino}
\begin{definicion} Sia $X$ uno spazio topologico e $I = [0,1]$. Un \textit{cammino} in $X$ \`e un'applicazione continua $\alpha: I \to X$.
\begin{itemize}
    \item $\alpha(0) = x_0 \in X$ si dice PUNTO INIZIALE di $\alpha$.
    \item $\alpha(1) = x_1 \in X$ si dice PUNTO FINALE di $\alpha$.
    \item $\alpha(I) \subseteq X$ si dice SOSTEGNO di $\alpha$.
\end{itemize}
\end{definicion}
\textbf{Osservazione:} Un'applicazione continua $\alpha: [a,b] \to X$ in senso stretto NON \`E UN CAMMINO. Tuttavia, l'applicazione $\varphi: [0,1] \to [a,b]$ definita da $\varphi(t) = (1-t)a + tb$ \`e un omeomorfismo. Pertanto, la composizione $\alpha \circ \varphi: [0,1] \to X$ \`e un cammino valido.

\textbf{NON CONFONDERE UN CAMMINO CON IL SUO SOSTEGNO.} In particolare, cammini diversi possono avere lo stesso sostegno.

\textbf{Esempio:} In $X = \mathbb{R}^2$, consideriamo il cammino $\alpha: I \to \mathbb{R}^2$ definito da:
\[ \alpha(t) = \left( \cos\left(\frac{\pi}{2}t\right), \sin\left(\frac{\pi}{2}t\right) \right) \]
Un altro cammino, con lo stesso sostegno, \`e $\beta: I \to \mathbb{R}^2$ definito da:
\[ \beta(t) = (t, \sqrt{1-t^2}) \]
Allora, chiaramente $\alpha \neq \beta$, pur descrivendo geometricamente lo stesso arco.

\textbf{Curva di Peano:} Un cammino pu\`o essere molto diverso dal nostro concetto intuitivo di "curva". G. Peano (1890) prov\`o che il quadrato $I^2 = I \times I$ \`e immagine continua di $I$. Cio\`e, $\exists \alpha: I \to I^2$ continua e suriettiva. Pi\`u in generale, $I^n$ \`e immagine continua di $I$.

\begin{definicion} Sia $X$ uno spazio topologico e $\alpha: I \to X$ un cammino. L'applicazione $\alpha^{-1}: I \to X$ si chiama \textit{cammino inverso} di $\alpha$, definita come:
\[ \alpha^{-1}(t) = \alpha(1-t) \]
Chiaramente, $\alpha^{-1}(0) = \alpha(1)$, $\alpha^{-1}(1) = \alpha(0)$ e $\alpha^{-1}(I) = \alpha(I)$.
\end{definicion}

\section{Prodotto Stellato}
\begin{notacion}[Lemma di Incollamento]
Se $X = A \cup B$ con $A, B$ entrambi chiusi (oppure entrambi aperti) in $X$, e $f|_A: A \to Y$, $g|_B: B \to Y$ sono continue con $f = g$ su $A \cap B$, allora l'applicazione $h$ definita da $f$ su $A$ e da $g$ su $B$ \`e continua. Lo useremo di continuo per mostrare la continuit\`a di applicazioni definite ``a pezzi''.
\end{notacion}

Siano $\alpha, \beta$ cammini in $X$ con $\alpha(1) = \beta(0)$. Si chiama \textbf{PRODOTTO STELLATO} (o concatenazione) di $\alpha$ e $\beta$, e si indica con $\alpha * \beta$, il cammino $\alpha * \beta: I \to X$ cos\`i definito:
\[ 
\alpha * \beta(t) = 
\begin{cases} 
\alpha(2t) & \text{se } 0 \leq t \leq 1/2 \\
\beta(2t-1) & \text{se } 1/2 \leq t \leq 1 
\end{cases}
\]

\textbf{Osservazioni:}
\begin{itemize}
    \item $\alpha * \beta$ \`e ben posto perch\'e nel punto di giuntura $t = 1/2$: 
    \[ \alpha\left(2 \cdot \frac{1}{2}\right) = \alpha(1) = \beta(0) = \beta\left(2 \cdot \frac{1}{2} - 1\right) \]
    \item La continuit\`a di $\alpha * \beta$ segue dal \textit{Lemma di Incollamento}.
    \item $\alpha * \beta(0) = \alpha(0)$, $\alpha * \beta(1) = \beta(1)$ e il sostegno \`e $\alpha(I) \cup \beta(I)$.
    \item Il "prodotto stellato" \textbf{non \`e associativo} rigorosamente come funzione. Infatti, se $\alpha(1) = \beta(0)$ e $\beta(1) = \gamma(0)$, in generale $\alpha * (\beta * \gamma) \neq (\alpha * \beta) * \gamma$. Ad esempio, in $t = 1/4$ si ha $(\alpha * (\beta * \gamma))(1/4) = \alpha(1/2)$, mentre $((\alpha * \beta) * \gamma)(1/4) = \alpha(1) = \beta(0)$: i due cammini percorrono le stesse tre tratte a ``velocit\`a'' diverse. Sono per\`o omotopi relativamente a $\partial I$ (vedi il teorema sulle propriet\`a della concatenazione, pi\`u avanti).
\end{itemize}

\section{Omotopia di Funzioni}
\begin{definicion} Siano $X, Y$ spazi topologici, $A \subseteq X$, e $f, g: X \to Y$ continue. $f$ e $g$ si dicono \textbf{omotope relativamente ad A} se $\exists H: X \times I \to Y$ continua, tale che:
\begin{enumerate}
    \item $H(x,0) = f(x) \quad \forall x \in X$
    \item $H(x,1) = g(x) \quad \forall x \in X$
    \item $H(x,t) = f(x) \quad \forall x \in A, \forall t \in I$
\end{enumerate}
\end{definicion}

In tal caso, scriveremo $H: f \simeq g \text{ rel } A$. Se $A = \emptyset$, si scrive semplicemente $f \simeq g$.

\textbf{Osservazione:} Se $H: f \simeq g \text{ rel } A \implies f|_A = g|_A$. Il viceversa in generale non vale.

\textbf{Definizione (Convesso):} $\Omega \subseteq \mathbb{R}^n$ si dice convesso se $\forall y_1, y_2 \in \Omega$, il segmento che li congiunge appartiene a $\Omega$: $\overline{y_1 y_2} = \{(1-t)y_1 + ty_2 \mid t \in I\} \subseteq \Omega$.

\textbf{Esercizi/Propriet\`a:}
\begin{enumerate}
    \item Siano $X$ sp. top., $f, g: X \to \Omega$ (con $\Omega \subseteq \mathbb{R}^n$ convesso) continue. Allora: 
    \[ f \simeq g \text{ rel } A \iff f|_A = g|_A \]
    \textit{Dimostrazione:} ($\Rightarrow$) \`e l'osservazione precedente. ($\Leftarrow$) Si pone $H(x,t) = (1-t)f(x) + tg(x)$. $H$ \`e continua e, per convessit\`a di $\Omega$, $H(x,t) \in \Omega$ (sta sul segmento da $f(x)$ a $g(x)$). Inoltre $H(x,0)=f(x)$, $H(x,1)=g(x)$ e, se $x \in A$, $f(x) = g(x)$ implica $H(x,t) = f(x)$ per ogni $t$.
    
    \item Siano $X, Y$ sp. top., $y_0, y_1 \in Y$, e $c_{y_0}, c_{y_1}: X \to Y$ le rispettive applicazioni costanti.
    \[ c_{y_0} \simeq c_{y_1} \iff \exists \alpha: I \to Y \text{ cammino con } \alpha(0)=y_0, \alpha(1)=y_1 \]
    \textit{Dimostrazione} (si suppone $X \neq \emptyset$). ($\Rightarrow$) Per ipotesi $\exists H: c_{y_0} \simeq c_{y_1}$. Fissando $\bar{x} \in X$, definiamo $\alpha: I \to Y$ come $\alpha(t) = H(\bar{x}, t)$: \`e continuo, $\alpha(0) = c_{y_0}(\bar x) = y_0$ e $\alpha(1) = y_1$, quindi \`e il cammino cercato.
    ($\Leftarrow$) Dato $\alpha$ da $y_0$ a $y_1$, basta porre $H(x,t) = \alpha(t)$ per ogni $x \in X$ (continua, perch\'e non dipende da $x$).
\end{enumerate}


\section{Omotopia e deformazioni su $S^n$}

\begin{ejemplo}
Consideriamo la sfera $n$-dimensionale:
\[S^n = \left\{ x = (x_1, \dots, x_{n+1}) \in \mathbb{R}^{n+1} \mid x_1^2 + \dots + x_{n+1}^2 = 1 \right\} \subset \mathbb{R}^{n+1} \setminus \{\bar{0}\}\]
e la \textbf{proiezione radiale}:
\Func{\phi}{\mathbb{R}^{n+1} \setminus \{\bar{0}\}}{\mathbb{R}^{n+1} \setminus \{\bar{0}\}}{x}{\frac{x}{\|x\|}}
Allora $\phi \simeq \mathrm{Id}_{\mathbb{R}^{n+1} \setminus \{\bar{0}\}}$ (e $\phi$ ha immagine in $S^n$).

\begin{figure}[H]
    \centering
    \begin{tikzpicture}[scale=1.5]
        \draw[thick] (0,0) circle (1);
        \filldraw (0,0) circle (1.2pt) node[below left] {$\bar{0}$};
        \coordinate (X) at (1.8,1.0);
        \coordinate (PX) at ($(0,0)!1cm/2.059cm!(X)$);
        \draw[dashed,gray] (0,0) -- (X);
        \draw[very thick,red] (PX) -- (X);
        \filldraw (X) circle (1.2pt) node[above right] {$x$};
        \filldraw[blue] (PX) circle (1.2pt) node[above left] {$\phi(x)$};
        \node[red,below right] at (1.3,0.6) {\small $H(x,s)$};
        \node at (-0.7,-0.8) {$S^n$};
    \end{tikzpicture}
    \caption{La proiezione radiale spinge $x$ lungo la semiretta uscente dall'origine fino a $S^n$; il segmento rosso \`e la traiettoria dell'omotopia.}
\end{figure}
\end{ejemplo}

\begin{proof}
Definiamo $H: (\mathbb{R}^{n+1} \setminus \{\bar{0}\}) \times I \to \mathbb{R}^{n+1} \setminus \{\bar{0}\}$ come
\[
H(x,s) = (1-s)\,x + s\,\phi(x) = \left( (1-s) + \frac{s}{\|x\|} \right) x .
\]
Il coefficiente $\lambda(x,s) = (1-s) + s/\|x\|$ \`e una combinazione convessa di $1$ e $1/\|x\|$, entrambi $> 0$: dunque $\lambda(x,s) > 0$ e quindi $H(x,s) = \lambda(x,s)\,x \neq \bar{0}$ (in altre parole, \emph{il segmento che unisce $x$ a $\phi(x)$ non passa per l'origine}, perch\'e giace sulla semiretta $\{\lambda x \mid \lambda > 0\}$). Cos\`i $H$ \`e ben definita e continua, e
$H(x,0) = x = \mathrm{Id}(x)$, $H(x,1) = \phi(x)$. Quindi $H: \mathrm{Id} \simeq \phi$.
\end{proof}

\begin{teo}
L'omotopia relativamente ad $A$ \`e una relazione di equivalenza (sull'insieme delle applicazioni continue $X \to Y$).
\end{teo}

\begin{proof}
Siano $X, Y$ spazi topologici e $A \subseteq X$.

\textbf{Riflessiva.} Sia $f: X \to Y$ continua. Poniamo $H: X \times I \to Y$, $H(x,s) = f(x)$. $H$ \`e continua e $H(x,0) = f(x)$, $H(x,1) = f(x)$ per ogni $x$, e $H(x,s) = f(x)$ per ogni $x \in A$ e $s \in I$. Dunque $H: f \simeq f$ rel $A$.

\textbf{Simmetrica.} Siano $f, g: X \to Y$ continue e $H: f \simeq g$ rel $A$. Poniamo $K: X \times I \to Y$, $K(x,s) = H(x, 1-s)$. $K$ \`e continua (omotopia) e
\begin{enumerate}
    \item[i)] $K(x,0) = H(x,1) = g(x)$ per ogni $x \in X$;
    \item[ii)] $K(x,1) = H(x,0) = f(x)$ per ogni $x \in X$;
    \item[iii)] per ogni $x \in A$ e $s \in I$: $K(x,s) = H(x,1-s) = f(x) = g(x)$.
\end{enumerate}
Dunque $K: g \simeq f$ rel $A$.

\textbf{Transitiva.} Siano $f, g, g': X \to Y$ continue con $H: f \simeq g$ rel $A$ e $K: g \simeq g'$ rel $A$. Poniamo
\[
W: X \times I \to Y, \qquad W(x,s) = \begin{cases} H(x, 2s) & 0 \leqslant s \leqslant \tfrac12, \\[1mm] K(x, 2s-1) & \tfrac12 \leqslant s \leqslant 1. \end{cases}
\]
Osserviamo che per $s = \frac12$: $H(x, 1) = g(x) = K(x, 0)$, quindi $W$ \`e ben posta e, per il Lemma di Incollamento applicato ai chiusi $C_1 = X \times [0, \frac12]$ e $C_2 = X \times [\frac12, 1]$, \`e continua (\`e un'omotopia). Inoltre:
\begin{enumerate}
    \item[i)] $W(x,0) = H(x,0) = f(x)$ per ogni $x$;
    \item[ii)] $W(x,1) = K(x,1) = g'(x)$ per ogni $x$;
    \item[iii)] per $x \in A$ e $s \in I$: $W(x,s) = H(x,2s) = f(x)$ se $s \leqslant \frac12$ e $W(x,s) = K(x,2s-1) = g(x) = f(x)$ se $s \geqslant \frac12$.
\end{enumerate}
Dunque $W: f \simeq g'$ rel $A$\hfill$\square$
\end{proof}

\begin{intuizione}
La riflessiva \`e ``non deformare''; la simmetrica \`e ``percorrere la deformazione al contrario''; la transitiva \`e ``fare una deformazione e poi l'altra'' (in met\`a tempo ciascuna). \`E la stessa idea con cui si concatenano i cammini.
\end{intuizione}


\begin{teo}[Compatibilità dell'omotopia con la composizione di funzioni]
Siano $X, Y, Z$ spazi topologici, $f, g: X \to Y$ continue, $f', g': Y \to Z$ continue, $A \subseteq X$ e $B \subseteq Y$ con $f(A) \subseteq B$. Allora:
\[
f \simeq g \text{ rel } A \quad \text{e} \quad f' \simeq g' \text{ rel } B \implies f' \circ f \simeq g' \circ g \text{ rel } A
\]
\end{teo}

\begin{proof}
Siano $H: f \simeq g \text{ rel } A$ e $K: f' \simeq g' \text{ rel } B$ le rispettive omotopie. Definiamo la mappa:
\Func{W}{X \times I}{Z}{(x,s)}{K(H(x,s), s)}

Poiché $W$ è la composizione di funzioni continue, $W$ è continua (e quindi un'omotopia). Verifichiamo le tre proprietà:
\begin{enumerate}[label=\roman*)]
    \item $W(x,0) = K(H(x,0), 0) = K(f(x), 0) = (f' \circ f)(x)$ per ogni $x \in X$.
    \item $W(x,1) = K(H(x,1), 1) = K(g(x), 1) = (g' \circ g)(x)$ per ogni $x \in X$.
    \item Per ogni $x \in A$ e ogni $s \in I$:
    Poiché $f(A) \subseteq B$ e $H$ è un'omotopia relativa ad $A$, si ha $H(x,s) = f(x) \in B$. Inoltre, essendo $K$ un'omotopia relativa a $B$, per ogni punto $y \in B$ vale $K(y,s) = f'(y)$. Segue che:
    \[
    W(x,s) = K(H(x,s), s) = K(f(x), s) = f'(f(x)) = (f' \circ f)(x)
    \]
\end{enumerate}
Pertanto $W: f' \circ f \simeq g' \circ g \text{ rel } A$.
\end{proof}

\vspace{1.5em}

\section{Omotopia e campi di vettori tangenti su $S^n$}

\begin{prop}
Siano $X$ uno spazio topologico, $n \geqslant 1$ e $f, g: X \to S^n$ mappe continue. Se $f(x) \neq -g(x)$ per ogni $x \in X$, allora $f \simeq g$.
\end{prop}

% Nota en español: Esta proposición asegura que si dos aplicaciones a la esfera nunca toman valores antípodas, son homotópicas.
% La idea geométrica es trazar el segmento directo en R^{n+1} entre f(x) y g(x) y proyectarlo radialmente sobre S^n.

\begin{proof}
Anzitutto, dimostriamo che il segmento $\overline{f(x)g(x)}$ non passa per l'origine $\bar{0}$ di $\mathbb{R}^{n+1}$. 

Infatti, se per assurdo esistessero $x \in X$ e $s \in [0,1]$ tali che $(1-s)f(x) + s\,g(x) = \bar{0}$, si avrebbe $(1-s)f(x) = -s\,g(x)$. Passando alle norme (e usando $\|f(x)\| = \|g(x)\| = 1$) si ottiene $1-s = s$, cio\`e $s = 1/2$, e quindi $f(x) = -g(x)$, il che contraddice l'ipotesi.

Definiamo quindi l'omotopia $H: X \times I \to S^n$ tramite la proiezione radicale $\phi$:
\Func{H}{X \times I}{S^n}{(x,s)}{\frac{(1-s)f(x) + s\,g(x)}{\|(1-s)f(x) + s\,g(x)\|} = \phi\big((1-s)f(x) + s\,g(x)\big)}

$H$ è ben posta e continua in quanto il denominatore è sempre non nullo. Inoltre:
\begin{enumerate}[label=\roman*)]
    \item $H(x,0) = \frac{f(x)}{\|f(x)\|} = f(x), \quad \forall x \in X$.
    \item $H(x,1) = \frac{g(x)}{\|g(x)\|} = g(x), \quad \forall x \in X$.
\end{enumerate}
Quindi $H: f \simeq g$.
\end{proof}

\begin{coro}
Sia $f: S^n \to S^n$ una mappa continua.
\begin{enumerate}
    \item Se $f(x) \neq x$ per ogni $x \in S^n$, allora $f \simeq -\mathrm{Id}_{S^n}$.
    \item Se $f(x) \neq -x$ per ogni $x \in S^n$, allora $f \simeq \mathrm{Id}_{S^n}$.
\end{enumerate}
\end{coro}

\begin{prop}
Se $n$ è dispari, allora $\mathrm{Id}_{S^n} \simeq -\mathrm{Id}_{S^n}$.
\end{prop}

\begin{observacion}
Vale anche l'implicazione contraria "$\impliedby$": $\mathrm{Id}_{S^n} \simeq -\mathrm{Id}_{S^n} \implies n$ è dispari (senza dimostrazione: richiede strumenti pi\`u fini, come il grado di un'applicazione o l'omologia, dove $\deg(-\mathrm{Id}_{S^n}) = (-1)^{n+1} \neq 1$ per $n$ pari). Questo fatto \`e usato nel Teorema della palla pelosa e nel fatto che $S^n$ non \`e contraibile per $n$ pari.
\end{observacion}

\begin{proof}
Sia $n$ dispari. Allora esiste $k \geqslant 1$ tale che $n = 2k - 1$. Possiamo identificare $\mathbb{R}^{n+1} = \mathbb{R}^{2k} \cong \mathbb{C}^k$, quindi:
\[
S^n = S^{2k-1} = \left\{ z = (z_1, \dots, z_k) \in \mathbb{C}^k \mid |z|^2 = 1 \right\}
\]
Definiamo l'omotopia $H: S^n \times I \to S^n$ ruotando le coordinate complesse:
\Func{H}{S^n \times I}{S^n}{(z,s)}{e^{i\pi s} z = \big(e^{i\pi s}z_1, \dots, e^{i\pi s}z_k\big)}
dove $e^{i\pi s} = \cos(\pi s) + i\sin(\pi s)$.

$H$ è ben definita e continua. Infatti, per ogni $z \in S^n$ e $s \in I$:
\[
|H(z,s)| = |e^{i\pi s} z| = |e^{i\pi s}| \cdot |z| = 1 \cdot 1 = 1
\]
il che mostra che l'immagine risiede in $S^n$. Inoltre:
\begin{enumerate}[label=\roman*)]
    \item $H(z,0) = e^0 z = z = \mathrm{Id}_{S^n}(z)$ per ogni $z \in S^n$.
    \item $H(z,1) = e^{i\pi} z = -z = -\mathrm{Id}_{S^n}(z)$ per ogni $z \in S^n$.
\end{enumerate}
Di conseguenza $\mathrm{Id}_{S^n} \simeq -\mathrm{Id}_{S^n}$.
\end{proof}

\vspace{1.5em}

\begin{definicion}
Un \textbf{campo di vettori tangenti} su $S^n$ è un'applicazione $V: S^n \to \mathbb{R}^{n+1}$ tale che:
\[
\langle V(x), x \rangle = 0, \quad \forall x \in S^n
\]
Un campo $V$ si dice \textbf{continuo} se l'applicazione $V: S^n \to \mathbb{R}^{n+1}$ è continua.
\end{definicion}

\begin{figure}[H]
    \centering
    \begin{tikzpicture}[scale=1.2]
        % Esfera
        \draw[thick] (0,0) circle (1.2);
        \draw[dashed] (-1.2,0) arc (180:0:1.2 and 0.4);
        \draw (1.2,0) arc (0:-180:1.2 and 0.4);
        
        % Punto y plano tangente
        \coordinate (P) at (0.5, 0.7);
        \filldraw[black] (P) circle (1.5pt) node[below left] {$x$};
        
        % Vector normal (radio x)
        \draw[->,>=stealth,thick,blue] (0,0) -- (P) node[midway,above left] {$x$};
        
        % Plano tangente en el punto
        \draw[gray!70, fill=gray!20, opacity=0.6] (0.1, 1.1) -- (1.3, 0.9) -- (0.9, 0.3) -- (-0.3, 0.5) -- cycle;
        
        % Vector tangente V(x)
        \draw[->,>=stealth,ultra thick,red] (P) -- (1.0, 1.0) node[above right] {$V(x)$};
    \end{tikzpicture}
    \caption{Campo di vettori tangenti $V(x) \perp x$ su $S^n$.}
\end{figure}

\begin{prop}
Se esiste un campo continuo di vettori tangenti mai nullo $V: S^n \to \mathbb{R}^{n+1}$ (ossia con $V(x) \neq \bar{0}, \, \forall x \in S^n$), allora $\mathrm{Id}_{S^n} \simeq -\mathrm{Id}_{S^n}$.
\end{prop}

\begin{proof}
Poiché $V(x) \neq \bar{0}$ per ogni $x \in S^n$, possiamo definire la mappa normalizzata:
\Func{\widetilde{V}}{S^n}{S^n}{x}{\frac{V(x)}{\|V(x)\|}}
$\widetilde{V}$ è ben posta e continua. 

Osserviamo che $\widetilde{V}(x) \neq x$ per ogni $x \in S^n$. Infatti, se per assurdo esistesse $\bar{x} \in S^n$ tale che $\widetilde{V}(\bar{x}) = \bar{x}$, avremmo:
\[
1 = \langle \bar{x}, \bar{x} \rangle = \langle \widetilde{V}(\bar{x}), \bar{x} \rangle = \left\langle \frac{V(\bar{x})}{\|V(\bar{x})\|}, \bar{x} \right\rangle = \frac{1}{\|V(\bar{x})\|} \langle V(\bar{x}), \bar{x} \rangle = 0
\]
il che è una contraddizione ($1 = 0$). Pertanto $\widetilde{V}(x) \neq x$ per ogni $x \in S^n$.

Per il Corollario 1, si ha che $\widetilde{V} \simeq -\mathrm{Id}_{S^n}$.

Con un argomento analogo (dimostrando che $\widetilde{V}(x) \neq -x$ per ogni $x \in S^n$):
\[
\langle \widetilde{V}(x), x \rangle = \frac{1}{\|V(x)\|} \langle V(x), x \rangle = 0 \neq -1 = \langle -x, x \rangle
\]
si ottiene, per il Corollario 2, che $\widetilde{V} \simeq \mathrm{Id}_{S^n}$.

Dalla proprietà transitiva della relazione di omotopia segue che:
\[
\mathrm{Id}_{S^n} \simeq \widetilde{V} \simeq -\mathrm{Id}_{S^n} \implies \mathrm{Id}_{S^n} \simeq -\mathrm{Id}_{S^n}
\]
\end{proof}

\begin{prop}
Se $n$ è dispari, allora esiste un campo continuo di vettori tangenti mai nullo $V: S^n \to \mathbb{R}^{n+1}$ (ossia $V(x) \neq \bar{0}$ per ogni $x \in S^n$).
\end{prop}

\begin{proof}
Sia $n$ dispari. Allora $n = 2k-1$ per qualche $k \geqslant 1$. Identifichiamo $\mathbb{R}^{n+1} = \mathbb{R}^{2k}$ come la coppia di vettori $(x_1, \dots, x_k, y_1, \dots, y_k) \in \mathbb{R}^k \times \mathbb{R}^k$. Pertanto:
\[
S^n = S^{2k-1} = \left\{ (x_1, \dots, x_k, y_1, \dots, y_k) \in \mathbb{R}^{2k} \;\middle|\; \sum_{i=1}^k x_i^2 + \sum_{i=1}^k y_i^2 = 1 \right\}
\]
Definiamo la mappa $V: S^n \to \mathbb{R}^{2k}$ mediante:
\Func{V}{S^n}{\mathbb{R}^{2k}}{(x_1, \dots, x_k, y_1, \dots, y_k)}{(-y_1, \dots, -y_k, x_1, \dots, x_k)}

Verifichiamo le proprietà del campo di vettori:
\begin{enumerate}[label=\roman*)]
    \item $V$ è continua poiché le sue componenti sono applicazioni lineari (quindi continue) delle coordinate.
    \item $V$ è ortogonale a $x$ per ogni $x \in S^n$:
    \[
    \langle x, V(x) \rangle = \sum_{i=1}^k x_i(-y_i) + \sum_{i=1}^k y_i(x_i) = -\sum_{i=1}^k x_i y_i + \sum_{i=1}^k x_i y_i = 0
    \]
    Quindi $V$ è un campo continuo di vettori tangenti su $S^n$.
    \item La norma del campo è pari alla norma di $x$:
    \[
    \|V(x)\|^2 = \sum_{i=1}^k (-y_i)^2 + \sum_{i=1}^k x_i^2 = \sum_{i=1}^k x_i^2 + \sum_{i=1}^k y_i^2 = \|x\|^2 = 1
    \]
    Di conseguenza $\|V(x)\| = 1$, il che implica che $V(x) \neq \bar{0}$ per ogni $x \in S^n$.
\end{enumerate}
\end{proof}

\begin{figure}[H]
    \centering
    \begin{tikzpicture}[scale=1.2]
        % S^1
        \draw[thick] (-2.5,0) circle (0.8);
        \node at (-2.5, -1.1) {$S^1$ (campo tangenti mai nullo)};
        \draw[->,>=stealth,red,thick] (-2.5,0.8) -- (-3.1,0.8);
        \draw[->,>=stealth,red,thick] (-2.5,-0.8) -- (-1.9,-0.8);
        \draw[->,>=stealth,red,thick] (-3.3,0) -- (-3.3,-0.6);
        \draw[->,>=stealth,red,thick] (-1.7,0) -- (-1.7,0.6);

        % S^2
        \draw[thick] (2.5,0) circle (0.8);
        \draw[dashed] (1.7,0) arc (180:0:0.8 and 0.25);
        \draw (2.5,0) arc (0:-180:0.8 and 0.25);
        \node at (2.5, -1.1) {$S^2$ (non ammette campi mai nulli)};
        \filldraw[red] (2.5,0.8) circle (1.5pt) node[above] {$\bar{0}$};
    \end{tikzpicture}
    \caption{Visualizzazione intuitiva del Teorema del Pettine (Hairy Ball Theorem).}
\end{figure}

\begin{observacion}
Ricapitolando i risultati ottenuti sulle sfere:
\[
n \text{ è dispari} \iff \mathrm{Id}_{S^n} \simeq -\mathrm{Id}_{S^n} \iff \exists\, V \text{ campo continuo di vettori tangenti mai nullo su } S^n
\]
In particolare, se $n$ è pari, ogni campo continuo di vettori tangenti $V$ su $S^n$ ammette almeno un punto di annullamento, ovvero esiste $x \in S^n$ tale che $V(x) = \bar{0}$ (Teorema della palla pelosa).
\end{observacion}

\vspace{1.5em}

\begin{prop}
Se $n$ è pari, allora per ogni applicazione continua $f: S^n \to S^n$ esiste un punto $x_0 \in S^n$ tale che $f(x_0) = x_0$ oppure $f(x_0) = -x_0$.
\end{prop}

% Explicación en español:
% Esta proposición asegura que cualquier mapa continuo en una esfera de dimensión par tiene al menos un punto fijo o un punto antipodal fijo.

\begin{proof}
Procediamo per assurdo. Supponiamo che per ogni $x \in S^n$ si abbia $f(x) \neq x$ e $f(x) \neq -x$.

Poiché $f(x) \neq -x$ per ogni $x \in S^n$, per il Corollario 2 della sezione precedente si ha che $f \simeq \mathrm{Id}_{S^n}$.
Analogamente, poiché $f(x) \neq x$ per ogni $x \in S^n$, per il Corollario 1 si ha che $f \simeq -\mathrm{Id}_{S^n}$.

Per la proprietà di simmetria e transitività dell'omotopia, si deduce che:
\[
\mathrm{Id}_{S^n} \simeq f \simeq -\mathrm{Id}_{S^n} \implies \mathrm{Id}_{S^n} \simeq -\mathrm{Id}_{S^n}
\]
Ciò implica che $n$ debba essere dispari, in contraddizione con l'ipotesi che $n$ sia pari.
\end{proof}

\vspace{2em}

\section{Spazi omotopicamente equivalenti}

\begin{definicion}
Due spazi topologici $X$ e $Y$ si dicono \textbf{topologicamente equivalenti} (o \textbf{omeomorfi}), e si scrive $X \cong Y$, se esiste un'applicazione $f: X \to Y$ continua, invertibile e con inversa $f^{-1}: Y \to X$ continua.

Equivalentemente, esistono $f: X \to Y$ e $g: Y \to X$ continue tali che:
\[
g \circ f = \mathrm{Id}_X \quad \text{e} \quad f \circ g = \mathrm{Id}_Y
\]
In tal caso $g = f^{-1}$.
\end{definicion}

\begin{definicion}
Due spazi topologici $X$ e $Y$ si dicono \textbf{omotopicamente equivalenti} (e si scrive $X \simeq Y$) se esistono mappe continue $f: X \to Y$ e $g: Y \to X$ tali che:
\[
g \circ f \simeq \mathrm{Id}_X \quad \text{e} \quad f \circ g \simeq \mathrm{Id}_Y
\]
\end{definicion}

\begin{observacion}
La relazione di equivalenza omotopica è un indebolimento del concetto di omeomorfismo:
\[
X \cong Y \implies X \simeq Y
\]
Il viceversa non vale in generale. Ad esempio, consideriamo $X = \mathbb{R}^n$ e $Y = \{y_0\}$. Poiché $\mathbb{R}^n$ non è in biezioni con un singolo punto per $n \geqslant 1$, $\mathbb{R}^n \not\cong \{y_0\}$. Tuttavia, sono omotopicamente equivalenti ($\mathbb{R}^n \simeq \{y_0\}$).

Infatti, definiamo le applicazioni continue:
\Func{f}{\mathbb{R}^n}{\{y_0\}}{x}{y_0}
\Func{g}{\{y_0\}}{\mathbb{R}^n}{y_0}{\bar{0}}

Analizziamo le composizioni:
\begin{itemize}
    \item $(f \circ g): \{y_0\} \to \{y_0\}$ invia $y_0 \mapsto \bar{0} \mapsto y_0$, dunque $f \circ g = \mathrm{Id}_{\{y_0\}}$, che è banalmente omotopa a $\mathrm{Id}_{\{y_0\}}$.
    \item $(g \circ f): \mathbb{R}^n \to \mathbb{R}^n$ invia $x \mapsto \bar{0}$. La mappa costante $g \circ f = \bar{0}$ è omotopa a $\mathrm{Id}_{\mathbb{R}^n}$ tramite l'omotopia lineare $H(x,s) = (1-s)x$.
\end{itemize}
Dunque $\mathbb{R}^n \simeq \{y_0\}$.
\end{observacion}

\begin{teo}
L'equivalenza omotopica tra spazi topologici è una relazione di equivalenza.
\end{teo}

\begin{proof}[Dimostrazione del Teorema]
Verifichiamo le tre proprietà per dimostrare che la relazione di equivalenza omotopica $\simeq$ è una relazione di equivalenza:

\begin{itemize}
    \item \textbf{Riflessività:} Per ogni spazio topologico $X$, si ha $X \simeq X$.
    Infatti, basta considerare l'applicazione identità $\mathrm{Id}_X: X \to X$. Risulta $\mathrm{Id}_X \circ \mathrm{Id}_X = \mathrm{Id}_X \simeq \mathrm{Id}_X$, quindi la proprietà segue direttamente dalla definizione.

    \item \textbf{Simmetria:} Siano $X, Y$ spazi topologici tali che $X \simeq Y$.
    Per definizione esistono due mappe continue $f: X \to Y$ e $g: Y \to X$ tali che $g \circ f \simeq \mathrm{Id}_X$ e $f \circ g \simeq \mathrm{Id}_Y$. Scambiando i ruoli di $f$ e $g$ si ottiene immediatamente che $Y \simeq X$.

    \item \textbf{Transitività:} Siano $X, Y, Z$ spazi topologici tali che $X \simeq Y$ e $Y \simeq Z$.
    
    Per ipotesi esistono mappe continue:
    \[
    X \xleftrightarrow[g]{f} Y \quad \text{con } g \circ f \simeq \mathrm{Id}_X \text{ e } f \circ g \simeq \mathrm{Id}_Y
    \]
    \[
    Y \xleftrightarrow[g']{f'} Z \quad \text{con } g' \circ f' \simeq \mathrm{Id}_Y \text{ e } f' \circ g' \simeq \mathrm{Id}_Z
    \]
    
    Consideriamo le composizioni $f' \circ f: X \to Z$ e $g \circ g': Z \to X$. Verifichiamo le due omotopie necessarie:
    
    \begin{enumerate}[label=\roman*)]
        \item Per la composizione $(g \circ g') \circ (f' \circ f)$:
        \[
        (g \circ g') \circ (f' \circ f) = g \circ (g' \circ f') \circ f \simeq g \circ \mathrm{Id}_Y \circ f = g \circ f \simeq \mathrm{Id}_X
        \]
        \item Per la composizione $(f' \circ f) \circ (g \circ g')$:
        \[
        (f' \circ f) \circ (g \circ g') = f' \circ (f \circ g) \circ g' \simeq f' \circ \mathrm{Id}_Y \circ g' = f' \circ g' \simeq \mathrm{Id}_Z
        \]
    \end{enumerate}
    Di conseguenza $X \simeq Z$.
\end{itemize}
\end{proof}

\vspace{1.5em}

\begin{ejemplo}
Analizziamo alcuni esempi fondamentali di equivalenza omotopica:

\begin{enumerate}
    \item \textbf{Spazi convessi:} Se $\Omega \subseteq \mathbb{R}^n$ è un insieme convesso non vuoto e $x_0 \in \Omega$, allora $\Omega \simeq \{x_0\}$.
    
    % Nota en español: Todo conjunto convexo en R^n es contraíble a cualquiera de sus puntos mediante la homotopía en línea recta H(x,s) = (1-s)x + s*x_0.

    \item \textbf{Spazio euclideo privato dell'origine e la sfera $n$-dimensionale:}
    \[
    \mathbb{R}^{n+1} \setminus \{\bar{0}\} \simeq S^n
    \]
    
    Consideriamo le mappe continue:
    \Func{f}{\mathbb{R}^{n+1} \setminus \{\bar{0}\}}{S^n}{x}{\frac{x}{\|x\|} = \phi(x)}
    e l'inclusione canonica $g = i: S^n \hookrightarrow \mathbb{R}^{n+1} \setminus \{\bar{0}\}$ con $g(x) = x$.

    Analizziamo le due composizioni:
    \begin{itemize}
        \item $(g \circ f)(x) = g(\phi(x)) = \phi(x)$. Come dimostrato in precedenza, $\phi \simeq \mathrm{Id}_{\mathbb{R}^{n+1} \setminus \{\bar{0}\}}$ tramite la proiezione radicale.
        \item $(f \circ g)(x) = f(x) = \frac{x}{\|x\|} = x = \mathrm{Id}_{S^n}(x)$ per ogni $x \in S^n$ (poiché $\|x\| = 1$). Pertanto $f \circ g = \mathrm{Id}_{S^n} \simeq \mathrm{Id}_{S^n}$.
    \end{itemize}
    Dunque $\mathbb{R}^{n+1} \setminus \{\bar{0}\} \simeq S^n$.

    \item \textbf{Palla unitaria piena e il suo bordo (sfera):}
    Sia $B^n = \{x \in \mathbb{R}^n \mid \|x\| \leqslant 1\}$ la palla chiusa e $\partial B^n = \{x \in \mathbb{R}^n \mid \|x\| = 1\} = S^{n-1}$ il suo bordo.
    
    Allora, $B^{n+1} \setminus \{\bar{0}\} \simeq \partial B^{n+1} = S^n$.

    Consideriamo le mappe continue:
    \Func{f}{B^{n+1} \setminus \{\bar{0}\}}{S^n}{x}{\frac{x}{\|x\|}}
    e l'inclusione $g = i: S^n \hookrightarrow B^{n+1} \setminus \{\bar{0}\}$ dato da $g(x) = x$. Con lo stesso argomento della proiezione radicale si conclude che $B^{n+1} \setminus \{\bar{0}\} \simeq S^n$.
\end{enumerate}
\end{ejemplo}

\vspace{1.5em}

\begin{observacion}
Dalla transitività dell'equivalenza omotopica, se due spazi $X$ e $Y$ sono entrambi omotopicamente equivalenti a uno stesso spazio $Z$, allora $X \simeq Y$.

In particolare, consideriamo la sfera senza un punto: $S^n \setminus \{P\}$. Dimostreremo che $S^n \setminus \{P\} \cong \mathbb{R}^n$ e di conseguenza $S^n \setminus \{P\} \simeq \mathbb{R}^n \simeq \{x_0\}$.
\end{observacion}

\begin{prop}
La sfera $n$-dimensionale privata del Polo Nord $N = (0, \dots, 0, 1) \in \mathbb{R}^{n+1}$ è omeomorfa allo spazio euclideo $\mathbb{R}^n$:
\[
S^n \setminus \{N\} \cong \mathbb{R}^n
\]
\end{prop}

\begin{proof}
Consideriamo la \textbf{proiezione stereografica} dal Polo Nord $N = (0, \dots, 0, 1)$ sull'iperpiano $x_{n+1} = 0$, identificato con $\mathbb{R}^n$:
\Func{\varphi}{S^n \setminus \{N\}}{\mathbb{R}^n}{(x_1, \dots, x_n, x_{n+1})}{\left( \frac{x_1}{1 - x_{n+1}}, \, \dots, \, \frac{x_n}{1 - x_{n+1}} \right)}

Geometricamente, $\varphi(x)$ rappresenta l'unica intersezione dell'iperpiano $x_{n+1} = 0$ con la semiretta con origine in $N$ e passante per il punto $x \in S^n \setminus \{N\}$.

\begin{figure}[H]
    \centering
    \begin{tikzpicture}[scale=1.4]
        \draw[gray!60, fill=gray!8] (-2.6,-0.5) -- (2.4,-0.5) -- (3.4,0.5) -- (-1.6,0.5) -- cycle;
        \node[gray] at (-2.0,-0.25) {$\mathbb{R}^n$};
        \draw[thick] (0,0) circle (1.2);
        \draw[dashed] (-1.2,0) arc (180:0:1.2 and 0.3);
        \draw (1.2,0) arc (0:-180:1.2 and 0.3);
        \coordinate (N) at (0,1.2);
        \filldraw[blue] (N) circle (1.5pt) node[above] {$N$};
        \coordinate (X) at (0.983,0.688);
        \filldraw (X) circle (1.3pt) node[above right] {$x$};
        \draw[red, thick] (N) -- (2.304,0);
        \filldraw[red] (2.304,0) circle (1.5pt) node[below right] {$\varphi(x)$};
        \filldraw (0,-1.2) circle (1pt) node[below] {$S$};
    \end{tikzpicture}
    \caption{Proiezione stereografica dal Polo Nord $N$ sull'iperpiano equatoriale $x_{n+1}=0$ ($\cong \mathbb{R}^n$): $\varphi(x)$ \`e l'intersezione della retta per $N$ e $x$ con l'iperpiano.}
\end{figure}

L'applicazione $\varphi$ è continua e invertibile. La sua inversa $\varphi^{-1}: \mathbb{R}^n \to S^n \setminus \{N\}$ è data da:
\Func{\varphi^{-1}}{\mathbb{R}^n}{S^n \setminus \{N\}}{(y_1, \dots, y_n)}{\left( \frac{2y_1}{\|y\|^2 + 1}, \, \dots, \, \frac{2y_n}{\|y\|^2 + 1}, \, \frac{\|y\|^2 - 1}{\|y\|^2 + 1} \right)}
dove $\|y\|^2 = y_1^2 + \dots + y_n^2$.

Poiché sia $\varphi$ che $\varphi^{-1}$ sono continue, $\varphi$ è un omeomorfismo. Pertanto $S^n \setminus \{N\} \cong \mathbb{R}^n$.
\end{proof}

\begin{observacion}
Consideriamo la proiezione stereografica dal Polo Sud $S = (0, 0, \dots, 0, -1) \in S^n$:
\[
\varphi(S) = (0, 0, \dots, 0) = \bar{0} \in \mathbb{R}^n
\]
Inoltre, l'equatore della sfera $S^n \cap \{x_{n+1} = 0\}$ equivale a $S^{n-1}$, e per ogni punto $x \in S^{n-1}$ vale $\varphi(x) = x$.

Di conseguenza:
\[
S^n \setminus \{N\} \cong \mathbb{R}^n \implies S^n \setminus \{N, S\} \cong \mathbb{R}^n \setminus \{\bar{0}\}
\]
E poiché $\mathbb{R}^n \setminus \{\bar{0}\} \simeq S^{n-1}$, si ottiene che la sfera privata dei due poli è omotopicamente equivalente alla sfera di dimensione inferiore:
\[
S^n \setminus \{N, S\} \simeq S^{n-1}
\]
\end{observacion}

\vspace{1.5em}

\section{Retratti e Retratti di Deformazione}

\begin{definicion}
Sia $Y$ uno spazio topologico e $X \subseteq Y$. $X$ si dice \textbf{retratto} di $Y$ se esiste un'applicazione continua $r: Y \to X$ tale che:
\[
r \circ i = \mathrm{Id}_X \quad (\text{cioè } r(x) = x, \; \forall x \in X)
\]
L'applicazione $r$ si dice \textbf{retrazione}.
\end{definicion}

% Nota en español: Una retracción 'r' es una proyección continua de todo el espacio Y sobre su subespacio X que deja fijos a todos los puntos de X.

\begin{ejemplo}
\begin{enumerate}
    \item Per ogni spazio topologico $Y$ e per ogni punto $y_0 \in Y$, il singoletto $\{y_0\}$ è un retratto di $Y$.
    Infatti, l'applicazione costante $r: Y \to \{y_0\}$ definita da $r(y) = y_0$ è continua e soddisfa $r(y_0) = y_0$.

    \item In generale, il fatto che $X \subseteq Y$ \textbf{non} implica che $X$ sia un retratto di $Y$.
\end{enumerate}
\end{ejemplo}

\begin{ejemplo}[Controesempio]
Sia $Y = I = [0,1]$ e $X = \mathring{I} = (0,1)$.

Supponiamo per assurdo che esista una retrazione $r: [0,1] \to (0,1)$. Allora $r$ è continua e suriettiva. Poiché $[0,1]$ è un insieme compatto e $r$ è continua, l'immagine $r([0,1]) = (0,1)$ dovrebbe essere un sottoinsieme compatto di $\mathbb{R}$. Tuttavia, l'intervallo aperto $(0,1)$ non è compatto, il che genera una contraddizione.

\textit{Nota generale:} Se $Y$ è compatto (o connesso) e $X \subseteq Y$ non è compatto (o non connesso), allora $X$ non può essere un retratto di $Y$. Ad esempio, l'interno della palla aperta $X = \mathring{B}^n$ non è un retratto della palla chiusa $Y = B^n$.
\end{ejemplo}

\begin{prop}
Siano $Y, Z$ spazi topologici, $X \subseteq Y$ ed $r: Y \to X$ una retrazione. Allora per ogni applicazione continua $f: X \to Z$, esiste un'estensione continua $\bar{f}: Y \to Z$ tale che $\bar{f}\vert_X = f$.
\end{prop}

\begin{proof}
Basta definire $\bar{f} = f \circ r: Y \to Z$. Poiché è composizione di mappe continue, $\bar{f}$ è continua, e per ogni $x \in X$ si ha $\bar{f}(x) = f(r(x)) = f(x)$.
\end{proof}

\begin{observacion}
Essere un retratto \textbf{non} implica l'equivalenza omotopica: $X$ retratto di $Y \centernot\implies X \simeq Y$.

Come controesempio, si consideri $Y = S^1$ e $X = \{x_0\} \subset S^1$. Il punto $\{x_0\}$ è un retratto di $S^1$, ma $S^1 \not\simeq \{x_0\}$ in quanto hanno gruppo fondamentale diverso ($\pi_1(S^1) \cong \mathbb{Z} \neq \{0\} = \pi_1(\{x_0\})$).

Per avere un'equivalenza omotopica occorre che $r \circ i = \mathrm{Id}_X$ e anche $i \circ r \simeq \mathrm{Id}_Y$ (retratto di deformazione).
\end{observacion}

\vspace{1.5em}

\begin{ejemplo}[Esempi di retratti di deformazione]
\begin{enumerate}
    \item In generale, un singoletto $\{y_0\}$ \textbf{non} è un retratto di deformazione di $Y$ (ad esempio $\{x_0\} \subset S^1$).
    \item $\{x_0\} \subset \mathbb{R}^n$ è un retratto di deformazione di $\mathbb{R}^n$. Più in generale, se $\Omega \subseteq \mathbb{R}^n$ è un insieme convesso, allora ogni punto $\{x_0\} \subset \Omega$ è un retratto di deformazione di $\Omega$ tramite $r(x) = x_0$.
    \item $S^n$ è un retratto di deformazione di $\mathbb{R}^{n+1} \setminus \{\bar{0}\}$ tramite la proiezione radicale:
    \Func{r}{\mathbb{R}^{n+1} \setminus \{\bar{0}\}}{S^n}{x}{\frac{x}{\|x\|}}
    \item $S^n = \partial B^{n+1} \setminus \{\bar{0}\}$ è un retratto di deformazione di $B^{n+1} \setminus \{\bar{0}\}$.
    \item Per ogni $\delta > 0$, la palla chiusa di raggio $\delta$, $B_\delta^n = \{x \in \mathbb{R}^n \mid \|x\| \leqslant \delta\}$, è un retratto di deformazione di $\mathbb{R}^n$.
\end{enumerate}
\end{ejemplo}

\begin{proof}[Dimostrazione dell'Esempio 5]
Definiamo la mappa $r: \mathbb{R}^n \to B_\delta^n$ come la proiezione radiale sulla palla:
\[
r(x) = 
\begin{cases}
    x & \text{se } \|x\| \leqslant \delta \\
    \delta \dfrac{x}{\|x\|} & \text{se } \|x\| \geqslant \delta
\end{cases}
\]

\begin{figure}[H]
    \centering
    \begin{tikzpicture}[scale=1.2]
        % Plano R^n
        \draw[gray!50, fill=gray!5] (-2.5,-1.5) -- (2.5,-1.5) -- (2.5,1.5) -- (-2.5,1.5) -- cycle;
        \node[above left] at (2.5,-1.5) {$\mathbb{R}^n$};

        % Bola B_\delta^n
        \draw[thick, fill=blue!10] (0,0) circle (0.9);
        \node at (0,0) {$B_\delta^n$};

        % Puntos y proyecciones
        \filldraw (1.8,0.8) circle (1.2pt) node[right] {$x$};
        \filldraw (0.72,0.32) circle (1.2pt) node[below right] {$r(x)$};
        \draw[->,>=stealth,red,thick] (1.8,0.8) -- (0.72,0.32);

        \draw[dashed] (0,0) -- (1.8,0.8);
    \end{tikzpicture}
    \caption{Proiezione radiale di $\mathbb{R}^n$ sulla palla $B_\delta^n$.}
\end{figure}

Sulla frontiera $\|x\| = \delta$, le due definizioni coincidono: $\delta \frac{x}{\|x\|} = \delta \frac{x}{\delta} = x$. Per il \textit{Lemma di Incollamento}, $r$ è continua ed è ben posta.

Inoltre, $r(x) = x$ per ogni $x \in B_\delta^n$, e l'omotopia lineare tra $i \circ r$ e $\mathrm{Id}_{\mathbb{R}^n}$ mostra che $i \circ r \simeq \mathrm{Id}_{\mathbb{R}^n}$. Di conseguenza, $B_\delta^n$ è un retratto di deformazione di $\mathbb{R}^n$.
\end{proof}

\begin{ejemplo}
$S^{n-1}$ è un retratto di deformazione di $S^n \setminus \{N, S\}$.

\begin{figure}[H]
    \centering
    \begin{tikzpicture}[scale=1.3]
        % Esfera
        \draw[thick] (0,0) circle (1.2);
        \draw[dashed] (-1.2,0) arc (180:0:1.2 and 0.3);
        \draw[thick, blue] (1.2,0) arc (0:-180:1.2 and 0.3);
        \node[left, blue] at (-1.2, 0) {$S^{n-1}$};

        % Polos
        \filldraw[red] (0,1.2) circle (1.5pt) node[above] {$N$};
        \filldraw[red] (0,-1.2) circle (1.5pt) node[below] {$S$};

        % Ejes
        \draw[->,>=stealth,gray] (0,-1.5) -- (0,1.5);
        \draw[->,>=stealth,gray] (-1.5,0) -- (1.5,0);

        % Punto x y proyeccion
        \coordinate (X) at (0.6, 0.7);
        \filldraw (X) circle (1.2pt) node[above right] {$x$};
        
        \draw[->,>=stealth,thick,orange] (X) -- (0.9, 0.2) node[right] {$\varphi(x)$};
        \draw[->,>=stealth,thick,purple] (0.9, 0.2) -- (1.15, 0) node[below] {$\frac{\varphi(x)}{\|\varphi(x)\|}$};
    \end{tikzpicture}
    \caption{Deformazione di $S^n \setminus \{N,S\}$ su $S^{n-1}$ tramite proiezione stereografica e normalizzazione.}
\end{figure}

Geometricamente, questo si ottiene combinando la proiezione stereografica $\varphi: S^n \setminus \{N\} \to \mathbb{R}^n$, che invia $S^n \setminus \{N, S\}$ su $\mathbb{R}^n \setminus \{\bar{0}\}$, con la proiezione radicale da $\mathbb{R}^n \setminus \{\bar{0}\}$ sull'equatore $S^{n-1}$ definita da $\frac{\varphi(x)}{\|\varphi(x)\|}$.
\end{ejemplo}

\begin{proof}[Continuazione della dimostrazione]
Definiamo la mappa di retrazione:
\Func{r}{S^n \setminus \{N,S\}}{S^{n-1} = S^n \cap \{x_{n+1} = 0\}}{x}{\frac{\varphi(x)}{\|\varphi(x)\|}}

$r$ è continua e ben posta, in quanto la proiezione stereografica $\varphi$ è continua e $\varphi(x) \neq \bar{0}$ per ogni $x \in S^n \setminus \{N,S\}$.
Inoltre, per ogni $x \in S^{n-1}$, si ha $\varphi(x) = x \implies \frac{\varphi(x)}{\|\varphi(x)\|} = \frac{x}{\|x\|} = x$. Quindi $r$ è effettivamente una retrazione.

% Nota en español: Para probar que S^{n-1} es un retrato de deformación fuerte de S^n \ {N,S}, definimos la homotopía lineal normalizada entre la identidad e i o r.

Consideriamo ora la composizione $i \circ r: S^n \setminus \{N,S\} \to S^n \setminus \{N,S\}$ definita da $x \mapsto \frac{\varphi(x)}{\|\varphi(x)\|}$. Definiamo l'omotopia:
\Func{H}{(S^n \setminus \{N,S\}) \times I}{S^n \setminus \{N,S\}}{(x,s)}{\frac{(1-s)x + s\,r(x)}{\|(1-s)x + s\,r(x)\|}}

Scriviamo $x = (x', x_{n+1})$ con $x' = (x_1, \dots, x_n)$. Poich\'e $x \neq N, S$, si ha $x_{n+1} \neq \pm 1$, quindi $x' \neq \bar{0}$. Inoltre $\varphi(x) = x'/(1 - x_{n+1})$ \`e un multiplo \emph{positivo} di $x'$, per cui
\[
r(x) = \frac{\varphi(x)}{\|\varphi(x)\|} = \frac{(x', 0)}{\|x'\|} .
\]
Verifichiamo che $H$ sia ben posta e continua, ossia che il denominatore non si annulli mai:
\[
(1-s)x + s\,r(x) = \Big( \underbrace{\big(1-s + \tfrac{s}{\|x'\|}\big)}_{>0}\, x',\;\; (1-s)\,x_{n+1} \Big) \neq \bar{0}
\]
perch\'e la prima componente \`e un multiplo positivo del vettore non nullo $x'$. Inoltre il punto $H(x,s)$ ha ancora la parte orizzontale non nulla, quindi non \`e $N$ n\'e $S$: $H$ prende valori in $S^n \setminus \{N,S\}$.

Pertanto $H$ è ben posta, continua e realizza l'omotopia $\mathrm{Id}_{S^n \setminus \{N,S\}} \simeq i \circ r$. Di conseguenza, $S^{n-1}$ è un retratto di deformazione di $S^n \setminus \{N,S\}$.
\end{proof}

\vspace{1.5em}

\begin{ejemplo}
Consideriamo la sfera $S^n$ e un punto fissato $x_0 \in S^n$.
Osserviamo che l'equatore rispetto a $x_0$ è dato da:
\[
S^{n-1} = \{x \in S^n \mid \langle x, x_0 \rangle = 0\} = S^n \cap x_0^\perp
\]
Fissiamo $\varepsilon \in (0, 1)$ e definiamo la fascia equatoriale:
\[
W_\varepsilon = \{x \in S^n \mid |\langle x, x_0 \rangle| < \varepsilon\}
\]

\begin{figure}[H]
    \centering
    \begin{tikzpicture}[scale=1.5]
        \begin{scope}
            \clip (0,0) circle (1.2);
            \fill[red!15] (-1.3,-0.4) rectangle (1.3,0.4);
        \end{scope}
        \draw[thick] (0,0) circle (1.2);
        \draw[blue, dashed] (-1.2,0) arc (180:0:1.2 and 0.3);
        \draw[blue, thick] (-1.2,0) arc (180:360:1.2 and 0.3);
        \draw[red, dashed] (-1.131,0.4) arc (180:0:1.131 and 0.28);
        \draw[red] (-1.131,0.4) arc (180:360:1.131 and 0.28);
        \draw[red, dashed] (-1.131,-0.4) arc (180:0:1.131 and 0.28);
        \draw[red] (-1.131,-0.4) arc (180:360:1.131 and 0.28);
        \draw[->,>=stealth,thick] (0,0) -- (0,1.55) node[above] {$x_0$};
        \filldraw (0,1.2) circle (1.2pt);
        \filldraw (0,-1.2) circle (1.2pt) node[below] {$-x_0$};
        \node[blue, right] at (1.2,0) {$S^{n-1}$};
        \node[red, right] at (1.1,0.55) {$W_\varepsilon$};
    \end{tikzpicture}
    \caption{La fascia equatoriale $W_\varepsilon = \{|\langle x, x_0\rangle| < \varepsilon\}$ (in rosa) contiene l'equatore $S^{n-1}$ e non contiene i poli $\pm x_0$.}
\end{figure}

Osserviamo che:
\begin{enumerate}[label=\roman*)]
    \item $S^{n-1} = S^n \cap x_0^\perp \subset W_\varepsilon$.
    \item $|\langle \pm x_0, x_0 \rangle| = 1 > \varepsilon \implies \pm x_0 \notin W_\varepsilon$.
\end{enumerate}
Allora $S^{n-1}$ è un retratto di deformazione di $W_\varepsilon$.
\end{ejemplo}

\begin{proof}
Definiamo l'applicazione di proiezione ortogonale normalizzata:
\Func{r}{W_\varepsilon}{S^{n-1}}{x}{\frac{x - \langle x, x_0 \rangle x_0}{\|x - \langle x, x_0 \rangle x_0\|}}

\begin{figure}[H]
    \centering
    \begin{tikzpicture}[scale=1.2]
        % Circulo
        \draw[thick] (0,0) circle (1.2);
        
        % Vector x_0 y x
        \draw[->,>=stealth,thick,blue] (0,0) -- (1.2,0) node[right] {$x_0$};
        \draw[->,>=stealth,thick] (0,0) -- (0.7,0.97) node[above right] {$x$};
        
        % Proyeccion
        \draw[dashed, gray] (0.7,0.97) -- (0,0.97);
        \draw[->,>=stealth,thick,red] (0,0) -- (0,0.97) node[left] {$x - \langle x, x_0 \rangle x_0$};
        \draw[->,>=stealth,ultra thick,purple] (0,0) -- (0,1.2) node[above left] {$r(x)$};
    \end{tikzpicture}
    \caption{Proiezione ortogonale di $x$ sull'equatore $S^{n-1}$.}
\end{figure}

Verifichiamo le proprietà:
\begin{enumerate}[label=\roman*)]
    \item Il denominatore non si annulla mai in $W_\varepsilon$: infatti, $x - \langle x, x_0 \rangle x_0 = \bar{0} \implies x = \langle x, x_0 \rangle x_0 \implies x \mathbin{\|} x_0 \implies x = \pm x_0$. Tuttavia, $\pm x_0 \notin W_\varepsilon$, quindi il denominatore è sempre strettamente positivo e $r$ è continua e ben posta.
    \item $r(x) \in S^{n-1}$, in quanto:
    \[
    \langle r(x), x_0 \rangle = \frac{1}{\|x - \langle x, x_0 \rangle x_0\|} \Big( \langle x, x_0 \rangle - \langle x, x_0 \rangle \underbrace{\langle x_0, x_0 \rangle}_{=1} \Big) = 0
    \]
    \item Per ogni $x \in S^{n-1}$, vale $\langle x, x_0 \rangle = 0$, quindi $r(x) = \frac{x}{\|x\|} = x$. Dunque $r$ è una retrazione.
\end{enumerate}

Mostriamo che $i \circ r \simeq \mathrm{Id}_{W_\varepsilon}$ tramite l'omotopia lineare normalizzata:
\Func{H}{W_\varepsilon \times I}{W_\varepsilon}{(x,s)}{\frac{(1-s)x + s\,r(x)}{\|(1-s)x + s\,r(x)\|}}

Scriviamo $a = \langle x, x_0 \rangle$, con $|a| < \varepsilon < 1$, e $b = \sqrt{1-a^2} > 0$. Allora
\[
x = a\,x_0 + b\,r(x), \qquad r(x) \perp x_0, \ \|r(x)\| = 1 ,
\]
e quindi
\[
v := (1-s)x + s\,r(x) = (1-s)a\,x_0 + \big( (1-s)b + s \big)\, r(x).
\]
Poich\'e $x_0 \perp r(x)$ e $(1-s)b + s > 0$ (perch\'e $b>0$), si ha $v \neq \bar{0}$: $H$ \`e ben posta e continua.

$H$ prende valori in $W_\varepsilon$: dalla formula,
$\|v\|^2 = (1-s)^2 a^2 + \big((1-s)b + s\big)^2 \geqslant (1-s)^2 (a^2 + b^2) = (1-s)^2$, quindi
\[
|\langle H(x,s), x_0 \rangle| = \frac{(1-s)\,|a|}{\|v\|} \leqslant |a| < \varepsilon .
\]
Inoltre $H$ lascia fisso ogni punto di $S^{n-1}$ (l\`\i\ $a=0$ e $r(x)=x$).

Le verifiche ai bordi confermano che:
\begin{itemize}
    \item $H(x,0) = \frac{x}{\|x\|} = x = \mathrm{Id}_{W_\varepsilon}(x)$.
    \item $H(x,1) = \frac{r(x)}{\|r(x)\|} = r(x) = (i \circ r)(x)$.
\end{itemize}
Di conseguenza $\mathrm{Id}_{W_\varepsilon} \simeq i \circ r$, dimostrando che $S^{n-1}$ è un retratto di deformazione di $W_\varepsilon$.
\end{proof}

\vspace{2em}

\section{Spazi Contraibili}

\begin{definicion}
Uno spazio topologico $Y$ si dice \textbf{contraibile} a un suo punto $y_0 \in Y$ se l'applicazione identità $\mathrm{Id}_Y$ è omotopa alla mappa costante $c_{y_0}: Y \to Y$ (definita da $c_{y_0}(y) = y_0$ per ogni $y \in Y$):
\[
\mathrm{Id}_Y \simeq c_{y_0}
\]
\end{definicion}

% Nota en español: Un espacio contraíble es aquel que se puede deformar de manera continua hasta reducirse a un único punto.

\begin{teo}
Uno spazio topologico $Y$ è contraibile a un punto $y_0 \in Y$ se e solo se per ogni spazio topologico $X$, due qualsiasi mappe continue $f, g: X \to Y$ sono omotope ($f \simeq g$).
\end{teo}

\begin{proof}[Dimostrazione del Teorema sulla Contraibilità]
\begin{itemize}
    \item[$\impliedby$)] Banale. Basta scegliere $X = Y$ e prendere $f = \mathrm{Id}_Y$ e $g = c_{y_0}$.
    \item[$\implies$)] Siano $X$ uno spazio topologico e $f, g: X \to Y$ due mappe continue qualsiasi.
    
    Poiché $Y$ è contraibile a $y_0$, si ha $c_{y_0} \simeq \mathrm{Id}_Y$. Applicando la proprietà di composizione dell'omotopia:
    \[
    f \simeq f \quad \text{e} \quad c_{y_0} \simeq \mathrm{Id}_Y \implies c_{y_0} \circ f \simeq \mathrm{Id}_Y \circ f = f
    \]
    \[
    g \simeq g \quad \text{e} \quad c_{y_0} \simeq \mathrm{Id}_Y \implies c_{y_0} \circ g \simeq \mathrm{Id}_Y \circ g = g
    \]
    
    Osserviamo che per ogni $x \in X$:
    \[
    (c_{y_0} \circ f)(x) = c_{y_0}(f(x)) = y_0 \quad \text{e} \quad (c_{y_0} \circ g)(x) = c_{y_0}(g(x)) = y_0
    \]
    Quindi $c_{y_0} \circ f = c_{y_0} \circ g = c_{y_0}$ (la mappa costante da $X$ in $Y$ che invia tutto in $y_0$). Per transitività e simmetria:
    \[
    f \simeq c_{y_0} \circ f = c_{y_0} \circ g \simeq g \implies f \simeq g
    \]
\end{itemize}
\end{proof}

\begin{coro}
Se $Y$ è contraibile a un punto $y_0 \in Y$, allora $Y$ è contraibile a ogni suo punto $y \in Y$ (ovvero $c_y \simeq \mathrm{Id}_Y, \, \forall y \in Y$).
\end{coro}
\begin{proof}
Per il teorema, due qualsiasi mappe continue $X \to Y$ sono omotope; in particolare (con $X = Y$) $c_y \simeq \mathrm{Id}_Y$.
\end{proof}

\begin{coro}
Se $Y$ è contraibile, allora $Y$ è connesso per archi.
\end{coro}
\begin{proof}
Per il corollario precedente $c_{y_0} \simeq c_{y}$ per ogni $y \in Y$ (come mappe $Y \to Y$). Per l'esercizio sulle mappe costanti, esiste un cammino da $y_0$ a $y$. Quindi ogni punto si collega a $y_0$.
\end{proof}

\begin{observacion}
Il viceversa del Corollario 2 \textbf{non} vale in generale: connesso per archi $\centernot\implies$ contraibile.

Come controesempio, per ogni $n$ pari, la sfera $S^n$ è connessa per archi, ma la mappa identità $\mathrm{Id}_{S^n}$ non è omotopa a $-\mathrm{Id}_{S^n}$ (e quindi la sfera non è contraibile, $S^n \not\simeq \{x_0\}$).
\end{observacion}

\vspace{1.5em}

\begin{teo}
Per uno spazio topologico $Y$, le seguenti affermazioni sono equivalenti:
\begin{enumerate}
    \item $Y$ è contraibile.
    \item $Y$ ammette un suo punto $y_0$ come retratto di deformazione.
    \item $Y$ è omotopicamente equivalente a un punto (singoletto).
\end{enumerate}
\end{teo}

\begin{proof}
\begin{itemize}
    \item[$(1) \implies (2)$:] Sia $y_0 \in Y$. Per ipotesi, $c_{y_0} \simeq \mathrm{Id}_Y$.
    Sia $r: Y \to \{y_0\}$ l'unica retrazione $r(y) = y_0$. Si ha $r \circ i = \mathrm{Id}_{\{y_0\}}$.
    Inoltre, $(i \circ r)(y) = i(y_0) = y_0 = c_{y_0}(y)$. Poiché $c_{y_0} \simeq \mathrm{Id}_Y$, segue che $i \circ r \simeq \mathrm{Id}_Y$.
    Quindi $\{y_0\}$ è un retratto di deformazione di $Y$.

    \item[$(2) \implies (3)$:] Segue direttamente dalla definizione di retratto di deformazione (che fornisce un'equivalenza omotopica con il sottoinsieme $\{y_0\}$).

    \item[$(3) \implies (1)$:] Se $Y \simeq \{*\}$, esistono mappe $f: Y \to \{*\}$ e $g: \{*\} \to Y$ con $g \circ f \simeq \mathrm{Id}_Y$. Posto $y_1 = g(*)$, si ha $(g \circ f)(y) = y_1 = c_{y_1}(y)$, quindi $c_{y_1} \simeq \mathrm{Id}_Y$: $Y$ è contraibile a $y_1$, e dunque (corollario precedente) a ogni suo punto.
\end{itemize}
\end{proof}

\vspace{1.5em}

\begin{ejemplo}
\begin{enumerate}
    \item $\mathbb{R}^n$ e qualsiasi insieme convesso $\Omega \subseteq \mathbb{R}^n$ sono contraibili.
    
    \item \textbf{Insiemi stellati:} Ricordiamo che un sottoinsieme $S \subseteq \mathbb{R}^n$ si dice \textbf{stellato rispetto a un punto} $x_0 \in S$ se per ogni $x \in S$ il segmento $\overline{x\, x_0}$ è interamente contenuto in $S$:
    \[
    \overline{x\, x_0} = \{(1-s)x + s\,x_0 \mid s \in [0,1]\} \subseteq S
    \]

    \begin{figure}[H]
        \centering
        \begin{tikzpicture}[scale=1.1]
            % Dominio estrellado
            \draw[thick, fill=gray!10] 
                (0,0) -- (1.5,0.5) -- (0.8,1.5) -- (2.2,2.0) -- (0.5,2.5) 
                -- (-1,1.8) -- (-2,2.2) -- (-1.5,0.8) -- (-2,-0.5) -- (-0.5,-0.2) -- cycle;
            
            % Punto centro x_0
            \coordinate (X0) at (0, 1.1);
            \filldraw[blue] (X0) circle (1.5pt) node[below right] {$x_0$};

            % Puntos en la frontera y segmentos
            \coordinate (P1) at (2.2, 2.0);
            \coordinate (P2) at (-2, 2.2);
            \coordinate (P3) at (-1.5, 0.8);

            \draw[dashed, red] (X0) -- (P1);
            \draw[dashed, red] (X0) -- (P2);
            \draw[dashed, red] (X0) -- (P3);

            \filldraw (P1) circle (1pt) node[above right] {$x$};
        \end{tikzpicture}
        \caption{Insieme $S \subset \mathbb{R}^n$ stellato rispetto a $x_0$.}
    \end{figure}

    \begin{prop}
    Se $S \subseteq \mathbb{R}^n$ è stellato rispetto a $x_0$, allora $S$ è contraibile.
    \end{prop}

    \begin{proof}
    Definiamo la retrazione $r: S \to \{x_0\}$ con $r(x) = x_0$.
    Mostriamo che $i \circ r \simeq \mathrm{Id}_S$ tramite l'omotopia lineare:
    \Func{H}{S \times I}{S}{(x,s)}{(1-s)x + s\,x_0}
    
    Poiché $S$ è stellato rispetto a $x_0$, $H(x,s) \in S$ per ogni $x \in S$ e $s \in I$. Inoltre $H(x,0) = x$ e $H(x,1) = x_0$. Dunque $\{x_0\}$ è un retratto di deformazione di $S$, e quindi $S$ è contraibile a $x_0$.
    \end{proof}
\end{enumerate}
\end{ejemplo}

\begin{prop}
Se $X$ e $Y$ sono spazi topologici contraibili, allora il loro prodotto cartesiano $X \times Y$ è contraibile.
\end{prop}

\begin{proof}
Siano $x_0 \in X$ e $y_0 \in Y$ i punti rispetto a cui $X$ e $Y$ si contraggono.
Per ipotesi, esistono le omotopie $H: c_{x_0} \simeq \mathrm{Id}_X$ su $X \times I$ e $K: c_{y_0} \simeq \mathrm{Id}_Y$ su $Y \times I$.

Definiamo la mappa prodotto per l'omotopia:
\Func{W}{(X \times Y) \times I}{X \times Y}{\big((x,y), s\big)}{\big(H(x,s), K(y,s)\big)}

$W$ è continua in quanto le sue componenti $H$ e $K$ sono continue. Verifichiamo le condizioni agli estremi:
\begin{enumerate}[label=\roman*)]
    \item $W((x,y), 0) = (H(x,0), K(y,0)) = (x_0, y_0) = c_{(x_0,y_0)}(x,y)$.
    \item $W((x,y), 1) = (H(x,1), K(y,1)) = (x,y) = \mathrm{Id}_{X \times Y}(x,y)$.
\end{enumerate}
Pertanto $W: c_{(x_0,y_0)} \simeq \mathrm{Id}_{X \times Y}$, ovvero $X \times Y$ è contraibile.
\end{proof}

\begin{ejemplo}
Se $S \subset \mathbb{R}^n$ è un insieme stellato e $\Omega \subset \mathbb{R}^m$ è convesso, il cilindro/prodotto $S \times \Omega$ è uno spazio contraibile.
\end{ejemplo}

\begin{prop}
Siano $X, Y$ spazi topologici. Se $Y$ è contraibile, allora $X \times Y \simeq X$.
\end{prop}

\begin{proof}
Sia $y_0 \in Y$ tale che $c_{y_0} \simeq \mathrm{Id}_Y$ tramite l'omotopia $H: Y \times I \to Y$.
Consideriamo le proiezioni e inclusioni continue:
\Func{f}{X \times Y}{X}{(x,y)}{x}
\Func{g}{X}{X \times Y}{x}{(x, y_0)}

Analizziamo le due composizioni:
\begin{itemize}
    \item $(f \circ g)(x) = f(x, y_0) = x = \mathrm{Id}_X(x)$, quindi $f \circ g = \mathrm{Id}_X \simeq \mathrm{Id}_X$.
    \item $(g \circ f)(x,y) = g(x) = (x, y_0)$. Mostriamo che $g \circ f \simeq \mathrm{Id}_{X \times Y}$ costruendo l'omotopia:
    \Func{K}{(X \times Y) \times I}{X \times Y}{\big((x,y), s\big)}{\big(x, H(y,s)\big)}
    $K$ è continua e verifica $K((x,y),0) = (x, y_0) = (g \circ f)(x,y)$ e $K((x,y),1) = (x, y) = \mathrm{Id}_{X \times Y}(x,y)$.
\end{itemize}
Di conseguenza $X \times Y \simeq X$.
\end{proof}

\section{Gruppo Fondamentale e Omotopia di Cammini}

Un \textbf{cappio} (o cammino chiuso) in $X$ con punto base $x_0$ \`e un cammino $\sigma: I \to X$ con $\sigma(0) = \sigma(1) = x_0$. Indichiamo con $\Lambda_{x_0} = \Lambda_{x_0}(X)$ l'insieme dei cappi in $x_0$. La relazione di omotopia relativamente al bordo $\partial I = \{0,1\}$ \`e una relazione di equivalenza su $\Lambda_{x_0}$ (riflessiva, simmetrica e transitiva: si compongono/invertono le omotopie). Poniamo:
\[
\Pi_1(X, x_0) = \Lambda_{x_0} \big/ \!\simeq \text{ rel } \partial I
\]
Gli elementi di $\Pi_1(X, x_0)$ sono le classi di omotopia dei cammini chiusi in $X$ con base in $x_0$, relativamente a $\partial I$.

\begin{figure}[H]
    \centering
    \begin{tikzpicture}[scale=1.2]
        % Plano X
        \draw[gray!50, fill=gray!5] (-2,-1) -- (2.5,-1) -- (3.5,1) -- (-1,1) -- cycle;
        \node[above left] at (3.5,1) {$X$};

        % Punto base x_0
        \coordinate (X0) at (0,-0.2);
        \filldraw[blue] (X0) circle (1.5pt) node[below] {$x_0$};

        % Lazos
        \draw[thick, red] (X0) to[out=0, in=60] (0.8, 0.6) to[out=120, in=180] (X0);
        \node[red, right] at (0.8, 0.6) {$\sigma$};

        \draw[thick, ForestGreen] (X0) to[out=-30, in=30] (1.2, 0) to[out=90, in=150] (X0);
        \node[ForestGreen, right] at (1.2, 0) {$\tau$};

        \draw[thick, Purple] (X0) to[out=120, in=180] (-0.8, 0.5) to[out=-90, in=0] (X0);
    \end{tikzpicture}
    \caption{Cappi in $X$ con punto base $x_0$.}
\end{figure}

Per ogni $\sigma \in \Lambda_{x_0}$, la classe di omotopia è indicata con:
\[
[\sigma] = \{ \tau \in \Lambda_{x_0} \mid \sigma \simeq \tau \text{ rel } \partial I \}
\]

\begin{observacion}
Siano $\sigma, \tau \in \Lambda_{x_0}$, ovvero $\sigma(1) = \tau(0) = x_0$.
Si ha che la concatenazione $\sigma * \tau \in \Lambda_{x_0}$, l'inverso $\sigma^{-1} \in \Lambda_{x_0}$ e il cammino costante $c_{x_0} \in \Lambda_{x_0}$.
\end{observacion}

\vspace{1.5em}

\begin{teo}
Sia $X$ uno spazio topologico.
\begin{enumerate}[label=\alph*)]
    \item \textbf{Compatibilità dell'omotopia rel $\partial I$ con l'operazione di concatenazione $*$:}
    Siano $\alpha, \beta, \alpha', \beta'$ cammini in $X$ tali che $\alpha(1) = \beta(0)$ e $\alpha'(1) = \beta'(0)$.
    Se $\alpha \simeq \alpha' \text{ rel } \partial I$ e $\beta \simeq \beta' \text{ rel } \partial I$, allora:
    \[
    \alpha * \beta \simeq \alpha' * \beta' \text{ rel } \partial I
    \]

    \item \textbf{Associatività a meno di omotopia:}
    Siano $\alpha, \beta, \gamma$ cammini in $X$ con $\alpha(1) = \beta(0)$ e $\beta(1) = \gamma(0)$. Allora:
    \[
    (\alpha * \beta) * \gamma \simeq \alpha * (\beta * \gamma) \text{ rel } \partial I
    \]

    \item \textbf{Elemento neutro a meno di omotopia:}
    Per ogni cammino $\alpha$ in $X$ con $\alpha(0) = x_0$ e $\alpha(1) = x_1$:
    \[
    c_{x_0} * \alpha \simeq \alpha \text{ rel } \partial I \quad \text{e} \quad \alpha * c_{x_1} \simeq \alpha \text{ rel } \partial I
    \]

    \item \textbf{Elemento inverso a meno di omotopia:}
    Per ogni cammino $\alpha$ in $X$ con $\alpha(0) = x_0$ e $\alpha(1) = x_1$:
    \[
    \alpha * \alpha^{-1} \simeq c_{x_0} \text{ rel } \partial I \quad \text{e} \quad \alpha^{-1} * \alpha \simeq c_{x_1} \text{ rel } \partial I
    \]
\end{enumerate}
\end{teo}

\begin{proof}
Dimostriamo le varie parti del teorema:

\textbf{Dimostrazione di a):}
Siano $H: \alpha \simeq \alpha' \text{ rel } \partial I$ e $K: \beta \simeq \beta' \text{ rel } \partial I$ le due omotopie relative al bordo. Definiamo la mappa $W: I \times I \to X$:
\[
W(t,s) = 
\begin{cases}
    H(2t, s) & t \leqslant \frac{1}{2} \\
    K(2t-1, s) & t \geqslant \frac{1}{2}
\end{cases}
\]

Verifichiamo la continuità e ben posizione per $t = \frac{1}{2}$:
\[
H(2 \cdot \tfrac{1}{2}, s) = H(1, s) = \alpha(1) = \beta(0) = K(0, s) = K(2 \cdot \tfrac{1}{2} - 1, s)
\]
Quindi $W$ è ben definita e continua per il Lemma di Incollamento.

Verifichiamo le condizioni ai bordi:
\begin{enumerate}[label=\roman*)]
    \item Per $s = 0$:
    \[
    W(t,0) = 
    \begin{cases}
        H(2t,0) = \alpha(2t) & t \leqslant \frac{1}{2} \\
        K(2t-1,0) = \beta(2t-1) & t \geqslant \frac{1}{2}
    \end{cases}
    = (\alpha * \beta)(t)
    \]
    \item Per $s = 1$:
    \[
    W(t,1) = 
    \begin{cases}
        H(2t,1) = \alpha'(2t) & t \leqslant \frac{1}{2} \\
        K(2t-1,1) = \beta'(2t-1) & t \geqslant \frac{1}{2}
    \end{cases}
    = (\alpha' * \beta')(t)
    \]
    \item Per $t = 0$ e $t = 1$, poiché $H$ e $K$ sono omotopie relative a $\partial I$:
    \[
    W(0,s) = H(0,s) = \alpha(0) = (\alpha * \beta)(0)
    \]
    \[
    W(1,s) = K(1,s) = \beta(1) = (\alpha * \beta)(1)
    \]
\end{enumerate}
Dunque $W: \alpha * \beta \simeq \alpha' * \beta' \text{ rel } \partial I$.

\vspace{1em}

\textbf{Dimostrazione di b) (Associatività):}
Calcoliamo le definizioni esplicite delle due riparametrizzazioni dei cammini:
\[
((\alpha * \beta) * \gamma)(t) = 
\begin{cases}
    (\alpha * \beta)(2t) = \alpha(4t) & 0 \leqslant t \leqslant \frac{1}{4} \\
    (\alpha * \beta)(2t) = \beta(4t-1) & \frac{1}{4} \leqslant t \leqslant \frac{1}{2} \\
    \gamma(2t-1) & \frac{1}{2} \leqslant t \leqslant 1
\end{cases}
\]
e
\[
(\alpha * (\beta * \gamma))(t) = 
\begin{cases}
    \alpha(2t) & 0 \leqslant t \leqslant \frac{1}{2} \\
    (\beta * \gamma)(2t-1) = \beta(4t-2) & \frac{1}{2} \leqslant t \leqslant \frac{3}{4} \\
    (\beta * \gamma)(2t-1) = \gamma(4t-3) & \frac{3}{4} \leqslant t \leqslant 1
\end{cases}
\]

Per costruire l'omotopia $H: I \times I \to X$ tra le due riparametrizzazioni, suddividiamo il quadrato $I \times I$ mediante due rette nel piano $(t,s)$:
la retta che unisce $(\frac{1}{4},0)$ a $(\frac{1}{2},1)$, di equazione $s = 4t-1$, e la retta che unisce $(\frac{1}{2},0)$ a $(\frac{3}{4},1)$, di equazione $s = 4t-2$.

\begin{figure}[H]
    \centering
    \begin{tikzpicture}[scale=2]
        % Cuadrado I x I
        \draw[thick] (0,0) -- (2,0) -- (2,1.5) -- (0,1.5) -- cycle;
        \node[below left] at (0,0) {$0$};
        \node[below] at (0.5,0) {$1/4$};
        \node[below] at (1,0) {$1/2$};
        \node[below] at (1.5,0) {$3/4$};
        \node[below] at (2,0) {$1$};

        \node[above] at (1,1.5) {$1/2$};
        \node[above] at (1.5,1.5) {$3/4$};

        % Divisiones
        \draw[thick, blue] (0.5,0) -- (1,1.5);
        \draw[thick, blue] (1,0) -- (1.5,1.5);

        % Regiones
        \node at (0.4, 0.75) {$A \quad (\alpha)$};
        \node at (1.0, 0.75) {$B \quad (\beta)$};
        \node at (1.6, 0.75) {$C \quad (\gamma)$};

        % Ejes
        \draw[->,>=stealth] (0,0) -- (2.3,0) node[right] {$t$};
        \draw[->,>=stealth] (0,0) -- (0,1.7) node[above] {$s$};
    \end{tikzpicture}
    \caption{Suddivisione del quadrato $I \times I$ per l'omotopia di associatività.}
\end{figure}

Definiamo le tre regioni del quadrato $I \times I$:
\[
A = \{ (t,s) \in I \times I \mid s \geqslant 4t-1 \}
\]
\[
B = \{ (t,s) \in I \times I \mid 4t-2 \leqslant s \leqslant 4t-1 \}
\]
\[
C = \{ (t,s) \in I \times I \mid s \leqslant 4t-2 \}
\]

Definiamo l'omotopia $H: I \times I \to X$ sulle tre regioni:
\[
H(t,s) = 
\begin{cases}
    \alpha\left( \dfrac{4t}{s+1} \right) & (t,s) \in A \\[1em]
    \beta(4t - s - 1) & (t,s) \in B \\[1em]
    \gamma\left( \dfrac{4t - s - 2}{2 - s} \right) & (t,s) \in C
\end{cases}
\]

Verifichiamo la continuità sulle linee di separazione (Lemma di Incollamento):
\begin{itemize}
    \item Sulla retta $s = 4t-1$ (confine tra $A$ e $B$):
    \[
    \alpha\left(\frac{4t}{(4t-1)+1}\right) = \alpha(1) = \beta(0) = \beta(4t - (4t-1) - 1)
    \]
    \item Sulla retta $s = 4t-2$ (confine tra $B$ e $C$):
    \[
    \beta(4t - (4t-2) - 1) = \beta(1) = \gamma(0) = \gamma\left(\frac{4t - (4t-2) - 2}{2 - (4t-2)}\right)
    \]
\end{itemize}

Controlliamo le condizioni per $s = 0$ e $s = 1$:
\begin{itemize}
    \item Per $s = 0$, le regioni corrispondono a $t \leqslant 1/4$, $1/4 \leqslant t \leqslant 1/2$ e $t \geqslant 1/2$:
    \[
    H(t,0) = 
    \begin{cases}
        \alpha(4t) & t \leqslant 1/4 \\
        \beta(4t-1) & 1/4 \leqslant t \leqslant 1/2 \\
        \gamma(2t-1) & t \geqslant 1/2
    \end{cases}
    = ((\alpha * \beta) * \gamma)(t)
    \]
    \item Per $s = 1$, le regioni corrispondono a $t \leqslant 1/2$, $1/2 \leqslant t \leqslant 3/4$ e $t \geqslant 3/4$:
    \[
    H(t,1) = 
    \begin{cases}
        \alpha(2t) & t \leqslant 1/2 \\
        \beta(4t-2) & 1/2 \leqslant t \leqslant 3/4 \\
        \gamma(4t-3) & t \geqslant 3/4
    \end{cases}
    = (\alpha * (\beta * \gamma))(t)
    \]
\end{itemize}
Inoltre $H(0,s) = \alpha(0)$ e $H(1,s) = \gamma(1)$ per ogni $s \in I$.
Pertanto $H$ è continua e dimostra l'associatività a meno di omotopia: $(\alpha * \beta) * \gamma \simeq \alpha * (\beta * \gamma) \text{ rel } \partial I$.
\end{proof}

\textbf{Dimostrazione di c) (Elemento neutro):}
Dimostriamo che $c_{x_0} * \alpha \simeq \alpha \text{ rel } \partial I$ (la dimostrazione per $\alpha * c_{x_1} \simeq \alpha$ è del tutto analoga).

Consideriamo la retta nel piano $(t,s)$ che unisce i punti $(\frac{1}{2}, 0)$ e $(1,1)$, avente equazione $s = 2t - 1$.

\begin{figure}[H]
    \centering
    \begin{tikzpicture}[scale=2]
        % Cuadrado I x I
        \draw[thick] (0,0) -- (1.6,0) -- (1.6,1.2) -- (0,1.2) -- cycle;
        \node[below left] at (0,0) {$0$};
        \node[below] at (0.8,0) {$1/2$};
        \node[below] at (1.6,0) {$1$};

        % Linea de separación s = 2t - 1
        \draw[thick, blue] (0.8,0) -- (1.6,1.2);

        % Regiones A y B
        \node at (0.5, 0.6) {$A \quad (\alpha)$};
        \node at (1.3, 0.4) {$B \quad (c_{x_1})$};

        % Ejes
        \draw[->,>=stealth] (0,0) -- (1.9,0) node[right] {$t$};
        \draw[->,>=stealth] (0,0) -- (0,1.4) node[above] {$s$};
    \end{tikzpicture}
    \caption{Divisione del quadrato $I \times I$ per l'omotopia dell'elemento neutro.}
\end{figure}

Definiamo le regioni $A$ e $B$:
\[
A = \{ (t,s) \in I \times I \mid s \geqslant 2t - 1 \}
\]
\[
B = \{ (t,s) \in I \times I \mid s \leqslant 2t - 1 \}
\]

Definiamo l'omotopia $H: I \times I \to X$:
\[
H(t,s) = 
\begin{cases}
    \alpha\left( \dfrac{2t}{s+1} \right) & (t,s) \in A \\[1em]
    x_1 & (t,s) \in B
\end{cases}
\]

Verifichiamo la continuità sulla linea di giunzione $s = 2t - 1$ ($2t = s + 1$):
\[
\alpha\left( \frac{2t}{s+1} \right) = \alpha(1) = x_1
\]
Per il Lemma di Incollamento, $H$ è ben posta e continua.

Verifichiamo le condizioni ai bordi:
\begin{enumerate}[label=\roman*)]
    \item Per $s = 0$:
    \[
    H(t,0) = 
    \begin{cases}
        \alpha(2t) & t \leqslant \frac{1}{2} \\
        x_1 & t \geqslant \frac{1}{2}
    \end{cases}
    = (\alpha * c_{x_1})(t)
    \]
    \item Per $s = 1$:
    \[
    H(t,1) = \alpha\left( \frac{2t}{2} \right) = \alpha(t)
    \]
    \item Per $t = 0$ e $t = 1$:
    \[
    H(0,s) = \alpha(0) = x_0 \quad \text{e} \quad H(1,s) = x_1
    \]
\end{enumerate}
Quindi $H: \alpha * c_{x_1} \simeq \alpha \text{ rel } \partial I$.

\vspace{1.5em}

\textbf{Dimostrazione di d) (Elemento inverso):}
Dimostriamo che $\alpha * \alpha^{-1} \simeq c_{x_0} \text{ rel } \partial I$.

Consideriamo le rette $s = 2t$ (tra $(0,0)$ e $(\frac{1}{2},1)$) e $s = 2 - 2t$ (tra $(\frac{1}{2},1)$ e $(1,0)$).

\begin{figure}[H]
    \centering
    \begin{tikzpicture}[scale=2]
        % Cuadrado I x I
        \draw[thick] (0,0) -- (1.6,0) -- (1.6,1.2) -- (0,1.2) -- cycle;
        \node[below left] at (0,0) {$(0,0)$};
        \node[below] at (1.6,0) {$(1,0)$};
        \node[above] at (0.8,1.2) {$(1/2,1)$};

        % Triangulo interior B
        \draw[thick, blue] (0,0) -- (0.8,1.2) -- (1.6,0);

        % Regiones
        \node at (0.3, 0.7) {$A$};
        \node at (0.8, 0.4) {$B \quad (\alpha(s))$};
        \node at (1.3, 0.7) {$C$};

        % Ejes
        \draw[->,>=stealth] (0,0) -- (1.8,0) node[right] {$t$};
        \draw[->,>=stealth] (0,0) -- (0,1.4) node[above] {$s$};
    \end{tikzpicture}
    \caption{Suddivisione per l'omotopia dell'elemento inverso.}
\end{figure}

Definiamo le regioni del quadrato:
\[
A = \{ (t,s) \in I \times I \mid 2t \leqslant s \}
\]
\[
B = \{ (t,s) \in I \times I \mid s \leqslant 2t \text{ e } s \leqslant 2 - 2t \}
\]
\[
C = \{ (t,s) \in I \times I \mid 2t \geqslant 2 - s \}
\]

Definiamo l'omotopia $H: I \times I \to X$:
\[
H(t,s) = 
\begin{cases}
    \alpha(2t) & (t,s) \in A \\
    \alpha(s) & (t,s) \in B \\
    \alpha^{-1}(2t-1) = \alpha(2 - 2t) & (t,s) \in C
\end{cases}
\]

Sui confini della regione $B$:
\begin{itemize}
    \item Per $s = 2t$: $\alpha(2t) = \alpha(s)$.
    \item Per $s = 2 - 2t$: $\alpha(s) = \alpha(2 - 2t) = \alpha^{-1}(2t - 1)$.
\end{itemize}
Quindi $H$ è continua per il Lemma di Incollamento.

Verifichiamo le condizioni ai bordi:
\begin{enumerate}[label=\roman*)]
    \item Per $s = 0$: $H(t,0) = \alpha(0) = x_0 = c_{x_0}(t)$ per ogni $t \in I$.
    \item Per $s = 1$:
    \[
    H(t,1) = 
    \begin{cases}
        \alpha(2t) & t \leqslant \frac{1}{2} \\
        \alpha^{-1}(2t-1) & t \geqslant \frac{1}{2}
    \end{cases}
    = (\alpha * \alpha^{-1})(t)
    \]
    \item Per $t = 0$ e $t = 1$:
    \[
    H(0,s) = \alpha(0) = x_0 \quad \text{e} \quad H(1,s) = \alpha^{-1}(1) = \alpha(0) = x_0
    \]
\end{enumerate}
Pertanto $H: \alpha * \alpha^{-1} \simeq c_{x_0} \text{ rel } \partial I$\hfill$\square$

\vspace{2em}

\begin{coro}[Struttura del Gruppo Fondamentale]
Sia $X$ uno spazio topologico e $x_0 \in X$. L'insieme delle classi di omotopia $\Pi_1(X, x_0)$ munito dell'operazione di prodotto tra classi:
\Func{\cdot}{\Pi_1(X, x_0) \times \Pi_1(X, x_0)}{\Pi_1(X, x_0)}{([\sigma], [\tau])}{[\sigma * \tau]}
è un \textbf{gruppo}, detto \textbf{Gruppo Fondamentale} di $X$ con punto base $x_0$ (o \textbf{Gruppo di Poincaré}).
\end{coro}

\begin{proof}
Dimostriamo che l'operazione è ben definita e verifica gli assiomi di gruppo:
\begin{enumerate}
    \item \textbf{Buona definizione:} Siano $[\sigma] = [\sigma']$ e $[\tau] = [\tau']$, ossia $\sigma \simeq \sigma' \text{ rel } \partial I$ e $\tau \simeq \tau' \text{ rel } \partial I$. Per la proprietà a) del teorema precedente, $\sigma * \tau \simeq \sigma' * \tau' \text{ rel } \partial I$, il che implica che $[\sigma * \tau] = [\sigma' * \tau']$.
    
    \item \textbf{Associatività:} Per ogni $[\sigma], [\tau], [\gamma] \in \Pi_1(X, x_0)$:
    \[
    ([\sigma] \cdot [\tau]) \cdot [\gamma] = [\sigma * \tau] \cdot [\gamma] = [(\sigma * \tau) * \gamma] \stackrel{\text{b)}}{=} [\sigma * (\tau * \gamma)] = [\sigma] \cdot ([\tau] \cdot [\gamma])
    \]
    
    \item \textbf{Elemento neutro:} La classe del cammino costante $[c_{x_0}]$ è l'elemento neutro del gruppo:
    \[
    [c_{x_0}] \cdot [\sigma] = [c_{x_0} * \sigma] \stackrel{\text{c)}}{=} [\sigma] \quad \text{e} \quad [\sigma] \cdot [c_{x_0}] = [\sigma * c_{x_0}] \stackrel{\text{c)}}{=} [\sigma]
    \]
    
    \item \textbf{Elemento inverso:} Per ogni $[\sigma] \in \Pi_1(X, x_0)$, esiste la classe del cammino percorso in senso opposto $[\sigma^{-1}] \in \Pi_1(X, x_0)$, che funge da inverso:
    \[
    [\sigma] \cdot [\sigma^{-1}] = [\sigma * \sigma^{-1}] \stackrel{\text{d)}}{=} [c_{x_0}] \quad \text{e} \quad [\sigma^{-1}] \cdot [\sigma] = [\sigma^{-1} * \sigma] \stackrel{\text{d)}}{=} [c_{x_0}]
    \]
\end{enumerate}
\end{proof}

\section{Ruolo del Punto Base e Invarianza del Gruppo Fondamentale}

\begin{ejemplo}
Sia $X = S^1 \cup \{x_0\}$ (consideriamo questo spazio con la topologia naturale indotta da $\mathbb{R}^2$).

\begin{figure}[H]
    \centering
    \begin{tikzpicture}[scale=1.2]
        % Plano X
        \draw[gray!50, fill=gray!5] (-2,-0.8) -- (2.5,-0.8) -- (3.5,0.8) -- (-1,0.8) -- cycle;
        \node[above left] at (3.5,0.8) {$X$};

        % Circulo S^1
        \draw[thick] (0,0) circle (0.6);
        \node at (-0.3, 0.4) {$S^1$};
        \filldraw (0,-0.6) circle (1pt) node[below] {$x_1$};

        % Punto aislado x_0
        \filldraw[blue] (2,0) circle (1.5pt) node[below right] {$x_0$};
    \end{tikzpicture}
    \caption{Spazio $X = S^1 \cup \{x_0\}$ con un punto isolato $x_0$.}
\end{figure}

Chiaramente, per il punto isolato $x_0$ si ha $\Lambda_{x_0} = \{c_{x_0}\}$, quindi:
\[
\Pi_1(X, x_0) = \{[c_{x_0}]\} = 0
\]
Invece, per qualsiasi punto $x_1 \in S^1$, si ha:
\[
\Pi_1(X, x_1) \neq 0
\]
poiché non tutti i cappi contenuti in $S^1$ sono omotopi al cammino costante $c_{x_1}$. Questo mostra che, in generale, il gruppo fondamentale dipende dalla scelta del punto base se lo spazio non è connesso per archi. (Il fatto che $\Pi_1(S^1) \neq 0$ sarà dimostrato rigorosamente pi\`u avanti, con Seifert--Van Kampen: $\Pi_1(S^1) \cong \mathbb{Z}$.)
\end{ejemplo}

\vspace{1.5em}

\begin{teo}
Siano $X$ uno spazio topologico e $x_0, x_1 \in X$. Se esiste un cammino $\alpha: I \to X$ con $\alpha(0) = x_0$ e $\alpha(1) = x_1$, allora l'applicazione:
\Func{\alpha_*}{\Pi_1(X, x_0)}{\Pi_1(X, x_1)}{[\sigma]}{[\alpha^{-1} * \sigma * \alpha]}
è un \textbf{isomorfismo di gruppi}.
\end{teo}

\begin{proof}
Dimostriamo le tre proprietà dell'applicazione $\alpha_*$:

\begin{figure}[H]
    \centering
    \begin{tikzpicture}[scale=1.2]
        % Plano X
        \draw[gray!50, fill=gray!5] (-2,-1) -- (3,-1) -- (4,1) -- (-1,1) -- cycle;
        \node[above left] at (4,1) {$X$};

        % Puntos x_0 y x_1
        \coordinate (X0) at (-0.5,-0.3);
        \coordinate (X1) at (2.5,0.2);
        \filldraw[blue] (X0) circle (1.5pt) node[below] {$x_0$};
        \filldraw[blue] (X1) circle (1.5pt) node[below] {$x_1$};

        % Lazo \sigma en x_0
        \draw[thick, red] (X0) to[out=130, in=50] (-0.5,0.6) to[out=-50, in=-130] (X0);
        \node[red, above] at (-0.5,0.6) {$\sigma$};

        % Cammino \alpha de x_0 a x_1
        \draw[thick, ForestGreen, ->, >=stealth] (X0) to[out=20, in=160] (2.5,0.2);
        \node[ForestGreen, above] at (1.0, 0.2) {$\alpha$};

        % Cammino \alpha^{-1}
        \draw[thick, Purple, ->, >=stealth] (X1) to[out=-140, in=-40] (X0);
        \node[Purple, below] at (1.0, -0.3) {$\alpha^{-1}$};
    \end{tikzpicture}
    \caption{Coniugazione di un cappio $\sigma$ in $x_0$ tramite il cammino $\alpha$.}
\end{figure}

\begin{itemize}
    \item \textbf{Buona definizione:} Per ogni $\sigma \in \Lambda_{x_0}$, la concatenazione $\alpha^{-1} * \sigma * \alpha$ è un cappio centrato in $x_1$, quindi $\alpha^{-1} * \sigma * \alpha \in \Lambda_{x_1}$.
    Inoltre, se $[\sigma] = [\sigma']$, ossia $\sigma \simeq \sigma' \text{ rel } \partial I$, allora poichè $\alpha \simeq \alpha \text{ rel } \partial I$ e $\alpha^{-1} \simeq \alpha^{-1} \text{ rel } \partial I$, si ha:
    \[
    \alpha^{-1} * \sigma * \alpha \simeq \alpha^{-1} * \sigma' * \alpha \text{ rel } \partial I \implies [\alpha^{-1} * \sigma * \alpha] = [\alpha^{-1} * \sigma' * \alpha]
    \]
    Quindi $\alpha_*$ non dipende dal rappresentante scelto.

    \item \textbf{Iniettività:} Siano $\sigma, \tau \in \Lambda_{x_0}$ tali che $\alpha_*([\sigma]) = \alpha_*([\tau])$.
    Allora $[\alpha^{-1} * \sigma * \alpha] = [\alpha^{-1} * \tau * \alpha]$, ossia $\alpha^{-1} * \sigma * \alpha \simeq \alpha^{-1} * \tau * \alpha \text{ rel } \partial I$.
    
    Concatenando a sinistra per $\alpha$ e a destra per $\alpha^{-1}$:
    \[
    \alpha * (\alpha^{-1} * \sigma * \alpha) * \alpha^{-1} \simeq \alpha * (\alpha^{-1} * \tau * \alpha) * \alpha^{-1} \text{ rel } \partial I
    \]
    Usando l'associatività e i risultati sugli elementi inversi e neutri ($\alpha * \alpha^{-1} \simeq c_{x_0}$):
    \[
    c_{x_0} * \sigma * c_{x_0} \simeq c_{x_0} * \tau * c_{x_0} \implies \sigma \simeq \tau \text{ rel } \partial I \implies [\sigma] = [\tau]
    \]
    Pertanto $\alpha_*$ è iniettiva.

    \item \textbf{Suriettività:} Sia $[\eta] \in \Pi_1(X, x_1)$ un qualsiasi elemento con $\eta \in \Lambda_{x_1}$.
    Poniamo $\sigma = \alpha * \eta * \alpha^{-1} \in \Lambda_{x_0}$. Calcolando l'immagine tramite $\alpha_*$:
    \[
    \alpha_*([\sigma]) = [\alpha^{-1} * (\alpha * \eta * \alpha^{-1}) * \alpha] = [(\alpha^{-1} * \alpha) * \eta * (\alpha^{-1} * \alpha)]
    \]
    Poiché $\alpha^{-1} * \alpha \simeq c_{x_1} \text{ rel } \partial I$, si ottiene:
    \[
    c_{x_1} * \eta * c_{x_1} \simeq \eta \text{ rel } \partial I \implies \alpha_*([\sigma]) = [\eta]
    \]
    Quindi $\alpha_*$ è suriettiva.

    \item \textbf{Omomorfismo di gruppi:} Siano $\sigma, \tau \in \Lambda_{x_0}$:
    \[
    \alpha_*([\sigma]) \cdot \alpha_*([\tau]) = [\alpha^{-1} * \sigma * \alpha] \cdot [\alpha^{-1} * \tau * \alpha] = [\alpha^{-1} * \sigma * (\alpha * \alpha^{-1}) * \tau * \alpha]
    \]
    Essendo $\alpha * \alpha^{-1} \simeq c_{x_0} \text{ rel } \partial I$:
    \[
    [\alpha^{-1} * \sigma * c_{x_0} * \tau * \alpha] = [\alpha^{-1} * (\sigma * \tau) * \alpha] = \alpha_*([\sigma * \tau]) = \alpha_*([\sigma] \cdot [\tau])
    \]
\end{itemize}
Dunque $\alpha_*$ è un isomorfismo di gruppi.
\end{proof}

\begin{coro}
Sia $X$ uno spazio topologico connesso per archi. Allora il gruppo fondamentale $\Pi_1(X, x)$ è isomorfo per tutti i punti $x \in X$ (ed è quindi indipendente dal punto base a meno di isomorfismo, potendo scriverlo semplicemente $\Pi_1(X)$).
\end{coro}

\begin{observacion}
Se $X$ è uno spazio topologico e $x_0 \in X$, la condizione $\Pi_1(X, x_0) = 0$ equivale a dire che ogni cappio in $x_0$ è omotopo al cammino costante $c_{x_0}$ relativamente al bordo $\partial I$. Ciò equivale a dire che per ogni coppia di cammini $\alpha, \beta: I \to X$ con $\alpha(0) = \beta(0) = x_0$ e $\alpha(1) = \beta(1)$, si ha $\alpha \simeq \beta \text{ rel } \partial I$.
\end{observacion}

\begin{proof}[Dimostrazione dell'equivalenza]
($\Leftarrow$) Basta applicare l'ipotesi con $\alpha = \sigma$ (un cappio in $x_0$) e $\beta = c_{x_0}$.

($\Rightarrow$) Siano $\alpha, \beta$ con $\alpha(0) = \beta(0) = x_0$ e $\alpha(1) = \beta(1) = x_1$. Allora $\alpha * \beta^{-1} \in \Lambda_{x_0}$, e poich\'e $\Pi_1(X, x_0) = 0$ si ha $\alpha * \beta^{-1} \simeq c_{x_0}$ rel $\partial I$. Usando compatibilit\`a, associativit\`a, inverso e neutro:
\[
\alpha \simeq \alpha * c_{x_1} \simeq \alpha * (\beta^{-1} * \beta) \simeq (\alpha * \beta^{-1}) * \beta \simeq c_{x_0} * \beta \simeq \beta \quad \text{rel } \partial I. \qedhere
\]
\end{proof}

\section{Invarianza Topologica del Gruppo Fondamentale}

\begin{prop}
Siano $X, Y$ spazi topologici, $x_0 \in X$ e $f: X \to Y$ un'applicazione continua.
Allora l'applicazione indotta da $f$:
\Func{f_*}{\Pi_1(X, x_0)}{\Pi_1(Y, f(x_0))}{[\sigma]}{[f \circ \sigma]}
è un \textbf{omomorfismo di gruppi} (più precisamente indicato con $f_{*x_0}$).
\end{prop}

\begin{proof}
Verifichiamo le due proprietà:
\begin{itemize}
    \item \textbf{Buona definizione:} Per ogni $\sigma \in \Lambda_{x_0}(X)$, la composizione $f \circ \sigma$ è un cappio in $Y$ centrato in $f(x_0)$, ossia $f \circ \sigma \in \Lambda_{f(x_0)}(Y)$.
    Se $[\sigma] = [\sigma']$, si ha $\sigma \simeq \sigma' \text{ rel } \partial I$. Sia $H: \sigma \simeq \sigma' \text{ rel } \partial I$. Allora $f \circ H: I \times I \to Y$ \`e continua, $(f\circ H)(t,0) = f(\sigma(t))$, $(f \circ H)(t,1) = f(\sigma'(t))$ e $(f \circ H)(0,s) = f(x_0) = (f\circ H)(1,s)$. Dunque $f \circ \sigma \simeq f \circ \sigma' \text{ rel } \partial I$ (\`e un caso particolare del teorema di compatibilit\`a con la composizione, con la omotopia costante $f \simeq f$), quindi $[f \circ \sigma] = [f \circ \sigma']$.

    \item \textbf{Omomorfismo:} Siano $\sigma, \tau \in \Lambda_{x_0}(X)$. Osserviamo che per ogni $t \in I$:
    \[
    (f \circ (\sigma * \tau))(t) = 
    \begin{cases}
        f(\sigma(2t)) & t \leqslant 1/2 \\
        f(\tau(2t-1)) & t \geqslant 1/2
    \end{cases}
    = ((f \circ \sigma) * (f \circ \tau))(t)
    \]
    Quindi $f \circ (\sigma * \tau) = (f \circ \sigma) * (f \circ \tau)$. Di conseguenza:
    \[
    f_*([\sigma] \cdot [\tau]) = f_*([\sigma * \tau]) = [f \circ (\sigma * \tau)] = [(f \circ \sigma) * (f \circ \tau)] = f_*([\sigma]) \cdot f_*([\tau])
    \]
\end{itemize}
\end{proof}

\vspace{1em}

\begin{prop}
\begin{enumerate}[label=\alph*)]
    \item Siano $X, Y, Z$ spazi topologici, $x_0 \in X$, $f: X \to Y$ e $g: Y \to Z$ continue. Allora:
    \[
    (g \circ f)_* = g_* \circ f_* \quad \text{su } \Pi_1(X, x_0)
    \]
    \item L'applicazione identità $\mathrm{Id}_X: X \to X$ induce l'identità sul gruppo fondamentale:
    \[
    (\mathrm{Id}_X)_* = \mathrm{Id}_{\Pi_1(X, x_0)}
    \]
\end{enumerate}
\end{prop}

\begin{proof}
\begin{itemize}
    \item[a)] Per ogni $[\sigma] \in \Pi_1(X, x_0)$:
    \[
    (g \circ f)_*([\sigma]) = [(g \circ f) \circ \sigma] = [g \circ (f \circ \sigma)] = g_*([f \circ \sigma]) = g_*(f_*([\sigma])) = (g_* \circ f_*)([\sigma])
    \]
    \item[b)] Per ogni $[\sigma] \in \Pi_1(X, x_0)$:
    \[
    (\mathrm{Id}_X)_*([\sigma]) = [\mathrm{Id}_X \circ \sigma] = [\sigma] = \mathrm{Id}_{\Pi_1(X, x_0)}([\sigma])
    \]
\end{itemize}
\end{proof}

\vspace{1.5em}

\begin{teo}[Invarianza Topologica]
Siano $X, Y$ spazi topologici, $x_0 \in X$ e $f: X \to Y$ un \textbf{omeomorfismo}. Poniamo $y_0 = f(x_0) \in Y$.
Allora l'applicazione indotta:
\[
f_*: \Pi_1(X, x_0) \longrightarrow \Pi_1(Y, y_0)
\]
è un \textbf{isomorfismo di gruppi}.
\end{teo}

\begin{proof}
Poiché $f$ è un omeomorfismo, esiste l'inversa continua $f^{-1}: Y \to X$ tale che $f^{-1} \circ f = \mathrm{Id}_X$ e $f \circ f^{-1} = \mathrm{Id}_Y$.
Applicando le proprietà di composizione per le mappe indotte:
\[
(f^{-1} \circ f)_* = (f^{-1})_* \circ f_* = (\mathrm{Id}_X)_* = \mathrm{Id}_{\Pi_1(X, x_0)}
\]
\[
(f \circ f^{-1})_* = f_* \circ (f^{-1})_* = (\mathrm{Id}_Y)_* = \mathrm{Id}_{\Pi_1(Y, y_0)}
\]
Quindi $f_*$ è un omomorfismo invertibile, la cui inversa è $(f_*)^{-1} = (f^{-1})_*$. Pertanto $f_*$ è un isomorfismo di gruppi\hfill$\square$
\end{proof}

\begin{coro}
Siano $X, Y$ spazi topologici connessi per archi. Se $X \cong Y$ (sono omeomorfi), allora $\Pi_1(X) \cong \Pi_1(Y)$.
\end{coro}

\begin{observacion}
Il viceversa non vale in generale. Ad esempio, consideriamo $X = \mathbb{R}^n$ e $Y = \{y_0\}$. Chiaramente $\mathbb{R}^n \not\cong \{y_0\}$ per $n \geqslant 1$. Tuttavia, $\Pi_1(\mathbb{R}^n) \cong \Pi_1(\{y_0\}) = 0$.
\end{observacion}

\vspace{2em}

\section{Invarianza Omotopica del Gruppo Fondamentale}

\begin{lema}
Siano $X, Y$ spazi topologici e $f_1, f_2: X \to Y$ mappe continue tali che $H: f_1 \simeq f_2$ (con $H: X \times I \to Y$ l'omotopia).
Fissato $x_0 \in X$, sia $\alpha: I \to Y$ il cammino in $Y$ definito da $\alpha(s) = H(x_0, s)$, che connette $f_1(x_0)$ ad $f_2(x_0)$ (ossia $\alpha(0) = f_1(x_0)$ e $\alpha(1) = f_2(x_0)$).

Allora vale la seguente relazione tra gli omomorfismi indotti:
\[
\alpha_* \circ (f_1)_* = (f_2)_*
\]

\begin{figure}[H]
    \centering
    \begin{tikzcd}[scale=1.3]
        \Pi_1(X, x_0) \arrow[r, "(f_1)_*"] \arrow[dr, "(f_2)_*"'] & \Pi_1(Y, f_1(x_0)) \arrow[d, "\alpha_*"] \\
        & \Pi_1(Y, f_2(x_0))
    \end{tikzcd}
    \caption{Diagramma commutativo per l'omomorfismo indotto tra mappe omotope.}
\end{figure}
\end{lema}

\begin{proof}[Dimostrazione (Facoltativa)]
Dobbiamo dimostrare che per ogni $\sigma \in \Lambda_{x_0}(X)$ si ha $(f_2)_*([\sigma]) = (\alpha_* \circ (f_1)_*)([\sigma])$, ossia $[f_2 \circ \sigma] = [\alpha^{-1} * (f_1 \circ \sigma) * \alpha]$\hfill(*)

Sia $H: X \times I \to Y$ l'omotopia tra $f_1$ e $f_2$.
Notiamo che il cammino coniugato si scrive come:
\[
(\alpha^{-1} * (f_1 \circ \sigma) * \alpha)(t) = 
\begin{cases}
    \alpha^{-1}(4t) = \alpha(1-4t) = H(x_0, 1-4t) & t \leqslant 1/4 \\
    (f_1 \circ \sigma)(4t-1) = H(\sigma(4t-1), 0) & 1/4 \leqslant t \leqslant 1/2 \\
    \alpha(2t-1) = H(x_0, 2t-1) & t \geqslant 1/2
\end{cases}
\]

D'altra parte, il cammino $f_2 \circ \sigma$ verifica $f_2(\sigma(t)) = H(\sigma(t), 1)$ per ogni $t \in I$. Ponendo $y_0 = f_2(x_0)$, si ha anche la decomposizione con cappi costanti:
\[
f_2 \circ \sigma \simeq (c_{y_0} * (f_2 \circ \sigma)) * c_{y_0} \text{ rel } \partial I
\]
dove
\[
(c_{y_0} * (f_2 \circ \sigma) * c_{y_0})(t) = 
\begin{cases}
    y_0 = f_2(x_0) = H(x_0, 1) & t \leqslant 1/4 \\
    H(\sigma(4t-1), 1) & 1/4 \leqslant t \leqslant 1/2 \\
    H(x_0, 1) & t \geqslant 1/2
\end{cases}
\]

Definiamo ora la mappa di omotopia $K: I \times I \to Y$:
\[
K(t,s) = 
\begin{cases}
    H(x_0, 1 - 4t(1-s)) & t \leqslant 1/4 \\[0.5em]
    H(\sigma(4t-1), s) & 1/4 \leqslant t \leqslant 1/2 \\[0.5em]
    H(x_0, 1 - 2(1-t)(1-s)) & t \geqslant 1/2
\end{cases}
\]

Controlliamo la continuità sulle giunzioni $t = 1/4$ e $t = 1/2$:
\begin{itemize}
    \item Per $t = 1/4$: $H(x_0, 1 - 1(1-s)) = H(x_0, s) = H(\sigma(0), s)$.
    \item Per $t = 1/2$: $H(\sigma(1), s) = H(x_0, s) = H(x_0, 1 - 2(1/2)(1-s)) = H(x_0, s)$.
\end{itemize}
Per il Lemma di Incollamento, $K$ è ben definita e continua.

Verifichiamo le condizioni ai bordi per la variabile $s$:
\begin{enumerate}[label=\roman*)]
    \item Per $s = 0$: $K(t,0) = (\alpha^{-1} * (f_1 \circ \sigma) * \alpha)(t)$.
    \item Per $s = 1$: $K(t,1) = (c_{y_0} * (f_2 \circ \sigma) * c_{y_0})(t)$.
    \item Per ogni $s \in I$: $K(0,s) = H(x_0, 1) = f_2(x_0) = y_0$ e $K(1,s) = H(x_0, 1) = y_0$.
\end{enumerate}
Pertanto $K: \alpha^{-1} * (f_1 \circ \sigma) * \alpha \simeq c_{y_0} * (f_2 \circ \sigma) * c_{y_0} \text{ rel } \partial I$. Poich\'e, per la propriet\`a dell'elemento neutro, $c_{y_0} * (f_2 \circ \sigma) * c_{y_0} \simeq f_2 \circ \sigma \text{ rel } \partial I$, per transitivit\`a si ottiene $\alpha^{-1} * (f_1 \circ \sigma) * \alpha \simeq f_2 \circ \sigma \text{ rel } \partial I$, cio\`e (*), ossia $\alpha_* \circ (f_1)_* = (f_2)_*$\hfill$\square$
\end{proof}

\vspace{1.5em}

\begin{teo}[Invarianza Omotopica del Gruppo Fondamentale]
Siano $X, Y$ spazi topologici omotopicamente equivalenti ($X \simeq Y$), e siano $f: X \to Y$ e $g: Y \to X$ mappe continue tali che $g \circ f \simeq \mathrm{Id}_X$ e $f \circ g \simeq \mathrm{Id}_Y$.
Allora per ogni $x_0 \in X$, l'applicazione indotta:
\[
f_*: \Pi_1(X, x_0) \longrightarrow \Pi_1(Y, f(x_0))
\]
è un \textbf{isomorfismo di gruppi}.
\end{teo}

\begin{proof}
\textbf{Passo 1: $f_*$ \`e iniettiva (e $g_*$ \`e suriettiva).} Sia $H: \mathrm{Id}_X \simeq g \circ f$ e sia $\alpha(s) = H(x_0, s)$, cammino da $x_0$ a $(g \circ f)(x_0)$. Applicando il Lemma a $f_1 = \mathrm{Id}_X$ e $f_2 = g \circ f$ si ottiene il diagramma commutativo

\begin{center}
\begin{tikzcd}
\Pi_1(X, x_0) \arrow[r, "(\mathrm{Id}_X)_*"] \arrow[dr, "(g \circ f)_*"'] & \Pi_1(X, x_0) \arrow[d, "\alpha_*"] \\
 & \Pi_1(X, (g \circ f)(x_0))
\end{tikzcd}
\end{center}

cio\`e $(g \circ f)_* = \alpha_* \circ (\mathrm{Id}_X)_* = \alpha_*$, che \`e un isomorfismo (Teorema del cambio di punto base). Poich\'e $(g \circ f)_* = g_* \circ f_*$, ne segue che $f_*: \Pi_1(X, x_0) \to \Pi_1(Y, f(x_0))$ \`e \textbf{iniettiva} e che $g_*: \Pi_1(Y, f(x_0)) \to \Pi_1(X, (g\circ f)(x_0))$ \`e \textbf{suriettiva}.

\textbf{Passo 2: $g_*$ \`e iniettiva.} Il ruolo di $f$ e $g$ \`e simmetrico: da $f \circ g \simeq \mathrm{Id}_Y$ e dal Passo 1 (applicato a $g$, con punto base $y_1 = f(x_0)$ e omotopia $K: \mathrm{Id}_Y \simeq f \circ g$) segue che $g_*: \Pi_1(Y, y_1) \to \Pi_1(X, g(y_1))$ \`e iniettiva.

\textbf{Passo 3: conclusione.} Per i Passi 1 e 2, $g_*: \Pi_1(Y, f(x_0)) \to \Pi_1(X, (g\circ f)(x_0))$ \`e un isomorfismo. Da $g_* \circ f_* = \alpha_*$ si ricava
\[
f_* = (g_*)^{-1} \circ \alpha_* ,
\]
composizione di due isomorfismi: dunque $f_*$ \`e un isomorfismo di gruppi\hfill$\square$
\end{proof}

\begin{observacion}[Sull'enunciato negli appunti]
Negli appunti il teorema \`e enunciato per ogni $x_0 \in g(Y)$, e la dimostrazione usa un $y_0 \in Y$ con $g(y_0) = x_0$: \emph{per quell'ipotesi l'argomento degli appunti \`e corretto}. Infatti, per $x_0 = g(y_0)$, sia $K: \mathrm{Id}_Y \simeq f \circ g$ e $\beta(s) = K(y_0, s)$ (da $y_0$ a $f(g(y_0)) = f(x_0)$): per il Lemma, $f_{*x_0} \circ g_{*y_0} = (f \circ g)_* = \beta_* \circ (\mathrm{Id}_Y)_* = \beta_*$ \`e un isomorfismo, quindi $f_{*x_0}$ \`e suriettiva (ed \`e iniettiva per il Passo 1). (Nella trascrizione automatica l'ipotesi $x_0 \in g(Y)$ si era persa, e nella precedente revisione l'argomento era stato giudicato erroneo per questo motivo: la valutazione era ingiusta.) La dimostrazione data qui sopra, invece, prova l'enunciato pi\`u forte, valido per \emph{ogni} $x_0 \in X$, anche fuori da $g(Y)$.
\end{observacion}

\begin{coro}
Siano $X, Y$ spazi topologici connessi per archi. Se $X \simeq Y$ (sono omotopicamente equivalenti), allora $\Pi_1(X) \cong \Pi_1(Y)$.
\end{coro}

\begin{observacion}
Il viceversa non è vero in generale. Ad esempio, prendendo $X = S^2$ e $Y = \{y_0\}$, si ha $\Pi_1(S^2) = 0 = \Pi_1(\{y_0\})$, ma $S^2 \not\simeq \{y_0\}$ in quanto $S^2$ non è uno spazio contraibile ($\mathrm{Id}_{S^2} \not\simeq -\mathrm{Id}_{S^2}$).
\end{observacion}

\vspace{1.5em}

\begin{ejemplo}[Calcolo dei Gruppi Fondamentali per Spazi Notevoli]
\begin{enumerate}
    \item Per ogni spazio topologico contraibile $X$, si ha $\Pi_1(X) = 0$.
    Infatti, $X \simeq \{y_0\} \implies \Pi_1(X) \cong \Pi_1(\{y_0\}) = 0$.
    Esempi tipici: $\mathbb{R}^n$, insiemi convessi $\Omega \subseteq \mathbb{R}^n$ e insiemi stellati $S \subseteq \mathbb{R}^n$.

    \item Per lo spazio euclideo bucato nell'origine:
    \[
    \Pi_1(\mathbb{R}^n \setminus \{\bar{0}\}) \cong \Pi_1(S^{n-1})
    \]
    poiché $S^{n-1}$ è un retratto di deformazione di $\mathbb{R}^n \setminus \{\bar{0}\}$.

    \item Per la sfera privata di due punti distinti $p \neq q$:
    \[
    \Pi_1(S^n \setminus \{p,q\}) \cong \Pi_1(S^{n-1})
    \]
    in quanto $S^n \setminus \{p,q\} \cong S^n \setminus \{N, S\} \simeq S^{n-1}$.
\end{enumerate}
\end{ejemplo}

\section{Il Gruppo Fondamentale dello Spazio Prodotto}

\begin{teo}
Siano $X, Y$ spazi topologici, $x_0 \in X$ e $y_0 \in Y$.
Allora il gruppo fondamentale del prodotto cartesiano $X \times Y$ è isomorfo al prodotto diretto dei gruppi fondamentali:
\[
\Pi_1(X \times Y, (x_0, y_0)) \cong \Pi_1(X, x_0) \times \Pi_1(Y, y_0)
\]
\end{teo}

\begin{proof}
Siano $p_1: X \times Y \to X$ e $p_2: X \times Y \to Y$ le proiezioni canoniche definite da $p_1(x,y) = x$ e $p_2(x,y) = y$.
Essendo $p_1$ e $p_2$ continue, esse inducono gli omomorfismi di gruppi $p_{1*}: \Pi_1(X \times Y, (x_0, y_0)) \to \Pi_1(X, x_0)$ e $p_{2*}: \Pi_1(X \times Y, (x_0, y_0)) \to \Pi_1(Y, y_0)$.

Definiamo l'applicazione:
\Func{\Phi}{\Pi_1(X \times Y, (x_0, y_0))}{\Pi_1(X, x_0) \times \Pi_1(Y, y_0)}{[\eta]}{(p_{1*}([\eta]), \, p_{2*}([\eta])) = ([p_1 \circ \eta], \, [p_2 \circ \eta])}

Verifichiamo che $\Phi$ sia un isomorfismo di gruppi:

\begin{itemize}
    \item \textbf{$\Phi$ è un omomorfismo:} Per ogni $[\eta], [\eta'] \in \Pi_1(X \times Y, (x_0, y_0))$:
    \[
    \begin{aligned}
        \Phi([\eta] \cdot [\eta']) &= (p_{1*}([\eta] \cdot [\eta']), \, p_{2*}([\eta] \cdot [\eta'])) \\
        &= (p_{1*}([\eta]) \cdot p_{1*}([\eta']), \, p_{2*}([\eta]) \cdot p_{2*}([\eta'])) \\
        &= (p_{1*}([\eta]), \, p_{2*}([\eta])) \cdot (p_{1*}([\eta']), \, p_{2*}([\eta'])) = \Phi([\eta]) \cdot \Phi([\eta'])
    \end{aligned}
    \]

    \item \textbf{$\Phi$ è suriettiva:} Siano $\sigma \in \Lambda_{x_0}(X)$ e $\tau \in \Lambda_{y_0}(Y)$.
    Definiamo il cappio $\eta: I \to X \times Y$ tramite $\eta(t) = (\sigma(t), \tau(t))$. Poiché le sue componenti sono continue, $\eta$ è continua e $\eta \in \Lambda_{(x_0, y_0)}(X \times Y)$.
    Calcolando l'immagine:
    \[
    \Phi([\eta]) = ([p_1 \circ \eta], [p_2 \circ \eta]) = ([\sigma], [\tau])
    \]
    Quindi $\Phi$ è suriettiva.

    \item \textbf{$\Phi$ è iniettiva:} Siano $\eta, \eta' \in \Lambda_{(x_0, y_0)}(X \times Y)$ tali che $\Phi([\eta]) = \Phi([\eta'])$.
    Allora $(p_{1*}([\eta]), p_{2*}([\eta])) = (p_{1*}([\eta']), p_{2*}([\eta'])) \implies [p_1 \circ \eta] = [p_1 \circ \eta']$ e $[p_2 \circ \eta] = [p_2 \circ \eta']$.
    
    Quindi esistono due omotopie relative al bordo:
    \[
    H: p_1 \circ \eta \simeq p_1 \circ \eta' \text{ rel } \partial I \quad \text{e} \quad K: p_2 \circ \eta \simeq p_2 \circ \eta' \text{ rel } \partial I
    \]
    Poniamo l'omotopia prodotto $W: I \times I \to X \times Y$ definita da $W(t,s) = (H(t,s), K(t,s))$.
    $W$ è continua e soddisfa:
    \begin{enumerate}[label=\roman*)]
        \item $W(t,0) = (H(t,0), K(t,0)) = ((p_1 \circ \eta)(t), (p_2 \circ \eta)(t)) = \eta(t)$.
        \item $W(t,1) = (H(t,1), K(t,1)) = ((p_1 \circ \eta')(t), (p_2 \circ \eta')(t)) = \eta'(t)$.
        \item $W(0,s) = (x_0, y_0) = W(1,s)$ per ogni $s \in I$.
    \end{enumerate}
    Pertanto $W: \eta \simeq \eta' \text{ rel } \partial I \implies [\eta] = [\eta']$, dimostrando l'iniettività.
\end{itemize}
In conclusione, $\Phi$ è un isomorfismo di gruppi\hfill$\square$
\end{proof}

\vspace{1.5em}

\begin{ejemplo}
Calcolo del gruppo fondamentale per prodotti topologici notevoli:

\begin{enumerate}
    \item \textbf{Cilindro vuoto:} $X = S^1 \times I$ oppure $S^1 \times \mathbb{R}$.
    \[
    \Pi_1(S^1 \times I) \cong \Pi_1(S^1) \times \Pi_1(I) \cong \mathbb{Z} \times \{0\} \cong \mathbb{Z}
    \]
    Geometricamente, $\mathbb{Z}$ conta il numero di avvolgimenti (con segno) che un cappio compie attorno a $S^1$.

    \item \textbf{Toro bidimensionale (Superficie torica):} $X = T^2 = S^1 \times S^1$.
    \[
    \Pi_1(T^2) = \Pi_1(S^1 \times S^1) \cong \Pi_1(S^1) \times \Pi_1(S^1) \cong \mathbb{Z} \times \mathbb{Z} = \mathbb{Z}^2
    \]

    \item \textbf{Toro pieno:} $X = S^1 \times B^2$.
    Poiché il disco chiuso $B^2$ è convesso (e quindi contraibile, $\Pi_1(B^2) = 0$):
    \[
    \Pi_1(S^1 \times B^2) \cong \Pi_1(S^1) \times \Pi_1(B^2) \cong \mathbb{Z} \times \{0\} \cong \mathbb{Z}
    \]
\end{enumerate}

\begin{figure}[H]
    \centering
    \begin{subfigure}[b]{0.3\textwidth}
        \centering
        \begin{tikzpicture}[scale=0.8]
            \draw[thick] (0,0) ellipse (0.8cm and 0.3cm);
            \draw[thick] (0,2) ellipse (0.8cm and 0.3cm);
            \draw[thick] (-0.8,0) -- (-0.8,2);
            \draw[thick] (0.8,0) -- (0.8,2);
        \end{tikzpicture}
        \caption{Cilindro $S^1 \times I$}
    \end{subfigure}
    \hfill
    \begin{subfigure}[b]{0.3\textwidth}
        \centering
        \begin{tikzpicture}[scale=0.8]
            \draw[thick] (0,0) ellipse (1.2cm and 0.5cm);
            \draw (-0.5,0) arc (180:360:0.5cm and 0.2cm);
            \draw (-0.35,-0.1) arc (180:0:0.35cm and 0.15cm);
        \end{tikzpicture}
        \caption{Toro $T^2 = S^1 \times S^1$}
    \end{subfigure}
    \hfill
    \begin{subfigure}[b]{0.3\textwidth}
        \centering
        \begin{tikzpicture}[scale=0.8]
            \draw[thick, fill=gray!20] (0,0) ellipse (1.2cm and 0.5cm);
            \draw (-0.5,0) arc (180:360:0.5cm and 0.2cm);
            \draw (-0.35,-0.1) arc (180:0:0.35cm and 0.15cm);
        \end{tikzpicture}
        \caption{Toro pieno $S^1 \times B^2$}
    \end{subfigure}
    \caption{Rappresentazione grafica degli spazi prodotto.}
\end{figure}
\end{ejemplo}

\vspace{2em}

\section{Il Toro $n$-Dimensionale}

\begin{definicion}[\textbf{IMPORTANTE PER L'ESAME}]
Si chiama \textbf{toro $n$-dimensionale} lo spazio topologico quoziente:
\[
T^n = \mathbb{R}^n \big/ \mathbb{Z}^n
\]
dotato della topologia quoziente indotta dalla topologia naturale su $\mathbb{R}^n$.

Un ruolo fondamentale è svolto dalla \textbf{proiezione canonica suriettiva}:
\Func{\pi}{\mathbb{R}^n}{T^n = \mathbb{R}^n / \mathbb{Z}^n}{x}{[x]_{\mathbb{Z}^n} = x + \mathbb{Z}^n}

Per ogni $x, y \in \mathbb{R}^n$, la relazione di equivalenza è data da:
\[
x \sim_\pi y \iff \pi(x) = \pi(y) \iff x - y \in \mathbb{Z}^n \iff x_i - y_i \in \mathbb{Z}, \quad \forall i = 1, \dots, n
\]
\end{definicion}

\begin{prop}
Sia $Q = I^n = [0,1]^n$ il cubo unitario $n$-dimensionale.
La restrizione della proiezione canonica $\pi|_Q: Q \to T^n$ è continua e suriettiva, e identifica $T^n \cong Q / \sim$.
\end{prop}

\begin{proof}
Dimostriamo che $\pi(Q) = T^n$:
\begin{itemize}
    \item [$\subseteq$)] Banale poiché $Q \subset \mathbb{R}^n \implies \pi(Q) \subseteq \pi(\mathbb{R}^n) = T^n$.
    \item [$\supseteq$)] Sia $y = (y_1, \dots, y_n) \in \mathbb{R}^n$. Per ogni $i = 1, \dots, n$, definiamo $x_i = y_i - \lfloor y_i \rfloor$, dove $\lfloor y_i \rfloor$ indica la parte intera.
    Allora $x_i \in [0,1]$ per ogni $i$, il che implica $x = (x_1, \dots, x_n) \in Q = I^n$.
    Inoltre $y - x = (\lfloor y_1 \rfloor, \dots, \lfloor y_n \rfloor) \in \mathbb{Z}^n$, quindi $\pi(x) = \pi(y)$.
\end{itemize}
Pertanto $\pi(Q) = T^n$, ovvero la mappa è suriettiva.

\textbf{Identificazione $T^n \cong Q/\!\sim$.} Su $Q$ la relazione indotta da $\pi$ \`e $x \sim y \iff x - y \in \mathbb{Z}^n$: poich\'e $x,y \in [0,1]^n$, ci\`o accade solo se $x = y$ oppure se $x,y$ stanno su facce opposte del cubo con le restanti coordinate uguali (ad esempio, per $n=2$: $(0,y) \sim (1,y)$ e $(x,0) \sim (x,1)$). Inoltre $\pi|_Q$ \`e una applicazione continua e suriettiva da uno spazio compatto ($Q$) a uno spazio di Hausdorff ($T^n \cong S^1 \times \dots \times S^1$, vedi il teorema seguente); dunque \`e chiusa, e quindi una applicazione quoziente. Ne segue che $\pi|_Q$ induce un omeomorfismo $Q/\!\sim\ \cong T^n$\hfill$\square$
\end{proof}

\begin{coro}
Il toro $n$-dimensionale $T^n$ è uno spazio topologico \textbf{compatto}.
Infatti, $T^n = \pi(Q)$ è l'immagine continua di $Q = [0,1]^n$, che è compatto in $\mathbb{R}^n$ per il Teorema di Heine-Borel.
\end{coro}

\vspace{1.5em}

\begin{teo}
Il toro $n$-dimensionale è omeomorfo al prodotto cartesiano di $n$ circonferenze:
\[
T^n \cong \underbrace{S^1 \times \dots \times S^1}_{n \text{ volte}}
\]
\end{teo}

\begin{proof}
Analizziamo prima il caso $n = 1$:
Consideriamo la mappa esponenziale $p: \mathbb{R} \to S^1$ data da $p(t) = e^{2\pi i t} = (\cos(2\pi t), \sin(2\pi t))$.
$p$ è continua, suriettiva e aperta (\`e aperta perch\'e l'immagine di un intervallo aperto di lunghezza $\leqslant 1$ \`e un arco aperto di $S^1$, e ogni aperto di $\mathbb{R}$ \`e unione di tali intervalli).
Due punti $t, t' \in \mathbb{R}$ soddisfano $p(t) = p(t') \iff e^{2\pi i t} = e^{2\pi i t'} \iff t - t' \in \mathbb{Z}$.
Per il Teorema Fondamentale dell'Omeomorfismo Quoziente, $p$ induce un omeomorfismo $\tilde{p}: \mathbb{R}/\mathbb{Z} \xrightarrow{\cong} S^1$.

In generale per $n \geqslant 1$, definiamo la mappa:
\Func{\Phi}{\mathbb{R}^n}{\underbrace{S^1 \times \dots \times S^1}_{n \text{ volte}}}{(x_1, \dots, x_n)}{(p(x_1), \dots, p(x_n))}
$\Phi$ è ben posta, continua, suriettiva e aperta.
Inoltre, $\Phi(x) = \Phi(y) \iff x - y \in \mathbb{Z}^n$.
Pertanto, $\Phi$ induce un omeomorfismo sul quoziente:
\[
\widetilde{\Phi}: \mathbb{R}^n / \mathbb{Z}^n \xrightarrow{\;\cong\;} S^1 \times \dots \times S^1
\]
\end{proof}

\begin{coro}
Il gruppo fondamentale del toro $n$-dimensionale è:
\[
\Pi_1(T^n) \cong \Pi_1(S^1 \times \dots \times S^1) \cong \underbrace{\mathbb{Z} \times \dots \times \mathbb{Z}}_{n \text{ volte}} = \mathbb{Z}^n
\]
\end{coro}

\begin{figure}[H]
    \centering
    \begin{tikzpicture}[scale=1.5]
        % Cuadrado Q
        \draw[thick] (0,0) -- (2,0) node[midway, below] {$(x,0)$} -- (2,2) node[midway, right] {$(1,y)$} -- (0,2) node[midway, above] {$(x,1)$} -- (0,0) node[midway, left] {$(0,y)$};
        
        % Flechas de identificación
        \draw[thick, red, ->, >=stealth] (0,0) -- (0,1.2);
        \draw[thick, red, ->, >=stealth] (2,0) -- (2,1.2);
        \draw[thick, blue, ->, >=stealth] (0,0) -- (1.2,0);
        \draw[thick, blue, ->, >=stealth] (0,2) -- (1.2,2);

        \node at (1,1) {$Q = I \times I$};

        % Proceso de enrollado
        \draw[->, >=stealth, ForestGreen, ultra thick] (2.3, 1) -- (3.3, 1) node[midway, above] {\tiny incollo i lati};

        % Cilindro intermedio
        \draw[thick] (4,0.5) ellipse (0.3cm and 0.6cm);
        \draw[thick] (5.5,0.5) ellipse (0.3cm and 0.6cm);
        \draw[thick] (4,1.1) -- (5.5,1.1);
        \draw[thick] (4,-0.1) -- (5.5,-0.1);

        \draw[->, >=stealth, ForestGreen, ultra thick] (5.8, 1) -- (6.6, 1);

        % Toro
        \draw[thick] (7.5,0.8) ellipse (0.6cm and 0.3cm);
        \draw (7.25,0.8) arc (180:360:0.25cm and 0.1cm);
    \end{tikzpicture}
    \caption{Costruzione del toro $T^2$ identificando i lati opposti del quadrato $Q = I \times I$.}
\end{figure}

\begin{ejercicio*}[\textbf{DOMANDA D'ESAME}]
\textbf{Quali sono i modelli topologici del toro $T^n$?}
\begin{enumerate}
    \item Il quoziente dello spazio euclideo: $\mathbb{R}^n / \mathbb{Z}^n$.
    \item Il prodotto cartesiano di circonferenze: $S^1 \times \dots \times S^1$.
    \item Il quoziente del cubo unitario con le facce opposte identificate: $Q/\!\sim\ = I^n/\!\sim$.
\end{enumerate}
\end{ejercicio*}

\section{Spazi Semplicemente Connessi}

\begin{definicion}
Uno spazio topologico $X$ si dice \textbf{semplicemente connesso} se:
\begin{enumerate}
    \item $X$ è connesso per archi, e
    \item $\Pi_1(X) = 0$ (ossia ha gruppo fondamentale banale).
\end{enumerate}
\end{definicion}

\begin{observacion}
La condizione $\Pi_1(X, x_0) = 0$ equivale a dire che per ogni coppia di cammini $\alpha, \beta: I \to X$ con $\alpha(0) = \beta(0) = x_0$ e $\alpha(1) = \beta(1)$, si ha $\alpha \simeq \beta \text{ rel } \partial I$.
\end{observacion}

\begin{prop}
Se $X$ è uno spazio topologico contraibile, allora $X$ è semplicemente connesso.
\end{prop}

\begin{proof}
Se $X$ è contraibile, allora $X$ è connesso per archi. Inoltre $X \simeq \{x_0\}$, per cui $\Pi_1(X) \cong \Pi_1(\{x_0\}) = 0$. Pertanto $X$ è semplicemente connesso.
\end{proof}

\begin{observacion}
Il viceversa della proposizione precedente \textbf{non} vale: semplicemente connesso $\centernot\implies$ contraibile.

Come controesempio, si consideri la sfera 2-dimensionale $S^2$. Si ha che $S^2$ è semplicemente connessa ($\Pi_1(S^2) = 0$), ma $S^2$ non è contraibile in quanto $\mathrm{Id}_{S^2} \not\simeq -\mathrm{Id}_{S^2}$.
\end{observacion}

\vspace{1.5em}


\vspace{1.5em}

\begin{lema}[Numero di Lebesgue, per l'intervallo]
Sia $\{A, B\}$ un ricoprimento aperto di $I = [0,1]$. Allora esiste $\delta > 0$ tale che ogni sottoinsieme di $I$ di diametro $< \delta$ \`e contenuto in $A$ oppure in $B$. In particolare esiste una suddivisione $0 = t_0 < t_1 < \dots < t_k = 1$ tale che ciascun intervallo $[t_{r-1}, t_r]$ \`e contenuto in $A$ oppure in $B$.
\end{lema}

\begin{proof}
Se $A = I$ o $B = I$ non c'\`e nulla da dimostrare. Altrimenti la funzione $\ell(t) = \max\{ d(t, I \setminus A),\, d(t, I \setminus B) \}$ \`e continua su $I$ e strettamente positiva (se $t \in A$, che \`e aperto, allora $d(t, I\setminus A) > 0$; analogamente per $B$). Per compattezza ha minimo $\delta > 0$. Se $\mathrm{diam}(T) < \delta$ e $t \in T$, supponiamo ad esempio $\ell(t) = d(t, I \setminus A) \geqslant \delta$: allora $T \subseteq \{u : |u - t| < \delta\} \subseteq A$. La suddivisione si ottiene prendendo $k$ con $1/k < \delta$ e $t_r = r/k$.
\end{proof}

\begin{teo}
La sfera $S^n$ è semplicemente connessa se e solo se $n \geqslant 2$.
\end{teo}

\begin{proof}
\begin{itemize}
    \item[$\implies$)] Mostriamo che per $n = 0, 1$ la sfera non \`e semplicemente connessa:
    \begin{itemize}
        \item Per $n = 0$: $S^0 = \{-1, 1\}$ non è connessa per archi.
        \item Per $n = 1$: $\Pi_1(S^1) \cong \mathbb{Z} \neq 0$ (dimostrato pi\`u avanti con Seifert--Van Kampen).
    \end{itemize}

    \item[$\impliedby$)] Sia $n \geqslant 2$.

    \textbf{Connessione per archi.} Per $x, y \in S^n$ con $y \neq -x$ il segmento $\overline{xy}$ non passa per l'origine, quindi $t \mapsto \frac{(1-t)x + ty}{\|(1-t)x + ty\|}$ \`e un cammino in $S^n$ da $x$ a $y$. Se $y = -x$ si sceglie $z \in S^n \setminus \{x, -x\}$ (esiste perch\'e $n \geqslant 2$) e si concatenano un cammino da $x$ a $z$ e uno da $z$ a $-x$.

    \textbf{Setup.} Siano $p \neq q$ in $S^n$, $U = S^n \setminus \{p\}$, $V = S^n \setminus \{q\}$. Sono aperti, $U \cong V \cong \mathbb{R}^n$ (proiezione stereografica), quindi contraibili e con $\Pi_1 = 0$. Inoltre $W = U \cap V = S^n \setminus \{p,q\} \cong \mathbb{R}^n \setminus \{\bar 0\}$ \`e connesso per archi perch\'e $n \geqslant 2$. Fissiamo $x_0 \in W$ e un cappio $\sigma$ in $x_0$: dobbiamo mostrare che $\sigma \simeq c_{x_0}$ rel $\partial I$.

    \textbf{Suddivisione.} $A = \sigma^{-1}(U)$ e $B = \sigma^{-1}(V)$ sono un ricoprimento aperto di $I$. Per il Lemma di Lebesgue esiste una suddivisione di $I$ in intervalli ciascuno inviato da $\sigma$ in $U$ oppure in $V$. Raggruppando intervalli consecutivi inviati nello stesso aperto, otteniamo $0 = t_0 < t_1 < \dots < t_h = 1$ con $\delta_j = [t_{j-1}, t_j]$ tali che $\sigma(\delta_j)$ sta in $U$ e $\sigma(\delta_{j+1})$ in $V$ o viceversa (alternati). Posto $x_j = \sigma(t_j)$, per $1 \leqslant j \leqslant h-1$ si ha
    \[
    x_j \in \sigma(\delta_j) \cap \sigma(\delta_{j+1}) \subseteq U \cap V = W .
    \]

    \textbf{Cammini di collegamento.} Poich\'e $W$ \`e connesso per archi, per $1 \leqslant j \leqslant h-1$ esiste $\tau_j: I \to W$ con $\tau_j(0) = x_0$, $\tau_j(1) = x_j$. Poniamo $\tau_0 = \tau_h = c_{x_0}$ (si ha $x_0 = \sigma(0) = \sigma(1)$).

    \begin{center}
    \begin{tikzpicture}[scale=1.3]
        \draw[thick] (0,0) circle (1.8cm);
        \node[above right] at (1.2,1.3) {$S^n$};
        \filldraw[blue] (-1.7, 0) circle (1.5pt) node[left] {$x_0$};
        \draw[thick, red] (-1.7,0) to[out=60, in=180] (-0.5,0.7) to[out=0, in=120] (0.8,0.5) to[out=-60, in=20] (0.5,-1.2) to[out=200, in=-60] (-1.7,0);
        \filldraw (-0.5,0.7) circle (1pt) node[above] {$x_1$};
        \filldraw (0.8,0.5) circle (1pt) node[right] {$x_2$};
        \draw[ForestGreen, ->, >=stealth] (-1.7,0) to[out=20, in=220] (-0.5,0.7);
        \node[ForestGreen, below] at (-1.1,0.3) {$\tau_1$};
        \draw[Purple, ->, >=stealth] (-1.7,0) to[out=-20, in=200] (0.8,0.5);
        \node[Purple, below] at (-0.3, 0.1) {$\tau_2$};
    \end{tikzpicture}
    \end{center}

    \textbf{Decomposizione.} Sia $\sigma_j = \sigma|_{\delta_j} \circ \varphi_j$, con $\varphi_j(s) = (1-s)t_{j-1} + s t_j$: \`e un cammino da $x_{j-1}$ a $x_j$ con sostegno $\sigma(\delta_j)$. La concatenazione $\sigma_1 * \dots * \sigma_h$ \`e $\sigma \circ \varphi$ per una opportuna riparametrizzazione $\varphi: I \to I$ (lineare a tratti, con $\varphi(0)=0$, $\varphi(1)=1$), e $\sigma \circ \varphi \simeq \sigma$ rel $\partial I$ tramite $(t,s) \mapsto \sigma((1-s)\varphi(t) + s t)$. Dunque
    \[
    \sigma \simeq \sigma_1 * \dots * \sigma_h \simeq (\tau_0 * \sigma_1 * \tau_1^{-1}) * (\tau_1 * \sigma_2 * \tau_2^{-1}) * \dots * (\tau_{h-1} * \sigma_h * \tau_h^{-1}) \quad \text{rel } \partial I ,
    \]
    dove si sono inseriti i cammini $\tau_j^{-1} * \tau_j \simeq c_{x_j}$ e usate le propriet\`a di neutro e associativit\`a.

    \textbf{Conclusione.} Ogni blocco $\lambda_j = \tau_{j-1} * \sigma_j * \tau_j^{-1}$ \`e un cappio in $x_0$ contenuto interamente in $U$ oppure in $V$ (i $\tau$ stanno in $W \subseteq U \cap V$). Poich\'e $\Pi_1(U) = \Pi_1(V) = 0$, ciascun $\lambda_j \simeq c_{x_0}$ rel $\partial I$. Quindi $\sigma \simeq c_{x_0} * \dots * c_{x_0} \simeq c_{x_0}$ rel $\partial I$, cio\`e $\Pi_1(S^n) = 0$\hfill$\square$
\end{itemize}
\end{proof}

\vspace{2em}

\section{Esempio di Spazio Topologico con Gruppo Fondamentale Non Abeliano}

\begin{definicion}[Gruppo Libero su $m$ Generatori]
Sia $\Sigma = \{a_1, \dots, a_m\}$ un insieme finito (alfabeto di generatori).
Consideriamo i simboli e i loro rispettivi inversi $\{a_1, a_1^{-1}, \dots, a_m, a_m^{-1}\}$ unitamente al simbolo d'identità $e$, ed imponiamo le regole di riduzione:
\[
a_i e = a_i = e a_i, \quad a_i a_i^{-1} = e = a_i^{-1} a_i, \quad a_i^{-1} e = a_i^{-1} = e a_i^{-1}, \quad \forall i = 1, \dots, m
\]

Sia $k \in \mathbb{Z}$. Per ogni $i = 1, \dots, m$, la potenza $a_i^k$ è definita da:
\[
a_i^k = 
\begin{cases}
    \underbrace{a_i \cdots a_i}_{k \text{ volte}} & \text{se } k > 0 \\[1em]
    e & \text{se } k = 0 \\[0.5em]
    \underbrace{a_i^{-1} \cdots a_i^{-1}}_{|k| \text{ volte}} & \text{se } k < 0
\end{cases}
\]
In particolare, vale la proprietà usuale delle potenze $a_i^k \cdot a_i^h = a_i^{k+h}$ per ogni $k, h \in \mathbb{Z}$.
\end{definicion}

%%%
\section{Struttura del Gruppo Libero Generato da un Insieme Finito}

\begin{definicion}
Si chiama \textbf{gruppo libero generato dall'insieme} $\Sigma = \{a_1, \dots, a_m\}$ l'insieme $F_m$ costituito da tutte le sequenze finite (parole ridotte) degli elementi $a_i$ e delle loro inverse $a_i^{-1}$:
\[
F_m = \left\{ a_{i_1}^{k_1} a_{i_2}^{k_2} \cdots a_{i_r}^{k_r} \;\middle|\; r \in \mathbb{N}, \; i_j \in \{1, \dots, m\}, \; k_j \in \mathbb{Z} \setminus \{0\}, \; i_j \neq i_{j+1} \right\}
\]
con l'operazione di \textbf{concatenazione seguita da riduzione} (si cancellano le coppie $a_i^{k} a_i^{-k}$ e si fondono le potenze adiacenti della stessa lettera, $a_i^k a_i^h = a_i^{k+h}$); l'elemento neutro \`e la parola vuota $e$:
\Func{\cdot}{F_m \times F_m}{F_m}{(a_{i_1}^{k_1} \cdots a_{i_r}^{k_r}, \, a_{j_1}^{h_1} \cdots a_{j_s}^{h_s})}{a_{i_1}^{k_1} \cdots a_{i_r}^{k_r} \cdot a_{j_1}^{h_1} \cdots a_{j_s}^{h_s}}
\end{definicion}

\begin{observacion}
Analizziamo la struttura di $F_m$ nei casi a bassa dimensione:
\begin{enumerate}
    \item \textbf{Caso $m=1$ ($F_1$):} $\Sigma = \{a\}$. Ogni elemento si scrive come $a^k$ con $k \in \mathbb{Z}$.
    Essendo $a^k \cdot a^h = a^{k+h}$, si ha l'isomorfismo di gruppi abeliani:
    \[
    (F_1, \cdot) \cong (\mathbb{Z}, +)
    \]
    
    \item \textbf{Caso $m=2$ ($F_2$):} $\Sigma = \{a, b\}$. Gli elementi si scrivono come alternanze di potenze di $a$ e $b$ (ad esempio $a^{k_1} b^{h_1} \cdots a^{k_r} b^{h_r}$).
    Poiché $ab \neq ba$, $F_2$ \textbf{non è un gruppo commutativo}.
    
    \item \textbf{Caso generale $m \geqslant 2$ ($F_m$):} Poich\'e $F_2$ \`e (isomorfo a) un sottogruppo di $F_m$ per ogni $m \geqslant 2$, e un sottogruppo di un gruppo abeliano \`e abeliano, il gruppo libero $F_m$ \textbf{non è mai abeliano} per $m \geqslant 2$.
\end{enumerate}
\end{observacion}

\vspace{1.5em}


\begin{definicion}[Prodotto libero]
Siano $G, H$ gruppi. Poniamo $S(G,H)$ l'insieme delle ``parole'' formate con elementi di $G$ e $H$, cio\`e
\[
w \in S(G,H) \iff w = b_1 b_2 \cdots b_n \ \text{ con } n \in \mathbb{N},\ b_i \in G \cup H \ \ \forall i
\]
(dove la parola vuota corrisponde a $n = 0$). In $S(G,H)$ consideriamo la relazione di equivalenza $\sim$, secondo cui $w \sim w'$ se una \`e ottenuta dall'altra mediante un numero finito di operazioni del tipo:
\begin{itemize}
    \item rimuovere $e_G$ o $e_H$, o inserirli;
    \item sostituire una coppia di elementi consecutivi, entrambi di $G$ o entrambi di $H$, con il risultato in $G$ o in $H$ del loro prodotto, o viceversa.
\end{itemize}
Si chiama \textbf{prodotto libero} di $G$ e $H$ l'insieme $G * H = S(G,H)/\!\sim$, munito dell'operazione $[w][w'] = [ww']$.
\end{definicion}

\begin{observacion}
$e_G \sim e_H \sim (\text{parola vuota})$ e $[e_G] = [e_H]$ \`e l'elemento neutro di $G * H$. Inoltre, per ogni $w = b_1 \cdots b_n \in S(G,H)$: $[w]^{-1} = [b_n^{-1} \cdots b_1^{-1}]$.

Equivalentemente (forma normale): ogni classe contiene una \emph{unica} parola ridotta, cio\`e della forma $g_1 h_1 g_2 h_2 \cdots$ con lettere non banali alternativamente di $G \setminus \{e_G\}$ e $H \setminus \{e_H\}$. Cos\`i si pu\`o pensare a $G * H$ come all'insieme delle parole ridotte, con il prodotto dato da concatenazione seguita da riduzione. $G$ e $H$ sono sottogruppi di $G * H$.
\end{observacion}

\vspace{1.5em}
\section{Proprietà del Prodotto Libero di Gruppi}

\begin{prop}
Siano $G, H, K$ gruppi. Il prodotto libero di gruppi soddisfa le seguenti proprietà:
\begin{enumerate}
    \item $G * H \cong H * G$ (commutatività) e $G * (H * K) \cong (G * H) * K$ (associatività).
    Ha quindi senso scrivere il prodotto libero di $p$ gruppi $G_1 * G_2 * \dots * G_p$.
    \item $G * 0 \cong G \cong 0 * G$.
    \item Se $G \neq 0$, allora non esiste alcun gruppo $H$ tale che $G * H = 0$.
    \item Se $G \neq 0$ e $H \neq 0$, allora $G * H$ ha un numero infinito di elementi.
    Siano $a \in G \setminus \{e_G\}$ e $b \in H \setminus \{e_H\}$. In $G * H$ si possono formare le seguenti parole ridotte distinte e di lunghezza arbitrariamente grande:
    \[
    [a], \quad [b], \quad [ab], \quad [ba], \quad [aba], \quad [abab], \quad \dots
    \]
    \item Il gruppo libero su $n$ generatori $F_n$ è isomorfo al prodotto libero di $n$ copie di $\mathbb{Z}$:
    \[
    F_n \cong \underbrace{\mathbb{Z} * \dots * \mathbb{Z}}_{n \text{ volte}}
    \]
\end{enumerate}
\end{prop}

\vspace{1.5em}


\section{Illustrazione Intuitiva del Gruppo Fondamentale}

\begin{ejemplo}[Gruppo fondamentale di $S^1$]
Consideriamo la circonferenza $S^1$, un punto base $x_0 \in S^1$ e sia $\sigma$ il cappio antiorario che compie un unico giro completo su $S^1$ partendo e ritornando in $x_0$.

\begin{figure}[H]
    \centering
    \begin{tikzpicture}[scale=1.2]
        \draw[thick] (0,0) circle (1.2cm);
        \draw[thick, ->, >=stealth] (0.1, 1.2) -- (-0.1, 1.2);
        \node at (1.4, 0.8) {$S^1$};
        \filldraw[blue] (-1.2,0) circle (1.5pt) node[left] {$x_0$};
    \end{tikzpicture}
    \caption{Lazo generatore $\sigma$ su $S^1$ con orientamento antiorario.}
\end{figure}

Le classi di omotopia distinte in $\Pi_1(S^1, x_0)$ corrispondono al numero di giri compiuti dal cappio:
\[
[c_{x_0}], \quad [\sigma], \quad [\sigma^{-1}], \quad [\sigma * \sigma] = [\sigma]^2, \quad \dots
\]
(Questa \`e solo l'intuizione; la dimostrazione rigorosa segue dal teorema di Seifert--Van Kampen, pi\`u avanti.) Pertanto, il gruppo fondamentale di $S^1$ è generato dall'unica classe $[\sigma]$ ed è isomorfo al gruppo libero $F_1$:
\[
\Pi_1(S^1, x_0) = \left\{ [\sigma]^k \;\middle|\; k \in \mathbb{Z} \right\} \cong F_1 \cong \mathbb{Z}
\]
\end{ejemplo}

\begin{ejemplo}[Gruppo fondamentale dell'unione di due circonferenze tangenti $S^1 \vee S^1$]
Consideriamo lo spazio $X = S^1 \vee S^1$ formato da due circonferenze tangenti esternamente in un punto base $x_0$.
Siano $\sigma$ e $\tau$ i due cappi orientati in senso antiorario su ciascuna delle due circonferenze.

\begin{figure}[H]
    \centering
    \begin{tikzpicture}[scale=1.0]
        \draw[thick] (-0.9,0) circle (0.9cm);
        \draw[thick] (0.9,0) circle (0.9cm);
        \filldraw[blue] (0,0) circle (2pt) node[below] {$x_0$};
        
        \draw[thick, ->, >=stealth] (-0.9, 0.9) -- (-1.0, 0.9);
        \node at (-0.9, 1.2) {$\sigma$};
        
        \draw[thick, ->, >=stealth] (0.9, 0.9) -- (0.8, 0.9);
        \node at (0.9, 1.2) {$\tau$};
    \end{tikzpicture}
    \caption{Rosa a due foglie $X = S^1 \vee S^1$ con i cappi generatori $\sigma$ e $\tau$.}
\end{figure}

In $\Pi_1(X, x_0)$, le concatenazioni di $\sigma$ e $\tau$ generano classi distinte a due a due:
\[
[c_{x_0}], \quad [\sigma], \quad [\tau], \quad [\sigma]^k, \quad [\tau]^h, \quad [\sigma * \tau], \quad [\tau * \sigma], \quad [\sigma * \tau^{-1}], \quad \dots
\]
Le classi di equivalenza sono le parole ridotte del tipo $[\sigma]^{k_1} [\tau]^{h_1} \cdots [\sigma]^{k_r} [\tau]^{h_r}$ con $k_i, h_i \in \mathbb{Z}$.
Di conseguenza, il gruppo fondamentale non è abeliano ed è isomorfo al gruppo libero su due generatori:
\[
\Pi_1(S^1 \vee S^1, x_0) \cong F_2
\]

Generalizzando a una rosa di $m$ circonferenze $X_m = \bigvee_{i=1}^m \mathscr{C}_i$, si ottiene:
\[
\Pi_1(X_m) \cong F_m
\]
\end{ejemplo}

\vspace{1.5em}


\section{Teoremi di Seifert-Van Kampen}

\begin{teo}[\textbf{I Teorema di Seifert-Van Kampen}]
Siano $X$ uno spazio topologico connesso per archi, $U, V \subseteq X$ aperti connessi per archi tali che $X = U \cup V$, e sia $W = U \cap V$ connesso per archi.

Consideriamo le inclusioni $i_1: W \hookrightarrow U$ e $i_2: W \hookrightarrow V$, e i rispettivi omomorfismi indotti $i_{1*}: \Pi_1(W) \to \Pi_1(U)$ e $i_{2*}: \Pi_1(W) \to \Pi_1(V)$.

Sia $N$ il più piccolo sottogruppo normale di $\Pi_1(U) * \Pi_1(V)$ contenente l'insieme di elementi del tipo:
\[
\left\{ (i_{1*}([\sigma]))^{-1} \cdot i_{2*}([\sigma]) \;\middle|\; [\sigma] \in \Pi_1(W) \right\}
\]
Allora il gruppo fondamentale di $X$ è dato dal quoziente:
\[
\Pi_1(X) \cong \frac{\Pi_1(U) * \Pi_1(V)}{N}
\]
\end{teo}

\begin{coro}
Nelle ipotesi del I Teorema di Seifert-Van Kampen, se in più l'intersezione $W = U \cap V$ è \textbf{semplicemente connessa} ($\Pi_1(W) = 0$), allora:
\[
\Pi_1(X) \cong \Pi_1(U) * \Pi_1(V)
\]
\end{coro}

\begin{coro}
Per ogni $n \geqslant 2$, la sfera $n$-dimensionale $S^n$ è semplicemente connessa:
\[
\Pi_1(S^n) = 0
\]
\end{coro}

\begin{proof}
Siano $p, q \in S^n$ con $p \neq q$. Definiamo gli aperti $U = S^n \setminus \{p\}$ e $V = S^n \setminus \{q\}$.
Allora $X = U \cup V$ e $W = U \cap V = S^n \setminus \{p,q\}$.

Gli aperti $U$ e $V$ sono omeomorfi a $\mathbb{R}^n$, quindi sono contrattili e semplicemente connessi ($\Pi_1(U) = \Pi_1(V) = 0$). Inoltre, per $n \geqslant 2$, $W \cong \mathbb{R}^n \setminus \{\bar{0}\} \simeq S^{n-1}$ è connesso per archi.

Applicando il I Teorema di Seifert-Van Kampen:
\[
\Pi_1(S^n) \cong \frac{\Pi_1(U) * \Pi_1(V)}{N} \cong \frac{0 * 0}{N} = 0
\]
\end{proof}

\begin{notacion}[\textbf{AVVISO PER L'ESAME}]
\textit{Se alla prova d'esame si richiede la dimostrazione che $S^n$ è semplicemente connessa per $n \geqslant 2$, occorre riportare la dimostrazione diretta dettagliata con il Lemma del numero di Lebesgue (del 06/10/2026)}.
\end{notacion}

\vspace{1.5em}

\begin{teo}[\textbf{II Teorema di Seifert-Van Kampen}]
Siano $X$ uno spazio topologico connesso per archi, $U, V \subseteq X$ aperti connessi per archi tali che $X = U \cup V$, e supponiamo che l'intersezione $W = U \cap V$ sia l'unione disgiunta di due aperti semplicemente connessi.

Se $V$ è semplicemente connesso, allora:
\[
\Pi_1(X) \cong \Pi_1(U) * \mathbb{Z}
\]
\end{teo}

\begin{coro}
$\Pi_1(S^1) \cong \mathbb{Z}$.
\end{coro}

\begin{proof}
Siano $p, q \in S^1$ due punti distinti ($p \neq q$). Definiamo $U = S^1 \setminus \{p\}$ e $V = S^1 \setminus \{q\}$.
Allora $S^1 = U \cup V$ e l'intersezione $W = U \cap V = S^1 \setminus \{p, q\}$ è l'unione disgiunta di due intervalli aperti (ciascuno semplice connesso e connesso per archi).

\begin{figure}[H]
    \centering
    \begin{tikzpicture}[scale=1.2]
        \draw[thick] (0,0) circle (1.2cm);
        \node[above right] at (0.85, 0.85) {$S^1$};
        \filldraw[blue] (-1.2, 0) circle (1.5pt) node[left] {$p$};
        \filldraw[red] (1.2, 0) circle (1.5pt) node[right] {$q$};
    \end{tikzpicture}
    \caption{Decomposizione di $S^1$ in due aperti $U, V$ con intersezione $W$ disgiunta.}
\end{figure}

Inoltre $U \cong V \cong \mathbb{R}$ sono connessi per archi e $V$ \`e semplicemente connesso ($\Pi_1(V) = 0$).
Dal II Teorema di Seifert-Van Kampen segue direttamente:
\[
\Pi_1(S^1) = \Pi_1(U) * \mathbb{Z} = 0 * \mathbb{Z} = \mathbb{Z}
\]
\end{proof}

\vspace{1.5em}

\begin{lema}
Siano $X$ uno spazio topologico, $x_0 \in X$ e $J$ un arco aperto, cio\`e un sottospazio $J \cong\ ]-1,1[$ tramite un omeomorfismo $\psi: J \to\ ]-1,1[$ con $\psi(x_0) = 0$ (quindi $x_0 \in J$). Posto $Y = X \cup J$, supponiamo $X \cap J = \{x_0\}$ e che $X$ e $J$ siano chiusi in $Y$.
Allora $X$ \`e un \textbf{retratto di deformazione} di $Y$ e, di conseguenza, l'inclusione induce un isomorfismo
\[
\Pi_1(X, x_0) \cong \Pi_1(X \cup J, x_0).
\]
\end{lema}

\begin{proof}
Sia $r: Y \to X$ definita da $r(y) = y$ su $X$ e $r(y) = x_0$ su $J$. Sono compatibili su $X \cap J = \{x_0\}$, e $X, J$ sono chiusi: per il Lemma di Incollamento $r$ \`e continua. Chiaramente $r \circ i = \mathrm{Id}_X$.

Definiamo $H: Y \times I \to Y$ ponendo $H(y,s) = y$ se $y \in X$, e $H(y,s) = \psi^{-1}\big((1-s)\,\psi(y)\big)$ se $y \in J$. Sulla intersezione $X \cap J$ si ha $H(x_0, s) = \psi^{-1}(0) = x_0$ in entrambe le definizioni, quindi $H$ \`e continua (Lemma di Incollamento su $X \times I$ e $J \times I$, chiusi). Inoltre $H(y,0) = y = \mathrm{Id}_Y(y)$ e $H(y,1) = r(y) = (i\circ r)(y)$. Dunque $\mathrm{Id}_Y \simeq i \circ r$, cio\`e $X$ \`e un retratto di deformazione di $Y$ (e $H$ fissa $X$ puntualmente). Per l'invarianza omotopica del gruppo fondamentale, $i_*: \Pi_1(X,x_0) \to \Pi_1(Y,x_0)$ \`e un isomorfismo\hfill$\square$
\end{proof}

\begin{observacion}[Nota di correzione]
Negli appunti originali il lemma era enunciato con ``$I'$ intervallo aperto di $\mathbb{R}^n$'' e la omotopia scritta come $(1-s)y + s\,r(y)$; ma nell'applicazione alla rosa a foglie $I'$ \`e un \emph{arco di circonferenza}, per il quale il segmento retto non resta in $Y$. La versione qui sopra (con l'omeomorfismo $\psi$) \`e quella che si applica correttamente.
\end{observacion}

\vspace{1.5em}
\begin{coro}
Sia $X_n = \bigcup_{i=1}^n \mathscr{C}_i$ la rosa a $n$ foglie (unione puntata di $n$ circonferenze). Allora:
\[
\Pi_1(X_n) \cong F_n \quad (\text{gruppo libero su } n \text{ generatori})
\]
\end{coro}

\begin{proof}
Procediamo per induzione su $n \geqslant 1$:
\begin{figure}[H]
    \centering
    \begin{tikzpicture}[scale=1.0]
        % Rosa a 3 foglie
        \draw[thick] (-0.8,0.8) circle (0.6cm);
        \draw[thick] (0.8,0.8) circle (0.6cm);
        \draw[thick] (-0.8,-0.8) circle (0.6cm);
        \filldraw[blue] (0,0) circle (2pt) node[below] {$x_0$};
        \node at (-1.4,1.2) {$\mathscr{C}_1$};
        \node at (1.4,1.2) {$\mathscr{C}_2$};
        \node at (-1.4,-1.2) {$\mathscr{C}_n$};
        \draw[dashed] (0,-0.4) -- (0,-1.2);
    \end{tikzpicture}
    \caption{Rosa a $n$ foglie $X_n$.}
\end{figure}

Per $n=1$, $X_1 = S^1 \implies \Pi_1(X_1) \cong \mathbb{Z} = F_1$.
Assumiamo il risultato vero per $n-1$. Fissiamo $q \in \mathscr{C}_n \setminus \{x_0\}$ e poniamo $U = X_n \setminus \{q\}$, $V = \mathscr{C}_n \setminus \{x_0\}$.
Allora $U$ e $V$ sono aperti in $X_n$, $X_n = U \cup V$ e $W = U \cap V = \mathscr{C}_n \setminus \{q, x_0\}$ è un'unione disgiunta di due aperti semplicemente connessi.

$V$ è un intervallo aperto, quindi semplicemente connesso.
Inoltre $U = X_{n-1} \cup J$, dove $J = \mathscr{C}_n \setminus \{q\}$ \`e un arco aperto che interseca $X_{n-1}$ solo in $x_0$ (entrambi sono chiusi in $U$: $X_{n-1}$ \`e compatto, e $J$ \`e chiuso in $U$ perch\'e $\overline{J} = \mathscr{C}_n$ incontra $X_{n-1}$ solo in $x_0 \in J$). Quindi $X_{n-1}$ \`e un retratto di deformazione di $U$.
Per il Lemma precedente, $\Pi_1(U) \cong \Pi_1(X_{n-1})$.

Sfruttando il II Teorema di Seifert-Van Kampen:
\[
\Pi_1(X_n) = \Pi_1(U) * \mathbb{Z} = \Pi_1(X_{n-1}) * \mathbb{Z} \cong F_{n-1} * \mathbb{Z} = F_n
\]
\end{proof}
\vspace{1.5em}

\begin{ejercicio*}
Calcolare il gruppo fondamentale dello spazio $X = \mathbb{R}^3 \setminus S^1$.

\begin{proof}[Svolgimento]
Rappresentiamo $S^1$ nel piano $z=0$ come $S^1 = \{(x,y,z) \in \mathbb{R}^3 \mid z=0, \, x^2+y^2=1\}$.
Consideriamo il disco chiuso $\overline{B^2} = \{(x,y,z) \in \mathbb{R}^3 \mid z=0, \, x^2+y^2 \leqslant 1\}$.

Decomponiamo $X = \mathbb{R}^3 \setminus S^1$ nell'unione di due aperti $U$ e $V$:
\begin{itemize}
    \item $U = \mathbb{R}^3 \setminus \overline{B^2}$ è un aperto semplicemente connesso ($\Pi_1(U) = 0$): intuitivamente, il disco chiuso si pu\`o ``schiacciare'' in un punto, per cui $U$ ha lo stesso tipo di omotopia di $\mathbb{R}^3 \setminus \{\text{pt}\} \simeq S^2$, che \`e semplicemente connessa.
    \item $V = \{(x,y,z) \in \mathbb{R}^3 \mid x^2+y^2 < 1\}$ (la parte interna di un cilindro infinito attorno all'asse $z$ con direttrice $S^1$). $V$ è contrattile, quindi semplicemente connesso ($\Pi_1(V) = 0$).
\end{itemize}

L'intersezione $W = U \cap V$ si scompone nell'unione disgiunta di due aperti semplicemente connessi:
\[
A_1 = \left\{(x,y,z) \in \mathbb{R}^3 \;\middle|\; z > 0, \, x^2+y^2 < 1\right\}
\]
\[
A_2 = \left\{(x,y,z) \in \mathbb{R}^3 \;\middle|\; z < 0, \, x^2+y^2 < 1\right\}
\]

\begin{figure}[H]
    \centering
    \begin{tikzpicture}[scale=1.0]
        % Cilindro V
        \draw[thick, gray!60] (-1.2,-2) -- (-1.2,2);
        \draw[thick, gray!60] (1.2,-2) -- (1.2,2);
        
        % Disco B^2
        \fill[red!30, opacity=0.6] (0,0) ellipse (1.2cm and 0.4cm);
        \draw[thick, red] (0,0) ellipse (1.2cm and 0.4cm);
        \node[red, left] at (-1.2,0) {$S^1$};

        % Regioni A_1 y A_2
        \node[ForestGreen] at (0,1.2) {$A_1 \subset W$};
        \node[ForestGreen] at (0,-1.2) {$A_2 \subset W$};

        % Ejes
        \draw[->, >=stealth] (-2.2,-1) -- (2.2,-1) node[right] {$x$};
        \draw[->, >=stealth] (0,-2.2) -- (0,2.2) node[above] {$z$};
    \end{tikzpicture}
    \caption{Decomposizione dello spazio $X = \mathbb{R}^3 \setminus S^1$ per l'applicazione di Seifert-Van Kampen.}
\end{figure}

Siamo nelle ipotesi del II Teorema di Seifert-Van Kampen:
\[
\Pi_1(\mathbb{R}^3 \setminus S^1) \cong \Pi_1(U) * \mathbb{Z} = 0 * \mathbb{Z} = \mathbb{Z}
\]
\end{proof}
\end{ejercicio*}

\vspace{2em}



\section{Applicazioni ed Esempi del Gruppo Fondamentale}

\begin{lema}
La circonferenza $S^1 = \partial B^2$ \textbf{non è un retratto} del disco chiuso $B^2$.
\end{lema}

\begin{proof}
Supponiamo per assurdo che esista una retrazione $r: B^2 \to S^1$, ossia un'applicazione continua tale che $r \circ i = \mathrm{Id}_{S^1}$, dove $i: S^1 \hookrightarrow B^2$ è l'inclusione canonica.

Passando agli omomorfismi indotti sul gruppo fondamentale:
\begin{tikzcd}[scale=1.2]
    \Pi_1(S^1) \arrow[r, "i_*"] \arrow[rr, bend right=30, "(\mathrm{Id}_{S^1})_*"'] & \Pi_1(B^2) \arrow[r, "r_*"] & \Pi_1(S^1)
\end{tikzcd}

Si ha che $(r \circ i)_* = r_* \circ i_* = (\mathrm{Id}_{S^1})_* = \mathrm{Id}_{\Pi_1(S^1)}$.
Sostituendo i gruppi fondamentali noti ($\Pi_1(S^1) \cong \mathbb{Z}$ e $\Pi_1(B^2) = 0$ in quanto $B^2$ è convesso):
\[
\Pi_1(S^1) = \mathrm{Id}_*(\Pi_1(S^1)) = r_*(i_*(\Pi_1(S^1))) \subseteq r_*(\Pi_1(B^2)) = r_*(0) = 0
\]
Otteniamo quindi che $\mathbb{Z} = 0$, il che è una palese contraddizione\hfill$\square$
\end{proof}

\vspace{1.5em}

\begin{teo}[Teorema del Punto Fisso di Brouwer in dimensione 2]
Ogni applicazione continua $f: B^2 \to B^2$ ammette almeno un punto fisso, ossia esiste $\bar{x} \in B^2$ tale che $f(\bar{x}) = \bar{x}$.
\end{teo}

\begin{proof}
Supponiamo per assurdo che $f(x) \neq x$ per ogni $x \in B^2$.
Possiamo allora definire un'applicazione $r: B^2 \to \partial B^2 = S^1$ costruita geometricamente come segue: per ogni $x \in B^2$, si considera la semiretta uscente da $f(x)$ e passante per $x$, e si definisce $r(x)$ come il punto di intersezione di tale semiretta con la circonferenza di bordo $S^1$.

\begin{figure}[H]
    \centering
    \begin{tikzpicture}[scale=1.3]
        % Disco B^2
        \draw[thick] (0,0) circle (1.5cm);
        \node[above right] at (1.1,1.1) {$B^2$};
        
        % Punti f(x) e x
        \filldraw[blue] (-0.6, -0.3) circle (1.5pt) node[below] {$f(x)$};
        \filldraw[red] (0.2, 0.1) circle (1.5pt) node[above] {$x$};
        
        % Semiretta e intersezione r(x)
        \draw[thick, ForestGreen, ->, >=stealth] (-0.6,-0.3) -- (1.3416, 0.6708);
        \filldraw[ForestGreen] (1.3416, 0.6708) circle (1.5pt) node[right] {$r(x)$};
    \end{tikzpicture}
    \caption{Costruzione della retrazione $r(x)$ a partire da una mappa priva di punti fissi.}
\end{figure}

L'applicazione $r$ così definita è continua. Per verificarlo esplicitamente, posto $u = x - f(x) \neq \bar{0}$, il punto $r(x) = x + \lambda(x)\,u$ \`e quello con $\lambda \geqslant 0$ tale che $\|x + \lambda u\|^2 = 1$, cio\`e $\|u\|^2\lambda^2 + 2\langle x,u\rangle \lambda + \|x\|^2 - 1 = 0$, e quindi
\[
\lambda(x) = \frac{-\langle x,u\rangle + \sqrt{\langle x,u\rangle^2 + \|u\|^2\,(1-\|x\|^2)}}{\|u\|^2},
\]
che \`e continua in $x$ (il discriminante \`e $\geqslant 0$ perch\'e $\|x\| \leqslant 1$, e $u \neq \bar 0$). Se $\|x\| = 1$ si ha $\langle x, u\rangle = 1 - \langle x, f(x)\rangle \geqslant 0$ (perch\'e $\|f(x)\| \leqslant 1$), quindi $\lambda(x) = 0$ e $r(x) = x$. Come visto sopra, $r(x) = x$ per ogni $x \in S^1 = \partial B^2$ (il punto $x$ giace gi\`a sul bordo).

Pertanto $r: B^2 \to S^1$ sarebbe una retrazione continua dal disco chiuso alla sua circonferenza di bordo. Ma questo contraddice il Lemma precedente, che stabilisce che $S^1$ non è un retratto di $B^2$.
Di conseguenza, deve esistere almeno un punto fisso $\bar{x} \in B^2$ tale che $f(\bar{x}) = \bar{x}$\hfill$\square$
\end{proof}

\vspace{2em}


\section{Gruppo Fondamentale del Toro Privato di un Punto}

\begin{prop}
Sia $T^2 = S^1 \times S^1$ il toro bidimensionale e $p \in T^2$.
Il gruppo fondamentale del toro privato di un punto è il gruppo libero su due generatori:
\[
\Pi_1(T^2 \setminus \{p\}) \cong F_2
\]
\end{prop}

\begin{proof}
Rappresentiamo il toro come il quoziente del quadrato unitario $T^2 = Q / \sim$, dove $Q = [0,1] \times [0,1]$ e la relazione di equivalenza identifica i lati opposti.
Sia $p \in \mathring{Q}$ un punto interno del quadrato (la scelta di $p$ \`e irrilevante: le traslazioni di $T^2$ sono omeomorfismi e agiscono transitivamente sui punti).

Consideriamo la proiezione radiale dal punto $p$ verso il bordo $\partial Q$:
\Func{r}{Q \setminus \{p\}}{\partial Q}{x}{\text{punto di intersezione della semiretta uscente da } p \text{ passante per } x \text{ con } \partial Q}

\begin{figure}[H]
    \centering
    \begin{tikzpicture}[scale=1.2]
        % Cuadrado Q
        \draw[thick] (0,0) -- (2.5,0) node[midway, below] {$a$} -- (2.5,2.5) node[midway, right] {$b$} -- (0,2.5) node[midway, above] {$a$} -- (0,0) node[midway, left] {$b$};
        
        % Flechas de identificación de lados
        \draw[thick, ->, >=stealth] (1.1,0) -- (1.3,0);
        \draw[thick, ->, >=stealth] (1.1,2.5) -- (1.3,2.5);
        \draw[thick, ->>, >=stealth] (0,1.1) -- (0,1.3);
        \draw[thick, ->>, >=stealth] (2.5,1.1) -- (2.5,1.3);

        % Punto p interno y proyectado
        \filldraw[red] (1.2, 1.2) circle (1.5pt) node[below left] {$p$};
        \filldraw[blue] (1.85, 1.65) circle (1.5pt) node[above left] {$x$};
        \draw[dashed, ForestGreen, ->, >=stealth] (1.2,1.2) -- (2.5, 2.1);
        \filldraw[ForestGreen] (2.5,2.1) circle (1.5pt) node[right] {$r(x)$};
    \end{tikzpicture}
    \caption{Proiezione radiale $r: Q \setminus \{p\} \to \partial Q$.}
\end{figure}

L'applicazione $r$ è ben definita e continua.
Inoltre, per ogni $x \in \partial Q$ si ha $r(x) = x$, quindi $r$ è una retrazione di $Q \setminus \{p\}$ su $\partial Q$.

Costruiamo la seguente omotopia $H: (Q \setminus \{p\}) \times I \to Q \setminus \{p\}$ data da:
\[
H(x, s) = (1-s)x + s\, r(x)
\]
$H$ \`e continua (e ha valori in $Q \setminus \{p\}$: il segmento da $x$ a $r(x)$ giace sulla semiretta uscente da $p$ che passa per $x$, quindi non contiene $p$; inoltre $Q$ \`e convesso) e soddisfa $H(x,0) = x = \mathrm{Id}_{Q \setminus \{p\}}(x)$ e $H(x,1) = r(x) = (i \circ r)(x)$.
Ciò prova che $\partial Q$ è un \textbf{retratto di deformazione} di $Q \setminus \{p\}$.

Poich\'e $H$ lascia fisso ogni punto di $\partial Q$, \`e compatibile con le identificazioni dei lati opposti, quindi passa al quoziente (si usa che $\pi \times \mathrm{Id}_I$ \`e ancora una applicazione quoziente, perch\'e $I$ \`e localmente compatto). Si ottiene una omotopia $\widetilde{H}$ su $T^2 \setminus \{p\} = (Q \setminus \{p\})/\!\sim$ e una retrazione
\[
\widetilde{r}: T^2 \setminus \{p\} \longrightarrow \partial Q / \sim ,
\]
cio\`e $\partial Q/\!\sim$ \`e un retratto di deformazione del toro bucato.

Nel quoziente i quattro lati si identificano a due a due (lato $a$ con lato $a$, lato $b$ con lato $b$) e i quattro vertici diventano un unico punto: $\partial Q / \sim$ \`e quindi formato da due circonferenze ($a$ e $b$) incollate in un punto, ossia la rosa a due foglie $S^1 \vee S^1$.

Sfruttando l'invarianza omotopica del gruppo fondamentale:
\[
\Pi_1(T^2 \setminus \{p\}) \cong \Pi_1(\partial Q / \sim) \cong \Pi_1(S^1 \vee S^1) \cong F_2
\]
\end{proof}

\vspace{2em}


    \chapter{Rivestimenti e spazi proiettivi}

\section{Spazi di Rivestimento (Rivestimenti)}

\begin{definicion}
Siano $\widetilde{X}, X$ spazi topologici, con $\widetilde{X}$ connesso per archi, e $p: \widetilde{X} \to X$ un'applicazione continua.

La coppia $(\widetilde{X}, p)$ si chiama \textbf{spazio di rivestimento} di $X$ se per ogni $x \in X$ esiste un intorno aperto $U$ di $x$ in $X$ tale che la sua controimmagine $p^{-1}(U)$ sia l'unione disgiunta di aperti in $\widetilde{X}$:
\[
p^{-1}(U) = \bigsqcup_{\alpha} \widetilde{U}_\alpha
\]
dove per ogni $\alpha$, la restrizione $p|_{\widetilde{U}_\alpha}: \widetilde{U}_\alpha \to U$ è un \textbf{omeomorfismo}.
\end{definicion}

\begin{notacion}
\begin{itemize}
    \item $p$ si dice \textbf{applicazione di rivestimento}.
    \item $U$ si dice \textbf{intorno ben coperto} di $x$.
    \item Ciascun aperto $\widetilde{U}_\alpha$ si dice \textbf{foglio} su $U$.
    \item La controimmagine di un punto $p^{-1}(x)$ si chiama \textbf{fibra} sopra $x$.
\end{itemize}
\end{notacion}

\begin{figure}[H]
    \centering
    \begin{tikzpicture}[scale=1.2]
        % Espacio revestimento \widetilde{X}
        \draw[gray!40, fill=gray!5] (-2,1.8) -- (2,1.8) -- (3,3) -- (-1,3) -- cycle;
        \node[above left] at (3,3) {$\widetilde{X}$};

        % Hojas \widetilde{U}_\alpha
        \draw[thick, blue] (-0.8, 2.6) to[out=20, in=160] (0.8, 2.6);
        \node[blue, right] at (0.8, 2.6) {$\widetilde{U}_1$};
        \draw[thick, blue] (-0.8, 2.2) to[out=20, in=160] (0.8, 2.2);
        \node[blue, right] at (0.8, 2.2) {$\widetilde{U}_2$};

        % Proyección p
        \draw[->, >=stealth, thick] (0, 1.6) -- (0, 0.6) node[midway, right] {$p$};

        % Espacio base X
        \draw[gray!40, fill=gray!5] (-2,-0.8) -- (2,-0.8) -- (3,0.4) -- (-1,0.4) -- cycle;
        \node[above left] at (3,0.4) {$X$};

        % Intorno U y punto x
        \draw[thick, ForestGreen] (-0.8, -0.2) to[out=20, in=160] (0.8, -0.2);
        \node[ForestGreen, right] at (0.8, -0.2) {$U$};
        \filldraw[black] (0, -0.2) circle (1.5pt) node[below] {$x$};
    \end{tikzpicture}
    \caption{Struttura locale di uno spazio di rivestimento $p: \widetilde{X} \to X$.}
\end{figure}

\begin{prop}[Proprietà di base dei rivestimenti]
Sia $(\widetilde{X}, p)$ uno spazio di rivestimento di $X$.
\begin{enumerate}
    \item \textbf{$p$ è suriettiva:}
    Per ogni $x \in X$, sia $U$ un intorno ben coperto di $x$. Fissato un foglio $\widetilde{U}_\alpha$, poichè $p|_{\widetilde{U}_\alpha}: \widetilde{U}_\alpha \to U$ è un omeomorfismo, esiste $\widetilde{x} \in \widetilde{U}_\alpha$ tale che $p(\widetilde{x}) = x$.

    \item \textbf{$p$ è un omeomorfismo locale:}
    Per ogni $\widetilde{x} \in \widetilde{X}$, posto $x = p(\widetilde{x}) \in X$, sia $U$ un intorno ben coperto di $x$.
    Siccome $\widetilde{x} \in p^{-1}(x) \subset p^{-1}(U) = \bigsqcup \widetilde{U}_\alpha$, esiste un unico foglio $\widetilde{U}_\alpha$ contenente $\widetilde{x}$.
    L'aperto $\widetilde{U}_\alpha$ è un intorno aperto di $\widetilde{x}$ in $\widetilde{X}$ e $p|_{\widetilde{U}_\alpha}: \widetilde{U}_\alpha \to U$ è un omeomorfismo su $U = p(\widetilde{U}_\alpha)$.

    \begin{observacion}
    Ogni omeomorfismo è un omeomorfismo locale, ma il viceversa non vale in generale.
    Controesempio: La mappa $p: \mathbb{R} \to S^1$ data da $t \mapsto e^{2\pi i t}$ è un omeomorfismo locale (ed è un rivestimento), ma non è un omeomorfismo globale poichè non è iniettiva.
    \end{observacion}

    \item \textbf{La fibra $p^{-1}(x)$ è uno spazio discreto (formato da punti isolati):}
    Sia $\widetilde{x} \in p^{-1}(x)$ e sia $U$ un intorno ben coperto di $x$.
    Allora $\widetilde{x} \in \widetilde{U}_\alpha$ per un unico foglio, ed essendo $\widetilde{U}_\alpha$ aperto in $\widetilde{X}$ e $p|_{\widetilde{U}_\alpha}$ iniettiva, si ha $p^{-1}(x) \cap \widetilde{U}_\alpha = \{\widetilde{x}\}$.
    Questo prova che ogni punto della fibra è isolato.

    \item \textbf{$p$ è un'applicazione aperta:}
    Sia $\widetilde{A}$ un aperto di $\widetilde{X}$. Mostriamo che $p(\widetilde{A})$ è aperto in $X$.
    Sia $x \in p(\widetilde{A})$ e sia $U$ un intorno ben coperto di $x$.
    Esiste $\widetilde{x} \in \widetilde{A}$ tale che $p(\widetilde{x}) = x$. Poiché $\widetilde{x} \in p^{-1}(U)$, esiste un foglio $\widetilde{U}_\alpha$ contenente $\widetilde{x}$.
    
    L'intersezione $\widetilde{A} \cap \widetilde{U}_\alpha$ è un aperto di $\widetilde{U}_\alpha$.
    Essendo $p|_{\widetilde{U}_\alpha}: \widetilde{U}_\alpha \to U$ un omeomorfismo (e quindi una mappa aperta), l'immagine $p(\widetilde{A} \cap \widetilde{U}_\alpha)$ è un aperto di $U$, e quindi aperto anche in $X$.
    Inoltre $x \in p(\widetilde{A} \cap \widetilde{U}_\alpha) \subseteq p(\widetilde{A})$, il che dimostra che $p(\widetilde{A})$ è un intorno aperto di ogni suo punto, ossia $p(\widetilde{A})$ è aperto in $X$.
\end{enumerate}
\end{prop}

\begin{prop}[Aperti nel rivestimento e nella base]
Sia $p: \widetilde{X} \to X$ una applicazione continua, \textbf{aperta} e suriettiva (ad esempio una applicazione di rivestimento, per le propriet\`a 1 e 4 sopra). Per ogni $A \subseteq X$:
\[
A \text{ è aperto in } X \iff p^{-1}(A) \text{ è aperto in } \widetilde{X}.
\]
\end{prop}

\begin{proof}
($\Rightarrow$) \`E la continuit\`a di $p$. ($\Leftarrow$) Se $p^{-1}(A)$ \`e aperto, poich\'e $p$ \`e aperta e suriettiva, $A = p(p^{-1}(A))$ \`e aperto. \qedhere
\end{proof}

\begin{observacion}
L'ipotesi ``$p$ aperta'' \`e essenziale: la sola continuit\`a e suriettivit\`a non bastano (ad esempio l'identit\`a $\mathrm{id}: (\mathbb{R}, \text{discreta}) \to \mathbb{R}$ \`e continua e suriettiva, ma $p^{-1}(A)$ \`e sempre aperto, anche se $A$ non lo \`e). Negli appunti originali la proposizione era enunciata senza questa ipotesi, che per\`o veniva usata nella dimostrazione.
\end{observacion}

\begin{coro}[Teorema dell'Omeomorfismo Quoziente per i Rivestimenti]
Se $(\widetilde{X}, p)$ è uno spazio di rivestimento di $X$, allora la mappa indotta sul quoziente
\[
\widetilde{p}: \widetilde{X} / \sim_p \longrightarrow X
\]
è un \textbf{omeomorfismo}, dove $\widetilde{x} \sim_p \widetilde{y} \iff p(\widetilde{x}) = p(\widetilde{y})$. In altri termini $X$ \`e il quoziente di $\widetilde{X}$ rispetto alla relazione indotta dalle fibre di $p$.
\end{coro}
\begin{proof}
Per la proposizione, $p$ \`e una applicazione quoziente (continua, suriettiva e con la topologia di $X$ determinata da $p^{-1}$). Un'applicazione quoziente induce un omeomorfismo dal quoziente per la relazione delle sue fibre.
\end{proof}

\begin{prop}[Corrispondenza tra punti della fibra e fogli]
Sia $p: \widetilde{X} \to X$ un rivestimento, $x \in X$ e $U$ un intorno ben coperto di $x$. Esiste una corrispondenza biunivoca tra i punti della fibra $p^{-1}(x)$ e i fogli $\widetilde{U}_\alpha$ sopra $U$.
\end{prop}

\begin{proof}
Sia $p^{-1}(U) = \bigsqcup_\alpha \widetilde{U}_\alpha$. Ad ogni $\widetilde{x} \in p^{-1}(x) \subset p^{-1}(U)$ associamo l'unico foglio che lo contiene (unico perch\'e i fogli sono disgiunti). Viceversa, dato un foglio $\widetilde{U}_\alpha$, la restrizione $p|_{\widetilde{U}_\alpha}: \widetilde{U}_\alpha \to U$ \`e un omeomorfismo, quindi esiste un unico punto $\widetilde{x} \in \widetilde{U}_\alpha$ con $p(\widetilde{x}) = x$. Le due costruzioni sono una l'inversa dell'altra\hfill$\square$
\end{proof}

\section{Esempi di Spazi di Rivestimento}

\begin{ejemplo}[Rivestimento della circonferenza $S^1$ tramite la retta reale $\mathbb{R}$]
Consideriamo l'applicazione mappa esponenziale:
\Func{p}{\mathbb{R}}{S^1}{t}{e^{2\pi i t} = (\cos(2\pi t), \sin(2\pi t))}

Allora la coppia $(\mathbb{R}, p)$ è uno \textbf{spazio di rivestimento} di $S^1$ ad \textbf{infinite foglie}.

\begin{proof}
Dimostriamo le proprietà richieste:

\begin{itemize}
    \item \textbf{$p$ è continua, suriettiva ed aperta:}
    La continuità segue dalla continuità delle componenti trigonometriche. La suriettivit\`a \`e nota; l'apertura segue dal fatto che l'immagine di un intervallo aperto di lunghezza $\leqslant 1$ \`e un arco aperto di $S^1$ e ogni aperto di $\mathbb{R}$ \`e unione di tali intervalli.

    \item \textbf{Costruzione degli intorni ben coperti:}
    Fissiamo un punto $x \in S^1$ e scegliamo un punto della fibra $\widetilde{x} \in p^{-1}(x)$, con $\widetilde{x} = t \in \mathbb{R}$ tale che $x = e^{2\pi i t}$.

    Scegliamo $a, b \in \mathbb{R}$ tali che $a < t < b$ e $b - a < 1$.
    L'intervallo aperto $]a, b[$ è un intorno aperto di $t$ in $\mathbb{R}$. Poiché $p$ è un'applicazione aperta, l'immagine:
    \[
    p(]a, b[) = \wideparen{p(a)p(b)}
    \]
    è un arco aperto di $S^1$ contenente $x = p(t)$.

    \item \textbf{Verifica dell'omeomorfismo locale:}
    La restrizione $p|_{]a,b[}: ]a,b[ \to \wideparen{p(a)p(b)}$ è una bigezione continua ed aperta, quindi è un \textbf{omeomorfismo}.

    Infatti, per la continuità di $p$, la restrizione $p|_{]a,b[}$ è continua. Essendo $]a,b[$ aperto in $\mathbb{R}$ e $p$ una mappa aperta, $p|_{]a,b[}$ è aperta.
    Per l'iniettività, siano $t, t' \in ]a,b[$ tali che $p(t) = p(t')$.
    Allora $e^{2\pi i t} = e^{2\pi i t'} \implies 2\pi(t - t') \in 2\pi\mathbb{Z} \implies t - t' \in \mathbb{Z}$.
    Poiché $|t - t'| < b - a < 1$, l'unica possibilità è $t - t' = 0 \implies t = t'$.

    \item \textbf{Scomposizione della controimmagine in fogli disgiunti:}
    Osserviamo che la controimmagine dell'arco aperto $U = \wideparen{p(a)p(b)}$ si scompone nell'unione disgiunta:
    \[
    p^{-1}(U) = p^{-1}(p(]a,b[)) = \coprod_{z \in \mathbb{Z}} ]a+z, b+z[
    \]
    Per ogni $z \in \mathbb{Z}$, la traslazione $t_z: ]a,b[ \to ]a+z, b+z[$ definita da $t_z(t) = t + z$ è un omeomorfismo.

    Pertanto, la restrizione di $p$ su ciascun foglio $]a+z, b+z[$ si esprime come la composizione dell'omeomorfismo di restrizione e della traslazione:
    \[
    p|_{]a+z, b+z[} = p|_{]a,b[} \circ t_{-z}
    \]
    Essendo composizione di omeomorfismi, $p|_{]a+z, b+z[}: ]a+z, b+z[ \to U$ è un omeomorfismo per ogni $z \in \mathbb{Z}$.
\end{itemize}
\end{proof}

\begin{observacion}
L'applicazione $p: \mathbb{R} \to S^1$ data da $t \mapsto e^{2\pi i t}$ è anche un \textbf{omomorfismo di gruppi} dal gruppo additivo $(\mathbb{R}, +)$ al gruppo moltiplicativo $(S^1, \cdot)$.
Infatti, per ogni $t, t' \in \mathbb{R}$:
\[
p(t + t') = e^{2\pi i (t + t')} = e^{2\pi i t} \cdot e^{2\pi i t'} = p(t) \cdot p(t')
\]
\end{observacion}
\end{ejemplo}

\vspace{1.5em}

\begin{ejemplo}[Interpretazione Geometrica: L'Elica Cilindrica]
Sia $\mathcal{E} \subset \mathbb{R}^3$ l'elica cilindrica di equazioni parametriche:
\[
\begin{cases}
    x = \cos(2\pi t) \\
    y = \sin(2\pi t) \\
    z = t
\end{cases}
\quad t \in \mathbb{R}
\]

\begin{figure}[H]
    \centering
    \begin{tikzpicture}[scale=1.2]
        % elica (proiezione obliqua): x = cos(2 pi t), y = sin(2 pi t), z = t
        \draw[domain=-1:2.2, samples=200, smooth, variable=\t, blue, thick]
            plot ({0.9*cos(360*\t)}, {0.35*sin(360*\t) + 1.1*\t});
        \node[blue, right] at (1.0,2.3) {$\mathcal{E}$};
        % circonferenza base
        \draw[thick, ForestGreen] (0,-1.9) ellipse (0.9cm and 0.35cm);
        \node[ForestGreen, right] at (0.9,-1.9) {$S^1$};
        \draw[dashed, gray] (0.9,-1.0) -- (0.9,-1.9);
        \draw[dashed, gray] (-0.9, 0.0) -- (-0.9,-1.9);
        \draw[->, >=stealth, thick, Purple] (-1.6, 0.8) -- (-1.6, -1.6) node[midway, left] {$p$};
    \end{tikzpicture}
    \caption{L'elica cilindrica $\mathcal{E}$ si proietta sulla circonferenza $S^1$ ``avvolgendosi'' infinite volte: ogni punto di $S^1$ ha infinite controimmagini (una per ogni giro).}
\end{figure}

Come insieme di punti, l'elica $\mathcal{E}$ è omeomorfa alla retta reale $\mathbb{R}$ tramite l'omeomorfismo $f: \mathbb{R} \to \mathcal{E}$ dato da $t \mapsto (\cos(2\pi t), \sin(2\pi t), t)$.

La proiezione verticale $p: \mathcal{E} \to S^1$ definita da $(\cos(2\pi t), \sin(2\pi t), t) \mapsto (\cos(2\pi t), \sin(2\pi t))$ corrisponde alla mappa esponenziale $t \mapsto e^{2\pi i t}$, confermando geometricamente che l'elica «srotola» la circonferenza in un rivestimento ad infinite foglie.
\end{ejemplo}

\vspace{2em}

\begin{ejemplo}[Rivestimento del Toro $n$-Dimensionale $T^n$]
Consideriamo l'applicazione proiettiva:
\Func{\overline{p}}{\mathbb{R}^n}{\underbrace{S^1 \times \dots \times S^1}_{n \text{ volte}} = T^n}{(t_1, \dots, t_n)}{(p(t_1), \dots, p(t_n))}
dove $p: \mathbb{R} \to S^1$ è la mappa esponenziale $p(t) = e^{2\pi i t}$.

Allora $(\mathbb{R}^n, \overline{p})$ è uno \textbf{spazio di rivestimento} del toro $n$-dimensionale $T^n$.

\begin{proof}
Sia $(z_1, \dots, z_n) \in S^1 \times \dots \times S^1 = T^n$.
Per ogni $r = 1, \dots, n$, scegliamo un punto $t_r \in p^{-1}(z_r)$ ed intervalli $]a_r, b_r[$ con $a_r < t_r < b_r$ e $b_r - a_r < 1$.

Consideriamo il prodotto cartesiano degli archi aperti su $T^n$:
\[
V = \overline{p}\left( ]a_1, b_1[ \times \dots \times ]a_n, b_n[ \right) = \widearc{p(a_1)p(b_1)} \times \dots \times \widearc{p(a_n)p(b_n)}
\]
$V$ è un intorno aperto ben coperto di $(z_1, \dots, z_n)$ in $T^n$.

La controimmagine $\overline{p}^{-1}(V)$ in $\mathbb{R}^n$ si decompone nell'unione disgiunta di cubi aperti traslati tramite il reticolo $\mathbb{Z}^n$:
\[
\overline{p}^{-1}(V) = \bigsqcup_{h = (h_1, \dots, h_n) \in \mathbb{Z}^n} \left( ]a_1 + h_1, b_1 + h_1[ \times \dots \times ]a_n + h_n, b_n + h_n[ \right)
\]

Per ogni $h \in \mathbb{Z}^n$, la traslazione $t_h: \prod_{r=1}^n ]a_r, b_r[ \to \prod_{r=1}^n ]a_r + h_r, b_r + h_r[$ data da $t_h(x) = x + h$ è un omeomorfismo.
Inoltre, la restrizione della proiezione sul singolo blocco traslato soddisfa:
\[
\overline{p}|_{\prod_{r=1}^n ]a_r + h_r, b_r + h_r[} = \overline{p}|_{\prod_{r=1}^n ]a_r, b_r[} \circ t_{-h}
\]
Trattandosi di composizione di omeomorfismi, ciascuna restrizione è un omeomorfismo sul prodotto aperto $V$.
Ciò prova che $(\mathbb{R}^n, \overline{p})$ è un rivestimento di $T^n$\hfill$\square$
\end{proof}
\end{ejemplo}

\vspace{2em}

\section{Lo Spazio Proiettivo Reale e il suo Rivestimento tramite $S^n$}

\begin{definicion}[Spazio Proiettivo Reale $\mathbb{R}P^n$]
Consideriamo lo spazio vettoriale $\mathbb{R}^{n+1} \setminus \{\bar{0}\}$.
Definiamo la seguente relazione di equivalenza $\sim$ su $\mathbb{R}^{n+1} \setminus \{\bar{0}\}$:
\[
x \sim y \iff \exists t \in \mathbb{R} \setminus \{0\} \text{ tale che } y = t x
\]
Lo \textbf{spazio proiettivo reale} di dimensione $n$ è lo spazio topologico quoziente:
\[
\mathbb{R}P^n = \left( \mathbb{R}^{n+1} \setminus \{\bar{0}\} \right) \big/ \sim
\]
ovvero l'insieme di tutte le rette passanti per l'origine in $\mathbb{R}^{n+1}$ dotate della topologia quoziente. Indichiamo con
\Func{\pi}{\mathbb{R}^{n+1} \setminus \{\bar{0}\}}{\mathbb{R}P^n}{x}{[x]_\sim = \{t x \mid t \neq 0\}}
la proiezione canonica (continua per definizione di topologia quoziente).
\end{definicion}

\section{Proprietà dello Spazio Proiettivo Reale $\mathbb{R}P^n$}

\begin{prop}
$\mathbb{R}P^n$ è uno spazio connesso (e connesso per archi).
\end{prop}

\begin{proof}
Essendo $\mathbb{R}P^n = \pi\left(\mathbb{R}^{n+1} \setminus \{\bar{0}\}\right)$ l'immagine continua dello spazio $\mathbb{R}^{n+1} \setminus \{\bar{0}\}$, ed essendo quest'ultimo connesso per archi per ogni $n \geqslant 1$ (e $\pi$ continua), ne consegue che $\mathbb{R}P^n$ è connesso per archi.
\end{proof}

\begin{prop}
$\mathbb{R}P^n$ è uno spazio di Hausdorff ($T_2$).
\end{prop}

\begin{proof}
Sia $\mathcal{C}$ un cono aperto di vertice $\bar{0}$ in $\mathbb{R}^{n+1} \setminus \{\bar{0}\}$.
Osserviamo che per ogni $x \in \mathcal{C}$, l'intera retta vettoriale privata dell'origine è contenuta in $\mathcal{C}$, ossia $\pi^{-1}(\pi(\mathcal{C})) = \mathcal{C}$.
Definiamo $\overline{\mathcal{C}} = \pi(\mathcal{C})$. Per la definizione di topologia quoziente, un insieme $A \subseteq \mathbb{R}P^n$ è aperto se e solo se $\pi^{-1}(A)$ è un aperto di $\mathbb{R}^{n+1} \setminus \{\bar{0}\}$.
Essendo $\pi^{-1}(\overline{\mathcal{C}}) = \mathcal{C}$ un cono aperto, $\overline{\mathcal{C}}$ è un aperto di $\mathbb{R}P^n$.

\begin{figure}[H]
    \centering
    \begin{tikzpicture}[scale=1.0]
        % Conos abiertos en R^{n+1}
        \draw[thick, ForestGreen, fill=ForestGreen!10] (0,0) -- (1,2) -- (-0.5,2) -- cycle;
        \draw[thick, ForestGreen, fill=ForestGreen!10] (0,0) -- (-1,-2) -- (0.5,-2) -- cycle;
        
        \draw[thick, blue, fill=blue!10] (0,0) -- (2,0.8) -- (2,-0.2) -- cycle;
        \draw[thick, blue, fill=blue!10] (0,0) -- (-2,-0.8) -- (-2,0.2) -- cycle;

        \filldraw[red] (0,0) circle (1.5pt) node[below left] {$\bar{0}$};
        
        % Ejes/Rettas
        \draw[thick] (-0.8,1.6) -- (0.8,-1.6) node[below right] {$[x]_\sim$};
        \draw[thick] (1.6,0.4) -- (-1.6,-0.4) node[below left] {$[y]_\sim$};
    \end{tikzpicture}
    \caption{Separazione di due rette distinte tramite coni aperti disgiunti in $\mathbb{R}^{n+1}$.}
\end{figure}

Siano ora $[x]_\sim, [y]_\sim \in \mathbb{R}P^n$ due punti distinti con $[x]_\sim \neq [y]_\sim$, il che equivale a dire che le rette vettoriali ordinate da $x$ e $y$ sono distinte.
Possiamo scegliere due coni aperti $\mathcal{C}_1$ e $\mathcal{C}_2$ di assi $[x]_\sim$ e $[y]_\sim$ rispettivamente, tali che $\mathcal{C}_1 \cap \mathcal{C}_2 = \varnothing$. Esplicitamente: siano $u = x/\|x\|$, $w = y/\|y\|$; poich\'e le rette sono distinte, $d = \min\{\|u - w\|, \|u + w\|\} > 0$. Posto $\delta = d/2$ e
\[
\mathcal{C}_1 = \Big\{ v \neq \bar 0 \;:\; \min\big\|\tfrac{v}{\|v\|} \mp u\big\| < \delta \Big\}, \qquad
\mathcal{C}_2 = \Big\{ v \neq \bar 0 \;:\; \min\big\|\tfrac{v}{\|v\|} \mp w\big\| < \delta \Big\},
\]
i due insiemi sono aperti, invarianti per moltiplicazione per scalari non nulli (coni), contengono rispettivamente le rette di $x$ e $y$, e sono disgiunti: se $v$ stesse in entrambi, per la disuguaglianza triangolare $\|\pm u \mp w\| < 2\delta = d$, assurdo.

Ma allora le loro immagini $\overline{\mathcal{C}}_1 = \pi(\mathcal{C}_1)$ e $\overline{\mathcal{C}}_2 = \pi(\mathcal{C}_2)$ sono intorni aperti disgiunti di $[x]_\sim$ e $[y]_\sim$ in $\mathbb{R}P^n$, ovvero $\overline{\mathcal{C}}_1 \cap \overline{\mathcal{C}}_2 = \varnothing$.
Questo dimostra che $\mathbb{R}P^n$ soddisfa l'assioma $T_2$\hfill$\square$
\end{proof}

\vspace{1.5em}

\begin{prop}
Consideriamo la sfera unitaria $S^n \subset \mathbb{R}^{n+1} \setminus \{\bar{0}\}$ e la restrizione della proiezione canonica:
\Func{p = \pi|_{S^n}}{S^n}{\mathbb{R}P^n}{x}{[x]_\sim}
Allora $(S^n, p)$ è uno \textbf{spazio di rivestimento a 2 foglie} di $\mathbb{R}P^n$.
\end{prop}

\begin{coro}
Lo spazio proiettivo reale è omeomorfo al quoziente della sfera rispetto all'azione dei punti antipodali:
\[
\mathbb{R}P^n \cong S^n \big/ \{ \pm \mathrm{Id}_{S^n} \}
\]
\end{coro}

\begin{proof}
Sia $p: S^n \to \mathbb{R}P^n$ la mappa di rivestimento a 2 foglie. Due punti $x, x' \in S^n$ soddisfano $p(x) = p(x')$ se e solo se $x' = \pm x$ (punti antipodali).
Dato che la relazione indotta dal rivestimento identifica ogni punto con il suo antipodo ($x \sim_p x' \iff x' = \pm x$), dal Teorema dell'Omeomorfismo Quoziente segue che:
\[
\mathbb{R}P^n = S^n \big/ \sim_p \;\;\cong\;\; S^n \big/ \{ \pm \mathrm{Id}_{S^n} \}
\]
\end{proof}

\begin{coro}
Lo spazio proiettivo reale $\mathbb{R}P^n$ è uno spazio topologico \textbf{compatto}.
\end{coro}

\begin{proof}
$\mathbb{R}P^n = p(S^n)$ è l'immagine continua della sfera unitaria $S^n$, la quale è un sottoinsieme chiuso e limitato di $\mathbb{R}^{n+1}$, e quindi compatta per il Teorema di Heine-Borel.
\end{proof}

\vspace{2em}

\section{Esempi Notevoli di Spazi Proiettivi Reali}

\begin{ejemplo}[Caso $n=1$: $\mathbb{R}P^1 \cong S^1$]
Dimostriamo che la retta proiettiva reale è omeomorfa alla circonferenza $S^1$.

\begin{proof}
Rappresentiamo la circonferenza come l'insieme dei numeri complessi di modulo unitario $S^1 = \{z \in \mathbb{C} \mid |z| = 1\}$. Ogni elemento $z \in S^1$ si scrive nella forma $z = e^{i\theta}$ con $\theta \in \mathbb{R}$.

Consideriamo l'applicazione di elevamento al quadrato $f: S^1 \to S^1$ definita da:
\[
f(z) = z^2 = e^{2i\theta}
\]
$f$ è ben posta, continua e suriettiva.

\begin{itemize}
    \item \textbf{$f$ è un'applicazione aperta:}
    Sia $\widearc{a,b}$ un arco aperto su $S^1$. L'immagine $f(\widearc{a,b})$ corrisponde all'arco duplicato $\widearc{a^2, b^2}$. Se la lunghezza dell'arco iniziale supera $\pi$, l'immagine ricopre l'intera circonferenza $S^1$ (che è aperta), altrimenti produce un arco aperto.

    \item \textbf{Relazione di equivalenza indotta:}
    Per ogni $z, z' \in S^1$, si ha:
    \[
    f(z) = f(z') \iff z^2 = (z')^2 \iff z' = \pm z
    \]
    Quindi la relazione indotta da $f$ identifica esattamente i punti antipodali di $S^1$.
\end{itemize}

Poich\'e $f$ \`e continua, aperta e suriettiva, essa \`e una applicazione quoziente, quindi induce un omeomorfismo
\[
\widetilde{f}: S^1 \big/ \sim_f \;=\; S^1 \big/ \{ \pm \mathrm{Id}_{S^1} \} \xrightarrow{\ \cong\ } S^1 .
\]
D'altra parte, per il corollario sul rivestimento $S^n \to \mathbb{R}P^n$ (con $n=1$), $S^1 / \{\pm \mathrm{Id}\} \cong \mathbb{R}P^1$. Componendo, $\mathbb{R}P^1 \cong S^1$\hfill$\square$
\end{proof}
\end{ejemplo}

\vspace{1.5em}

\begin{ejemplo}[Caso $n=2$: Il Piano Proiettivo Reale $\mathbb{R}P^2$]
Analizziamo la rappresentazione topologica del piano proiettivo reale $\mathbb{R}P^2 = S^2 / \{ \pm \mathrm{Id} \}$.

\begin{figure}[H]
    \centering
    \begin{tikzpicture}[scale=1.3]
        % emisfero superiore
        \draw[thick] (-1.5,0) arc (180:0:1.5);
        \draw[thick] (-1.5,0) arc (180:360:1.5 and 0.4);
        \draw[thick, dashed] (-1.5,0) arc (180:0:1.5 and 0.4);
        \node at (0,0.8) {$S^2_+$};
        \draw[->, >=stealth, thick] (0,-0.7) -- (0,-1.4) node[midway, right] {$(x,y,z) \mapsto (x,y)$};
        % disco
        \begin{scope}[shift={(0,-3.0)}]
            \draw[thick, fill=gray!10] (0,0) circle (1.5cm);
            \node at (0,0) {$\overline{B^2}$};
            \draw[thick, red, ->, >=stealth] (0:1.5) arc (0:100:1.5);
            \draw[thick, red] (100:1.5) arc (100:180:1.5);
            \draw[thick, red, ->, >=stealth] (180:1.5) arc (180:280:1.5);
            \draw[thick, red] (280:1.5) arc (280:360:1.5);
            \node[red] at (50:1.8) {$a$};
            \node[red] at (230:1.8) {$a$};
            \filldraw[blue] (0:1.5) circle (1.5pt) node[right] {$A'$};
            \filldraw[blue] (180:1.5) circle (1.5pt) node[left] {$A$};
        \end{scope}
    \end{tikzpicture}
    \caption{$S^2_+ \cong \overline{B^2}$ tramite la proiezione ortogonale; sul bordo si identificano i punti antipodali (i due semicerchi, entrambi percorsi in senso antiorario, sono identificati: parola $aa$).}
\end{figure}

Consideriamo l'emisfero superiore chiuso della sfera 2-dimensionale:
\[
S^2_+ = \left\{ (x,y,z) \in \mathbb{R}^3 \;\middle|\; x^2 + y^2 + z^2 = 1, \; z \geqslant 0 \right\}
\]

Ogni punto dell'emisfero inferiore $S^2_-$ ha un unico punto antipodale nell'emisfero superiore $S^2_+$. L'unica sovrapposizione avviene sull'equatore $z=0$, dove i punti antipodali si trovano sul bordo $\partial S^2_+ \cong S^1$. Inoltre $p|_{S^2_+}: S^2_+ \to \mathbb{R}P^2$ \`e continua e suriettiva da uno spazio compatto a uno spazio di Hausdorff, quindi \`e chiusa e dunque una applicazione quoziente; le sue fibre sono i singoli punti di $S^2_+$ con $z>0$ e le coppie di punti antipodali dell'equatore.

L'applicazione di proiezione ortogonale $f: S^2_+ \to \overline{B^2}$ definita da $f(x,y,z) = (x,y)$ è un \textbf{omeomorfismo} tra l'emisfero superiore $S^2_+$ ed il disco chiuso bidimensionale $\overline{B^2}$.

Pertanto, il piano proiettivo reale $\mathbb{R}P^2$ è topologicamente equivalente al \textbf{disco chiuso $\overline{B^2}$ con i punti antipodali del bordo $\partial B^2 = S^1$ identificati a due a due}:
\[
\mathbb{R}P^2 \cong \overline{B^2} \big/ \sim_{\text{antipodale su } \partial B^2}
\]
\end{ejemplo}

\section{Modelli topologici di $\mathbb{R}P^n$}

\begin{teo}
Lo spazio proiettivo reale $\mathbb{R}P^2$ è omeomorfo al quoziente del disco chiuso bidimensionale:
\[
\mathbb{R}P^2 \cong \overline{B^2} \big/ \sim
\]
dove $\sim$ è la relazione di equivalenza su $\overline{B^2}$, non banale solo sul bordo $\partial \overline{B^2} \cong S^1$, che identifica a due a due i punti antipodali di $\partial \overline{B^2}$.
\end{teo}

\vspace{1em}

\begin{prop}[Generalizzazione in dimensione $n$]
Per ogni $n \geqslant 1$, lo spazio proiettivo reale $\mathbb{R}P^n$ soddisfa:
\[
\mathbb{R}P^n = S^n \big/ \{\pm \mathrm{Id}\} = S^n \big/ \sim \;\;\cong\;\; \overline{B^n} \big/ \sim
\]
dove $\sim$ è la relazione di equivalenza su $\overline{B^n}$, non banale solo sul bordo $\partial \overline{B^n} = S^{n-1}$, che ne identifica a due a due i punti antipodali.
\end{prop}

\begin{proof}[Dimostrazione (per ogni $n \geqslant 1$, in particolare $n=2$ del teorema precedente)]
Sia $S^n_+ = \{x \in S^n \mid x_{n+1} \geqslant 0\}$. Ogni retta per l'origine incontra $S^n_+$ (se $x_{n+1} < 0$ si usa $-x$), e la incontra in due punti (antipodali) solo se sta nell'iperpiano $x_{n+1} = 0$, cio\`e se i punti sono sul bordo $\partial S^n_+ = S^{n-1}$. Dunque $p|_{S^n_+}: S^n_+ \to \mathbb{R}P^n$ \`e continua e suriettiva, con fibre ``banali'' fuori dal bordo e coppie antipodali sul bordo. Essendo $S^n_+$ compatto e $\mathbb{R}P^n$ di Hausdorff, $p|_{S^n_+}$ \`e chiusa, quindi quoziente: $\mathbb{R}P^n \cong S^n_+/\!\sim$. Infine la proiezione $(x_1, \dots, x_{n+1}) \mapsto (x_1, \dots, x_n)$ \`e un omeomorfismo $S^n_+ \to \overline{B^n}$ (inversa $y \mapsto (y, \sqrt{1-\|y\|^2})$) che manda $\partial S^n_+$ in $\partial \overline{B^n} = S^{n-1}$ e rispetta l'identificazione antipodale\hfill$\square$
\end{proof}

\begin{ejercicio*}[\textbf{DOMANDA D'ESAME}]
\textbf{Quali sono i modelli topologici di $\mathbb{R}P^n$?}
\begin{enumerate}
    \item Il quoziente dello spazio vettoriale: $\left(\mathbb{R}^{n+1} \setminus \{\bar{0}\}\right) / \sim$.
    \item Il quoziente della sfera unitaria per l'azione antipodale: $S^n / \{\pm \mathrm{Id}_{S^n}\}$.
    \item Il quoziente del disco chiuso con identificazione antipodale sul bordo: $\overline{B^n} / \sim$.
\end{enumerate}
\end{ejercicio*}

\vspace{2em}

\section{Dimostrazione: $(S^n, p)$ \`e un rivestimento di $\mathbb{R}P^n$}

\begin{teo}
La coppia $(S^n, p)$, con $p = \pi|_{S^n}: S^n \to \mathbb{R}P^n$, è uno \textbf{spazio di rivestimento a 2 foglie} dello spazio proiettivo reale $\mathbb{R}P^n$.
\end{teo}

\begin{proof}
Dimostriamo le proprietà fondamentali della mappa di rivestimento $p$:

\begin{itemize}
    \item \textbf{$p$ è continua:}
    La mappa $p = \pi|_{S^n}$ è la restrizione alla sfera della proiezione canonica $\pi: \mathbb{R}^{n+1} \setminus \{\bar{0}\} \to \mathbb{R}P^n$, che è continua per definizione di topologia quoziente.

    \item \textbf{$p$ è suriettiva:}
    Sia $[y]_\sim \in \mathbb{R}P^n$, con $y \in \mathbb{R}^{n+1} \setminus \{\bar{0}\}$.
    Normalizziamo il vettore ponendo $x = \frac{y}{\|y\|} \in S^n$.
    Allora l'immagine tramite $p$ è $p(x) = \pi(x) = \pi\left(\frac{y}{\|y\|}\right) = \pi(y) = [y]_\sim$.

    \item \textbf{Caratterizzazione della fibra:}
    Per ogni $x, x' \in S^n$, mostriamo che $p(x) = p(x') \iff x' = \pm x$.
    \begin{itemize}
        \item[$\impliedby$)] Banale: se $x' = \pm x \implies [x']_\sim = [x]_\sim \implies p(x') = p(x)$.
        \item[$\implies$)] Se $p(x) = p(x')$, allora $[x]_\sim = [x']_\sim$, il che implica che esiste $t \neq 0$ tale che $x' = t x$.
        Passando alle norme euclidee:
        \[
        1 = \|x'\| = \|t x\| = |t| \cdot \|x\| = |t| \implies t = \pm 1 \implies x' = \pm x
        \]
    \end{itemize}
    Di conseguenza, per ogni $[y]_\sim \in \mathbb{R}P^n$, la fibra $p^{-1}([y]_\sim)$ contiene esattamente due punti antipodali:
    \[
    p^{-1}([y]_\sim) = \left\{ \frac{y}{\|y\|}, \, -\frac{y}{\|y\|} \right\}
    \]
\end{itemize}

\begin{figure}[H]
    \centering
    \begin{tikzpicture}[scale=1.2]
        % Sfera S^n
        \draw[thick] (0,0) circle (1.5cm);
        \draw[dashed] (-1.5,0) arc (180:0:1.5cm and 0.4cm);
        \draw (-1.5,0) arc (180:360:1.5cm and 0.4cm);
        \node[above right] at (1.2,1.2) {$S^n$};

        % Punti antipodali
        \filldraw[blue] (0.9, 0.7) circle (1.5pt) node[above right] {$x$};
        \filldraw[blue] (-0.9, -0.7) circle (1.5pt) node[below left] {$-x$};
        
        % Retta per l'origine
        \draw[thick, gray, dashed] (-1.8,-1.4) -- (1.8,1.4) node[right] {$[y]_\sim$};
        \filldraw[red] (0,0) circle (1.5pt) node[below] {$\bar{0}$};
    \end{tikzpicture}
    \caption{Intersezione della retta vettoriale $[y]_\sim$ con la sfera $S^n$ nei punti antipodali $\pm x$.}
\end{figure}

\begin{itemize}
    \item \textbf{Costruzione degli intorni ben coperti.} Sia $[y]_\sim \in \mathbb{R}P^n$ e $x = y/\|y\|$. Poniamo
    \[
    V = \{ [v]_\sim \;:\; \langle v, x \rangle \neq 0 \} \subseteq \mathbb{R}P^n .
    \]
    La condizione $\langle v,x\rangle \neq 0$ \`e invariante per $v \mapsto tv$, e $\pi^{-1}(V) = \{v : \langle v, x\rangle \neq 0\}$ \`e aperto; quindi $V$ \`e un intorno aperto di $[y]_\sim$ (contiene $[x]_\sim$). La sua controimmagine in $S^n$ \`e l'unione disgiunta di due emisferi aperti
    \[
    p^{-1}(V) = \widetilde{U}_+ \sqcup \widetilde{U}_-, \qquad \widetilde{U}_\pm = \{ z \in S^n : \pm\langle z, x \rangle > 0 \},
    \]
    antipodali tra loro ($\widetilde{U}_- = -\widetilde{U}_+$), aperti in $S^n$ e disgiunti.

    Mostriamo che $p|_{\widetilde{U}_+}: \widetilde{U}_+ \to V$ \`e un omeomorfismo (per $\widetilde{U}_-$ \`e identico).
    \begin{enumerate}[label=\roman*)]
        \item \`E continua (restrizione di $p$) e \textbf{iniettiva}: se $p(z) = p(z')$ con $z, z' \in \widetilde{U}_+$, allora $z' = \pm z$; ma $-z \in \widetilde{U}_-$, quindi $z' = z$.
        \item \`E \textbf{suriettiva}: se $[v]_\sim \in V$, allora $\langle v, x \rangle \neq 0$ e il punto $z = \operatorname{sgn}\langle v,x\rangle \, v/\|v\|$ sta in $\widetilde{U}_+$ e $p(z) = [v]_\sim$.
        \item L'\textbf{inversa} $V \to \widetilde{U}_+$ \`e continua. Sia $s: \pi^{-1}(V) \to \widetilde{U}_+$ data da $s(v) = \operatorname{sgn}\langle v,x\rangle\, v/\|v\|$: \`e continua (il segno \`e localmente costante su $\pi^{-1}(V)$) e costante sulle fibre di $\pi$, perch\'e $s(tv) = s(v)$ per ogni $t \neq 0$. Poich\'e $\pi|_{\pi^{-1}(V)}: \pi^{-1}(V) \to V$ \`e una applicazione quoziente (restrizione di un quoziente a un aperto saturo), $s$ induce una applicazione continua $V \to \widetilde{U}_+$, che \`e proprio $(p|_{\widetilde{U}_+})^{-1}$.
    \end{enumerate}
    Quindi ciascun $\widetilde{U}_\pm$ \`e un foglio omeomorfo a $V$ tramite $p$. (Inoltre $S^n$ \`e connessa per archi per $n \geqslant 1$, come richiesto dalla definizione di rivestimento.)
\end{itemize}
In conclusione, $(S^n, p)$ è uno spazio di rivestimento a 2 foglie di $\mathbb{R}P^n$\hfill$\square$
\end{proof}

% ======== NUOVO (pp. 89-92 degli appunti, 13/10/2026) ========
\section{Altri esempi di rivestimenti e sollevamenti}

\begin{ejemplo}[Rivestimenti del toro]
\begin{enumerate}
    \item L'applicazione
    \Func{p}{\mathbb{R} \times S^1}{S^1 \times S^1 = T^2}{(x,z)}{(e^{2\pi i x},\, z)}
    \`e un'applicazione di rivestimento del toro $T^2$ tramite il \emph{cilindro} $\mathbb{R} \times S^1$ (si ``avvolge'' soltanto la prima circonferenza).
    \item Anche l'applicazione
    \Func{p}{\mathbb{R}^2 = \mathbb{R} \times \mathbb{R}}{\mathbb{R} \times S^1}{(x,y)}{(x,\, e^{2\pi i y})}
    \`e un rivestimento (del cilindro).
\end{enumerate}
\end{ejemplo}

\begin{observacion}
Il rivestimento di uno spazio $X$ \textbf{non \`e unico}: sia $\mathbb{R}^2$ sia il cilindro $\mathbb{R} \times S^1$ sono rivestimenti di $T^2$ (il primo tramite $(x,y) \mapsto (e^{2\pi i x}, e^{2\pi i y})$, il secondo come sopra). Il rivestimento $\mathbb{R}^2 \to T^2$ ``srotola'' entrambe le circonferenze, quello dal cilindro ne srotola una sola.
\end{observacion}

\begin{figure}[H]
    \centering
    \begin{tikzpicture}[scale=1.1]
        % cilindro
        \draw[thick] (-1.2,0) ellipse (0.35 and 0.8);
        \draw[thick] (-1.2,0.8) -- (1.6,0.8);
        \draw[thick] (-1.2,-0.8) -- (1.6,-0.8);
        \draw[thick, dashed] (1.6,0) ellipse (0.35 and 0.8);
        \node at (0.2,0) {$\mathbb{R} \times S^1$};
        \draw[->, >=stealth, thick, Purple] (0.2,-1.1) -- (0.2,-2.0) node[midway, right] {$p$};
        % toro (schematico)
        \begin{scope}[shift={(0.2,-3.4)}]
            \draw[thick] (0,0) ellipse (1.8 and 0.9);
            \draw[thick] (-0.75,0.05) to[bend right=40] (0.75,0.05);
            \draw[thick] (-0.55,-0.02) to[bend left=40] (0.55,-0.02);
            \node at (2.6,0) {$T^2$};
        \end{scope}
    \end{tikzpicture}
    \caption{Il cilindro $\mathbb{R}\times S^1$ si avvolge sul toro $T^2$ lungo la direzione di $\mathbb{R}$.}
\end{figure}

\section{Sollevamenti}

\begin{definicion}[Sollevamento]
Sia $(\widetilde{X}, p)$ uno spazio di rivestimento di $X$, $Y$ uno spazio topologico e $f: Y \to X$ continua. Un'applicazione continua $\widetilde{f}: Y \to \widetilde{X}$ si dice \textbf{sollevamento} di $f$ se $p \circ \widetilde{f} = f$, cio\`e se il seguente diagramma commuta:
\begin{center}
\begin{tikzcd}
 & \widetilde{X} \arrow[d, "p"] \\
Y \arrow[ur, dashed, "\widetilde{f}"] \arrow[r, "f"'] & X
\end{tikzcd}
\end{center}
\end{definicion}

\begin{intuizione}
Sollevare $f$ significa ``scegliere con continuit\`a'' una controimmagine, nella fibra di $p$, del punto $f(y)$ per ogni $y \in Y$. Il teorema seguente dice che, una volta scelta la controimmagine in un punto, la scelta in tutti gli altri punti (se $Y$ \`e connesso) \`e obbligata.
\end{intuizione}

\begin{teo}[Unicit\`a del sollevamento]
Sia $(\widetilde{X}, p)$ uno spazio di rivestimento di $X$, $Y$ uno spazio topologico \textbf{connesso} e $f: Y \to X$ continua. Siano $\widetilde{f}_1, \widetilde{f}_2: Y \to \widetilde{X}$ due sollevamenti di $f$ (cio\`e $p \circ \widetilde{f}_1 = p \circ \widetilde{f}_2 = f$). Se esiste $y_0 \in Y$ tale che $\widetilde{f}_1(y_0) = \widetilde{f}_2(y_0)$, allora $\widetilde{f}_1 = \widetilde{f}_2$.
\end{teo}

\begin{proof}
Poniamo $Z = \{ y \in Y \mid \widetilde{f}_1(y) = \widetilde{f}_2(y) \}$. Vogliamo mostrare che $Z$ \`e non vuoto, aperto e chiuso: essendo $Y$ connesso, seguir\`a $Z = Y$.

\begin{itemize}
    \item \textbf{$Z \neq \varnothing$:} $y_0 \in Z$ per ipotesi.
    \item \textbf{$Z$ \`e aperto.} Sia $y \in Z$ e sia $U$ un intorno aperto ben coperto di $f(y) = p(\widetilde{f}_1(y)) \in X$. Sia $W$ l'unico foglio su $U$ contenente $\widetilde{f}_1(y) = \widetilde{f}_2(y)$. Poniamo $V = \widetilde{f}_1^{-1}(W) \cap \widetilde{f}_2^{-1}(W)$: \`e un aperto di $Y$ che contiene $y$. Se $y' \in V$, allora $\widetilde{f}_1(y'), \widetilde{f}_2(y') \in W$ e $p(\widetilde{f}_1(y')) = f(y') = p(\widetilde{f}_2(y'))$; poich\'e $p|_W: W \to U$ \`e un omeomorfismo (in particolare iniettiva), $\widetilde{f}_1(y') = \widetilde{f}_2(y')$, ossia $y' \in Z$. Dunque $V \subseteq Z$ e $Z$ \`e intorno di ogni suo punto.
    \item \textbf{$Z$ \`e chiuso} (cio\`e $Y \setminus Z$ \`e aperto). Sia $y \notin Z$: $\widetilde{f}_1(y) \neq \widetilde{f}_2(y)$, ma $p(\widetilde{f}_1(y)) = f(y) = p(\widetilde{f}_2(y))$, cio\`e i due punti stanno nella stessa fibra. Sia $U$ un intorno aperto ben coperto di $f(y)$ e siano $W_1, W_2$ i fogli su $U$ contenenti rispettivamente $\widetilde{f}_1(y)$ e $\widetilde{f}_2(y)$. Poich\'e $p$ \`e iniettiva su ciascun foglio e i due punti sono distinti con la stessa immagine, appartengono a fogli \emph{distinti}, quindi $W_1 \cap W_2 = \varnothing$. Posto $V = \widetilde{f}_1^{-1}(W_1) \cap \widetilde{f}_2^{-1}(W_2)$, aperto contenente $y$, per ogni $y' \in V$ si ha $\widetilde{f}_1(y') \in W_1$ e $\widetilde{f}_2(y') \in W_2$, che sono disgiunti: $\widetilde{f}_1(y') \neq \widetilde{f}_2(y')$. Dunque $V \subseteq Y \setminus Z$.
\end{itemize}
Per la connessione di $Y$ si conclude $Z = Y$, cio\`e $\widetilde{f}_1 = \widetilde{f}_2$\hfill$\square$
\end{proof}

\begin{observacion}
L'ipotesi di connessione \`e essenziale: se $Y = \{a\} \sqcup \{b\}$ \`e discreto e $f$ costante su $X = S^1$, si possono sollevare $a$ e $b$ in due punti diversi di una stessa fibra di $\mathbb{R} \to S^1$ senza alcun legame tra loro.
\end{observacion}


% ======== NUOVO (pp. 93-100 degli appunti) ========
\begin{coro}
Sia $(\widetilde{X}, p)$ uno spazio di rivestimento di $X$ e sia $h: \widetilde{X} \to \widetilde{X}$ continua con $p \circ h = p$. Se esiste $\widetilde{x}_0 \in \widetilde{X}$ con $h(\widetilde{x}_0) = \widetilde{x}_0$, allora $h = \mathrm{Id}_{\widetilde{X}}$.
\end{coro}

\begin{proof}
$\widetilde{X}$ \`e connesso per archi per definizione. Le applicazioni $h$ e $\mathrm{Id}_{\widetilde{X}}$ sono entrambe sollevamenti di $p: \widetilde{X} \to X$ (infatti $p \circ h = p = p \circ \mathrm{Id}_{\widetilde{X}}$) e coincidono in $\widetilde{x}_0$. Per il teorema di unicit\`a del sollevamento (con $Y = \widetilde{X}$ connesso), $h = \mathrm{Id}_{\widetilde{X}}$\hfill$\square$
\end{proof}

\begin{intuizione}
Una ``simmetria'' del rivestimento (un omeomorfismo $h$ di $\widetilde{X}$ che rispetta le fibre) che fissa anche un solo punto \`e necessariamente l'identit\`a: \`e rigidamente determinata dal valore in un punto.
\end{intuizione}

\section{Sollevamento di cammini e di omotopie}

\begin{teo}[Sollevamento dei cammini]
Sia $(\widetilde{X}, p)$ uno spazio di rivestimento di $X$, $x_0 \in X$ e $\widetilde{x}_0 \in p^{-1}(x_0)$. Se $\alpha: I \to X$ \`e un cammino con $\alpha(0) = x_0$, allora esiste un \textbf{unico} cammino $\widetilde{\alpha}: I \to \widetilde{X}$ tale che
\[
p \circ \widetilde{\alpha} = \alpha \qquad \text{e} \qquad \widetilde{\alpha}(0) = \widetilde{x}_0 .
\]
($\widetilde{\alpha}$ \`e detto \textbf{sollevamento} di $\alpha$ \emph{uscente da} $\widetilde{x}_0$.)
\end{teo}

\begin{figure}[H]
    \centering
    \begin{tikzpicture}[scale=1.2]
        \draw[gray!60, fill=gray!8] (-0.4,1.6) -- (3.6,1.6) -- (4.2,2.8) -- (0.2,2.8) -- cycle;
        \node[gray] at (3.9,3.0) {$\widetilde{X}$};
        \draw[gray!60] (-0.4,0) -- (3.6,0);
        \node[gray] at (3.9,0.2) {$X$};
        \filldraw (0.6,2.1) circle (1.3pt) node[left] {$\widetilde{x}_0$};
        \draw[thick, red] (0.6,2.1) to[out=20, in=200] (1.8,2.4) to[out=20, in=180] (3.0,2.0);
        \node[red, above] at (1.9,2.45) {$\widetilde{\alpha}$};
        \filldraw (0.6,0) circle (1.3pt) node[below] {$x_0$};
        \draw[thick, blue] (0.6,0) to[out=20, in=200] (1.8,0.3) to[out=20, in=180] (3.0,-0.1);
        \node[blue, below] at (1.9,0.1) {$\alpha$};
        \draw[->, >=stealth, thick, Purple] (-0.8,1.5) -- (-0.8,0.4) node[midway, left] {$p$};
        \draw[dashed, gray] (0.6,2.1) -- (0.6,0);
    \end{tikzpicture}
    \caption{Il cammino $\alpha$ in $X$ si solleva, in modo unico, a un cammino $\widetilde{\alpha}$ in $\widetilde{X}$ che parte dal punto prescelto $\widetilde{x}_0$ della fibra.}
\end{figure}

\begin{proof}
\textbf{Unicit\`a.} Segue dal teorema di unicit\`a del sollevamento, con $Y = I$ connesso: due sollevamenti di $\alpha$ uscenti da $\widetilde{x}_0$ coincidono in $0$, dunque coincidono.

\textbf{Esistenza.}
\begin{itemize}
    \item[\textbf{I passo.}] Supponiamo che esista un intorno aperto ben coperto $U$ di $x_0$ in $X$ con $\alpha(I) \subseteq U$. Sia $W$ l'unico foglio su $U$ contenente $\widetilde{x}_0$; allora $p|_W: W \to U$ \`e un omeomorfismo. Poniamo
    \[
    \widetilde{\alpha} = (p|_W)^{-1} \circ \alpha : I \to W \subseteq \widetilde{X} .
    \]
    \`E continua (composizione di continue), $p \circ \widetilde{\alpha} = p \circ (p|_W)^{-1} \circ \alpha = \alpha$ e $\widetilde{\alpha}(0) = (p|_W)^{-1}(x_0) = \widetilde{x}_0$.

    \item[\textbf{II passo.}] Supponiamo che esistano due aperti ben coperti $U_1, U_2$ di $X$ e $t_1 \in\, ]0,1[$ tali che, posto $I_1 = [0,t_1]$, $I_2 = [t_1,1]$, risulti $\alpha(I_1) \subseteq U_1$ e $\alpha(I_2) \subseteq U_2$. Siano $\alpha_i = \alpha|_{I_i}$. Per il I passo (applicato a $\alpha_1$ e, poi, a $\alpha_2$) esiste un unico sollevamento $\widetilde{\alpha}_1: I_1 \to \widetilde{X}$ di $\alpha_1$ uscente da $\widetilde{x}_0$, ed esiste un unico sollevamento $\widetilde{\alpha}_2: I_2 \to \widetilde{X}$ di $\alpha_2$ uscente da $\widetilde{\alpha}_1(t_1)$. Poich\'e $\widetilde{\alpha}_1(t_1) = \widetilde{\alpha}_2(t_1)$, per il Lemma di Incollamento ($I_1, I_2$ chiusi) l'applicazione
    \[
    \widetilde{\alpha}(t) = \begin{cases} \widetilde{\alpha}_1(t) & t \in I_1 \\ \widetilde{\alpha}_2(t) & t \in I_2 \end{cases}
    \]
    \`e continua, $\widetilde{\alpha}(0) = \widetilde{x}_0$ e $p \circ \widetilde{\alpha} = \alpha$ su $I_1$ e su $I_2$, quindi su tutto $I$.

    \item[\textbf{III passo (caso generale).}] Gli aperti ben coperti di $X$ ricoprono $X$, quindi $\{\alpha^{-1}(U) \mid U \text{ ben coperto}\}$ \`e un ricoprimento aperto del compatto $I$. Per il Lemma del numero di Lebesgue esistono $0 = t_0 < t_1 < \dots < t_k = 1$ e aperti ben coperti $U_1, \dots, U_k$ tali che, posto $I_i = [t_{i-1}, t_i]$, si abbia $\alpha(I_i) \subseteq U_i$ per ogni $i$. Si procede come nel II passo, sollevando in successione $\alpha|_{I_1}$ (a partire da $\widetilde{x}_0$), poi $\alpha|_{I_2}$ (a partire dall'estremo finale del sollevamento precedente), e cos\`i via, e si incollano i sollevamenti ottenuti. Si ottiene $\widetilde{\alpha}$ con $p \circ \widetilde{\alpha} = \alpha$ e $\widetilde{\alpha}(0) = \widetilde{x}_0$\hfill$\square$
\end{itemize}
\end{proof}

\begin{observacion}
\begin{enumerate}
    \item Se $\widetilde{\alpha}$ \`e il sollevamento di $\alpha$ (uscente da $\widetilde{x}_0$), allora $\widetilde{\alpha}^{-1}$ \`e un sollevamento di $\alpha^{-1}$. Infatti, per ogni $t \in I$: $(p \circ \widetilde{\alpha}^{-1})(t) = p(\widetilde{\alpha}(1-t)) = \alpha(1-t) = \alpha^{-1}(t)$.
    \item Sia $\sigma: I \to X$ un \textbf{cappio} in $x_0$ e sia $\widetilde{x}_0 \in p^{-1}(x_0)$. Se $\widetilde{\sigma}$ \`e l'unico sollevamento di $\sigma$ uscente da $\widetilde{x}_0$, allora in generale $\widetilde{\sigma}$ \textbf{non \`e un cappio}; tuttavia il suo punto finale sta nella fibra: $\widetilde{\sigma}(1) \in p^{-1}(x_0)$, perch\'e $p(\widetilde{\sigma}(1)) = \sigma(1) = x_0$. (Esempio: il cappio $t \mapsto e^{2\pi i t}$ in $S^1$ si solleva a $t \mapsto t$ in $\mathbb{R}$, che va da $0$ a $1 \neq 0$.)
\end{enumerate}
\end{observacion}

\begin{teo}[Sollevamento dell'omotopia]
Sia $(\widetilde{X}, p)$ uno spazio di rivestimento di $X$, $x_0 \in X$, $\widetilde{x}_0 \in p^{-1}(x_0)$, $\alpha: I \to X$ un cammino con $\alpha(0) = x_0$ e $\widetilde{\alpha}$ il suo sollevamento uscente da $\widetilde{x}_0$. Se $H: I \times I \to X$ \`e un'omotopia con $H(t,0) = \alpha(t)$ per ogni $t \in I$, allora esiste un \textbf{unico} sollevamento continuo $\widetilde{H}: I \times I \to \widetilde{X}$ di $H$ (cio\`e $p \circ \widetilde{H} = H$) tale che $\widetilde{H}(t,0) = \widetilde{\alpha}(t)$ per ogni $t \in I$.
\end{teo}

\begin{proof}
\textbf{Unicit\`a.} Il dominio $I \times I$ \`e connesso e $\widetilde{H}$ \`e determinata dalle condizioni $p \circ \widetilde{H} = H$ e $\widetilde{H}(0,0) = \widetilde{\alpha}(0) = \widetilde{x}_0$: segue dal teorema di unicit\`a del sollevamento.

\textbf{Esistenza} (\emph{dimostrazione facoltativa}). Poich\'e $I \times I$ \`e compatto e $H$ \`e continua, $H(I \times I)$ \`e coperto da un numero finito di aperti ben coperti; per il Lemma di Lebesgue (applicato al ricoprimento $\{H^{-1}(U)\}$ di $I \times I$) esistono suddivisioni $0 = a_0 < \dots < a_k = 1$ e $0 = b_0 < \dots < b_h = 1$ tali che, posto $R_{ij} = [a_{i}, a_{i+1}] \times [b_j, b_{j+1}]$, ciascun $H(R_{ij})$ \`e contenuto in un aperto ben coperto $U_{ij}$ di $X$.

Sollevamo $H$ un rettangolo alla volta, in ordine lessicografico ($R_{00}, R_{10}, \dots$, poi la riga successiva). Supponiamo di aver gi\`a definito il sollevamento $\widetilde{H}$ sull'unione $A$ dei rettangoli precedenti; l'intersezione $A \cap R_{ij}$ \`e l'unione del lato sinistro e del lato inferiore di $R_{ij}$ (o solo uno dei due, o vuota per il primo rettangolo), cio\`e un insieme \textbf{connesso} $L$ che contiene l'angolo in basso a sinistra. Allora:
\begin{itemize}
    \item $\widetilde{H}(L) \subseteq p^{-1}(U_{ij}) = \bigsqcup_\beta W_\beta$; essendo $L$ connesso e i fogli $W_\beta$ aperti e disgiunti, $\widetilde{H}(L)$ \`e contenuto in \emph{un solo} foglio $W$ (per il primo rettangolo si sceglie il foglio contenente $\widetilde{x}_0$).
    \item Si definisce $\widetilde{H}$ su $R_{ij}$ ponendo $\widetilde{H} = (p|_W)^{-1} \circ H$. Su $L$ questa formula coincide con il sollevamento gi\`a definito (per l'unicit\`a del sollevamento su $L$ connesso, dato che entrambe sollevano $H|_L$ e coincidono nell'angolo).
\end{itemize}
Per il Lemma di Incollamento (unione finita di chiusi) $\widetilde{H}$ \`e continua su tutto $I \times I$, $p \circ \widetilde{H} = H$ e $\widetilde{H}(t,0) = \widetilde{\alpha}(t)$ (perch\'e sul lato inferiore $\widetilde{H}$ solleva $\alpha$ a partire da $\widetilde{x}_0$, e vale l'unicit\`a del sollevamento dei cammini)\hfill$\square$
\end{proof}

\begin{observacion}[Importante]
Se $t_0 \in I$ e $H(t_0, s) = \alpha(t_0)$ per ogni $s \in I$ (cio\`e se l'omotopia lascia fisso il punto $\alpha(t_0)$ durante tutta la deformazione), allora anche $\widetilde{H}(t_0, s) = \widetilde{\alpha}(t_0)$ per ogni $s \in I$: il sollevamento ``non si muove'' sopra i punti fissati.
\end{observacion}

\begin{proof}
Consideriamo i cammini $\widetilde{\beta}(s) = \widetilde{H}(t_0, s)$ e il cammino costante $c_{\widetilde{\alpha}(t_0)}$. Si ha $\widetilde{\beta}(0) = \widetilde{H}(t_0, 0) = \widetilde{\alpha}(t_0) = c_{\widetilde{\alpha}(t_0)}(0)$. Inoltre $p \circ \widetilde{\beta}(s) = H(t_0, s) = \alpha(t_0)$ e $p \circ c_{\widetilde{\alpha}(t_0)}(s) = p(\widetilde{\alpha}(t_0)) = \alpha(t_0)$. Quindi sono due sollevamenti dello stesso cammino costante, che coincidono in $s = 0$: per l'unicit\`a del sollevamento $\widetilde{\beta} = c_{\widetilde{\alpha}(t_0)}$, ossia $\widetilde{H}(t_0, s) = \widetilde{\alpha}(t_0)$ per ogni $s$\hfill$\square$
\end{proof}

\begin{prop}[Omotopia di cammini e sollevamenti]
Sia $(\widetilde{X}, p)$ uno spazio di rivestimento di $X$, $x_0 \in X$, $\widetilde{x}_0 \in p^{-1}(x_0)$, $\alpha, \beta: I \to X$ cammini con $\alpha(0) = \beta(0) = x_0$, e siano $\widetilde{\alpha}, \widetilde{\beta}$ i rispettivi sollevamenti uscenti da $\widetilde{x}_0$. Allora
\[
\alpha \simeq \beta \text{ rel } \partial I \implies \widetilde{\alpha} \simeq \widetilde{\beta} \text{ rel } \partial I .
\]
In particolare $\widetilde{\alpha}(1) = \widetilde{\beta}(1)$: i sollevamenti di cammini omotopi hanno lo \emph{stesso punto finale}.
\end{prop}

\begin{proof}
Sia $H: \alpha \simeq \beta \text{ rel } \partial I$. Per il teorema di sollevamento dell'omotopia esiste un'unica $\widetilde{H}: I \times I \to \widetilde{X}$ con $p \circ \widetilde{H} = H$ e $\widetilde{H}(t,0) = \widetilde{\alpha}(t)$. Poich\'e $H(0,s) = x_0$ e $H(1,s) = \alpha(1)$ per ogni $s$, l'osservazione precedente d\`a $\widetilde{H}(0,s) = \widetilde{x}_0$ e $\widetilde{H}(1,s) = \widetilde{\alpha}(1)$ per ogni $s$.

Poniamo $\widetilde{\beta}'(t) = \widetilde{H}(t,1)$. Allora $\widetilde{\beta}'(0) = \widetilde{H}(0,1) = \widetilde{x}_0 = \widetilde{\beta}(0)$ e $p \circ \widetilde{\beta}'(t) = H(t,1) = \beta(t) = p \circ \widetilde{\beta}(t)$ per ogni $t$. Per l'unicit\`a del sollevamento, $\widetilde{\beta}' = \widetilde{\beta}$. Dunque $\widetilde{H}$ \`e un'omotopia da $\widetilde{\alpha}$ a $\widetilde{\beta}$ relativa a $\partial I$ (gli estremi restano fissi: $\widetilde{H}(0,s) = \widetilde{x}_0$, $\widetilde{H}(1,s) = \widetilde{\alpha}(1)$)\hfill$\square$
\end{proof}


% ======== NUOVO (pp. 101-104) ========
\section{Rivestimenti e gruppo fondamentale}

\begin{coro}
Sia $(\widetilde{X}, p)$ uno spazio di rivestimento di $X$, $x_0 \in X$ e $\widetilde{x}_0 \in p^{-1}(x_0)$. Allora l'omomorfismo
\[
p_*: \Pi_1(\widetilde{X}, \widetilde{x}_0) \longrightarrow \Pi_1(X, x_0)
\]
\`e \textbf{iniettivo}.
\end{coro}

\begin{proof}
Siano $[\widetilde{\sigma}], [\widetilde{\tau}] \in \Pi_1(\widetilde{X}, \widetilde{x}_0)$ con $p_*[\widetilde{\sigma}] = p_*[\widetilde{\tau}]$, cio\`e $[p \circ \widetilde{\sigma}] = [p \circ \widetilde{\tau}]$, quindi $p \circ \widetilde{\sigma} \simeq p \circ \widetilde{\tau}$ rel $\partial I$. Ora $\widetilde{\sigma}$ e $\widetilde{\tau}$ sono i sollevamenti, uscenti da $\widetilde{x}_0$, rispettivamente di $p \circ \widetilde{\sigma}$ e $p \circ \widetilde{\tau}$. Per la proposizione sull'omotopia di cammini e sollevamenti, $\widetilde{\sigma} \simeq \widetilde{\tau}$ rel $\partial I$, cio\`e $[\widetilde{\sigma}] = [\widetilde{\tau}]$\hfill$\square$
\end{proof}

\begin{observacion}
Di conseguenza $\Pi_1(X, x_0)$ contiene un sottogruppo $H = p_*\big(\Pi_1(\widetilde{X}, \widetilde{x}_0)\big)$ isomorfo a $\Pi_1(\widetilde{X}, \widetilde{x}_0)$.
\end{observacion}

\begin{prop}
Se $(\widetilde{X}, p)$ \`e uno spazio di rivestimento di $X$ e $X$ \`e \textbf{semplicemente connesso}, allora $p: \widetilde{X} \to X$ \`e un omeomorfismo.
\end{prop}

\begin{proof}
$p$ \`e continua, suriettiva e aperta; basta provare che \`e iniettiva (una bigezione continua e aperta \`e un omeomorfismo). Siano $\widetilde{x}_0, \widetilde{x}_1 \in \widetilde{X}$ con $p(\widetilde{x}_0) = p(\widetilde{x}_1) =: x_0$. Poich\'e $\widetilde{X}$ \`e connesso per archi, esiste un cammino $\widetilde{\alpha}: I \to \widetilde{X}$ con $\widetilde{\alpha}(0) = \widetilde{x}_0$ e $\widetilde{\alpha}(1) = \widetilde{x}_1$. Posto $\alpha = p \circ \widetilde{\alpha}$, si ha $\alpha(0) = \alpha(1) = x_0$, cio\`e $\alpha \in \Lambda_{x_0}(X)$. Poich\'e $X$ \`e semplicemente connesso, $\alpha \simeq c_{x_0}$ rel $\partial I$. Ma $\widetilde{\alpha}$ \`e il sollevamento di $\alpha$ uscente da $\widetilde{x}_0$ e $c_{\widetilde{x}_0}$ \`e il sollevamento di $c_{x_0}$ uscente da $\widetilde{x}_0$; per la proposizione precedente $\widetilde{\alpha} \simeq c_{\widetilde{x}_0}$ rel $\partial I$, quindi hanno lo stesso punto finale:
\[
\widetilde{x}_1 = \widetilde{\alpha}(1) = c_{\widetilde{x}_0}(1) = \widetilde{x}_0 . \qquad \square
\]
\end{proof}

\begin{observacion}
Uno spazio semplicemente connesso non ammette rivestimenti non banali (cio\`e, con pi\`u di un foglio).
\end{observacion}

\begin{definicion}[Rivestimento universale]
Uno spazio di rivestimento $(\widetilde{X}, p)$ di $X$ si dice \textbf{rivestimento universale} se $\widetilde{X}$ \`e semplicemente connesso.
\end{definicion}

\begin{ejemplo}
Si ha la catena di rivestimenti $\mathbb{R}^2 \to \mathbb{R} \times S^1 \to T^2$: il toro \`e rivestito dal cilindro e il cilindro da $\mathbb{R}^2$, ma $\mathbb{R}^2$ (semplicemente connesso) non ammette rivestimenti non banali. Sono rivestimenti universali: $\mathbb{R} \to S^1$, $\mathbb{R}^n \to T^n$ e $S^n \to \mathbb{R}P^n$ per $n \geqslant 2$.
\end{ejemplo}

\begin{teo}[Risultato centrale]
Sia $(\widetilde{X}, p)$ uno spazio di rivestimento di $X$, $x_0 \in X$, $\widetilde{x}_0 \in p^{-1}(x_0)$ e $H = p_*\big(\Pi_1(\widetilde{X}, \widetilde{x}_0)\big) \leqslant \Pi_1(X, x_0)$. Allora esiste una corrispondenza biunivoca tra la fibra $p^{-1}(x_0)$ e l'insieme quoziente (dei laterali destri)
\[
H \backslash \Pi_1(X, x_0) = \big\{ Hg \mid g \in \Pi_1(X, x_0) \big\}.
\]
\end{teo}

\begin{proof}
Dato $\sigma \in \Lambda_{x_0}(X)$, sia $\widetilde{\sigma}$ il suo sollevamento uscente da $\widetilde{x}_0$ (unico per il teorema di sollevamento dei cammini). Definiamo
\[
f: \Pi_1(X, x_0) \to p^{-1}(x_0), \qquad [\sigma] \mapsto \widetilde{\sigma}(1).
\]
\begin{enumerate}
    \item \textbf{$f$ \`e ben definita.} Se $[\sigma] = [\tau]$, allora $\sigma \simeq \tau$ rel $\partial I$ e, per la proposizione sull'omotopia di cammini e sollevamenti, $\widetilde{\sigma}(1) = \widetilde{\tau}(1)$. Inoltre $\widetilde{\sigma}(1) \in p^{-1}(x_0)$ perch\'e $p(\widetilde{\sigma}(1)) = \sigma(1) = x_0$.
    \item \textbf{$f$ \`e suriettiva.} Sia $\widetilde{x}_1 \in p^{-1}(x_0)$. Poich\'e $\widetilde{X}$ \`e connesso per archi esiste $\widetilde{\alpha}$ da $\widetilde{x}_0$ a $\widetilde{x}_1$. Posto $\alpha = p \circ \widetilde{\alpha}$, si ha $\alpha(0) = p(\widetilde{x}_0) = x_0$ e $\alpha(1) = p(\widetilde{x}_1) = x_0$, quindi $[\alpha] \in \Pi_1(X, x_0)$ e $\widetilde{\alpha}$ \`e il suo sollevamento uscente da $\widetilde{x}_0$: $f([\alpha]) = \widetilde{\alpha}(1) = \widetilde{x}_1$.
    \item \textbf{$f$ \`e costante sui laterali destri di $H$,} cio\`e $[\tau] \in H[\sigma] \Rightarrow f([\tau]) = f([\sigma])$. Sia $[\tau] = h[\sigma]$ con $h \in H$, quindi $h = p_*[\widetilde{\gamma}]$ per un cappio $\widetilde{\gamma}$ in $\widetilde{x}_0$. Allora
    \[
    [\tau] = p_*[\widetilde{\gamma}]\,[\sigma] = [(p \circ \widetilde{\gamma}) * \sigma].
    \]
    Il sollevamento di $(p \circ \widetilde{\gamma}) * \sigma$ uscente da $\widetilde{x}_0$ \`e $\widetilde{\gamma} * \widetilde{\sigma}$: infatti $\widetilde{\gamma} * \widetilde{\sigma}$ \`e ben definito (il punto finale di $\widetilde{\gamma}$ \`e $\widetilde{x}_0$, che \`e l'inizio di $\widetilde{\sigma}$), parte da $\widetilde{x}_0$ e $p \circ (\widetilde{\gamma} * \widetilde{\sigma}) = (p\circ\widetilde{\gamma}) * (p \circ \widetilde{\sigma}) = (p \circ \widetilde{\gamma}) * \sigma$. Per l'unicit\`a del sollevamento, $f([\tau]) = (\widetilde{\gamma} * \widetilde{\sigma})(1) = \widetilde{\sigma}(1) = f([\sigma])$.
    \item \textbf{Iniettivit\`a al quoziente.} Siano $\sigma, \tau \in \Lambda_{x_0}(X)$ con $f([\sigma]) = f([\tau])$, cio\`e $\widetilde{\sigma}(1) = \widetilde{\tau}(1)$. Allora $\widetilde{\gamma} = \widetilde{\sigma} * \widetilde{\tau}^{-1}$ \`e un cappio in $\widetilde{x}_0$, quindi $h = p_*[\widetilde{\gamma}] \in H$ e
    \[
    h[\tau] = [(p \circ \widetilde{\sigma}) * (p \circ \widetilde{\tau}^{-1})][\tau] = [\sigma * \tau^{-1} * \tau] = [\sigma * c_{x_0}] = [\sigma],
    \]
    cio\`e $[\sigma] \in H[\tau]$.
\end{enumerate}
Quindi $f$ induce $F: H\backslash \Pi_1(X, x_0) \to p^{-1}(x_0)$, $H[\sigma] \mapsto \widetilde{\sigma}(1)$, ben definita per il punto 3, suriettiva per il punto 2 e iniettiva per il punto 4: \`e una bigezione\hfill$\square$
\end{proof}

\begin{coro}
Tutte le fibre di un rivestimento hanno la stessa cardinalit\`a (uguale all'indice $[\Pi_1(X,x_0) : H]$; in particolare, se $\widetilde{X}$ \`e universale, a $|\Pi_1(X, x_0)|$).
\end{coro}


% ======== NUOVO (pp. 105-106) ========
\begin{coro}
Se $(\widetilde{X}, p)$ \`e un rivestimento \textbf{universale} di $X$ e $x_0 \in X$, $\widetilde{x}_0 \in p^{-1}(x_0)$, l'applicazione
\[
F: \Pi_1(X, x_0) \longrightarrow p^{-1}(x_0), \qquad [\sigma] \longmapsto \widetilde{\sigma}(1)
\]
\`e una \textbf{bigezione}.
\end{coro}
\begin{proof}
Poich\'e $\Pi_1(\widetilde{X}, \widetilde{x}_0) = 0$, il sottogruppo $H = p_*\Pi_1(\widetilde{X}, \widetilde{x}_0)$ \`e banale e $H \backslash \Pi_1(X, x_0) = \Pi_1(X, x_0)$. Si applica il teorema centrale\hfill$\square$
\end{proof}

\begin{prop}
Per ogni $n \geqslant 2$: $\Pi_1(\mathbb{R}P^n) \cong \mathbb{Z}_2$.
\end{prop}
\begin{proof}
Per $n \geqslant 2$, $S^n$ \`e semplicemente connessa, quindi $p: S^n \to \mathbb{R}P^n$ \`e un rivestimento universale a due fogli. Per ogni $[y]_\sim \in \mathbb{R}P^n$ la fibra \`e $p^{-1}([y]_\sim) = \{ y/\|y\|,\ -y/\|y\| \}$, di due elementi. Per il corollario precedente $\Pi_1(\mathbb{R}P^n)$ ha esattamente due elementi, dunque \`e il gruppo (unico) di ordine $2$: $\Pi_1(\mathbb{R}P^n) \cong \mathbb{Z}_2$\hfill$\square$
\end{proof}

\begin{intuizione}
In $\mathbb{R}P^n$ ($n \geqslant 2$) un cappio che sollevato in $S^n$ va da un punto al suo antipodo \emph{non} si pu\`o contrarre; ma percorso due volte torna al punto di partenza e diventa contraibile. Per questo $\Pi_1 = \mathbb{Z}_2$.
\end{intuizione}

\begin{teo}[Gruppo fondamentale della circonferenza]
$\Pi_1(S^1) \cong \mathbb{Z}$.
\end{teo}
\begin{proof}
Consideriamo il rivestimento universale $p: \mathbb{R} \to S^1$, $t \mapsto e^{2\pi i t}$. Fissiamo $x_0 = 1 \in S^1$, per cui $p^{-1}(1) = \mathbb{Z}$, e $\widetilde{x}_0 = 0 \in p^{-1}(1)$. Per il corollario precedente,
\[
F: \Pi_1(S^1, 1) \to p^{-1}(1) = \mathbb{Z}, \qquad [\sigma] \mapsto \widetilde{\sigma}(1)
\]
\`e una bigezione. Resta da mostrare che $F$ \`e un omomorfismo di gruppi: $F([\sigma][\tau]) = F([\sigma]) + F([\tau])$.

Siano $\sigma, \tau \in \Lambda_1(S^1)$ e $\widetilde{\sigma}, \widetilde{\tau}$ i sollevamenti uscenti da $0$; poniamo $m = \widetilde{\sigma}(1) \in \mathbb{Z}$ e $n = \widetilde{\tau}(1) \in \mathbb{Z}$. Definiamo
\[
\widetilde{\tau}_1(t) = m + \widetilde{\tau}(t), \qquad t \in I .
\]
\`E un cammino in $\mathbb{R}$ con $\widetilde{\tau}_1(0) = m = \widetilde{\sigma}(1)$ ed \`e ancora un sollevamento di $\tau$, perch\'e
\[
p(\widetilde{\tau}_1(t)) = e^{2\pi i (m + \widetilde{\tau}(t))} = \underbrace{e^{2\pi i m}}_{=1}\, e^{2\pi i \widetilde{\tau}(t)} = \tau(t).
\]
Dunque $\widetilde{\sigma} * \widetilde{\tau}_1$ \`e ben definito, parte da $0$ e $p \circ (\widetilde{\sigma} * \widetilde{\tau}_1) = (p \circ \widetilde{\sigma}) * (p \circ \widetilde{\tau}_1) = \sigma * \tau$: per l'unicit\`a del sollevamento, $\widetilde{\sigma * \tau} = \widetilde{\sigma} * \widetilde{\tau}_1$. Allora
\[
F([\sigma][\tau]) = (\widetilde{\sigma * \tau})(1) = \widetilde{\tau}_1(1) = m + \widetilde{\tau}(1) = m + n = F([\sigma]) + F([\tau]).
\]
Quindi $F$ \`e un isomorfismo di gruppi: $\Pi_1(S^1) \cong \mathbb{Z}$\hfill$\square$
\end{proof}

\begin{observacion}
Questa \`e una seconda dimostrazione di $\Pi_1(S^1) \cong \mathbb{Z}$, indipendente da Seifert--Van Kampen: usa il fatto che $\mathbb{R} \to S^1$ \`e un rivestimento universale. Pu\`o quindi essere usata (senza circolarit\`a) per tutti i risultati precedenti che la richiedevano, ad esempio $\Pi_1(S^n)=0$ solo per $n \geqslant 2$, il lemma sulla non-retrazione del disco e il teorema di Brouwer.
\end{observacion}

\begin{definicion}[Grado di un cappio]
Sia $p: (\mathbb{R}, 0) \to (S^1, 1)$ il rivestimento $t \mapsto e^{2\pi i t}$. Per $\sigma \in \Lambda_1(S^1)$ sia $\widetilde{\sigma}$ il suo sollevamento uscente da $\widetilde{x}_0 = 0$. Si chiama \textbf{grado} (o \textbf{indice di avvolgimento}) di $\sigma$ il numero intero
\[
\deg(\sigma) = \widetilde{\sigma}(1) \in \mathbb{Z}.
\]
\end{definicion}

\begin{intuizione}
$\deg(\sigma)$ conta quante volte (con segno: $+$ in senso antiorario) il cappio $\sigma$ gira attorno alla circonferenza.
\end{intuizione}

\begin{prop}
Per ogni $\sigma, \tau \in \Lambda_1(S^1)$: $\ \sigma \simeq \tau \text{ rel } \partial I \iff \deg(\sigma) = \deg(\tau)$.
\end{prop}
\begin{proof}
\[
\begin{aligned}
\sigma \simeq \tau \text{ rel } \partial I &\iff [\sigma] = [\tau] \overset{F \text{ biiettiva}}{\iff} F([\sigma]) = F([\tau]) \\
&\iff \widetilde{\sigma}(1) = \widetilde{\tau}(1) \iff \deg(\sigma) = \deg(\tau). \qquad \square
\end{aligned}
\]
\end{proof}

\begin{observacion}
Per ogni $n \in \mathbb{Z}$ sia $\omega_n(t) = e^{2\pi i n t}$. Allora $\omega_n \in \Lambda_1(S^1)$ e $\deg(\omega_n) = n$ (il suo sollevamento uscente da $0$ \`e $\widetilde{\omega}_n(t) = nt$). Quindi ogni classe di $\Pi_1(S^1)$ \`e rappresentata da un solo $\omega_n$: $\Pi_1(S^1) = \{ [\omega_n] \mid n \in \mathbb{Z} \}$.
\end{observacion}


% ======== NUOVO (pp. 107-108) ========
\section{Applicazioni di $\Pi_1(S^1)$}

\begin{teo}[Teorema fondamentale dell'algebra]
Sia $n \geqslant 1$ e sia $f(z) = a_0 + a_1 z + \dots + a_{n-1} z^{n-1} + z^n$ un polinomio monico a coefficienti $a_i \in \mathbb{C}$. Allora esiste $z_0 \in \mathbb{C}$ tale che $f(z_0) = 0$.
\end{teo}

\begin{intuizione}
Se $f$ non avesse radici, potremmo definire una deformazione continua dal cappio costante al cappio $e^{2\pi i n t}$, che gira $n$ volte attorno alla circonferenza: ma i due cappi hanno gradi diversi ($0$ e $n$), assurdo. L'idea \`e guardare la direzione di $f(z)$ sulla circonferenza di raggio $r$: per $r = 0$ \`e costante, per $r$ grande $f(z) \approx z^n$ gira $n$ volte.
\end{intuizione}

\begin{proof}
Supponiamo per assurdo che $f(z) \neq 0$ per ogni $z \in \mathbb{C}$. Definiamo
\[
G: I \times [0, +\infty[ \longrightarrow S^1, \qquad G(t,r) = \frac{f(r e^{2\pi i t})}{|f(r e^{2\pi i t})|} \cdot \frac{|f(r)|}{f(r)} .
\]
$G$ \`e ben definita e continua perch\'e $f$ non si annulla mai (e ha modulo $1$ come prodotto di due numeri complessi di modulo $1$). Poniamo
\[
H: I \times I \to S^1, \qquad H(t,s) = \begin{cases} G\!\left(t, \dfrac{s}{1-s}\right) & s < 1 \\[2mm] e^{2\pi i n t} & s = 1. \end{cases}
\]
\textbf{$H$ \`e continua.} Sull'insieme $s<1$ \`e composizione di funzioni continue. Rimane la continuit\`a nei punti con $s = 1$: basta provare che $G(t,r) \to e^{2\pi i n t}$ per $r \to +\infty$ \emph{uniformemente in $t \in I$}. Scriviamo, per $|z| \geqslant 1$, $f(z) = z^n (1 + g(z))$ con
\[
g(z) = \frac{a_{n-1}}{z} + \dots + \frac{a_0}{z^n}, \qquad |g(z)| \leqslant \frac{C}{|z|}, \quad C = |a_0| + \dots + |a_{n-1}|.
\]
Per $r$ grande $|g(re^{2\pi i t})| \leqslant C/r < 1/2$ (uniformemente in $t$), quindi
\[
G(t,r) = e^{2\pi i n t} \cdot \frac{1 + g(re^{2\pi i t})}{|1 + g(re^{2\pi i t})|} \cdot \frac{|1 + g(r)|}{1 + g(r)} ,
\]
dove si \`e usato che $f(re^{2\pi i t})/|f(re^{2\pi i t})| = e^{2\pi i n t}(1+g)/|1+g|$ e che $r^n>0$, per cui $|f(r)|/f(r) = |1+g(r)|/(1+g(r))$. I due ultimi fattori tendono a $1$ per $r \to \infty$ uniformemente in $t$, poich\'e $g \to 0$ uniformemente. Quindi $H$ \`e continua e dunque \`e un'omotopia.

\textbf{Estremi.}
\begin{itemize}
    \item $H(t,0) = G(t,0) = \dfrac{f(0)}{|f(0)|} \cdot \dfrac{|f(0)|}{f(0)} = 1 = c_1(t)$ per ogni $t$.
    \item $H(t,1) = e^{2\pi i n t} = \omega_n(t)$.
    \item $H(0,s) = G(0,r) = \dfrac{f(r)}{|f(r)|} \cdot \dfrac{|f(r)|}{f(r)} = 1$ e, analogamente, $H(1,s) = G(1,r) = 1$ (perch\'e $e^{2\pi i \cdot 1} = 1$).
\end{itemize}
Dunque $H: c_1 \simeq \omega_n$ rel $\partial I$. Per la proposizione sul grado, $0 = \deg(c_1) = \deg(\omega_n) = n$, il che contraddice $n \geqslant 1$. Pertanto esiste $z_0 \in \mathbb{C}$ con $f(z_0) = 0$\hfill$\square$
\end{proof}


% ======== NUOVO (pp. 109-112) ========
\section{Rivestimenti definiti da gruppi propriamente discontinui}

Sia $\widetilde{X}$ uno spazio topologico e $G$ un gruppo di omeomorfismi di $\widetilde{X}$ (con l'operazione di composizione). Definiamo su $\widetilde{X}$ la relazione
\[
\widetilde{x} \sim_G \widetilde{y} \iff \exists f \in G : \ \widetilde{y} = f(\widetilde{x}).
\]
\`E una relazione di equivalenza perch\'e $G$ \`e un gruppo:
\begin{itemize}
    \item \emph{riflessiva}, perch\'e $\mathrm{Id}_{\widetilde{X}} \in G$ e $\widetilde{x} = \mathrm{Id}(\widetilde{x})$;
    \item \emph{simmetrica}, perch\'e se $f \in G$ e $\widetilde{y} = f(\widetilde{x})$, allora $f^{-1} \in G$ e $\widetilde{x} = f^{-1}(\widetilde{y})$;
    \item \emph{transitiva}, perch\'e se $f, g \in G$, $\widetilde{y} = f(\widetilde{x})$, $\widetilde{z} = g(\widetilde{y})$, allora $g \circ f \in G$ e $\widetilde{z} = (g \circ f)(\widetilde{x})$.
\end{itemize}
Poniamo $X = \widetilde{X}/G = \widetilde{X}/\!\sim_G$, con la topologia quoziente, e indichiamo con
\[
\pi: \widetilde{X} \to X = \widetilde{X}/G, \qquad \widetilde{x} \mapsto [\widetilde{x}]_{\sim_G}
\]
la proiezione. Quindi $A \subseteq X$ \`e aperto $\iff$ $\pi^{-1}(A)$ \`e aperto in $\widetilde{X}$.
Dato $\widetilde{x} \in \widetilde{X}$, la classe $\pi^{-1}(\pi(\widetilde{x})) = \{ f(\widetilde{x}) \mid f \in G \} = G\widetilde{x}$ si chiama \textbf{orbita} di $\widetilde{x}$ rispetto a $G$, e l'applicazione
\[
\varphi: G \times \widetilde{X} \to \widetilde{X}, \qquad (f, \widetilde{x}) \mapsto f(\widetilde{x})
\]
si chiama \textbf{azione} di $G$ su $\widetilde{X}$.

\begin{prop}
$\pi: \widetilde{X} \to \widetilde{X}/G$ \`e un'applicazione aperta.
\end{prop}
\begin{proof}
Sia $\widetilde{A}$ aperto in $\widetilde{X}$. Allora
\[
\pi^{-1}(\pi(\widetilde{A})) = \{ \widetilde{y} \mid \exists \widetilde{x} \in \widetilde{A},\ \pi(\widetilde{y}) = \pi(\widetilde{x}) \} = \{ \widetilde{y} \mid \exists \widetilde{x} \in \widetilde{A},\ \exists f \in G,\ \widetilde{y} = f(\widetilde{x}) \} = \bigcup_{f \in G} f(\widetilde{A}),
\]
che \`e aperto perch\'e ogni $f \in G$ \`e un omeomorfismo, e quindi $f(\widetilde{A})$ \`e aperto. Per definizione di topologia quoziente, $\pi(\widetilde{A})$ \`e aperto\hfill$\square$
\end{proof}

\begin{observacion}
Gli esempi visti si interpretano tutti come quozienti per un gruppo di omeomorfismi:
\begin{enumerate}
    \item $p: \mathbb{R} \to S^1 = \mathbb{R}/\mathbb{Z}$, $t \mapsto e^{2\pi i t}$. Infatti $p(t) = p(t') \iff t - t' \in \mathbb{Z} \iff \exists h \in \mathbb{Z}: t = t' + h = t_h(t')$, dove $t_h: \mathbb{R} \to \mathbb{R}$, $x \mapsto x + h$ \`e un omeomorfismo (traslazione), e $t_h \circ t_{h'} = t_{h+h'}$. Dunque $S^1 = \mathbb{R}/G$ con $G = \{ t_h \mid h \in \mathbb{Z} \} \cong \mathbb{Z}$.
    \item $p: \mathbb{R}^n \to T^n = \mathbb{R}^n/\mathbb{Z}^n$: $p(x) = p(y) \iff \exists z \in \mathbb{Z}^n : y = x + z = t_z(x)$, con $t_z(x) = x + z$ omeomorfismo di $\mathbb{R}^n$. Quindi $T^n = \mathbb{R}^n/G$ con $G = \{ t_z \mid z \in \mathbb{Z}^n \} \cong \mathbb{Z}^n$.
    \item Per ogni $n \geqslant 2$ (e anche per $n = 1$), $p: S^n \to \mathbb{R}P^n = S^n/\{\pm \mathrm{Id}\}$ con $G = \{ \mathrm{Id}, -\mathrm{Id} \} \cong \mathbb{Z}_2$.
\end{enumerate}
\end{observacion}

\begin{definicion}[Azione propriamente discontinua]
Sia $\widetilde{X}$ uno spazio topologico e $G$ un gruppo di omeomorfismi di $\widetilde{X}$. Si dice che l'azione di $G$ su $\widetilde{X}$ \`e \textbf{propriamente discontinua} se
\[
\forall \widetilde{x} \in \widetilde{X} \ \ \exists \widetilde{U} \text{ intorno aperto di } \widetilde{x} \ : \ f(\widetilde{U}) \cap \widetilde{U} = \varnothing \quad \forall f \in G,\ f \neq \mathrm{Id}_{\widetilde{X}}.
\]
\end{definicion}

\begin{intuizione}
Ogni punto ha un intorno che i ``movimenti'' non identici di $G$ spostano completamente fuori da s\'e stesso (nessun $f \neq \mathrm{Id}$ ha punti fissi, nemmeno ``quasi'' fissi). Es.: la traslazione $t \mapsto t+1$ di $\mathbb{R}$ sposta l'intervallo $]x-\tfrac13, x+\tfrac13[$ fuori da s\'e stesso; la rotazione di un angolo $\theta \neq 0$ attorno all'origine \emph{non} \`e propriamente discontinua (l'origine \`e fissa).
\end{intuizione}

\begin{prop}
L'azione di $G$ \`e propriamente discontinua se e solo se per ogni $\widetilde{x} \in \widetilde{X}$ esiste un intorno aperto $\widetilde{U}$ di $\widetilde{x}$ tale che
\[
f(\widetilde{U}) \cap g(\widetilde{U}) = \varnothing \qquad \forall f \neq g \in G.
\]
\end{prop}
\begin{proof}
($\Leftarrow$) \`E banale (si prende $g = \mathrm{Id}$).
($\Rightarrow$) Sia $\widetilde{U}$ un intorno aperto di $\widetilde{x}$ con $f(\widetilde{U}) \cap \widetilde{U} = \varnothing$ per ogni $f \neq \mathrm{Id}$. Supponiamo per assurdo che esistano $f \neq g$ in $G$ con $f(\widetilde{U}) \cap g(\widetilde{U}) \neq \varnothing$. Applicando l'omeomorfismo $g^{-1} \in G$:
\[
\varnothing \neq g^{-1}\big( f(\widetilde{U}) \cap g(\widetilde{U}) \big) = (g^{-1} \circ f)(\widetilde{U}) \cap \widetilde{U}.
\]
Poich\'e $g^{-1} \circ f \in G$, ci\`o forza $g^{-1} \circ f = \mathrm{Id}$, cio\`e $f = g$: assurdo\hfill$\square$
\end{proof}

\begin{teo}
Sia $\widetilde{X}$ uno spazio topologico connesso per archi e $G$ un gruppo di omeomorfismi di $\widetilde{X}$ ad azione propriamente discontinua. Allora $\pi: \widetilde{X} \to X = \widetilde{X}/G$ \`e un'applicazione di rivestimento.
\end{teo}
\begin{proof}
$\pi$ \`e continua e suriettiva (proiezione al quoziente) e $\widetilde{X}$ \`e connesso per archi. Sia $x_0 = G\widetilde{x}_0 \in X$ con $\widetilde{x}_0 \in \widetilde{X}$. Per ipotesi esiste un intorno aperto $\widetilde{U}$ di $\widetilde{x}_0$ con $f(\widetilde{U}) \cap g(\widetilde{U}) = \varnothing$ per ogni $f \neq g$ in $G$. Poniamo $U = \pi(\widetilde{U})$: \`e aperto perch\'e $\pi$ \`e aperta, e contiene $x_0 = \pi(\widetilde{x}_0)$.

\textbf{$(\pi|_{\widetilde{U}})_\#: \widetilde{U} \to U$ \`e un omeomorfismo.} \`E continua, suriettiva per costruzione, aperta (perch\'e $\pi$ \`e aperta e $\widetilde{U}$ \`e aperto), ed \`e \emph{iniettiva}: se $\widetilde{y}, \widetilde{y}' \in \widetilde{U}$ e $\pi(\widetilde{y}) = \pi(\widetilde{y}')$, allora esiste $f \in G$ con $\widetilde{y}' = f(\widetilde{y})$, cio\`e $\widetilde{y}' \in \widetilde{U} \cap f(\widetilde{U})$, e quindi $f = \mathrm{Id}$ (altrimenti l'intersezione sarebbe vuota), cio\`e $\widetilde{y}' = \widetilde{y}$.

\textbf{La controimmagine di $U$ \`e unione disgiunta di fogli.} Poich\'e $\pi \circ g = \pi$ per ogni $g \in G$, si ha
\[
\pi^{-1}(U) = \pi^{-1}(\pi(\widetilde{U})) = \bigcup_{g \in G} g(\widetilde{U}),
\]
e l'unione \`e disgiunta per la scelta di $\widetilde{U}$. Ogni $g(\widetilde{U})$ \`e aperto in $\widetilde{X}$ e
\[
(\pi|_{g(\widetilde{U})})_\# = (\pi|_{\widetilde{U}})_\# \circ g^{-1}
\]
\`e composizione di omeomorfismi, dunque un omeomorfismo $g(\widetilde{U}) \to U$ per ogni $g \in G$. Pertanto $U$ \`e un intorno ben coperto di $x_0$ e $\pi$ \`e un'applicazione di rivestimento\hfill$\square$
\end{proof}

\begin{observacion}
$\mathbb{R} \to S^1 = \mathbb{R}/\mathbb{Z}$, $\mathbb{R}^n \to T^n = \mathbb{R}^n/\mathbb{Z}^n$ e $S^n \to \mathbb{R}P^n = S^n/\{\pm \mathrm{Id}\}$ sono tutti esempi di quozienti per gruppi di omeomorfismi ad azione propriamente discontinua (in particolare, questo d\`a una dimostrazione unificata del fatto che sono rivestimenti). Nel teorema precedente i fogli sopra $U$ sono in corrispondenza biunivoca con gli elementi di $G$.
\end{observacion}

\begin{figure}[H]
    \centering
    \begin{tikzpicture}[scale=1.1]
        \draw[thick] (-0.3,0) -- (7.3,0);
        \foreach \x in {0,1,2,3,4,5,6,7} {
            \filldraw (\x,0) circle (1.2pt);
        }
        \foreach \x in {0,1,2,3,4,5,6,7} {
            \draw[red, very thick] (\x-0.28,0) -- (\x+0.28,0);
        }
        \node[below] at (3,-0.1) {$\widetilde{U} + h = t_h(\widetilde{U})$, $h \in \mathbb{Z}$};
        \draw[->, >=stealth, thick, Purple] (3.5,-0.9) -- (3.5,-1.7) node[midway, right] {$\pi$};
        \draw[thick] (3.5,-2.4) circle (0.55);
        \draw[red, very thick] ({3.5+0.55*cos(-30)}, {-2.4+0.55*sin(-30)}) arc (-30:30:0.55);
        \node[right] at (4.1,-2.4) {$U$};
    \end{tikzpicture}
    \caption{I traslati $t_h(\widetilde{U})$ (in rosso) di un piccolo intorno sono disgiunti e vengono tutti ``avvolti'' sullo stesso intorno $U \subset S^1$.}
\end{figure}

\section{Gruppo fondamentale di un quoziente}

Sia $(\widetilde{X}, \pi)$ un rivestimento di $X = \widetilde{X}/G$ (con $G$ ad azione propriamente discontinua), $x_0 \in X$ e $\widetilde{x}_0 \in \pi^{-1}(x_0)$. Abbiamo la seguente corrispondenza biunivoca (teorema centrale sui rivestimenti, con $\pi^{-1}(x_0) = G\widetilde{x}_0$):
\[
\Pi_1(X, x_0) \big/ \pi_*\big(\Pi_1(\widetilde{X}, \widetilde{x}_0)\big) \ \longleftrightarrow\ \pi^{-1}(x_0) = G\widetilde{x}_0 \ \longleftrightarrow\ G, \qquad [\sigma] \mapsto \widetilde{\sigma}(1) = g(\widetilde{x}_0) \mapsto g .
\]
(L'applicazione $g \mapsto g(\widetilde{x}_0)$ da $G$ all'orbita \`e iniettiva: se $g(\widetilde{x}_0) = g'(\widetilde{x}_0)$ allora $g^{-1}g'$ ha un punto fisso e, per la propriet\`a di discontinuit\`a, $g^{-1}g' = \mathrm{Id}$.)

\begin{teo}
Se $\widetilde{X}$ \`e il rivestimento \textbf{universale} di $X = \widetilde{X}/G$, allora l'applicazione
\[
F: \Pi_1(X, x_0) \longrightarrow G, \qquad [\sigma] \longmapsto g, \ \text{ dove } \ \widetilde{\sigma}(1) = g(\widetilde{x}_0),
\]
\`e un \textbf{isomorfismo di gruppi}.
\end{teo}
\begin{proof}
Per il corollario sui rivestimenti universali $F$ \`e una bigezione. Basta provare che \`e un omomorfismo: per ogni $[\sigma], [\tau] \in \Pi_1(X, x_0)$, $F([\sigma][\tau]) = F([\sigma]) \circ F([\tau])$.

Sia $F([\sigma][\tau]) = F([\sigma * \tau]) = h \in G$, cio\`e $h(\widetilde{x}_0) = \widetilde{\sigma * \tau}(1)$. Siano $f, g \in G$ tali che $f(\widetilde{x}_0) = \widetilde{\sigma}(1)$ e $g(\widetilde{x}_0) = \widetilde{\tau}(1)$, dove $\widetilde{\sigma}, \widetilde{\tau}$ sono i sollevamenti di $\sigma, \tau$ uscenti da $\widetilde{x}_0$. Poniamo $\widetilde{\tau}_1 = f \circ \widetilde{\tau}$. Allora:
\begin{enumerate}
    \item $\widetilde{\tau}_1$ \`e un sollevamento di $\tau$ (perch\'e $\pi \circ f = \pi$, quindi $\pi \circ \widetilde{\tau}_1 = \pi \circ \widetilde{\tau} = \tau$) uscente da $f(\widetilde{x}_0) = \widetilde{\sigma}(1)$;
    \item perci\`o $\widetilde{\sigma} * \widetilde{\tau}_1$ \`e ben definito, parte da $\widetilde{x}_0$ e solleva $\sigma * \tau$; per l'unicit\`a del sollevamento, $\widetilde{\sigma} * \widetilde{\tau}_1 = \widetilde{\sigma * \tau}$.
\end{enumerate}
Quindi $h(\widetilde{x}_0) = (\widetilde{\sigma} * \widetilde{\tau}_1)(1) = \widetilde{\tau}_1(1) = f(\widetilde{\tau}(1)) = f(g(\widetilde{x}_0))$. Per l'iniettivit\`a di $g \mapsto g(\widetilde{x}_0)$ (vedi sopra), $h = f \circ g$, cio\`e $F([\sigma][\tau]) = F([\sigma]) \circ F([\tau])$\hfill$\square$
\end{proof}

\begin{coro}
$\Pi_1(T^n) \cong \mathbb{Z}^n$ e, per $n \geqslant 2$, $\Pi_1(\mathbb{R}P^n) \cong \mathbb{Z}_2$.
\end{coro}
\begin{proof}
$\mathbb{R}^n \to T^n = \mathbb{R}^n/\mathbb{Z}^n$ e $S^n \to \mathbb{R}P^n = S^n/\{\pm\mathrm{Id}\}$ ($n \geqslant 2$) sono rivestimenti universali dati da gruppi propriamente discontinui $G \cong \mathbb{Z}^n$ e $G \cong \mathbb{Z}_2$ rispettivamente\hfill$\square$
\end{proof}


% ======== NUOVO (pp. 113-114) ========
\section{Algoritmo per il calcolo di $\Pi_1$}

\textbf{Problema.} Dato uno spazio topologico $X$, determinare $\Pi_1(X)$.

\begin{figure}[H]
    \centering
    \begin{tikzpicture}[node distance=1.0cm and 1.2cm,
        box/.style={draw, rounded corners, align=center, minimum height=0.9cm, text width=4.2cm, fill=gray!8},
        dec/.style={draw, diamond, aspect=2.3, align=center, inner sep=1pt, fill=blue!6},
        fin/.style={draw, rounded corners, fill=green!12, minimum height=0.7cm},
        >=Stealth]
        \node[dec] (d1) {$X$ \`e connesso\\ per archi?};
        \node[box, right=2.6cm of d1] (no1) {\textbf{NO:} si applicano i passi successivi a ciascuna componente connessa per archi di $X$};
        \node[dec, below=1.2cm of d1] (d2) {$\Pi_1(X) = 0$?};
        \node[fin, right=2.6cm of d2] (f1) {FINE};
        \node[box, below=1.2cm of d2] (p3) {Trovare $Y$ con $\Pi_1(Y)$ noto, topologicamente o omotopicamente equivalente a $X$: allora $\Pi_1(X) = \Pi_1(Y)$};
        \node[fin, below=0.8cm of p3] (f2) {FINE};
        \draw[->] (d1) -- node[above] {NO} (no1);
        \draw[->] (d1) -- node[left] {S\`I} (d2);
        \draw[->] (d2) -- node[above] {S\`I} (f1);
        \draw[->] (d2) -- node[left] {NO} (p3);
        \draw[->] (p3) -- (f2);
    \end{tikzpicture}
    \caption{Schema dell'algoritmo per calcolare il gruppo fondamentale.}
\end{figure}

\begin{ejemplo}
Sia $X = \mathbb{R}^2 \setminus \{ x \in \mathbb{R}^2 \mid \|x\| = 1 \ \lor\ \|x\| = 2 \}$ (il piano privato delle due circonferenze centrate nell'origine di raggi $1$ e $2$). Calcoliamo $\Pi_1(X)$.

\begin{center}
\begin{tikzpicture}[scale=1.1]
    \fill[red!10] (0,0) circle (1);
    \fill[blue!10] (0,0) circle (2); \fill[red!10] (0,0) circle (1);
    \fill[green!10, even odd rule] (0,0) circle (3) (0,0) circle (2);
    \draw[thick, dashed] (0,0) circle (1);
    \draw[thick, dashed] (0,0) circle (2);
    \node at (0,0) {$A$};
    \node at (1.5,0) {$B$};
    \node at (2.5,0.0) {$C$};
\end{tikzpicture}
\end{center}

$X$ \textbf{non} \`e connesso per archi: ha tre componenti connesse per archi:
\[
A = \{\|x\| < 1\}, \qquad B = \{1 < \|x\| < 2\}, \qquad C = \{\|x\| > 2\}.
\]
\begin{itemize}
    \item $\Pi_1(A) = 0$, perch\'e $A$ \`e convesso (un disco aperto \`e contraibile).
    \item $B$ ha come retratto di deformazione la circonferenza $S^1 = \{\|x\| = 3/2\}$ (proiezione radiale), quindi $\Pi_1(B) = \Pi_1(S^1) \cong \mathbb{Z}$. In modo equivalente, $B \cong S^1 \times\, ]1,2[$ e $\Pi_1(S^1 \times\, ]1,2[) \cong \Pi_1(S^1) \times \Pi_1(]1,2[) = \mathbb{Z} \times 0 = \mathbb{Z}$.
    \item $C \cong S^1 \times\, ]0, +\infty[$ ha $S^1$ come retratto di deformazione, quindi $\Pi_1(C) \cong \mathbb{Z}$.
\end{itemize}
Dunque $\Pi_1(A) = 0$, $\Pi_1(B) \cong \Pi_1(C) \cong \mathbb{Z}$ (ciascuna componente va trattata separatamente, con il suo punto base).
\end{ejemplo}

\begin{ejemplo}
Sia $X = S^2 \setminus \gamma$, con $\gamma$ una circonferenza (ad esempio l'equatore). $X$ non \`e connesso per archi: ha due componenti connesse per archi $A$ e $B$ (le due calotte aperte). Ciascuna \`e omeomorfa a un disco aperto $\mathring{B}^2$ (per proiezione), quindi contraibile: $\Pi_1(A) = \Pi_1(B) = 0$.
\end{ejemplo}

\begin{ejemplo}
Sia $X = \mathbb{R}^3 \setminus r$, con $r$ una retta. $X$ \`e connesso per archi e $\Pi_1(X) \neq 0$: un cappio che gira attorno alla retta non si pu\`o contrarre al cappio costante (la retta \`e infinita, non c'\`e modo di ``sfilare'' il cappio). Per calcolarlo:
\begin{itemize}
    \item \textbf{Metodo 1.} Un cilindro infinito $\mathscr{C} = S^1 \times \mathbb{R}$ avente $r$ come asse (privato dell'asse) \`e un retratto di deformazione di $X$ (si spinge ogni punto radialmente sul cilindro), quindi $\Pi_1(X) = \Pi_1(S^1 \times \mathbb{R}) = \Pi_1(S^1) \times \Pi_1(\mathbb{R}) = \mathbb{Z}$.
    \item \textbf{Metodo 2.} Si pu\`o prendere un piano ortogonale a $r$ che interseca $r$ in un punto $p$: $\mathbb{R}^3 \setminus r$ si retrae per deformazione su $\mathbb{R}^2 \setminus \{p\}$. Quindi $\Pi_1(X) = \Pi_1(\mathbb{R}^2 \setminus \{p\}) = \Pi_1(S^1) = \mathbb{Z}$.
\end{itemize}
\end{ejemplo}

\begin{figure}[H]
    \centering
    \begin{tikzpicture}[scale=1.2]
        \draw[thick] (0,-1.6) -- (0,1.6) node[above] {$r$};
        \draw[blue, thick] (0,0.9) ellipse (0.9 and 0.3);
        \draw[blue, thick] (0,-0.9) ellipse (0.9 and 0.3);
        \draw[blue, thick] (-0.9,0.9) -- (-0.9,-0.9);
        \draw[blue, thick] (0.9,0.9) -- (0.9,-0.9);
        \node[blue] at (1.6,0) {cilindro $S^1\times\mathbb{R}$};
        \draw[red, thick, ->, >=stealth] (0,0) ellipse (0.5 and 0.18);
        \node[red, right] at (0.5,0.2) {cappio};
    \end{tikzpicture}
    \caption{$\mathbb{R}^3 \setminus r$ si retrae per deformazione su un cilindro che ha $r$ come asse; il cappio attorno a $r$ genera $\Pi_1 = \mathbb{Z}$.}
\end{figure}

\begin{ejemplo}
Sia $X = \mathbb{R}^2 \setminus (S^1 \cup \{p\})$, con $S^1$ la circonferenza unitaria e $p$ un punto. Ci sono due configurazioni possibili.

\begin{center}
\begin{tikzpicture}[scale=1.0]
    % p interno
    \begin{scope}
        \draw[thick, dashed] (0,0) circle (1);
        \fill[red!10] (0,0) circle (1);
        \filldraw (0.35,0.2) circle (1.5pt) node[right] {$p$};
        \node at (-0.5,-0.4) {$A$}; \node at (1.7,0.8) {$B$};
        \node[below] at (0,-1.3) {$p$ interno a $S^1$};
    \end{scope}
    % p esterno
    \begin{scope}[shift={(5.5,0)}]
        \draw[thick, dashed] (0,0) circle (1);
        \fill[red!10] (0,0) circle (1);
        \filldraw (2.0,0.2) circle (1.5pt) node[right] {$p$};
        \node at (0,0) {$A$}; \node at (1.2,1.2) {$B$};
        \node[below] at (0.8,-1.3) {$p$ esterno a $S^1$};
    \end{scope}
\end{tikzpicture}
\end{center}

\begin{enumerate}
    \item \textbf{$p$ interno alla circonferenza.} $X$ non \`e connesso per archi e ha due componenti $A$ (il disco aperto privato di $p$) e $B$ (la parte esterna al disco). $\Pi_1(A) \neq 0$ (un cappio attorno a $p$ non si contrae); $A$ ha $S^1$ come retratto di deformazione, quindi $\Pi_1(A) = \Pi_1(S^1) = \mathbb{Z}$; e $\Pi_1(B) = \mathbb{Z}$ (gi\`a visto).
    \item \textbf{$p$ esterno alla circonferenza.} Ancora due componenti: $A$ \`e il disco aperto, quindi $\Pi_1(A) = 0$; $B$ \`e l'esterno del disco privato di $p$ e $\Pi_1(B) \neq 0$ (un cappio attorno al disco, o attorno a $p$, non si contrae). Osserviamo che $S^1 \vee S^1$ \`e un retratto di deformazione di $B$ (il piano privato di un disco chiuso e di un punto ha lo stesso tipo di omotopia del piano privato di due punti), quindi
    \[
    \Pi_1(B) = \Pi_1(S^1 \vee S^1) = F_2 .
    \]
\end{enumerate}
\end{ejemplo}



    % ======== NUOVO: CAPITOLO 3 (pp. 116-120) ========
\chapter{I gruppi di omologia simpliciale}

\section{Simplessi}

\begin{definicion}[Insieme convesso]
Sia $V$ uno spazio vettoriale su $\mathbb{R}$. Un sottoinsieme $C \subseteq V$ si dice \textbf{convesso} se per ogni $x, y \in C$ il segmento
\[
\overline{x,y} = \{ (1-t)x + ty \mid t \in [0,1] \}
\]
\`e contenuto in $C$.
\end{definicion}

\begin{definicion}[Vettori c-indipendenti]
Dato $k \in \mathbb{N}$, i vettori $v_0, \dots, v_k \in V$ si dicono \textbf{c-indipendenti} (``convessamente indipendenti'', detti anche \emph{affinemente indipendenti}) se i $k$ vettori $v_1 - v_0, \dots, v_k - v_0$ sono linearmente indipendenti.
\end{definicion}

\begin{observacion}
La definizione non dipende dal vettore $v_0$ scelto come ``iniziale'': se si sceglie come origine un altro $v_j$, i vettori $v_i - v_j$ ($i \neq j$) sono combinazioni lineari invertibili dei $v_i - v_0$ (infatti $v_i - v_j = (v_i - v_0) - (v_j - v_0)$), quindi sono indipendenti se e solo se lo sono questi.
\end{observacion}

\begin{ejemplo}
\begin{enumerate}
    \item $k = 0$: un singolo vettore $v_0$ \`e (banalmente) c-indipendente, perch\'e la famiglia dei vettori $v_i - v_0$ con $i \geqslant 1$ \`e vuota, e la famiglia vuota \`e linearmente indipendente.
    \item $k = 1$: $v_0, v_1$ sono c-indipendenti $\iff$ $v_1 - v_0$ \`e linearmente indipendente $\iff$ $v_1 - v_0 \neq \bar{0}$ $\iff$ $v_1 \neq v_0$.
    \item $k = 2$: $v_0, v_1, v_2$ sono c-indipendenti $\iff$ $v_1 - v_0, v_2 - v_0$ sono linearmente indipendenti, cio\`e \textbf{non paralleli} $\iff$ $v_0, v_1, v_2$ \textbf{non allineati}.
    \item $k = 3$: $v_0, v_1, v_2, v_3$ sono c-indipendenti $\iff$ $v_1 - v_0, v_2 - v_0, v_3 - v_0$ sono linearmente indipendenti, cio\`e \textbf{non complanari} $\iff$ $v_0, v_1, v_2, v_3$ non sono complanari.
\end{enumerate}
\end{ejemplo}

\begin{observacion}[Nota di correzione]
Negli appunti, al punto 1, si afferma che per $k = 0$ un vettore ``non \`e mai c-indipendente perch\'e $v_0 - v_0 = \bar{0}$''. Questo non \`e coerente con la definizione (per $k=0$ non ci sono vettori $v_i - v_0$ con $i \geqslant 1$, non compare $v_0 - v_0$) e con il seguito: i $0$-simplessi sono proprio i punti. Qui si usa la convenzione standard: ogni punto \`e c-indipendente.
\end{observacion}

\begin{figure}[H]
    \centering
    \begin{tikzpicture}[scale=1.0]
        % non c-indipendenti (allineati) e c-indipendenti
        \begin{scope}
            \draw[->] (0,0) -- (2.2,0); \draw[->] (0,0) -- (0,1.8);
            \filldraw (0.3,0.3) circle (1.5pt) node[below] {$v_0$};
            \filldraw (1.0,0.7) circle (1.5pt) node[above] {$v_1$};
            \filldraw (1.7,1.1) circle (1.5pt) node[above] {$v_2$};
            \draw[dashed] (0.3,0.3) -- (1.7,1.1);
            \node[below] at (1.1,-0.4) {\small non c-indipendenti};
        \end{scope}
        \begin{scope}[shift={(4.2,0)}]
            \draw[->] (0,0) -- (2.2,0); \draw[->] (0,0) -- (0,1.8);
            \filldraw (0.4,0.3) circle (1.5pt) node[below left] {$v_0$};
            \filldraw (1.9,0.5) circle (1.5pt) node[right] {$v_1$};
            \filldraw (0.9,1.5) circle (1.5pt) node[above] {$v_2$};
            \draw[blue, ->, >=stealth] (0.4,0.3) -- (1.9,0.5) node[midway, below] {\scriptsize $v_1-v_0$};
            \draw[red, ->, >=stealth] (0.4,0.3) -- (0.9,1.5) node[midway, left] {\scriptsize $v_2-v_0$};
            \node[below] at (1.1,-0.4) {\small c-indipendenti};
        \end{scope}
    \end{tikzpicture}
    \caption{Tre punti sono c-indipendenti se e solo se non sono allineati, cio\`e se $v_1-v_0$ e $v_2-v_0$ non sono paralleli.}
\end{figure}

\begin{definicion}[$k$-piano generato]
Dati $k+1$ vettori c-indipendenti $v_0, \dots, v_k$, l'insieme
\[
\Pi = \left\{ v \in V \;\Big|\; v = \sum_{i=0}^k t_i v_i,\ t_i \in \mathbb{R},\ \sum_{i=0}^k t_i = 1 \right\}
\]
si chiama \textbf{$k$-piano generato} da $v_0, \dots, v_k$.
\end{definicion}

\begin{observacion}
Se $v \in \Pi$, con $v = \sum_{i=0}^k t_i v_i$ e $\sum t_i = 1$, allora
\[
v = t_0 v_0 + \sum_{i=1}^k t_i v_i = \Big(1 - \sum_{i=1}^k t_i\Big) v_0 + \sum_{i=1}^k t_i v_i = v_0 + \sum_{i=1}^k t_i (v_i - v_0).
\]
Quindi $\Pi$ \`e il piano $k$-dimensionale di $V$ passante per $v_0$ e parallelo ai vettori linearmente indipendenti $v_1 - v_0, \dots, v_k - v_0$. In generale $\Pi$ \emph{non} \`e un sottospazio vettoriale, ma un \textbf{sottospazio affine}. Le \emph{trasformazioni affini} di $V$ (le applicazioni $V \to V$ che conservano i sottospazi affini) sono tutte e sole le funzioni $f = \varphi \circ t_v$, con $t_v$ la traslazione $x \mapsto x + v$ e $\varphi \in GL(V)$.
\end{observacion}

\begin{prop}
Sia $V$ di dimensione $n$ e sia $\{e_1, \dots, e_n\}$ una base di $V$. Dati $k+1$ vettori c-indipendenti $v_0, \dots, v_k$ (con $k \leqslant n$), esiste una trasformazione affine $f$ di $V$ tale che $f(\Pi) = \mathcal{L}(e_1, \dots, e_k)$ (il sottospazio generato).
\end{prop}
\begin{proof}
Consideriamo la traslazione $t_{-v_0}$: per l'osservazione precedente, $t_{-v_0}(\Pi) = \mathcal{L}(v_1 - v_0, \dots, v_k - v_0)$, e $v_1 - v_0, \dots, v_k - v_0$ sono linearmente indipendenti. Per il teorema di completamento della base esistono $n-k$ vettori $e'_{k+1}, \dots, e'_n$ di $V$ tali che $\{ v_1 - v_0, \dots, v_k - v_0, e'_{k+1}, \dots, e'_n \}$ sia una base di $V$. Sia $\varphi: V \to V$ lineare tale che $\varphi(v_i - v_0) = e_i$ per $i = 1, \dots, k$ e $\varphi(e'_j) = e_j$ per $j = k+1, \dots, n$. Allora $\varphi \in GL(V)$ (manda una base in una base). Quindi $f = \varphi \circ t_{-v_0}$ \`e una trasformazione affine e
\[
f(\Pi) = \varphi\big( t_{-v_0}(\Pi) \big) = \varphi\big( \mathcal{L}(v_1 - v_0, \dots, v_k - v_0) \big) = \mathcal{L}(e_1, \dots, e_k) . \qquad \square
\]
\end{proof}

\begin{definicion}[Convesso generato]
Sia $V$ uno spazio vettoriale su $\mathbb{R}$ e $X \subseteq V$. Si chiama \textbf{convesso generato} da $X$ il pi\`u piccolo sottoinsieme convesso di $V$ contenente $X$. Dunque $C$ \`e il convesso generato da $X$ se e solo se:
\begin{enumerate}
    \item $C$ \`e convesso;
    \item $X \subseteq C$;
    \item per ogni convesso $C'$ di $V$ con $X \subseteq C'$ si ha $C \subseteq C'$.
\end{enumerate}
\end{definicion}

\begin{observacion}
$C = \bigcap_{C' \text{ convesso},\ X \subseteq C'} C'$ (un'intersezione di convessi \`e convessa).
\end{observacion}

\begin{teo}
Sia $V$ uno spazio vettoriale su $\mathbb{R}$, $k \in \mathbb{N}$ e $v_0, \dots, v_k$ vettori c-indipendenti. Allora il convesso generato da $\{v_0, \dots, v_k\}$ \`e
\[
C = \left\{ v \in V \;\Big|\; v = \sum_{i=0}^k t_i v_i,\ t_i \geqslant 0\ \forall i,\ \sum_{i=0}^k t_i = 1 \right\}.
\]
Inoltre i coefficienti $t_i$ sono univocamente determinati da $v$ (sono le \emph{coordinate baricentriche} di $v$).
\end{teo}

\begin{proof}
Dobbiamo provare che $C$ \`e convesso, che $v_i \in C$ per ogni $i$, e che ogni convesso $C'$ contenente tutti i $v_i$ contiene $C$.

\medskip
\textbf{1. $C$ \`e convesso.} Siano $x, y \in C$, $x = \sum a_i v_i$, $y = \sum b_i v_i$ con $a_i, b_i \geqslant 0$ e $\sum a_i = \sum b_i = 1$. Per ogni $t \in [0,1]$:
\[
(1-t)x + ty = \sum_{i=0}^k \big( (1-t)a_i + t b_i \big) v_i,
\]
con coefficienti $(1-t)a_i + tb_i \geqslant 0$ e $\sum_i \big((1-t)a_i + tb_i\big) = (1-t) + t = 1$. Quindi $(1-t)x + ty \in C$, cio\`e $\overline{x,y} \subseteq C$.

\textbf{2. $v_i \in C$.} Per ogni $i$, $v_i = \sum_j \delta_{ij} v_j$ (delta di Kronecker), con $\delta_{ij} \geqslant 0$ e $\sum_j \delta_{ij} = 1$.

\textbf{3. Minimalit\`a.} Sia $C'$ convesso con $v_0, \dots, v_k \in C'$. Dimostriamo che ogni $v = \sum t_i v_i$ con $t_i \geqslant 0$, $\sum t_i = 1$, appartiene a $C'$, per induzione sul numero $p$ dei coefficienti non nulli.
\begin{itemize}
    \item Se $p = 1$, $v = v_i \in C'$ per qualche $i$.
    \item Supponiamo la tesi vera fino a $p$ e proviamola per $p+1$. Senza perdita di generalit\`a i $p+1$ coefficienti non nulli sono i primi: $v = \sum_{i=0}^{p} t_i v_i$, $t_i > 0$, $\sum_{i=0}^p t_i = 1$. Poich\'e $p+1 \geqslant 2$ coefficienti sono non nulli, $t_p < 1$, e
    \[
    v = \sum_{i=0}^{p-1} t_i v_i + t_p v_p = (1 - t_p) \underbrace{\sum_{i=0}^{p-1} \frac{t_i}{1 - t_p}\, v_i}_{=: w} + t_p v_p .
    \]
    I coefficienti $t_i/(1-t_p)$ sono $> 0$ e hanno somma $\frac{1-t_p}{1-t_p} = 1$, quindi $w$ ha $p$ coefficienti non nulli e, per ipotesi induttiva, $w \in C'$. Allora $v = (1-t_p)w + t_p v_p \in \overline{w, v_p} \subseteq C'$ perch\'e $C'$ \`e convesso.
\end{itemize}
Dunque $C \subseteq C'$ e $C$ \`e il convesso generato.

\textbf{Unicit\`a dei coefficienti.} Se $v = \sum a_i v_i = \sum b_i v_i$ con $\sum a_i = \sum b_i = 1$, allora, poich\'e $\sum_i (a_i - b_i) = 0$,
\[
\bar{0} = \sum_{i=0}^k (a_i - b_i) v_i = \sum_{i=1}^k (a_i - b_i)(v_i - v_0) .
\]
I vettori $v_i - v_0$ sono linearmente indipendenti, quindi $a_i = b_i$ per $i = 1, \dots, k$, e allora $a_0 = 1 - \sum_{i \geqslant 1} a_i = 1 - \sum_{i \geqslant 1} b_i = b_0$\hfill$\square$
\end{proof}

\begin{observacion}
Il convesso generato $C$ \`e contenuto nel $k$-piano $\Pi$ generato dagli stessi vettori ($C \subseteq \Pi$): \`e la parte di $\Pi$ dove i coefficienti sono tutti non negativi.
\end{observacion}


% ======== NUOVO (pp. 121-124) ========
\subsection{Terminologia e notazioni}

Siano $v_0, \dots, v_k \in V$ vettori c-indipendenti. L'insieme convesso generato da $v_0, \dots, v_k$ si chiama \textbf{$k$-simplesso (chiuso)} di vertici $v_0, \dots, v_k$ e si indica con
\[
[v_0, \dots, v_k] = \left\{ v = \sum_{i=0}^k t_i v_i \;\Big|\; t_i \geqslant 0,\ \sum_{i=0}^k t_i = 1 \right\}.
\]
I vettori $v_0, \dots, v_k$ si dicono \textbf{vertici} del $k$-simplesso e $k$ si dice \textbf{dimensione} del simplesso.

\medskip
Per ogni $h \leqslant k$ e per ogni scelta di indici $j_0, \dots, j_h \in \{0, \dots, k\}$ distinti, i vettori $v_{j_0}, \dots, v_{j_h}$ sono c-indipendenti (sottofamiglia di una famiglia c-indipendente) e il simplesso $[v_{j_0}, \dots, v_{j_h}] \subseteq [v_0, \dots, v_k]$ si dice \textbf{faccia $h$-dimensionale chiusa} di $[v_0, \dots, v_k]$. In particolare, fissato $j \in \{0, \dots, k\}$,
\[
[v_0, \dots, \widehat{v_j}, \dots, v_k] = [v_0, \dots, v_{j-1}, v_{j+1}, \dots, v_k]
\]
(il cappuccio indica che $v_j$ \`e omesso) si dice \textbf{faccia $(k-1)$-dimensionale chiusa opposta a $v_j$}.

\medskip
L'insieme
\[
(s) = (v_0, \dots, v_k) = \left\{ v \in V \;\Big|\; v = \sum_{i=0}^k t_i v_i,\ t_i > 0\ \forall i,\ \sum_{i=0}^k t_i = 1 \right\}
\]
(con coefficienti \emph{strettamente} positivi) si chiama \textbf{$k$-simplesso aperto} determinato da $v_0, \dots, v_k$. Per ogni $h \leqslant k$ e $j_0, \dots, j_h$ distinti, $(v_{j_0}, \dots, v_{j_h})$ si dice \textbf{faccia $h$-dimensionale aperta} del $k$-simplesso (chiuso) $[s] = [v_0, \dots, v_k]$. Chiaramente $(s) \subseteq [s]$.

\begin{ejercicio*}
\begin{enumerate}
    \item Dimostrare che $(s)$ \`e convesso.
    \item Dimostrare che $[s]$ \`e l'unione \emph{disgiunta} delle sue facce aperte $(s')$.
\end{enumerate}
\end{ejercicio*}
\textit{Soluzione.} (1) Se $x = \sum a_i v_i$ e $y = \sum b_i v_i$ hanno coefficienti $a_i, b_i > 0$, allora $(1-t)x + ty = \sum ((1-t)a_i + t b_i) v_i$ ha coefficienti $(1-t)a_i + t b_i > 0$ per ogni $t \in [0,1]$ (perch\'e almeno uno tra $1-t$ e $t$ \`e positivo) e somma $1$.
(2) Se $v = \sum t_i v_i \in [s]$, sia $J = \{ i \mid t_i > 0 \}$: allora $v$ appartiene alla faccia aperta $(v_j)_{j \in J}$. Poich\'e i coefficienti sono univocamente determinati (teorema precedente), $J$ \`e univocamente determinato da $v$, quindi $v$ sta in \emph{una sola} faccia aperta: l'unione \`e disgiunta.

\begin{ejemplo}
\begin{enumerate}
    \item \textbf{$k = 0$.} Impropriamente si chiama $0$-simplesso il pi\`u piccolo convesso contenente $v_0$, cio\`e $\{v_0\}$. Osserviamo che $k = 0$ \`e l'unico caso in cui il simplesso aperto e quello chiuso coincidono: $[v_0] = (v_0) = \{v_0\}$.
    \item \textbf{$k = 1$.} Siano $v_0, v_1$ c-indipendenti ($v_0 \neq v_1$). Allora
    \[
    [v_0, v_1] = \{ v \mid v = (1-t_1)v_0 + t_1 v_1,\ t_1 \geqslant 0,\ 1 - t_1 \geqslant 0 \} = \{ (1-t)v_0 + t v_1 \mid 0 \leqslant t \leqslant 1 \} = \overline{v_0 v_1},
    \]
    il \textbf{segmento} di estremi $v_0, v_1$. Inoltre $(v_0, v_1) = \{ (1-t)v_0 + t v_1 \mid 0 < t < 1 \}$ e
    \[
    [v_0, v_1] = (v_0, v_1) \sqcup (v_0) \sqcup (v_1).
    \]
    \item \textbf{$k = 2$.} Se $v_0, v_1, v_2$ non sono allineati, $[v_0, v_1, v_2]$ \`e il \textbf{triangolo pieno} di vertici $v_0, v_1, v_2$.
    \item \textbf{$k = 3$.} Se $v_0, v_1, v_2, v_3$ non sono complanari, $[v_0, v_1, v_2, v_3]$ \`e il \textbf{tetraedro pieno} di vertici $v_0, v_1, v_2, v_3$.
\end{enumerate}
\end{ejemplo}

\begin{figure}[H]
    \centering
    \begin{tikzpicture}[scale=1.0]
        % 0-simplesso
        \begin{scope}
            \filldraw (0,0) circle (2pt) node[below] {$v_0$};
            \node at (0,-1.2) {$[v_0]$};
        \end{scope}
        % 1-simplesso
        \begin{scope}[shift={(2.6,0)}]
            \draw[very thick, blue] (0,0) -- (1.8,0.5);
            \filldraw (0,0) circle (2pt) node[below] {$v_0$};
            \filldraw (1.8,0.5) circle (2pt) node[above right] {$v_1$};
            \node at (0.9,-1.2) {$[v_0, v_1]$};
        \end{scope}
        % 2-simplesso
        \begin{scope}[shift={(6.0,0)}]
            \fill[blue!15] (0,0) -- (2,0.2) -- (0.9,1.7) -- cycle;
            \draw[thick, blue] (0,0) -- (2,0.2) -- (0.9,1.7) -- cycle;
            \filldraw (0,0) circle (2pt) node[below left] {$v_0$};
            \filldraw (2,0.2) circle (2pt) node[below right] {$v_1$};
            \filldraw (0.9,1.7) circle (2pt) node[above] {$v_2$};
            \node at (1.0,-1.2) {$[v_0, v_1, v_2]$};
        \end{scope}
        % 3-simplesso
        \begin{scope}[shift={(9.8,0)}]
            \fill[blue!12] (0,0) -- (2.1,0.1) -- (1.0,1.8) -- cycle;
            \draw[thick, blue] (0,0) -- (2.1,0.1) -- (1.0,1.8) -- cycle;
            \draw[thick, blue] (0.9,0.7) -- (0,0);
            \draw[thick, blue] (0.9,0.7) -- (2.1,0.1);
            \draw[thick, blue] (0.9,0.7) -- (1.0,1.8);
            \filldraw (0,0) circle (2pt) node[below left] {$v_0$};
            \filldraw (2.1,0.1) circle (2pt) node[below right] {$v_1$};
            \filldraw (1.0,1.8) circle (2pt) node[above] {$v_2$};
            \filldraw (0.9,0.7) circle (2pt) node[above right] {$v_3$};
            \node at (1.0,-1.2) {$[v_0, \dots, v_3]$};
        \end{scope}
    \end{tikzpicture}
    \caption{I simplessi di dimensione $0$, $1$, $2$, $3$: punto, segmento, triangolo pieno, tetraedro pieno.}
\end{figure}

\begin{teo}
Sia $[s] = [v_0, \dots, v_k]$ un $k$-simplesso e sia fissato $j \in \{0, \dots, k\}$. Allora
\[
v \in [s] \iff \exists\, w \in [v_0, \dots, \widehat{v_j}, \dots, v_k] \ : \ v \in \overline{w, v_j}.
\]
In parole: $[s]$ \`e l'unione dei segmenti che congiungono $v_j$ con i punti della faccia opposta (un ``cono'' di vertice $v_j$ sulla faccia opposta).
\end{teo}
\begin{proof}
($\Leftarrow$) \`E banale: $v_j \in [s]$, la faccia opposta $[v_0, \dots, \widehat{v_j}, \dots, v_k] \subseteq [s]$, quindi $w \in [s]$ e, essendo $[s]$ convesso, $\overline{w, v_j} \subseteq [s]$.

($\Rightarrow$) Senza perdita di generalit\`a $j = k$. Sia $v \in [s]$: $v = \sum_{i=0}^k t_i v_i$, $t_i \geqslant 0$, $\sum t_i = 1$. Se $v = v_k$, ogni $w$ va bene ($v \in \overline{w, v_k}$). Sia quindi $v \neq v_k$: allora $t_k \neq 1$, cio\`e $t_k < 1$, e
\[
v = \sum_{i=0}^{k-1} t_i v_i + t_k v_k = (1 - t_k) \underbrace{\sum_{i=0}^{k-1} \frac{t_i}{1 - t_k}\, v_i}_{=: w} + t_k v_k .
\]
I coefficienti $t_i/(1-t_k)$ sono $\geqslant 0$ e hanno somma $\frac{1-t_k}{1-t_k} = 1$, quindi $w \in [v_0, \dots, v_{k-1}] = [v_0, \dots, v_{k-1}, \widehat{v_k}]$ e $v = (1-t_k)w + t_k v_k \in \overline{w, v_k}$\hfill$\square$
\end{proof}

\begin{definicion}[Coordinate baricentriche e baricentro]
Sia $[s] = [v_0, \dots, v_k]$ un $k$-simplesso, cio\`e per ogni $v \in [s]$ si ha $v = \sum_{i=0}^k t_i v_i$ con $t_i \geqslant 0$, $\sum t_i = 1$. I coefficienti $(t_0, \dots, t_k)$ sono univocamente determinati e si chiamano \textbf{coordinate baricentriche} di $v$ in $[s]$. Il punto $b = \left( \frac{1}{k+1}, \dots, \frac{1}{k+1} \right)$ (cio\`e $b = \frac{1}{k+1} \sum_i v_i$) si chiama \textbf{baricentro} del $k$-simplesso $[s]$.
\end{definicion}

\begin{intuizione}
Le coordinate baricentriche dicono ``quanto peso'' si d\`a a ciascun vertice per ottenere il punto $v$: $v$ \`e il baricentro di masse $t_0, \dots, t_k$ poste sui vertici. Il baricentro $b$ \`e il caso in cui tutti i pesi sono uguali.
\end{intuizione}


% ======== NUOVO (pp. 125-126, 22/10/2026) ========
\subsection{Il baricentro di un simplesso}

\begin{ejemplo}
Sia $[s] = [v_0, \dots, v_k]$ e $b = \left( \frac{1}{k+1}, \dots, \frac{1}{k+1} \right)$ il suo baricentro.
\begin{itemize}
    \item $k = 0$: $[v_0] = (v_0) = \{v_0\} = \{b\}$.
    \item $k = 1$: $[v_0, v_1] = \overline{v_0 v_1}$ e $b = \frac12 v_0 + \frac12 v_1$: $b$ \`e il \textbf{punto medio} del segmento, che lo divide in due parti uguali.
    \item $k = 2$: $[v_0, v_1, v_2]$ e
    \[
    b = \tfrac13 v_0 + \tfrac13 v_1 + \tfrac13 v_2 = \tfrac13 v_0 + \tfrac23\left( \tfrac12 v_1 + \tfrac12 v_2 \right) .
    \]
    Quindi $b$ \`e il baricentro del triangolo, cio\`e il punto di intersezione delle sue \textbf{mediane}; $b$ divide ciascuna mediana in due segmenti che stanno in rapporto $1 : 2$.
    \item $k = 3$: $[v_0, v_1, v_2, v_3]$ e
    \[
    b = \tfrac14 v_0 + \tfrac14 v_1 + \tfrac14 v_2 + \tfrac14 v_3 = \tfrac14 v_0 + \tfrac34\left( \tfrac13 v_1 + \tfrac13 v_2 + \tfrac13 v_3 \right).
    \]
    Quindi $b$ \`e il baricentro del tetraedro, cio\`e il punto di intersezione dei segmenti che congiungono ciascun vertice con il baricentro della faccia ad esso opposta, e li divide in due segmenti in rapporto $1 : 3$.
\end{itemize}
\end{ejemplo}

\begin{figure}[H]
    \centering
    \begin{tikzpicture}[scale=1.4]
        \coordinate (v0) at (0,0);
        \coordinate (v1) at (3,0);
        \coordinate (v2) at (1.1,2.4);
        \coordinate (m0) at ($(v1)!0.5!(v2)$);
        \coordinate (m1) at ($(v0)!0.5!(v2)$);
        \coordinate (m2) at ($(v0)!0.5!(v1)$);
        \draw[thick, fill=blue!8] (v0) -- (v1) -- (v2) -- cycle;
        \draw[gray] (v0) -- (m0); \draw[gray] (v1) -- (m1); \draw[gray] (v2) -- (m2);
        \coordinate (bb) at (1.3667,0.8);
        \filldraw[red] (bb) circle (2pt) node[above right] {$b$};
        \filldraw (v0) circle (1.5pt) node[below left] {$v_0$};
        \filldraw (v1) circle (1.5pt) node[below right] {$v_1$};
        \filldraw (v2) circle (1.5pt) node[above] {$v_2$};
    \end{tikzpicture}
    \caption{Il baricentro $b$ di un triangolo \`e l'intersezione delle mediane e le divide in rapporto $1:2$.}
\end{figure}

\begin{prop}
Sia $[s] = [v_0, \dots, v_k]$ un $k$-simplesso e $b$ il suo baricentro. Allora $b$ \`e il punto di intersezione dei segmenti che congiungono ciascun vertice $v_j$ con il baricentro della faccia ad esso opposta, e divide ciascuno di questi segmenti in due parti che stanno in rapporto $1 : k$ (la parte verso la faccia \`e $1/(k+1)$ del segmento, la parte verso il vertice \`e $k/(k+1)$).
\end{prop}
\begin{proof}
Per simmetria basta considerare $v_0$. Il baricentro della faccia opposta $[v_1, \dots, v_k]$ \`e $b_0' = \frac{1}{k}(v_1 + \dots + v_k)$ e
\[
b = \frac{1}{k+1} v_0 + \dots + \frac{1}{k+1} v_k = \frac{1}{k+1} v_0 + \frac{k}{k+1} \left( \frac1k v_1 + \dots + \frac1k v_k \right) = \frac{1}{k+1} v_0 + \frac{k}{k+1}\, b_0' .
\]
Dunque $b$ sta sul segmento $\overline{v_0 b_0'}$ e $b - v_0 = \frac{k}{k+1}(b_0' - v_0)$: la distanza da $v_0$ \`e i $k/(k+1)$ del segmento, quella da $b_0'$ \`e $1/(k+1)$, quindi $|b b_0'| : |v_0 b| = 1 : k$. Lo stesso calcolo vale per ogni vertice $v_j$, e il punto $b$, avendo coordinate baricentriche $\frac{1}{k+1}$ tutte uguali, \`e lo stesso in ogni caso\hfill$\square$
\end{proof}


% ======== NUOVO (pp. 127-130) ========
\section{Complessi simpliciali}

\begin{definicion}[Complesso simpliciale]
Si chiama \textbf{complesso simpliciale} (euclideo finito) un qualsiasi insieme \emph{finito} $K$ di simplessi \emph{aperti} di $\mathbb{R}^n$ tale che:
\begin{enumerate}
    \item[1)] se $(s) \in K$, allora $K$ contiene ogni faccia aperta di $(s)$;
    \item[2)] per ogni $(s), (s') \in K$: $(s) = (s')$ oppure $(s) \cap (s') = \varnothing$.
\end{enumerate}
\end{definicion}

\begin{intuizione}
Un complesso simpliciale \`e un ``oggetto geometrico costruito incollando punti, segmenti, triangoli, tetraedri, \dots'' in modo coerente: (1) se c'\`e un triangolo ci sono anche i suoi lati e i suoi vertici; (2) due simplessi distinti non si sovrappongono (si incontrano al pi\`u lungo facce comuni, che sono a loro volta elementi di $K$).
\end{intuizione}

\begin{definicion}[Dimensione e sostegno]
La \textbf{dimensione} del complesso simpliciale $K$ \`e $\dim K = \max_{(s) \in K} \dim [s]$. Si chiama \textbf{sostegno} di $K$ l'insieme
\[
|K| = \bigcup_{(s) \in K} (s) \subseteq \mathbb{R}^n .
\]
\end{definicion}

\begin{observacion}
\textbf{Non confondere $K$ e $|K|$!} $K$ \`e un insieme \emph{finito} di simplessi aperti (dato combinatorio), mentre $|K|$ \`e un sottoinsieme di $\mathbb{R}^n$ (lo spazio topologico). Lo stesso sostegno pu\`o essere il sostegno di complessi diversi.
\end{observacion}

\begin{ejemplo}[Decomposizione simpliciale di un simplesso chiuso]
Sia $[s] = [v_0, \dots, v_k]$ un $k$-simplesso chiuso. Allora $K = \{ (s') \mid (s') \text{ faccia aperta di } [s] \}$ \`e un complesso simpliciale di dimensione $k$ e sostegno $|K| = [s]$ (per l'esercizio precedente $[s]$ \`e l'unione disgiunta delle sue facce aperte). Un $k$-simplesso ha $2^{k+1} - 1$ facce aperte non vuote:
\begin{itemize}
    \item $k = 1$: $[v_0, v_1] = \overline{v_0 v_1}$, $K = \{ (v_0), (v_1), (v_0, v_1) \}$ (3 elementi).
    \item $k = 2$: $[v_0, v_1, v_2]$, $K = \{ (v_0), (v_1), (v_2), (v_0, v_1), (v_0, v_2), (v_1, v_2), (v_0, v_1, v_2) \}$ (7 elementi).
    \item $k = 3$: $[v_0, v_1, v_2, v_3]$, $K$ ha $4$ vertici, $6$ spigoli, $4$ facce triangolari e $1$ tetraedro aperto: $15$ elementi.
\end{itemize}
Cos\`i, il segmento chiuso $[v_0, v_1]$ \`e il sostegno di $K = \{ (v_0), (v_1), (v_0,v_1) \}$.
\end{ejemplo}

\begin{figure}[H]
    \centering
    \begin{tikzpicture}[scale=1.0]
        \begin{scope}
            \draw[very thick, blue] (0,0) -- (2,0);
            \filldraw (0,0) circle (2pt) node[below] {$v_0$};
            \filldraw (2,0) circle (2pt) node[below] {$v_1$};
            \node at (1,-1.0) {$k=1$: 3 simplessi};
        \end{scope}
        \begin{scope}[shift={(4.2,0)}]
            \fill[blue!15] (0,0) -- (2,0) -- (0.8,1.6) -- cycle;
            \draw[thick, blue] (0,0) -- (2,0) -- (0.8,1.6) -- cycle;
            \filldraw (0,0) circle (2pt) node[below left] {$v_0$};
            \filldraw (2,0) circle (2pt) node[below right] {$v_1$};
            \filldraw (0.8,1.6) circle (2pt) node[above] {$v_2$};
            \node at (1,-1.0) {$k=2$: 7 simplessi};
        \end{scope}
        \begin{scope}[shift={(8.4,0)}]
            \fill[blue!12] (0,0) -- (2.1,0.1) -- (1.0,1.7) -- cycle;
            \draw[thick, blue] (0,0) -- (2.1,0.1) -- (1.0,1.7) -- cycle;
            \draw[thick, blue] (0.9,0.6) -- (0,0); \draw[thick, blue] (0.9,0.6) -- (2.1,0.1); \draw[thick, blue] (0.9,0.6) -- (1.0,1.7);
            \filldraw (0,0) circle (2pt); \filldraw (2.1,0.1) circle (2pt);
            \filldraw (1.0,1.7) circle (2pt); \filldraw (0.9,0.6) circle (2pt);
            \node at (1,-1.0) {$k=3$: 15 simplessi};
        \end{scope}
    \end{tikzpicture}
    \caption{Le decomposizioni simpliciali di un $k$-simplesso chiuso, $k = 1,2,3$.}
\end{figure}

\begin{ejemplo}[Altri esempi e controesempi]
\begin{enumerate}
    \item \textbf{Suddivisioni.} Scelto $v_2 \in (v_0, v_1)$, il segmento si pu\`o anche descrivere come $K' = \{ (v_0), (v_2), (v_1), (v_0, v_2), (v_2, v_1) \}$, che \`e un complesso simpliciale diverso da $K$ ma con lo stesso sostegno: $|K'| = |K| = [v_0, v_1]$.
    \item Un triangolo pieno $[v_0, v_1, v_2]$ con un cammino di segmenti $(v_1, v_3), (v_3, v_4)$ attaccato al vertice $v_1$ \`e un complesso simpliciale di dimensione $2$ (formato da $5$ vertici, i segmenti $(v_0,v_1), (v_0, v_2), (v_1, v_2), (v_1, v_3), (v_3, v_4)$ e il triangolo $(v_0, v_1, v_2)$).
    \item Un esagono vuoto (poligono di lati $(v_0,v_1), \dots, (v_7, v_0)$) con un triangolo pieno attaccato a un vertice \`e un complesso simpliciale di dimensione $2$.
    \item \textbf{Ogni grafo} (finito, con vertici e spigoli ``segmenti'') \`e un complesso simpliciale di dimensione $1$.
    \item \textbf{Controesempio 1.} L'insieme $\{ (v_0), (v_1), (v_2), (v_3), (v_0, v_1), (v_2, v_3) \}$, dove i due segmenti $(v_0,v_1)$ e $(v_2,v_3)$ si intersecano a croce nel loro interno, \textbf{non} \`e un complesso simpliciale: $(v_0, v_1) \neq (v_2, v_3)$ ma $(v_0, v_1) \cap (v_2, v_3) \neq \varnothing$ (viola la condizione 2).
    \item \textbf{Controesempio 2.} $\{ (v_0), (v_1), (v_2), (v_0, v_1), (v_1, v_2), (v_0, v_1, v_2) \}$ (manca il lato $(v_0, v_2)$) \textbf{non} \`e un complesso simpliciale: $(v_0, v_2)$ \`e una faccia aperta di $[v_0, v_1, v_2]$ ma non appartiene all'insieme (viola la condizione 1).
\end{enumerate}
\end{ejemplo}

\begin{figure}[H]
    \centering
    \begin{tikzpicture}[scale=1.0]
        % esempio 2
        \begin{scope}
            \fill[blue!15] (0,0) -- (1.4,0) -- (0.5,1.3) -- cycle;
            \draw[thick, blue] (0,0) -- (1.4,0) -- (0.5,1.3) -- cycle;
            \draw[thick, blue] (1.4,0) -- (2.4,0.7) -- (3.3,0.2);
            \foreach \p/\n/\d in {(0,0)/v_0/below left, (1.4,0)/v_1/below, (0.5,1.3)/v_2/above, (2.4,0.7)/v_3/above, (3.3,0.2)/v_4/below right}
                \filldraw \p circle (1.5pt) node[\d] {$\n$};
            \node at (1.6,-1.0) {complesso (dim. 2)};
        \end{scope}
        % controesempio 1
        \begin{scope}[shift={(5.5,0)}]
            \draw[thick, red] (-0.8,0) -- (0.8,0); \draw[thick, red] (0,-0.8) -- (0,0.8);
            \foreach \p/\n/\d in {(-0.8,0)/v_0/left, (0.8,0)/v_1/right, (0,-0.8)/v_2/below, (0,0.8)/v_3/above}
                \filldraw \p circle (1.5pt) node[\d] {$\n$};
            \node at (0,-1.5) {\textbf{non} complesso (cond. 2)};
        \end{scope}
        % controesempio 2
        \begin{scope}[shift={(9.5,0)}]
            \fill[red!10] (0,0) -- (1.4,0) -- (0.5,1.3) -- cycle;
            \draw[thick, red] (0,0) -- (1.4,0) -- (0.5,1.3);
            \draw[thick, red, dashed] (0.5,1.3) -- (0,0);
            \foreach \p/\n/\d in {(0,0)/v_0/below left, (1.4,0)/v_1/below right, (0.5,1.3)/v_2/above}
                \filldraw \p circle (1.5pt) node[\d] {$\n$};
            \node at (0.7,-1.0) {\textbf{non} complesso (cond. 1)};
        \end{scope}
    \end{tikzpicture}
    \caption{Un complesso simpliciale e due controesempi.}
\end{figure}


% ======== NUOVO (pp. 131-132) ========
\begin{ejemplo}[Altri esempi, continuazione]
\begin{enumerate}
    \item[7)] Nel caso del controesempio 1 (due segmenti che si incrociano) si ottiene un complesso simpliciale con lo stesso sostegno aggiungendo il punto di intersezione $v_4$ come nuovo vertice e dividendo i due segmenti:
    \[
    K = \{ (v_0), \dots, (v_4), (v_0, v_4), (v_1, v_4), (v_2, v_4), (v_3, v_4) \}.
    \]
    Invece, nel controesempio 2 (il triangolo senza un lato) \textbf{non esiste alcun complesso simpliciale con quel sostegno}: infatti il sostegno di un complesso simpliciale \`e unione di simplessi \emph{chiusi} (proposizione seguente), quindi \`e un insieme chiuso, mentre ``triangolo meno un lato aperto'' non \`e chiuso.
    \item[8)] Un triangolo pieno $[v_0, v_1, v_2]$ bucato trasversalmente da un segmento $(v_3, v_4)$ ha $(v_0, v_1, v_2) \cap (v_3, v_4) \neq \varnothing$, quindi non \`e un complesso simpliciale. Ne \`e il sostegno il complesso che si ottiene aggiungendo il punto di intersezione $v_5$ come vertice, suddividendo il triangolo in tre triangoli e il segmento in due:
    \begin{align*}
    K = \{ & (v_0), \dots, (v_5),\ (v_0,v_1), (v_0,v_2), (v_1,v_2), (v_0,v_5), (v_1,v_5), (v_2,v_5), (v_3,v_5), (v_4,v_5), \\
    & (v_0,v_1,v_5), (v_1,v_2,v_5), (v_0,v_2,v_5) \}.
    \end{align*}
    \item[9)] Ogni poligono \`e il sostegno di un complesso simpliciale (si triangola); di conseguenza ogni poliedro \`e il sostegno di un complesso simpliciale.
\end{enumerate}
\end{ejemplo}

\begin{figure}[H]
    \centering
    \begin{tikzpicture}[scale=1.1]
        % esempio 7: triangolo trafitto
        \begin{scope}
            \draw[thick, red] (0,0) -- (3,0) -- (1.4,2.3) -- cycle;
            \draw[thick, red] (1.1,1.0) -- (3.0,0.4);
            \draw[thick, red, dashed] (0.7,0.5) -- (1.1,1.0);
            \node[below] at (1.5,-0.3) {\textbf{non} complesso: $(v_0,v_1,v_2)\cap(v_3,v_4)\neq\varnothing$};
        \end{scope}
        \begin{scope}[shift={(6.2,0)}]
            \fill[blue!12] (0,0) -- (3,0) -- (1.4,2.3) -- cycle;
            \draw[thick, blue] (0,0) -- (3,0) -- (1.4,2.3) -- cycle;
            \coordinate (v5) at (1.45,0.85);
            \draw[thick, blue] (v5) -- (0,0); \draw[thick, blue] (v5) -- (3,0); \draw[thick, blue] (v5) -- (1.4,2.3);
            \draw[thick, blue] (v5) -- (0.5,2.0); \draw[thick, blue] (v5) -- (2.8,1.4);
            \filldraw[red] (v5) circle (2pt) node[below] {$v_5$};
            \node[below] at (1.5,-0.3) {complesso dopo la suddivisione};
        \end{scope}
    \end{tikzpicture}
    \caption{Per rendere ``simpliciale'' una configurazione in cui i simplessi si intersecano nel loro interno, si aggiunge il punto di intersezione come nuovo vertice e si suddivide.}
\end{figure}

\begin{prop}
Se $K$ \`e un complesso simpliciale, allora
\[
|K| = \bigcup_{(s) \in K} [s].
\]
In particolare $|K|$ \`e un sottoinsieme \textbf{chiuso} (e compatto) di $\mathbb{R}^n$, essendo unione finita di simplessi chiusi (che sono compatti).
\end{prop}
\begin{proof}
``$\subseteq$'': $|K| = \bigcup_{(s) \in K} (s) \subseteq \bigcup_{(s) \in K} [s]$, perch\'e $(s) \subseteq [s]$.
``$\supseteq$'': per ogni $(s) \in K$, $[s]$ \`e l'unione delle facce aperte $(s')$ di $(s)$ (esercizio precedente), che per la condizione 1) appartengono a $K$; dunque $[s] = \bigcup (s') \subseteq |K|$\hfill$\square$
\end{proof}

\begin{definicion}[$n$-scheletro]
Sia $K$ un complesso simpliciale, $m = \dim K$ e $n \in \{0, \dots, m\}$. Si chiama \textbf{$n$-scheletro} di $K$ il complesso simpliciale
\[
K^n = \{ (s) \in K \mid \dim (s) \leqslant n \}.
\]
\end{definicion}

\begin{observacion}
\begin{enumerate}
    \item $\dim K^n = n$.
    \item $K^m = K$.
    \item $K^0$ \`e l'insieme dei \textbf{vertici} di $K$.
\end{enumerate}
\end{observacion}

\begin{intuizione}
L'$n$-scheletro conserva tutti i simplessi di dimensione al pi\`u $n$ e ``butta via'' quelli pi\`u grandi: per un tetraedro pieno, $K^0$ sono i 4 vertici, $K^1$ lo scheletro dei 6 spigoli, $K^2$ la superficie (4 facce) e $K^3$ tutto il tetraedro. (\`E ancora un complesso simpliciale perch\'e le facce di un simplesso hanno dimensione minore.)
\end{intuizione}


% ======== NUOVO (pp. 133-134) ========
\begin{ejemplo}[Scheletri]
\begin{enumerate}
    \item Sia $K$ la decomposizione simpliciale del tetraedro $[v_0, v_1, v_2, v_3]$. Allora
    \begin{align*}
    K^0 &= \{ (v_0), (v_1), (v_2), (v_3) \}, \\
    K^1 &= K^0 \cup \{ (v_0,v_1), (v_0,v_2), (v_0,v_3), (v_1,v_2), (v_1,v_3), (v_2,v_3) \}, \\
    K^2 &= K^1 \cup \{ (v_0,v_1,v_2), (v_0,v_1,v_3), (v_0,v_2,v_3), (v_1,v_2,v_3) \}, \\
    K^3 &= K.
    \end{align*}
    \item Sia $K$ una triangolazione di un quadrato (con un vertice al centro e otto triangoli): $K^0$ sono i 9 vertici, $K^1$ \`e il grafo formato dai vertici e dagli spigoli, e $K = K^2$.
\end{enumerate}
\end{ejemplo}

\begin{figure}[H]
    \centering
    \begin{tikzpicture}[scale=1.0]
        \foreach \s/\name in {0/{$K^0$},4.2/{$K^1$},8.4/{$K^2 = K$}} {
            \begin{scope}[shift={(\s,0)}]
                \node at (1,-1.0) {\name};
            \end{scope}
        }
        % K^0
        \foreach \x in {0,1,2} \foreach \y in {0,1,2} { \filldraw (\x*1.0,\y*1.0) circle (1.6pt); }
        % K^1
        \begin{scope}[shift={(4.2,0)}]
            \draw[thick] (0,0) rectangle (2,2);
            \draw[thick] (1,0) -- (1,2); \draw[thick] (0,1) -- (2,1);
            \draw[thick] (0,0) -- (2,2); \draw[thick] (0,2) -- (2,0);
            \foreach \x in {0,1,2} \foreach \y in {0,1,2} { \filldraw (\x*1.0,\y*1.0) circle (1.6pt); }
        \end{scope}
        % K^2
        \begin{scope}[shift={(8.4,0)}]
            \fill[blue!12] (0,0) rectangle (2,2);
            \draw[thick] (0,0) rectangle (2,2);
            \draw[thick] (1,0) -- (1,2); \draw[thick] (0,1) -- (2,1);
            \draw[thick] (0,0) -- (2,2); \draw[thick] (0,2) -- (2,0);
            \foreach \x in {0,1,2} \foreach \y in {0,1,2} { \filldraw (\x*1.0,\y*1.0) circle (1.6pt); }
        \end{scope}
    \end{tikzpicture}
    \caption{Gli scheletri di una triangolazione del quadrato: vertici, grafo, e il complesso completo.}
\end{figure}

\begin{definicion}[Suddivisione]
Siano $K, K^*$ complessi simpliciali. Si dice che $K^*$ \`e una \textbf{suddivisione} di $K$ se $|K| = |K^*|$ e per ogni $(s) \in K^*$ esiste $(s') \in K$ tale che $(s) \subseteq (s')$.
\end{definicion}

\begin{intuizione}
$K^*$ ha lo stesso sostegno di $K$, ma i suoi simplessi sono ``pi\`u piccoli'': ciascun simplesso di $K^*$ sta dentro un simplesso di $K$. \`E una triangolazione pi\`u fine della stessa figura.
\end{intuizione}

\begin{ejemplo}
\begin{enumerate}
    \item Sia $K = \{ (v_0), (v_1), (v_0, v_1) \}$ (un segmento). Allora, con $v_2 \in (v_0, v_1)$,
    \[ K^* = \{ (v_0), (v_1), (v_2), (v_0, v_2), (v_2, v_1) \} \]
    \`e una suddivisione di $K$. Se prendessimo un altro punto $v_3 \in (v_0, v_1)$, avremmo un'altra suddivisione
    \[ K' = \{ (v_0), (v_3), (v_1), (v_0, v_3), (v_3, v_1) \} \]
    di $K$, \textbf{non confrontabile} con $K^*$ (n\'e $K'$ \`e suddivisione di $K^*$, n\'e viceversa).
    \item Sia $K = \{ \text{vertici } v_0, v_1, v_2,\ \text{lati},\ \text{triangolo} \}$ (un triangolo pieno). Aggiungendo i punti medi dei lati $v_3, v_4, v_5$ e dividendo in $4$ triangoli si ottiene una suddivisione $K^*$ di $K$; dividendo invece il triangolo in due con il segmento $(v_2, v_3)$, dove $v_3 \in (v_0, v_1)$, si ottiene un'altra suddivisione $K'$ di $K$, non confrontabile con $K^*$.
\end{enumerate}
\end{ejemplo}

\begin{figure}[H]
    \centering
    \begin{tikzpicture}[scale=1.1]
        \begin{scope}
            \fill[blue!10] (0,0) -- (3,0) -- (1.5,2.4) -- cycle;
            \draw[thick] (0,0) -- (3,0) -- (1.5,2.4) -- cycle;
            \filldraw (0,0) circle (1.6pt) node[below left] {$v_0$};
            \filldraw (3,0) circle (1.6pt) node[below right] {$v_1$};
            \filldraw (1.5,2.4) circle (1.6pt) node[above] {$v_2$};
            \node at (1.5,-1.0) {$K$};
        \end{scope}
        \begin{scope}[shift={(4.6,0)}]
            \fill[blue!10] (0,0) -- (3,0) -- (1.5,2.4) -- cycle;
            \draw[thick] (0,0) -- (3,0) -- (1.5,2.4) -- cycle;
            \draw[thick] (1.5,0) -- (0.75,1.2) -- (2.25,1.2) -- cycle;
            \foreach \p in {(0,0),(3,0),(1.5,2.4),(1.5,0),(0.75,1.2),(2.25,1.2)} \filldraw \p circle (1.6pt);
            \node at (1.5,-1.0) {$K^*$ (4 triangoli)};
        \end{scope}
        \begin{scope}[shift={(9.2,0)}]
            \fill[blue!10] (0,0) -- (3,0) -- (1.5,2.4) -- cycle;
            \draw[thick] (0,0) -- (3,0) -- (1.5,2.4) -- cycle;
            \draw[thick] (1.0,0) -- (1.5,2.4);
            \foreach \p in {(0,0),(3,0),(1.5,2.4),(1.0,0)} \filldraw \p circle (1.6pt);
            \node[below] at (1.0,0) {$v_3$};
            \node at (1.5,-1.0) {$K'$ (2 triangoli)};
        \end{scope}
    \end{tikzpicture}
    \caption{Due suddivisioni diverse e non confrontabili dello stesso triangolo $K$.}
\end{figure}


% ======== NUOVO (pp. 135-140) ========
\section{La suddivisione baricentrica}

\begin{definicion}[Ordine parziale in un complesso]
Sia $K$ un complesso simpliciale. Introduciamo in $K$ una relazione d'ordine parziale: per ogni $(s), (s') \in K$,
\[
(s) \leqslant (s') \iff (s) \subseteq [s'],
\]
cio\`e $(s)$ \`e una faccia aperta di $[s']$. \`E riflessiva, transitiva e antisimmetrica: se $(s) \leqslant (s')$ e $(s') \leqslant (s)$ allora $(s) \subseteq (s') \subseteq (s)$, cio\`e $(s) = (s')$. Due simplessi possono per\`o non essere confrontabili: ad esempio due vertici distinti.
\end{definicion}

Per ogni $(s) \in K$ indichiamo con $b(s)$ il \textbf{baricentro} di $(s)$.

\begin{teo}[Prima suddivisione baricentrica]
Dato un complesso simpliciale $K$ con $m = \dim K$, l'insieme
\[
K^{(1)} = \Big\{ \big( b(s_0), \dots, b(s_p) \big) \;\Big|\; p = 0, \dots, m,\ (s_0), \dots, (s_p) \in K,\ (s_i) < (s_{i+1})\ \forall i \Big\}
\]
(catene \emph{strettamente} crescenti di simplessi di $K$) \`e una suddivisione di $K$, detta \textbf{prima suddivisione baricentrica} di $K$. Si pone $K^{(2)} = (K^{(1)})^{(1)}, \dots, K^{(q)} = (K^{(q-1)})^{(1)}$.
\end{teo}

\begin{intuizione}
Per ottenere $K^{(1)}$ si mette un nuovo vertice nel baricentro di \emph{ogni} simplesso di $K$ (anche degli spigoli e dei triangoli) e poi si ``congiungono'' i baricentri ordinati per inclusione: baricentro di un vertice $\to$ baricentro di uno spigolo che lo contiene $\to$ baricentro di un triangolo che contiene lo spigolo $\to \dots$
\end{intuizione}

\begin{ejemplo}
\begin{enumerate}
    \item \textbf{Segmento.} $K = \{ (v_0), (v_1), (v_0, v_1) \}$. Il baricentro di $(v_0, v_1)$ \`e il punto medio $b$. Allora
    \[
    K^{(1)} = \{ (v_0), (v_1), (b), (v_0, b), (b, v_1) \}
    \]
    e, iterando, $(K^{(1)})^{(1)} = \{ (v_0), (v_1), (b), (b_1), (b_2), (v_0, b_1), (b_1, b), (b, b_2), (b_2, v_1) \}$, dove $b_1, b_2$ sono i punti medi di $(v_0, b)$ e $(b, v_1)$: il segmento viene diviso in $4$ parti uguali.
    \item \textbf{Triangolo.} Se $K$ \`e la decomposizione di un triangolo $[v_0, v_1, v_2]$, $K^{(1)}$ ha $7$ vertici ($3$ vertici, $3$ punti medi dei lati e il baricentro $b$ del triangolo) e $6$ triangoli: ciascuno ha vertici $\{$vertice, punto medio di un lato uscente da quel vertice, baricentro $b\}$.
\end{enumerate}
\end{ejemplo}

\begin{figure}[H]
    \centering
    \begin{tikzpicture}[scale=1.3]
        \begin{scope}
            \fill[blue!10] (0,0) -- (3,0) -- (1.4,2.4) -- cycle;
            \draw[thick] (0,0) -- (3,0) -- (1.4,2.4) -- cycle;
            \node at (1.5,-0.6) {$K$};
        \end{scope}
        \begin{scope}[shift={(4.5,0)}]
            \coordinate (A) at (0,0); \coordinate (B) at (3,0); \coordinate (C) at (1.4,2.4);
            \coordinate (mAB) at (1.5,0); \coordinate (mBC) at (2.2,1.2); \coordinate (mAC) at (0.7,1.2);
            \coordinate (G) at (1.4667,0.8);
            \fill[blue!10] (A) -- (B) -- (C) -- cycle;
            \draw[thick] (A) -- (B) -- (C) -- cycle;
            \draw[thick] (A) -- (mBC); \draw[thick] (B) -- (mAC); \draw[thick] (C) -- (mAB);
            \foreach \p in {A,B,C,mAB,mBC,mAC} \filldraw (\p) circle (1.5pt);
            \filldraw[red] (G) circle (2pt) node[above right] {$b$};
            \node at (1.5,-0.6) {$K^{(1)}$ (6 triangoli)};
        \end{scope}
    \end{tikzpicture}
    \caption{La prima suddivisione baricentrica di un triangolo: le tre mediane lo dividono in $6$ triangoli.}
\end{figure}

\begin{prop}
Sia $[s] = [v_0, \dots, v_k]$ un $k$-simplesso. Allora:
\begin{enumerate}
    \item[(a)] per ogni $u, v \in [s]$: $d(u,v) \leqslant \max_{i,j = 0, \dots, k} d(v_i, v_j)$;
    \item[(b)] $\operatorname{diam}[s] = \max_{i,j = 0, \dots, k} d(v_i, v_j)$.
\end{enumerate}
In parole: il diametro di un simplesso \`e la lunghezza del suo spigolo pi\`u lungo.
\end{prop}
\begin{proof}
(a) Siano $u, v \in [s]$, $v = \sum_i t_i v_i$ con $t_i \geqslant 0$, $\sum t_i = 1$. Allora
\[
\begin{aligned}
d(u,v) = \|u - v\| &= \Big\| \Big(\sum_i t_i\Big) u - \sum_i t_i v_i \Big\| = \Big\| \sum_i t_i (u - v_i) \Big\| \\
&\leqslant \sum_i t_i \|u - v_i\| \leqslant \max_\delta \|u - v_\delta\| = \max_\delta d(u, v_\delta).
\end{aligned}
\]
Ora, scrivendo $u = \sum_j a_j v_j$ con $a_j \geqslant 0$, $\sum a_j = 1$, per ogni $\delta$:
\[
d(u, v_\delta) = \Big\| \sum_j a_j v_j - \Big(\sum_j a_j\Big) v_\delta \Big\| = \Big\| \sum_j a_j (v_j - v_\delta) \Big\| \leqslant \sum_j a_j \|v_j - v_\delta\| \leqslant \max_j d(v_j, v_\delta).
\]
Quindi $d(u, v) \leqslant \max_\delta d(u, v_\delta) \leqslant \max_{j, \delta} d(v_j, v_\delta)$.

(b) Per (a), $\operatorname{diam}[s] = \sup_{u,v \in [s]} d(u,v) \leqslant \max_{i,j} d(v_i, v_j)$; d'altra parte i vertici stanno in $[s]$, quindi $\max_{i,j} d(v_i, v_j) \leqslant \operatorname{diam}[s]$\hfill$\square$
\end{proof}

\begin{definicion}[Maglia]
Sia $K$ un complesso simpliciale. Si chiama \textbf{maglia} di $K$ il numero
\[
\operatorname{maglia}(K) = \max_{(s) \in K} \operatorname{diam}[s].
\]
\end{definicion}

\begin{teo}
Sia $K$ un complesso simpliciale, $m = \dim K$ e $n \in \mathbb{N}$. Allora
\[
0 \leqslant \operatorname{maglia}\big(K^{(n)}\big) \leqslant \left( \frac{m}{m+1} \right)^{n} \operatorname{maglia}(K) \ \ \Big( \leqslant \frac{m}{m+n} \operatorname{maglia}(K) \Big).
\]
\end{teo}

\begin{observacion}
Negli appunti la stima \`e scritta nella forma $\frac{m}{m+n}\operatorname{maglia}(K)$ (``senza dimostrazione''). Questa forma \`e vera, ma segue dalla stima pi\`u forte $\big(\frac{m}{m+1}\big)^n$, perch\'e per la disuguaglianza di Bernoulli $\big(1 + \frac1m\big)^n \geqslant 1 + \frac nm = \frac{m+n}{m}$, cio\`e $\big(\frac{m}{m+1}\big)^n \leqslant \frac{m}{m+n}$. Qui si dimostra il caso $n=1$ (da cui il caso generale per induzione, applicando il caso $n = 1$ a $K^{(n-1)}$, che ha ancora dimensione $m$).
\end{observacion}

\begin{proof}[Dimostrazione del caso $n=1$]
Un simplesso di $K^{(1)}$ ha vertici $b(s_0), \dots, b(s_p)$ con $(s_0) < \dots < (s_p)$ in $K$. Per la proposizione precedente il suo diametro \`e il massimo delle distanze tra due vertici, quindi basta stimare $d(b(s_i), b(s_j))$ con $i < j$. Poich\'e $b(s_i) \in [s_i] \subseteq [s_j]$, per il punto (a) della proposizione (la dimostrazione mostra $d(u, w) \leqslant \max_\delta d(u, w_\delta)$) si ha
\[
d(b(s_j), b(s_i)) \leqslant \max_{\delta} d(b(s_j), w_\delta),
\]
dove $w_0, \dots, w_q$ sono i vertici di $(s_j)$ ($q = \dim (s_j) \leqslant m$). Per ogni $\delta$,
\[
\|b(s_j) - w_\delta\| = \Big\| \frac{1}{q+1} \sum_{l} (w_l - w_\delta) \Big\| \leqslant \frac{q}{q+1} \max_l \|w_l - w_\delta\| \leqslant \frac{q}{q+1} \operatorname{diam}[s_j] \leqslant \frac{m}{m+1} \operatorname{maglia}(K),
\]
(il termine con $l = \delta$ \`e nullo, rimangono $q$ addendi; e $\frac{q}{q+1} \leqslant \frac{m}{m+1}$ perch\'e $q \leqslant m$). Dunque $\operatorname{maglia}(K^{(1)}) \leqslant \frac{m}{m+1} \operatorname{maglia}(K)$\hfill$\square$
\end{proof}

\begin{coro}
$\displaystyle \lim_{n \to +\infty} \operatorname{maglia}\big(K^{(n)}\big) = 0$.
\end{coro}

\begin{intuizione}
Suddividendo baricentricamente un numero sufficiente di volte, \emph{tutti} i simplessi diventano piccoli quanto si vuole, pur mantenendo lo stesso sostegno. Questo \`e lo strumento che servir\`a per mostrare che l'omologia simpliciale non dipende dalla triangolazione scelta.
\end{intuizione}

\section{Propriet\`a topologiche dei complessi simpliciali}

Sia $[s] = [v_0, \dots, v_k]$ un $k$-simplesso in $\mathbb{R}^n$ e $\Pi$ il $k$-piano da esso generato, cos\`i che $[s] \subseteq \Pi$. Per la proposizione sulle trasformazioni affini esiste una trasformazione affine $f: \mathbb{R}^n \to \mathbb{R}^n$ con $f(\Pi) = \mathcal{L}(e_1, \dots, e_k)$, dove $\{e_1, \dots, e_n\}$ \`e la base canonica. Osserviamo che una trasformazione affine \`e un omeomorfismo (composizione di una traslazione e di un isomorfismo lineare), quindi $f([s]) \subseteq \mathcal{L}(e_1, \dots, e_k)$ \`e omeomorfo a $[s]$.

Poich\'e $f(v_0) = \bar{0}$ e $f(v_i) = e_i$ ($i = 1, \dots, k$) (la parte lineare di $f$ manda $v_i - v_0$ in $e_i$), e $f$ conserva le combinazioni affini, per $v = \sum_i t_i v_i \in [s]$ si ha $f(v) = \sum_{i=1}^k t_i e_i$ con $t_i \geqslant 0$ e $t_1 + \dots + t_k = 1 - t_0 \leqslant 1$. Quindi
\[
f([s]) = \Big\{ x \in \mathbb{R}^k \;\Big|\; x = \sum_{i=1}^k t_i e_i,\ t_i \geqslant 0,\ \sum_{i=1}^k t_i \leqslant 1 \Big\} = \Delta^k,
\]
il \textbf{simplesso standard} di $\mathbb{R}^k$ (l'insieme dei punti di $\mathcal{L}(e_1, \dots, e_k) \cong \mathbb{R}^k$ con coordinate $\geqslant 0$ e somma $\leqslant 1$). Per $k = 1$: $\mathbb{R} = \mathcal{L}(e_1)$ e $[\bar{0}, e_1] = \{ t e_1 \mid 0 \leqslant t \leqslant 1 \}$.

\begin{prop}
Ogni $k$-simplesso chiuso $[s] \subset \mathbb{R}^n$ \`e omeomorfo al simplesso standard $\Delta^k \subset \mathbb{R}^k$ (e in particolare \`e compatto e connesso per archi, perch\'e $\Delta^k$ \`e chiuso e limitato e convesso).
\end{prop}


% ======== NUOVO (pp. 141-144) ========
Per $k=2$ si ha $\mathbb{R}^2 = \mathcal{L}(e_1, e_2)$ e
\[
[\bar{0}, e_1, e_2] = \{ v \in \mathbb{R}^2 \mid v = t_1 e_1 + t_2 e_2,\ t_1, t_2 \geqslant 0,\ t_1 + t_2 \leqslant 1 \},
\]
il triangolo di vertici $(0,0), (1,0), (0,1)$; per $k = 3$, $[\bar{0}, e_1, e_2, e_3]$ \`e il tetraedro con $t_1, t_2, t_3 \geqslant 0$ e $t_1 + t_2 + t_3 \leqslant 1$.

\begin{figure}[H]
    \centering
    \begin{tikzpicture}[scale=1.5]
        \begin{scope}
            \draw[->] (-0.2,0) -- (1.5,0) node[right] {$t_1$};
            \draw[->] (0,-0.2) -- (0,1.5) node[above] {$t_2$};
            \fill[blue!15] (0,0) -- (1,0) -- (0,1) -- cycle;
            \draw[thick, blue] (0,0) -- (1,0) -- (0,1) -- cycle;
            \node[below] at (1,0) {$e_1$}; \node[left] at (0,1) {$e_2$};
            \node at (0.7,-0.7) {$\Delta^2$};
        \end{scope}
        \begin{scope}[shift={(3.6,0)}]
            \draw[->] (-0.2,0) -- (1.6,0) node[right] {$t_2$};
            \draw[->] (0,-0.2) -- (0,1.6) node[above] {$t_3$};
            \draw[->] (0.2,0.2) -- (-0.7,-0.7) node[below left] {$t_1$};
            \draw[thick, blue] (-0.5,-0.5) -- (1.2,0) -- (0,1.2) -- cycle;
            \draw[thick, blue] (0,0) -- (-0.5,-0.5); \draw[thick, blue] (0,0) -- (1.2,0); \draw[thick, blue] (0,0) -- (0,1.2);
            \draw[thick, blue] (-0.5,-0.5) -- (0,1.2);
            \node at (0.4,-1.0) {$\Delta^3$};
        \end{scope}
    \end{tikzpicture}
    \caption{I simplessi standard $\Delta^2 \subset \mathbb{R}^2$ e $\Delta^3 \subset \mathbb{R}^3$.}
\end{figure}

\begin{observacion}[Conseguenza]
$f([s])$ (e quindi anche $[s]$) \`e un sottoinsieme chiuso e limitato di $\mathbb{R}^n$, dunque \textbf{compatto}.
\end{observacion}

\begin{prop}
Sia $K$ un complesso simpliciale. Allora:
\begin{enumerate}
    \item[a)] $|K|$ \`e compatto;
    \item[b)] $\{ [s] \}_{(s) \in K}$ \`e un ricoprimento \textbf{finito} di \textbf{chiusi} di $|K|$.
\end{enumerate}
\end{prop}
\begin{proof}
Si ha $|K| = \bigcup_{(s) \in K} [s]$, unione finita di compatti (quindi compatta e chiusa) e ciascun $[s]$ \`e chiuso\hfill$\square$
\end{proof}

\begin{prop}
Per ogni spazio topologico $X$ e ogni applicazione $f: |K| \to X$:
\[
f \text{ \`e continua} \iff f|_{[s]} \text{ \`e continua per ogni } (s) \in K.
\]
\end{prop}
\begin{proof}
($\Rightarrow$) \`E banale (restrizione di una continua). ($\Leftarrow$) Segue dal punto b) della proposizione precedente e dal Lemma di Incollamento (ricoprimento finito di chiusi)\hfill$\square$
\end{proof}

\begin{intuizione}
Per controllare la continuit\`a di una funzione su un complesso basta controllarla simplesso per simplesso: \`e il principio che permette di definire applicazioni ``a pezzi'' (lineari su ciascun simplesso) e di ottenere automaticamente applicazioni continue.
\end{intuizione}

\section{Coordinate baricentriche di un complesso}

Siano $K$ un complesso simpliciale e $v \in K^0$ un vertice. Sia $w \in |K| = \bigsqcup_{(s) \in K} (s)$. Allora esiste un \textbf{unico} $(s) \in K$ con $w \in (s)$. Se $(s) = (v_0, \dots, v_k)$, allora $w = \sum_{i=0}^k t_i v_i$ con $t_i > 0$ per ogni $i$ e $\sum_{i=0}^k t_i = 1$. Poniamo
\[
t_v(w) = \begin{cases} 0 & \text{se } v \neq v_i \ \ \forall i = 0, \dots, k, \\ t_i & \text{se } v = v_i. \end{cases}
\]

\begin{definicion}
L'applicazione $t_v: |K| \to \mathbb{R}$, $w \mapsto t_v(w)$ si chiama \textbf{(funzione) coordinata baricentrica} di $K$ rispetto al vertice $v$.
\end{definicion}

\begin{prop}
$t_v: |K| \to \mathbb{R}$ \`e continua.
\end{prop}
\begin{proof}
Per la proposizione precedente basta verificare che $t_v|_{[s]}$ sia continua per ogni $(s) \in K$. Ora, per ogni $(s) \in K$:
\begin{itemize}
    \item se $v$ \textbf{non} \`e un vertice di $[s]$, allora $t_v|_{[s]} \equiv 0$ (ogni punto di $[s]$ sta in una faccia aperta i cui vertici sono vertici di $[s]$): \`e costante, dunque continua;
    \item se $v = v_i$ \`e un vertice di $[s]$, allora $t_v|_{[s]}$ \`e la $i$-esima coordinata baricentrica di $[s]$, che \`e continua perch\'e, via l'omeomorfismo affine $f: [s] \to \Delta^k$, \`e (a meno di un'identificazione) una funzione coordinata (lineare) su $\Delta^k$ (con $t_0 = 1 - t_1 - \dots - t_k$)\hfill$\square$
\end{itemize}
\end{proof}

\begin{ejemplo}
Sia $K$ la decomposizione del triangolo $[v_0, v_1, v_2]$.
\begin{itemize}
    \item Se $w \in (v_1, v_2)$, allora esiste un unico $t \in\, ]0,1[$ con $w = (1-t)v_1 + tv_2$, quindi $t_{v_0}(w) = 0$, $t_{v_1}(w) = 1-t$, $t_{v_2}(w) = t$.
    \item Per il vertice $v_0 \in (v_0) \in K$ si ha $v_0 = 1 \cdot v_0$, quindi $t_{v_0}(v_0) = 1$, $t_{v_1}(v_0) = 0$, $t_{v_2}(v_0) = 0$.
    \item Se $x \in (v_0, v_1, v_2)$, allora $x = (1 - t_1 - t_2)v_0 + t_1 v_1 + t_2 v_2$, e dunque $t_{v_0}(x) = 1 - t_1 - t_2$, $t_{v_1}(x) = t_1$, $t_{v_2}(x) = t_2$.
\end{itemize}
In tutti i casi: $v_0 = \sum_{i=0}^2 t_{v_i}(v_0) v_i$, $w = \sum_{i=0}^2 t_{v_i}(w) v_i$, $x = \sum_{i=0}^2 t_{v_i}(x) v_i$.
\end{ejemplo}

\begin{prop}
Sia $K$ un complesso simpliciale. Allora per ogni $x \in |K|$
\[
x = \sum_{v \in K^0} t_v(x) \cdot v, \qquad \text{con } t_v(x) \geqslant 0 \ \ \forall v \in K^0 \quad \text{e} \quad \sum_{v \in K^0} t_v(x) = 1 .
\]
\end{prop}
\begin{proof}
Sia $x \in |K| = \bigsqcup_{(s) \in K} (s)$; esiste un unico $(s) = (v_0, \dots, v_k) \in K$ con $x \in (s)$, e $x = \sum_{i=0}^k t_i v_i$ con $t_i > 0$, $\sum t_i = 1$. Per definizione $t_{v_i}(x) = t_i$ e $t_v(x) = 0$ per ogni vertice $v \in K^0$ diverso da $v_0, \dots, v_k$. Quindi
\[
\sum_{v \in K^0} t_v(x)\, v = \underbrace{\sum_{v \notin \{v_0, \dots, v_k\}} t_v(x)\, v}_{=0} + \sum_{i=0}^k t_{v_i}(x)\, v_i = \sum_{i=0}^k t_i v_i = x,
\]
e analogamente $t_v(x) \geqslant 0$ per ogni $v$ e $\sum_{v \in K^0} t_v(x) = \sum_{i=0}^k t_i = 1$\hfill$\square$
\end{proof}

\begin{intuizione}
Le funzioni $t_v$, $v \in K^0$, forniscono un sistema di ``coordinate globali'' su $|K|$ (le \emph{coordinate baricentriche} di $K$): un punto di $|K|$ \`e determinato dai pesi $(t_v(x))_{v \in K^0}$, non nulli solo sui vertici dell'unico simplesso aperto che lo contiene.
\end{intuizione}


% ======== NUOVO (pp. 145-148) ========
\begin{lema}
Per ogni $x, y \in |K|$: $\ x = y \iff t_v(x) = t_v(y)$ per ogni $v \in K^0$.
\end{lema}
\begin{proof}
($\Rightarrow$) \`E banale. ($\Leftarrow$) Per la proposizione precedente $x = \sum_{v \in K^0} t_v(x)\, v = \sum_{v \in K^0} t_v(y)\, v = y$\hfill$\square$
\end{proof}

\begin{prop}
$|K|$ \`e uno spazio di Hausdorff ($T_2$).
\end{prop}
\begin{proof}
Siano $x, y \in |K|$, $x \neq y$. Per il lemma esiste $v \in K^0$ con $t_v(x) \neq t_v(y)$; senza perdita di generalit\`a $t_v(x) < t_v(y)$, e sia $r \in \mathbb{R}$ con $t_v(x) < r < t_v(y)$. Poich\'e $t_v$ \`e continua,
\[
A = t_v^{-1}\big(\,]-\infty, r[\,\big), \qquad B = t_v^{-1}\big(\,]r, +\infty[\,\big)
\]
sono aperti di $|K|$, con $x \in A$, $y \in B$ e $A \cap B = \varnothing$\hfill$\square$
\end{proof}

(Naturalmente $|K| \subset \mathbb{R}^n$ \`e Hausdorff anche perch\'e sottospazio di $\mathbb{R}^n$; la dimostrazione precedente \`e utile perch\'e usa solo le coordinate baricentriche.)

\begin{definicion}[Stella di un vertice]
Sia $K$ un complesso simpliciale e $v \in K^0$. Si chiama \textbf{stella} di $v$ l'insieme
\[
\operatorname{st}(v) = \bigcup_{\substack{(s) \in K \\ v \in [s]}} (s),
\]
cio\`e l'unione di tutti i simplessi aperti di $K$ che hanno $v$ come vertice.
\end{definicion}

\begin{intuizione}
La stella di $v$ \`e l'``intorno naturale'' di $v$ nel complesso: i punti che si raggiungono da $v$ muovendosi dentro i simplessi (aperti) che contengono $v$ come vertice, \emph{senza} le facce opposte a $v$.
\end{intuizione}

\begin{ejemplo}
\begin{enumerate}
    \item $K = \{ (v_0), (v_1), (v_0, v_1) \}$ (un segmento). Allora $\operatorname{st}(v_0) = (v_0) \cup (v_0, v_1) = [v_0, v_1] \setminus \{v_1\}$.
    \item $K$ = decomposizione del triangolo $[v_0, v_1, v_2]$. Allora
    \[
    \operatorname{st}(v_0) = (v_0) \cup (v_0, v_1) \cup (v_0, v_2) \cup (v_0, v_1, v_2) = [v_0, v_1, v_2] \setminus [v_1, v_2],
    \]
    cio\`e il triangolo privato del lato opposto a $v_0$.
    \item Sia $K$ formato da un triangolo pieno $[v_0, v_1, v_2]$ a cui sono attaccati in $v_1$ due spigoli $(v_1, v_3)$ e $(v_1, v_4)$, con $v_4$ vertice di un secondo triangolo pieno $[v_4, v_5, v_6]$. Allora $\operatorname{st}(v_0) = [v_0, v_1, v_2] \setminus [v_1, v_2]$; $\operatorname{st}(v_3) = (v_3) \cup (v_1, v_3) = [v_1, v_3] \setminus \{v_1\}$; e $\operatorname{st}(v_1)$ \`e l'unione di $(v_1)$, degli spigoli $(v_0,v_1)$, $(v_1,v_2)$, $(v_1, v_3)$, $(v_1, v_4)$ e del triangolo aperto $(v_0, v_1, v_2)$.
\end{enumerate}
\end{ejemplo}

\begin{figure}[H]
    \centering
    \begin{tikzpicture}[scale=1.3]
        \begin{scope}
            \fill[blue!15] (0,0) -- (2.2,0) -- (0.9,1.7) -- cycle;
            \draw[thick, blue] (0,0) -- (2.2,0) -- (0.9,1.7) -- cycle;
            \draw[thick, dashed, red] (2.2,0) -- (0.9,1.7);
            \filldraw (0,0) circle (2pt) node[below left] {$v_0$};
            \filldraw[white, draw=black] (2.2,0) circle (2pt) node[below right] {$v_1$};
            \filldraw[white, draw=black] (0.9,1.7) circle (2pt) node[above] {$v_2$};
            \node at (1.1,-0.9) {$\operatorname{st}(v_0)$: il lato $[v_1,v_2]$ \`e escluso};
        \end{scope}
    \end{tikzpicture}
    \caption{La stella di $v_0$ in un triangolo: tutto tranne il lato opposto (cerchi vuoti e tratteggio indicano punti esclusi).}
\end{figure}

\begin{prop}
Sia $K$ un complesso simpliciale e $v \in K^0$. Allora:
\begin{enumerate}
    \item[(1)] $\operatorname{st}(v)$ \`e un aperto di $|K|$;
    \item[(2)] se $v' \in K^0$ e $v' \in \operatorname{st}(v)$, allora $v' = v$ (la stella di $v$ contiene un solo vertice: $v$ stesso).
\end{enumerate}
\end{prop}
\begin{proof}
(1) Si ha $\operatorname{st}(v) = t_v^{-1}\big(\,]0, +\infty[\,\big)$: infatti $t_v(w) > 0$ se e solo se $v$ \`e un vertice dell'unico simplesso aperto $(s)$ che contiene $w$, cio\`e $w \in \operatorname{st}(v)$. Essendo $t_v$ continua, $\operatorname{st}(v)$ \`e aperto.

(2) Se $v' \in \operatorname{st}(v)$, esiste un unico $(s) \in K$ con $v' \in (s)$ e $v \in [s]$. Poich\'e $v'$ \`e un vertice, l'unico simplesso aperto che lo contiene \`e $(v')$, quindi $(s) = (v')$, e $v \in [s] = [v'] = \{v'\}$, cio\`e $v = v'$\hfill$\square$
\end{proof}

\begin{prop}
$\{ \operatorname{st}(v) \}_{v \in K^0}$ \`e un ricoprimento aperto e finito di $|K|$.
\end{prop}
\begin{proof}
Ogni $w \in |K|$ sta in un simplesso aperto $(s)$, che ha almeno un vertice $v$; allora $w \in \operatorname{st}(v)$. I $\operatorname{st}(v)$ sono aperti per la proposizione precedente e sono in numero finito perch\'e $K^0$ \`e finito\hfill$\square$
\end{proof}


% ======== NUOVO (pp. 147-152): Richiami sui gruppi abeliani, commutatori, simplessi orientati ========
\section{Richiami sui gruppi abeliani}

Sia $(G, +)$ un gruppo abeliano. Per ogni $x \in G$ e $k \in \mathbb{Z}$ poniamo
\[
kx = \begin{cases} \underbrace{x + \dots + x}_{k} & \text{se } k > 0, \\ 0 & \text{se } k = 0, \\ \underbrace{(-x) + \dots + (-x)}_{|k|} & \text{se } k < 0. \end{cases}
\]

\begin{definicion}[Generatori, indipendenza, base]
\begin{itemize}
    \item Una famiglia $(x_i)_{i \in I} \subseteq G$ si dice \textbf{insieme di generatori} per $G$ se per ogni $x \in G$ esistono $F \subseteq I$ finito, $F = \{j_1, \dots, j_n\}$, e $k_1, \dots, k_n \in \mathbb{Z}$ tali che $x = \sum_{i=1}^n k_i x_{j_i}$.
    \item $(x_i)_{i \in I}$ si dice \textbf{libero} (o linearmente indipendente) se per ogni $F \subseteq I$ finito, $F = \{j_1, \dots, j_n\}$, e per ogni $k_1, \dots, k_n \in \mathbb{Z}$: $\sum_{i=1}^n k_i x_{j_i} = 0 \implies k_i = 0$ per ogni $i = 1, \dots, n$.
    \item $(x_i)_{i \in I}$ si dice \textbf{base} di $G$ se \`e un sistema di generatori ed \`e libero. Un gruppo che ammette una base si dice \textbf{libero}.
    \item $(G, +)$ si dice \textbf{finitamente generato} se $G$ ammette un sistema di generatori finito.
\end{itemize}
Se $(G, +)$ \`e libero, tutte le sue basi hanno la stessa cardinalit\`a, che prende il nome di \textbf{rango} di $G$.
\end{definicion}

\begin{ejemplo}
$\mathbb{Z}^n$ \`e libero di rango $n$ (base canonica); $\mathbb{Z}_2$ \`e finitamente generato ma \emph{non} libero (l'elemento $1$ genera, ma $2 \cdot 1 = 0$ con $2 \neq 0$).
\end{ejemplo}

\begin{definicion}[Gruppo ciclico]
$(G, +)$ si dice \textbf{ciclico} se \`e generato da un unico elemento $x \in G$. A meno di isomorfismi: esiste un unico gruppo ciclico infinito, $\mathbb{Z}$, ed esiste un unico gruppo ciclico finito di ordine $n$, $\mathbb{Z}_n$, per ogni $n \in \mathbb{N}$, $n \geqslant 1$.
\end{definicion}

\begin{definicion}[Elemento di torsione]
Sia $(G, +)$ un gruppo abeliano. Un elemento $x \in G$ si dice \textbf{elemento di torsione} se esiste $k \in \mathbb{N}$, $k \geqslant 1$, tale che $kx = 0$.
\end{definicion}

\begin{prop}
L'insieme $\Gamma = \{ x \in G \mid x \text{ \`e di torsione} \}$ \`e un sottogruppo di $G$ (detto \textbf{sottogruppo di torsione}).
\end{prop}
\begin{proof}
Basta mostrare che per ogni $x, y \in \Gamma$ si ha $x - y \in \Gamma$. Se $kx = 0$ e $hy = 0$ con $k, h \geqslant 1$, allora $kh \geqslant 1$ e $kh(x - y) = h(kx) - k(hy) = h \cdot 0 - k \cdot 0 = 0$\hfill$\square$
\end{proof}

\begin{definicion}
Se $\Gamma$ \`e finito, la sua cardinalit\`a si chiama \textbf{ordine} di $\Gamma$. Se $\Gamma = \{0\}$, il gruppo $(G, +)$ si dice \textbf{privo di torsione}.
\end{definicion}

\begin{prop}[Propriet\`a]
\begin{enumerate}
    \item[a)] $G$ libero $\implies$ $G$ privo di torsione.
    \item[b)] Se $G$ \`e finitamente generato, allora $G$ privo di torsione $\implies$ $G$ libero. (Negli appunti questo punto compare scritto come la stessa implicazione del punto a); la forma corretta e standard \`e il viceversa, valido per i gruppi finitamente generati.)
    \item[c)] $G = G_1 \oplus G_2 \implies G/G_1 \cong G_2$.
    \item[d)] Se $G$ \`e finitamente generato e $H \leqslant G$, allora $G/H$ \`e finitamente generato.
\end{enumerate}
\end{prop}

\begin{teo}[Teorema fondamentale sui gruppi abeliani di tipo finito]
Sia $(G, +)$ un gruppo abeliano finitamente generato. Allora esiste un sottogruppo \textbf{libero} $G_0$ di $G$ tale che
\[
G = G_0 \oplus \Gamma ,
\]
dove $\Gamma$ \`e il sottogruppo di torsione. Si chiama \textbf{rango} di $G$ il rango di $G_0$. Inoltre $\Gamma$ \`e somma diretta di un numero finito $k$ di gruppi ciclici $\Gamma_1, \dots, \Gamma_k$ di ordini $t_1, \dots, t_k$ rispettivamente (nella forma standard, con $t_i > 1$ e $t_1 \mid t_2 \mid \dots \mid t_k$). In conclusione
\[
G \cong \mathbb{Z}^r \oplus \mathbb{Z}_{t_1} \oplus \dots \oplus \mathbb{Z}_{t_k}, \qquad r = \operatorname{rango}(G).
\]
\end{teo}
\begin{observacion}
Il testo del teorema nel PDF si interrompe a fine pagina dopo ``$\Gamma_1, \dots, \Gamma_k$ di ordine $t_1, \dots, t_k$ rispettivamente, con \dots''; qui \`e stata completata con l'enunciato standard (fattori invarianti).
\end{observacion}

\subsection{Commutatore e abelianizzato}

\begin{definicion}[Commutatore]
Sia $(G, \cdot)$ un gruppo e $a, b \in G$. L'elemento $[a, b] = a b a^{-1} b^{-1}$ si chiama \textbf{commutatore} di $a$ e $b$. Esso ``stima di quanto $a$ e $b$ siano lontani dal commutare'', perch\'e $[a,b] = e \iff ab = ba$. Si chiama \textbf{commutatore} di $G$ l'insieme $[G, G] = \{ [a,b] \mid a, b \in G \}$, e \textbf{derivato} di $G$ il sottogruppo $G'$ generato da $[G, G]$.
\end{definicion}

\begin{observacion}
\begin{enumerate}
    \item Ogni $x \in G'$ \`e il prodotto di un numero finito di commutatori (e dei loro inversi, che sono ancora commutatori: $[a,b]^{-1} = [b,a]$).
    \item $G'$ \`e un sottogruppo \textbf{normale} di $G$ (il coniugato di un commutatore \`e un commutatore: $g[a,b]g^{-1} = [gag^{-1}, gbg^{-1}]$).
    \item $G$ \`e abeliano $\iff$ $[G, G] = \{e\}$ (cio\`e $G' = \{e\}$).
\end{enumerate}
\end{observacion}

\begin{prop}
$G/G'$ \`e un gruppo abeliano, detto \textbf{abelianizzato} di $G$. Vale
\[
G' = \bigcap_{\substack{H \trianglelefteq G \\ G/H \text{ abeliano}}} H ,
\]
quindi $G/G'$ \`e il \textbf{pi\`u grande quoziente abeliano} di $G$.
\end{prop}
\begin{intuizione}
Abelianizzare significa ``far commutare'' tutto, cio\`e identificare $ab$ con $ba$. Ad esempio, l'abelianizzato del gruppo libero $F_2 = \langle a, b \rangle$ \`e $\mathbb{Z}^2$: \`e questo il motivo per cui, a differenza di $\Pi_1$, i gruppi di omologia che costruiremo sono sempre abeliani.
\end{intuizione}

\section{Simplessi orientati}

Sia $[s] = [v_0, \dots, v_k]$ un $k$-simplesso e siano $(i_0, \dots, i_k)$ e $(j_0, \dots, j_k)$ due permutazioni di $\{0, \dots, k\}$. Poniamo
\[
(i_0, \dots, i_k) \sim (j_0, \dots, j_k) \iff (j_0, \dots, j_k) \text{ \`e una permutazione \emph{pari} di } (i_0, \dots, i_k),
\]
cio\`e si ottiene da essa con un numero pari di trasposizioni. Se $P$ \`e l'insieme di tutte le permutazioni, allora per ogni $k \geqslant 1$ il quoziente $P/\!\sim$ ha esattamente \textbf{2 elementi} (le permutazioni pari e le dispari).

\begin{definicion}[Simplesso orientato]
Si chiama \textbf{$k$-simplesso orientato} e si indica con $\langle s \rangle = \langle v_{j_0}, \dots, v_{j_k} \rangle$ un $k$-simplesso $[s] = [v_0, \dots, v_k]$ munito di una classe fissata in $P/\!\sim$. Ciascun elemento di $P/\!\sim$ si chiama \textbf{orientazione} del $k$-simplesso.
\end{definicion}

Se $(i_0, \dots, i_k)$ \`e un'altra permutazione di $\{0, \dots, k\}$, allora
\[
\langle v_{i_0}, \dots, v_{i_k} \rangle = \begin{cases} \ \ \langle v_{j_0}, \dots, v_{j_k} \rangle & \text{se } (i_0, \dots, i_k) \sim (j_0, \dots, j_k), \\ -\langle v_{j_0}, \dots, v_{j_k} \rangle & \text{se } (i_0, \dots, i_k) \not\sim (j_0, \dots, j_k). \end{cases}
\]

\begin{intuizione}
Orientare un simplesso significa scegliere un ``verso di percorrenza'': per un segmento, quale estremo \`e l'inizio ($\langle v_0, v_1 \rangle = - \langle v_1, v_0 \rangle$); per un triangolo, il verso orario o antiorario ($\langle v_0, v_1, v_2 \rangle = \langle v_1, v_2, v_0 \rangle = - \langle v_1, v_0, v_2 \rangle$); per un tetraedro, la ``chiralit\`a'' (destra o sinistra). Scambiare due vertici inverte l'orientazione.
\end{intuizione}

\begin{figure}[H]
    \centering
    \begin{tikzpicture}[scale=1.2]
        \begin{scope}
            \draw[thick, ->, >=stealth, blue] (0,0) -- (2,0);
            \filldraw (0,0) circle (1.6pt) node[below] {$v_0$};
            \filldraw (2,0) circle (1.6pt) node[below] {$v_1$};
            \node at (1,-0.9) {$\langle v_0, v_1 \rangle$};
        \end{scope}
        \begin{scope}[shift={(4.2,0)}]
            \draw[thick, ->, >=stealth, red] (2,0) -- (0,0);
            \filldraw (0,0) circle (1.6pt) node[below] {$v_0$};
            \filldraw (2,0) circle (1.6pt) node[below] {$v_1$};
            \node at (1,-0.9) {$\langle v_1, v_0 \rangle = -\langle v_0, v_1 \rangle$};
        \end{scope}
        \begin{scope}[shift={(8.8,0)}]
            \fill[blue!10] (0,0) -- (2,0) -- (1,1.7) -- cycle;
            \draw[thick, blue] (0,0) -- (2,0) -- (1,1.7) -- cycle;
            \draw[->, >=stealth, thick, blue] (1.35,0.55) arc (-20:290:0.35);
            \filldraw (0,0) circle (1.6pt) node[below left] {$v_0$};
            \filldraw (2,0) circle (1.6pt) node[below right] {$v_1$};
            \filldraw (1,1.7) circle (1.6pt) node[above] {$v_2$};
            \node at (1,-0.9) {$\langle v_0, v_1, v_2 \rangle$ (antiorario)};
        \end{scope}
    \end{tikzpicture}
    \caption{Orientazione di un $1$-simplesso e di un $2$-simplesso.}
\end{figure}


% ======== NUOVO (pp. 153-156) ========
\begin{observacion}
Per ogni $k \geqslant 1$: $-\langle v_{j_0}, v_{j_1}, \dots, v_{j_k} \rangle = \langle v_{j_1}, v_{j_0}, \dots, v_{j_k} \rangle$ (scambiare due vertici \`e una permutazione dispari: inverte l'orientazione).
\end{observacion}

\begin{ejemplo}
\begin{enumerate}
    \item $k = 0$: $[v_0] = \{v_0\} = \langle v_0 \rangle$ ha un'unica orientazione naturale.
    \item $k = 1$: $[v_0, v_1] = \overline{v_0 v_1}$. Le permutazioni di $\{0,1\}$ sono $(0,1)$ e $(1,0)$ e $(0,1) \not\sim (1,0)$. Quindi ci sono due $1$-simplessi orientati: $\langle v_0, v_1 \rangle$ (percorso da $v_0$ a $v_1$) e $\langle v_1, v_0 \rangle = -\langle v_0, v_1 \rangle$ (percorso nel verso opposto).
    \item $k = 2$: a $[v_0, v_1, v_2]$ corrispondono i due $2$-simplessi orientati $\langle v_0, v_1, v_2 \rangle$ e $-\langle v_0, v_1, v_2 \rangle = \langle v_1, v_0, v_2 \rangle$ (i due versi di percorrenza del bordo del triangolo).
\end{enumerate}
In generale, un $k$-simplesso orientato $\langle s \rangle$ \`e il corrispondente $k$-simplesso chiuso in cui, avendo specificato un ordinamento dei vertici (e quindi dei $k$ vettori linearmente indipendenti $v_{j_1} - v_{j_0}, \dots, v_{j_k} - v_{j_0}$), \`e specificata un'\textbf{orientazione del $k$-piano} $\Pi \supseteq [s]$.
\end{ejemplo}

\section{Il gruppo delle $p$-catene}

Sia $K$ un complesso simpliciale, $m = \dim K$, $0 \leqslant p \leqslant m$, e sia $l_p$ il numero di $p$-simplessi di $K$.

\begin{definicion}[$p$-catene]
Si chiama \textbf{$p$-catena} di $K$ ogni combinazione lineare formale degli $l_p$ $p$-simplessi orientati di $K$ a coefficienti interi. Quindi, se $\langle s_1 \rangle, \dots, \langle s_{l_p} \rangle$ sono gli $l_p$ $p$-simplessi di $K$, ciascuno con una orientazione specificata, allora
\[
c \text{ \`e una $p$-catena di } K \iff \exists\, a_1, \dots, a_{l_p} \in \mathbb{Z} : \ c = \sum_{i=1}^{l_p} a_i \langle s_i \rangle .
\]
Per ogni $a \in \mathbb{Z}$ e ogni $p$-simplesso orientato $\langle v_0, \dots, v_p \rangle$:
\[
a \langle v_0, \dots, v_p \rangle = \begin{cases} \langle v_0, \dots, v_p \rangle \text{ sommato } a \text{ volte} & \text{se } a > 0, \\ 0 & \text{se } a = 0, \\ -\langle v_0, \dots, v_p \rangle \text{ sommato } |a| \text{ volte} & \text{se } a < 0. \end{cases}
\]
\end{definicion}

Poniamo
\[
C_p(K, \mathbb{Z}) = \{ c \mid c \text{ \`e una $p$-catena di } K \text{ a coefficienti in } \mathbb{Z} \} \qquad \text{(gruppo delle $p$-catene a coefficienti in $\mathbb{Z}$)}.
\]
\`E un gruppo abeliano finitamente generato rispetto all'operazione: per $c = \sum a_i \langle s_i \rangle$ e $c' = \sum b_i \langle s_i \rangle$,
\[
c + c' = \sum_{i=1}^{l_p} (a_i + b_i) \langle s_i \rangle .
\]
L'insieme $\{ \langle s_i \rangle \}_{i = 1, \dots, l_p}$ \`e un sistema finito di generatori di $C_p(K, \mathbb{Z})$ e, anzi, una \textbf{base}: $C_p(K, \mathbb{Z}) \cong \mathbb{Z}^{l_p}$ \`e un gruppo abeliano \emph{libero di rango $l_p$}. (Convenzione: $\langle s \rangle$ con l'orientazione opposta \`e $-\langle s \rangle$, quindi per ogni simplesso si sceglie un'orientazione di riferimento e le altre non danno nuovi generatori.)

Analogamente $C_p(K, \mathbb{R}) = \{ c \mid c \text{ \`e una $p$-catena di } K \text{ a coefficienti reali} \}$, con somma definita come prima. Inoltre, per ogni $\lambda \in \mathbb{R}$ e $c = \sum a_i \langle s_i \rangle \in C_p(K, \mathbb{R})$ si pone $\lambda c = \sum (\lambda a_i) \langle s_i \rangle$: $(C_p(K, \mathbb{R}), +, \cdot)$ \`e uno \textbf{spazio vettoriale reale} di dimensione finita $l_p$, avente $\{ \langle s_i \rangle \}$ come base.

\medskip
\textbf{Convenzione.} Per ogni $p < 0$ e per ogni $p > m = \dim K$: $C_p(K, \mathbb{Z}) = 0$ (e $C_p(K, \mathbb{R}) = 0$).

\begin{ejemplo}
Sia $K$ la decomposizione del triangolo $[v_0, v_1, v_2]$: $l_0 = 3$, $l_1 = 3$, $l_2 = 1$, quindi $C_0 \cong \mathbb{Z}^3$ (base $\langle v_0 \rangle, \langle v_1 \rangle, \langle v_2 \rangle$), $C_1 \cong \mathbb{Z}^3$ (base $\langle v_0, v_1 \rangle, \langle v_1, v_2 \rangle, \langle v_0, v_2 \rangle$), $C_2 \cong \mathbb{Z}$ (base $\langle v_0, v_1, v_2 \rangle$), e $C_p = 0$ per $p \geqslant 3$.
\end{ejemplo}

\section{L'operatore bordo}

\begin{definicion}[Bordo di un simplesso orientato]
Sia $\langle s \rangle = \langle v_0, \dots, v_p \rangle$ un $p$-simplesso orientato. Si chiama \textbf{bordo} di $\langle s \rangle$ la $(p-1)$-catena
\[
\partial_p \langle s \rangle = \sum_{i=0}^p (-1)^i \langle v_0, \dots, \widehat{v_i}, \dots, v_p \rangle ,
\]
dove $\widehat{v_i}$ indica che il vertice $v_i$ \`e omesso. Se $c \in C_p(K)$, $c = \sum_{i=1}^{l_p} a_i \langle s_i \rangle$, si chiama \textbf{bordo} di $c$ la $(p-1)$-catena
\[
\partial_p c = \sum_{i=1}^{l_p} a_i\, \partial_p \langle s_i \rangle .
\]
\end{definicion}

\begin{intuizione}
Il bordo di un segmento orientato $\langle v_0, v_1 \rangle$ \`e $\langle v_1 \rangle - \langle v_0 \rangle$ (``punto finale meno punto iniziale''); il bordo di un triangolo orientato \`e la somma dei suoi tre lati percorsi coerentemente con l'orientazione; i segni $(-1)^i$ servono a far s\`i che i lati siano orientati ``in modo coerente''.
\end{intuizione}

\begin{ejemplo}
\begin{itemize}
    \item $\partial_1 \langle v_0, v_1 \rangle = \langle v_1 \rangle - \langle v_0 \rangle$.
    \item $\partial_2 \langle v_0, v_1, v_2 \rangle = \langle v_1, v_2 \rangle - \langle v_0, v_2 \rangle + \langle v_0, v_1 \rangle$, cio\`e la somma dei tre lati percorsi $v_0 \to v_1 \to v_2 \to v_0$ (perch\'e $-\langle v_0, v_2 \rangle = \langle v_2, v_0 \rangle$).
    \item $\partial_3 \langle v_0, v_1, v_2, v_3 \rangle = \langle v_1, v_2, v_3 \rangle - \langle v_0, v_2, v_3 \rangle + \langle v_0, v_1, v_3 \rangle - \langle v_0, v_1, v_2 \rangle$ (le quattro facce).
\end{itemize}
\end{ejemplo}

\begin{figure}[H]
    \centering
    \begin{tikzpicture}[scale=1.3]
        \fill[blue!10] (0,0) -- (2.6,0) -- (1.0,2.0) -- cycle;
        \draw[thick, ->, >=stealth, blue] (0,0) -- (1.3,0);
        \draw[thick, blue] (1.3,0) -- (2.6,0);
        \draw[thick, ->, >=stealth, blue] (2.6,0) -- (1.8,1.0);
        \draw[thick, blue] (1.8,1.0) -- (1.0,2.0);
        \draw[thick, ->, >=stealth, blue] (1.0,2.0) -- (0.5,1.0);
        \draw[thick, blue] (0.5,1.0) -- (0,0);
        \draw[->, >=stealth, thick, red] (1.45,0.55) arc (-30:300:0.28);
        \filldraw (0,0) circle (1.6pt) node[below left] {$v_0$};
        \filldraw (2.6,0) circle (1.6pt) node[below right] {$v_1$};
        \filldraw (1.0,2.0) circle (1.6pt) node[above] {$v_2$};
        \node[below] at (1.3,-0.15) {$\partial \langle v_0,v_1,v_2\rangle = \langle v_0,v_1\rangle + \langle v_1,v_2\rangle + \langle v_2,v_0\rangle$};
    \end{tikzpicture}
    \caption{Il bordo di un triangolo orientato \`e il cammino chiuso formato dai suoi tre lati, percorsi nel verso indotto dall'orientazione.}
\end{figure}


% ======== NUOVO (pp. 157-158) ========
Per costruzione, $\partial_p: C_p(K, \mathbb{Z}) \to C_{p-1}(K, \mathbb{Z})$ \`e un \textbf{omomorfismo di gruppi} e $\partial_p: C_p(K, \mathbb{R}) \to C_{p-1}(K, \mathbb{R})$ \`e un'\textbf{applicazione lineare}.

\begin{observacion}[Buona definizione]
Il bordo $\partial_p \langle v_0, \dots, v_p \rangle$ \`e stato scritto scegliendo un ordinamento dei vertici; bisogna controllare che cambiando orientazione si ottenga il segno opposto, in modo da avere $\partial_p(-\langle s \rangle) = -\partial_p \langle s \rangle$. Basta verificarlo per la trasposizione di due vertici \emph{adiacenti}, $v_j \leftrightarrow v_{j+1}$ (ogni permutazione \`e prodotto di tali trasposizioni): i termini del bordo che omettono un vertice diverso da $v_j, v_{j+1}$ contengono la coppia scambiata e quindi cambiano segno; i due termini che omettono $v_j$ o $v_{j+1}$ sono
\[
(-1)^j \langle \dots, v_{j+1}, \dots \rangle + (-1)^{j+1} \langle \dots, v_j, \dots \rangle \quad \longmapsto \quad (-1)^j \langle \dots, v_j, \dots \rangle + (-1)^{j+1} \langle \dots, v_{j+1}, \dots \rangle ,
\]
che \`e l'opposto del primo. Dunque il bordo cambia segno.
\end{observacion}

\begin{ejemplo}
\begin{enumerate}
    \item $p = 0$: $C_{-1}(K) = 0$, quindi $\partial_0: C_0(K) \to C_{-1}(K) = 0$ \`e l'applicazione nulla ($c \mapsto 0$).
    \item $p = 1$, $\langle s \rangle = \langle v_0, v_1 \rangle$: $\partial_1 \langle v_0, v_1 \rangle = (-1)^0 \langle v_1 \rangle + (-1)^1 \langle v_0 \rangle = \langle v_1 \rangle - \langle v_0 \rangle$; e $\partial_1 \langle v_1, v_0 \rangle = \langle v_0 \rangle - \langle v_1 \rangle$.
    \item $p = 2$, $\langle s \rangle = \langle v_0, v_1, v_2 \rangle$:
    \[
    \partial_2 \langle s \rangle = (-1)^0 \langle v_1, v_2 \rangle + (-1)^1 \langle v_0, v_2 \rangle + (-1)^2 \langle v_0, v_1 \rangle = \langle v_0, v_1 \rangle + \langle v_1, v_2 \rangle + \langle v_2, v_0 \rangle,
    \]
    perch\'e $-\langle v_0, v_2 \rangle = \langle v_2, v_0 \rangle$.
\end{enumerate}
\end{ejemplo}

\begin{lema}
$\partial^2 = 0$, cio\`e $\partial_{p-1} \circ \partial_p = 0$ per ogni $p$.
\end{lema}

\begin{intuizione}
``Il bordo di un bordo \`e vuoto'': il bordo di un triangolo \`e un cammino chiuso, e un cammino chiuso non ha bordo (ogni vertice compare una volta come punto finale e una come punto iniziale, con segni opposti). \`E l'identit\`a che permette di definire l'omologia.
\end{intuizione}

\begin{proof}
Poich\'e $\partial_p$ \`e un omomorfismo di gruppi, basta mostrare che $(\partial_{p-1} \circ \partial_p) \langle s \rangle = 0$ per ogni $p$-simplesso orientato $\langle s \rangle = \langle v_0, \dots, v_p \rangle$. Si ha
\[
\partial_{p-1} \partial_p \langle s \rangle = \sum_{i=0}^p (-1)^i\, \partial_{p-1} \langle v_0, \dots, \widehat{v_i}, \dots, v_p \rangle .
\]
Nel simplesso $\langle v_0, \dots, \widehat{v_i}, \dots, v_p \rangle$ i vertici con indice $j < i$ hanno posizione $j$, quelli con $j > i$ hanno posizione $j - 1$. Quindi
\[
\partial_{p-1} \langle v_0, \dots, \widehat{v_i}, \dots, v_p \rangle = \sum_{j < i} (-1)^j \langle \dots, \widehat{v_j}, \dots, \widehat{v_i}, \dots \rangle + \sum_{j > i} (-1)^{j-1} \langle \dots, \widehat{v_i}, \dots, \widehat{v_j}, \dots \rangle .
\]
Perci\`o
\[
\partial_{p-1} \partial_p \langle s \rangle = \sum_{j < i} (-1)^{i+j} \langle \dots, \widehat{v_j}, \dots, \widehat{v_i}, \dots \rangle + \sum_{i < j} (-1)^{i+j-1} \langle \dots, \widehat{v_i}, \dots, \widehat{v_j}, \dots \rangle .
\]
Nella seconda somma si scambiano i nomi degli indici ($i \leftrightarrow j$): diventa $\sum_{j < i} (-1)^{i+j-1} \langle \dots, \widehat{v_j}, \dots, \widehat{v_i}, \dots \rangle$. Allora le due somme hanno gli stessi addendi con segni opposti e
\[
\partial_{p-1} \partial_p \langle s \rangle = \sum_{j < i} \big( (-1)^{i+j} + (-1)^{i+j-1} \big) \langle \dots, \widehat{v_j}, \dots, \widehat{v_i}, \dots \rangle = 0 . \qquad \square
\]
\end{proof}

\begin{ejemplo}
Per il triangolo: $\partial_1 \partial_2 \langle v_0, v_1, v_2 \rangle = \partial_1 \big( \langle v_0, v_1 \rangle + \langle v_1, v_2 \rangle + \langle v_2, v_0 \rangle \big) = (\langle v_1 \rangle - \langle v_0 \rangle) + (\langle v_2 \rangle - \langle v_1 \rangle) + (\langle v_0 \rangle - \langle v_2 \rangle) = 0$.
\end{ejemplo}


% ======== NUOVO (pp. 159-160) ========
\section{I gruppi di omologia simpliciale}

\begin{definicion}[Cicli e bordi]
Una $p$-catena $c \in C_p(K)$ si dice
\begin{itemize}
    \item \textbf{$p$-ciclo} se $\partial_p c = 0$;
    \item \textbf{$p$-bordo} se esiste $c' \in C_{p+1}(K)$ tale che $c = \partial_{p+1} c'$.
\end{itemize}
\end{definicion}

\begin{observacion}
Per il lemma $\partial^2 = 0$, ogni $p$-bordo \`e anche un $p$-ciclo: se $c = \partial_{p+1} c'$ allora $\partial_p c = \partial_p \partial_{p+1} c' = 0$.
\end{observacion}

\begin{definicion}[Gruppi dei cicli e dei bordi]
Il sottogruppo
\[
Z_p(K) = \{ c \in C_p(K) \mid c \text{ \`e un $p$-ciclo} \} = \operatorname{Ker} \partial_p
\]
si chiama \textbf{gruppo dei $p$-cicli} di $K$. Il sottogruppo
\[
B_p(K) = \{ c \in C_p(K) \mid c \text{ \`e un $p$-bordo} \} = \operatorname{Im} \partial_{p+1}
\]
si chiama \textbf{gruppo dei $p$-bordi}.
\end{definicion}

\begin{observacion}
Per il lemma, $B_p(K) \subseteq Z_p(K)$; poich\'e i gruppi sono abeliani, $B_p(K) \trianglelefteq Z_p(K)$.
\end{observacion}

\begin{definicion}[Omologia simpliciale]
Sia $K$ un complesso simpliciale e $0 \leqslant p \leqslant m = \dim K$. Si chiama \textbf{$p$-esimo gruppo di omologia simpliciale} di $K$ il gruppo quoziente
\[
H_p(K) = \frac{Z_p(K)}{B_p(K)} .
\]
\end{definicion}

\begin{intuizione}
I $p$-cicli sono le ``$p$-superfici senza bordo'' (ad esempio un cammino chiuso per $p=1$); i $p$-bordi sono quelli che si ottengono come bordo di qualcosa di dimensione $p+1$ \emph{presente} in $K$ (cicli che si possono ``riempire''). $H_p(K)$ misura i cicli che \emph{non} si riempiono, cio\`e i ``buchi $p$-dimensionali'' di $K$: componenti connesse per $p=0$, buchi circolari per $p=1$, cavit\`a per $p=2$.
\end{intuizione}

\begin{observacion}
\begin{enumerate}
    \item[0)] $H_p(K) = 0 \iff Z_p(K) = B_p(K) \iff$ ogni $p$-ciclo \`e un $p$-bordo.
    \item[1)] $H_p(K)$ \`e definito per ogni $p \in \mathbb{Z}$, con $H_p(K) = 0$ per $p < 0$ e per $p > \dim K$ (perch\'e $C_p(K) = 0$).
    \item[2)] $C_p(K, \mathbb{Z})$ \`e un gruppo abeliano libero di tipo finito (di rango $l_p$). Ne segue che $Z_p(K, \mathbb{Z}) \leqslant C_p(K, \mathbb{Z})$ \`e un gruppo abeliano libero di tipo finito (un sottogruppo di un gruppo abeliano libero di tipo finito \`e libero di tipo finito). Allora $H_p(K, \mathbb{Z}) = Z_p(K, \mathbb{Z}) / B_p(K, \mathbb{Z})$ \`e abeliano, di tipo finito, ma \textbf{non necessariamente libero}! Per il teorema fondamentale sui gruppi abeliani di tipo finito,
    \[
    H_p(K, \mathbb{Z}) = \beta_p(K, \mathbb{Z}) \oplus \Gamma_p(K, \mathbb{Z}),
    \]
    dove $\beta_p(K, \mathbb{Z})$ \`e libero e $\Gamma_p(K, \mathbb{Z})$ \`e il sottogruppo di torsione.
\end{enumerate}
\end{observacion}

\begin{definicion}[Gruppo di Betti, torsione e numeri di Betti]
\begin{itemize}
    \item $\beta_p(K, \mathbb{Z})$ si chiama \textbf{$p$-esimo gruppo di Betti} di $K$ (da non confondere con $B_p(K)$, il gruppo dei bordi);
    \item $\Gamma_p(K, \mathbb{Z})$ si chiama \textbf{$p$-esimo gruppo di torsione} di $K$;
    \item $\operatorname{rango}\big(H_p(K, \mathbb{Z})\big) = \operatorname{rango}\big(\beta_p(K, \mathbb{Z})\big) = b_p(K)$ si chiama \textbf{$p$-esimo numero di Betti} di $K$.
\end{itemize}
\end{definicion}

\begin{intuizione}
Il numero di Betti $b_p(K)$ conta \emph{quanti} buchi $p$-dimensionali indipendenti ha $K$; la torsione $\Gamma_p$ rileva fenomeni pi\`u sottili (come il ``mezzo giro'' del piano proiettivo, dove $H_1(\mathbb{R}P^2) = \mathbb{Z}_2$).
\end{intuizione}


% ======== NUOVO (pp. 161-166, 29/10/2026) ========
\begin{observacion}[Omologia a coefficienti reali]
$C_p(K, \mathbb{R})$ \`e uno spazio vettoriale reale di dimensione finita $l_p$; quindi $Z_p(K, \mathbb{R}) \subseteq C_p(K, \mathbb{R})$ \`e un sottospazio, e $H_p(K, \mathbb{R}) = Z_p(K, \mathbb{R}) / B_p(K, \mathbb{R})$ \`e uno spazio vettoriale reale di dimensione finita. $H_p(K, \mathbb{Z})$ \`e ``pi\`u ricco'' di $H_p(K, \mathbb{R})$ (contiene anche la torsione, che i coefficienti reali non vedono). Si dimostra che
\[
\dim H_p(K, \mathbb{R}) = \operatorname{rango} H_p(K, \mathbb{Z}) = b_p(K).
\]
\end{observacion}

\section{Componenti connesse di $|K|$}

Sia $K$ un complesso simpliciale e $v, w \in K^0$ vertici. Poniamo
\[
v \sim w \iff \exists\, n \in \mathbb{N}_0,\ \exists\, v_1, \dots, v_n \in K^0 : \ \langle v, v_1 \rangle + \langle v_1, v_2 \rangle + \dots + \langle v_{n-1}, v_n \rangle + \langle v_n, w \rangle \in C_1(K),
\]
cio\`e: $v$ e $w$ si possono collegare con una \textbf{spezzata} di spigoli di $K$ ($v \to v_1 \to \dots \to v_n \to w$, ogni tratto \`e un $1$-simplesso di $K$; per $n = 0$ \`e un solo spigolo, e si conviene $v \sim v$).

\begin{figure}[H]
    \centering
    \begin{tikzpicture}[scale=1.2]
        \draw[thick, blue] (0,0) -- (1,0.8) -- (2.1,0.3);
        \draw[thick, blue, dotted] (2.1,0.3) -- (3.2,0.5);
        \draw[thick, blue] (3.2,0.5) -- (4.2,1.0) -- (5.2,0.2);
        \foreach \p/\n/\d in {(0,0)/v/below, (1,0.8)/v_1/above, (2.1,0.3)/v_2/below, (3.2,0.5)/v_{n-1}/above, (4.2,1.0)/v_n/above, (5.2,0.2)/w/below}
            \filldraw \p circle (1.6pt) node[\d] {$\n$};
    \end{tikzpicture}
    \caption{$v \sim w$ se si pu\`o andare da $v$ a $w$ seguendo spigoli del complesso.}
\end{figure}

$\sim$ \`e una \textbf{relazione di equivalenza} su $K^0$:
\begin{itemize}
    \item \emph{riflessiva}: $\langle v, v \rangle = 0 \in C_1(K)$ (cammino di lunghezza zero), quindi $v \sim v$;
    \item \emph{simmetrica}: se $v \sim w$ si percorre la spezzata al contrario: $\langle w, v_n \rangle + \dots + \langle v_1, v \rangle \in C_1(K)$ perch\'e $\langle a, b \rangle = -\langle b, a \rangle$;
    \item \emph{transitiva}: se $v \sim w$ e $w \sim w'$ si concatenano le due spezzate.
\end{itemize}

\begin{observacion}
Se $K = \{ (v_0), (v_1) \}$ (due punti isolati, un complesso simpliciale), allora $v_0 \not\sim v_1$: i due vertici non sono connessi tra loro.
\end{observacion}

\begin{definicion}
Sia $K$ un complesso simpliciale e $v \in K^0$. Poniamo
\[
K_v = \bigcup_{\substack{w \in K^0 \\ w \sim v}} \operatorname{st}(w) = \bigcup_{\substack{(s) \in K \\ (s) \text{ ha un vertice} \sim v}} (s).
\]
\end{definicion}

\begin{prop}[Propriet\`a di $K_v$]
\begin{enumerate}
    \item[1)] Se $(s) \in K$ e $v, w$ sono vertici di $(s)$, allora $v \sim w$.
    \item[2)] $|K| = \bigcup_{v \in K^0} K_v$.
    \item[3)] Se $(s) \subseteq K_v$, allora $[s] \subseteq K_v$.
    \item[4)] $K_v$ \`e un aperto di $|K|$.
    \item[5)] $K_v \neq \varnothing$.
    \item[6)] $K_v = K_{v'} \iff v \sim v'$.
    \item[7)] $K_v$ \`e connesso per archi.
    \item[8)] $K_v \cap K_{v'} \neq \varnothing \iff K_v = K_{v'}$.
    \item[9)] $K_v$ \`e chiuso in $|K|$.
\end{enumerate}
\end{prop}

\begin{proof}
\textbf{1)} Se $(s) \in K$, allora $[s] \subseteq |K|$ e i suoi vertici $v = v_i$, $w = v_j$ sono collegati dal segmento $\overline{v_i v_j}$, che \`e una faccia (spigolo) di $[s]$, quindi un $1$-simplesso di $K$ (condizione 1 di complesso): $v \sim w$.

\textbf{2)} ``$\supseteq$'' \`e banale; ``$\subseteq$'': $|K| = \bigcup_{v \in K^0} \operatorname{st}(v) \subseteq \bigcup_{v} K_v$.

\textbf{3)} Se $(s) \subseteq K_v = \bigcup_{w \sim v} \operatorname{st}(w)$, allora $(s)$ ha un vertice $w \sim v$. Per 1), ogni vertice di $(s)$ \`e equivalente a $w$, e quindi a $v$. Dunque ogni faccia di $[s]$ ha un vertice equivalente a $v$ e $[s] = \bigcup_{(s') \text{ faccia di } (s)} (s') \subseteq K_v$.

\textbf{4)} $K_v = \bigcup_{w \sim v} \operatorname{st}(w)$ \`e unione di aperti (le stelle sono aperte).

\textbf{5)} $v \in K_v$.

\textbf{6)} ($\Rightarrow$) $K_v = K_{v'}$: allora $v \in K_{v'} = \bigcup_{w \sim v'} \operatorname{st}(w)$, quindi esiste $(s) \in K$ con $v \in (s)$ e $(s)$ ha un vertice $\sim v'$. Ma $v \in K^0$ implica $(s) = (v) = \{v\}$, quindi $v \sim v'$. ($\Leftarrow$) Se $v \sim v'$: $x \in K_v \iff \exists w \in K^0: x \in \operatorname{st}(w),\ w \sim v \iff \exists w: x \in \operatorname{st}(w),\ w \sim v' \iff x \in K_{v'}$.

\textbf{7)} Basta mostrare che per ogni $x \in K_v$ esiste un arco che congiunge $x$ a $v$. Sia $x \in K_v$: esiste $(s) \in K$ con $x \in (s)$ e $(s)$ ha un vertice $w \sim v$. Per 3), $[s] \subseteq K_v$ e $[s]$ \`e convesso, quindi $\overline{x, w} \subseteq [s] \subseteq K_v$. Inoltre $w \sim v$ d\`a una spezzata $v \to v_1 \to \dots \to v_n \to w$ di spigoli, con $v_i \sim v$ per ogni $i$; ciascuno spigolo ha un vertice $\sim v$, quindi sta in $K_v$ (per 3)). Allora $\gamma = \overline{x, w} \cup \overline{w, v_n} \cup \dots \cup \overline{v_1, v}$ \`e una spezzata contenuta in $K_v$ che congiunge $x$ a $v$.

\textbf{8)} ($\Leftarrow$) \`E banale per 5). ($\Rightarrow$) Sia $x \in K_v \cap K_{v'}$ e $(s)$ l'unico simplesso aperto con $x \in (s)$. Da $x \in K_v$, $(s)$ ha un vertice $w \sim v$; da $x \in K_{v'}$, ha un vertice $w' \sim v'$. Per 1), $w \sim w'$, quindi $v \sim w \sim w' \sim v'$ e per 6) $K_v = K_{v'}$.

\textbf{9)} Per 2) e 8), $K_v = |K| \setminus \bigcup_{v' \not\sim v} K_{v'}$, complementare di un'unione di aperti (4), dunque chiuso\hfill$\square$
\end{proof}

\begin{prop}
$\{ K_v \}_{v \in K^0}$ \`e l'insieme delle \textbf{componenti connesse (per archi)} di $|K|$.
\end{prop}
\begin{proof}
Per ogni $v \in K^0$, $K_v$ \`e aperto, chiuso, non vuoto e connesso per archi; inoltre due insiemi $K_v$ distinti sono disgiunti (propriet\`a 8) e ricoprono $|K|$ (propriet\`a 2). Dunque sono esattamente le componenti connesse per archi\hfill$\square$
\end{proof}

\begin{coro}
$|K|$ \`e connesso $\iff$ $v \sim v'$ per ogni $v, v' \in K^0$.
\end{coro}
\begin{proof}
Per la proposizione le componenti connesse di $|K|$ sono i $K_v$. Dunque $|K|$ \`e connesso $\iff |K| = K_v = \bigcup_{w \sim v} \operatorname{st}(w)$ per un (ogni) $v$ $\iff$ ogni $w \in K^0$ \`e equivalente a $v$\hfill$\square$
\end{proof}

\begin{intuizione}
$|K|$ \`e connesso se e solo se il ``grafo'' formato dai vertici e dagli spigoli di $K$ \`e connesso. Le componenti connesse sono gli ``isolotti'' separati.
\end{intuizione}

Poniamo $K'_v = \{ (s) \in K \mid (s) \subseteq K_v \} \subseteq K$.

\begin{prop}
$K'_v$ \`e un complesso simpliciale (con sostegno $|K'_v| = K_v$).
\end{prop}
\begin{proof}
La condizione 2) di complesso (due simplessi aperti distinti sono disgiunti) \`e banale perch\'e $K'_v \subseteq K$. Per la condizione 1): se $(s) \in K'_v$ e $(s')$ \`e una faccia aperta di $(s)$, allora $(s') \subseteq [s] \subseteq K_v$ (per 3)), cio\`e $(s') \in K'_v$\hfill$\square$
\end{proof}


% ======== NUOVO (pp. 167-174) ========
\begin{teo}
Sia $K$ un complesso simpliciale e siano $K_1, \dots, K_n$ i complessi $K'_{v}$ associati alle componenti connesse di $|K|$ (cio\`e $|K_i|$ sono le componenti connesse di $|K|$). Allora, per ogni $p$,
\[
H_p(K) \cong \bigoplus_{i=1}^n H_p(K_i) .
\]
\end{teo}
\begin{proof}
Per ogni $p$ vale $C_p(K) = C_p(K_1) \oplus \dots \oplus C_p(K_n)$: ogni $p$-simplesso di $K$ \`e contenuto in uno e uno solo dei sottoinsiemi $|K_i|$, quindi ogni $p$-catena di $K$ si decompone in modo unico come somma di $p$-catene dei vari $K_i$. Il bordo di un simplesso di $K_i$ \`e una catena di $K_i$ (le facce di un simplesso stanno nella stessa componente), quindi, essendo $\partial_p$ un omomorfismo,
\[
B_p(K) = \partial_{p+1}\big(C_{p+1}(K)\big) = \partial_{p+1}\big(C_{p+1}(K_1) \oplus \dots \oplus C_{p+1}(K_n)\big) = B_p(K_1) \oplus \dots \oplus B_p(K_n) .
\]
Inoltre, se $c \in Z_p(K)$, $c = c_1 + \dots + c_n$ con $c_i \in C_p(K_i)$, e $0 = \partial_p c = \partial_p c_1 + \dots + \partial_p c_n$, con $\partial_p c_i \in C_{p-1}(K_i)$ in addendi diversi di una somma diretta; ne segue $\partial_p c_i = 0$, cio\`e $c_i \in Z_p(K_i)$ per ogni $i$. Dunque $Z_p(K) = Z_p(K_1) \oplus \dots \oplus Z_p(K_n)$ e
\[
H_p(K) = \frac{Z_p(K)}{B_p(K)} = \frac{Z_p(K_1) \oplus \dots \oplus Z_p(K_n)}{B_p(K_1) \oplus \dots \oplus B_p(K_n)} \cong \frac{Z_p(K_1)}{B_p(K_1)} \oplus \dots \oplus \frac{Z_p(K_n)}{B_p(K_n)} = H_p(K_1) \oplus \dots \oplus H_p(K_n). \qquad \square
\]
\end{proof}

\section{Omologia $0$-dimensionale}

\begin{teo}
Sia $K$ un complesso simpliciale e sia $n$ il numero delle componenti connesse di $|K|$. Allora
\[
H_0(K, \mathbb{Z}) \cong \mathbb{Z}^n \qquad \big( \text{e } H_0(K, \mathbb{R}) \cong \mathbb{R}^n \big).
\]
\end{teo}

\begin{intuizione}
$H_0$ \emph{conta le componenti connesse}: $b_0(K) = n$. \`E il primo caso concreto in cui un gruppo di omologia misura una caratteristica geometrica elementare.
\end{intuizione}

\begin{proof}
Per il teorema precedente basta supporre $|K|$ connesso e provare che $H_0(K, \mathbb{R}) \cong \mathbb{R}$ (e $H_0(K, \mathbb{Z}) \cong \mathbb{Z}$). Poich\'e $\partial_0 = 0$, $Z_0 = C_0$ e $H_0(K, \mathbb{R}) = C_0(K, \mathbb{R}) / B_0(K, \mathbb{R})$. Sia $K^0 = \{ v_1, \dots, v_{l_0} \}$ e consideriamo
\[
\varphi: C_0(K, \mathbb{R}) \to \mathbb{R}, \qquad c = \sum_{i=1}^{l_0} \alpha_i \langle v_i \rangle \longmapsto \sum_{i=1}^{l_0} \alpha_i
\]
(la somma dei coefficienti). \`E un omomorfismo, ed \`e suriettivo perch\'e $\varphi(k \langle v_i \rangle) = k$ per ogni $k \in \mathbb{R}$. Per il teorema di isomorfismo $C_0(K, \mathbb{R}) / \operatorname{Ker} \varphi \cong \mathbb{R}$. Mostriamo che $\operatorname{Ker} \varphi = B_0(K, \mathbb{R})$.

$(\subseteq)$ Sia $c = \sum \alpha_i \langle v_i \rangle \in \operatorname{Ker} \varphi$, cio\`e $\sum \alpha_i = 0$. Fissato $v_0 \in K^0$:
\[
c = \sum_i \alpha_i \langle v_i \rangle - \Big( \sum_i \alpha_i \Big) \langle v_0 \rangle = \sum_{i} \alpha_i \big( \langle v_i \rangle - \langle v_0 \rangle \big).
\]
Poich\'e $|K|$ \`e connesso, $v_i \sim v_0$ per ogni $i$: esiste una spezzata di spigoli da $v_0$ a $v_i$, cio\`e una $1$-catena $c_i' \in C_1(K, \mathbb{R})$ con $\partial_1 c'_i = \langle v_i \rangle - \langle v_0 \rangle$ (la somma lungo il cammino ``telescopizza''). Allora $c = \sum_i \alpha_i \partial_1 c'_i = \partial_1 \big( \sum_i \alpha_i c'_i \big) \in B_0(K, \mathbb{R})$.

$(\supseteq)$ Se $c \in B_0(K, \mathbb{R})$, $c = \partial_1 c_1$ con $c_1 = \sum_j a_j \langle v_j, w_j \rangle$, quindi $c = \sum_j a_j (\langle w_j \rangle - \langle v_j \rangle)$ e $\varphi(c) = \sum a_j - \sum a_j = 0$.

Dunque $H_0(K, \mathbb{R}) = C_0(K, \mathbb{R}) / B_0(K, \mathbb{R}) = C_0 / \operatorname{Ker} \varphi \cong \mathbb{R}$\hfill$\square$
\end{proof}

\begin{ejemplo}
Se $K$ \`e un triangolo pieno (connesso), $H_0(K) \cong \mathbb{Z}$; se $K = \{ (v_0), (v_1) \}$ (due punti) allora $H_0(K) \cong \mathbb{Z}^2$.
\end{ejemplo}

\section{Invarianza topologica dei gruppi di omologia simpliciale}

\begin{teo}
Siano $K, L$ complessi simpliciali. Se $|K| \overset{\text{top}}{\cong} |L|$ (omeomorfi), allora $H_p(K) \overset{\text{alg}}{\cong} H_p(L)$ per ogni $p$.
\end{teo}

\begin{observacion}
\`E il risultato che giustifica l'uso dell'omologia come invariante topologico: \emph{non dipende dalla triangolazione scelta}, ma solo dallo spazio $|K|$. Gli appunti ne danno lo \textbf{schema} della dimostrazione; per i dettagli si veda S.~R.~Munkres, \emph{Elements of Algebraic Topology}.
\end{observacion}

\textbf{Schema della dimostrazione.}

\medskip
\textbf{I passo: applicazioni simpliciali e omologia.} Siano $K, L$ complessi simpliciali. Un'applicazione $f: |K| \to |L|$ si dice \textbf{applicazione simpliciale} se manda $p$-simplessi di $K$ in $q$-simplessi di $L$, con $q \leqslant p$. In particolare, per ogni $v \in K^0$ si ha $f(v) \in L^0$ e $f$ si ottiene estendendo ``linearmente'' la sua azione sui vertici:
\[
f: |K| \to |L|, \qquad x = \sum_{v \in K^0} t_v(x)\, v \ \longmapsto\ \sum_{v \in K^0} t_v(x)\, f(v) .
\]
Se $f$ \`e simpliciale, essa induce un omomorfismo di gruppi
\[
f_{\#p}: C_p(K) \to C_p(L), \qquad \langle v_0, \dots, v_p \rangle \longmapsto \begin{cases} 0 & \text{se } f(v_0), \dots, f(v_p) \text{ non sono c-indip.,} \\ \langle f(v_0), \dots, f(v_p) \rangle & \text{se sono c-indipendenti,} \end{cases}
\]
esteso per linearit\`a a tutte le $p$-catene. Si dimostra che $f_\#$ commuta con il bordo: $\partial_p \circ f_{\#p} = f_{\#, p-1} \circ \partial_p$ (diagramma commutativo). Ne segue facilmente che $f_{\#p}$ manda $p$-cicli in $p$-cicli e $p$-bordi in $p$-bordi, quindi induce un omomorfismo di gruppi
\[
f_{*p}: H_p(K) \to H_p(L), \qquad [c] \longmapsto [f_{\#p}\, c] .
\]

\textbf{II passo: applicazioni contigue.} Due applicazioni simpliciali $f, g: |K| \to |L|$ si dicono \textbf{contigue} se per ogni $p$-simplesso $\langle v_0, \dots, v_p \rangle$ di $K$ i vertici $f(v_0), \dots, f(v_p), g(v_0), \dots, g(v_p)$ generano un simplesso di $L$ (di dimensione al pi\`u $2p+1$). Si dimostra che se $f, g$ sono continue e contigue allora $f_{*p} = g_{*p}$ per ogni $p$ (con un argomento di omotopia di catene).

\textbf{III passo: approssimazione simpliciale.} Sia $f: |K| \to |L|$ continua. Si dimostra che esiste una suddivisione baricentrica $K' = K^{(n)}$ di $K$ tale che per ogni $v \in (K')^0$ esiste $u \in L^0$ con $f(\operatorname{st}(v)) \subseteq \operatorname{st}(u)$. (Si usa il numero di Lebesgue per il ricoprimento $\{ f^{-1}(\operatorname{st}(u)) \}_{u \in L^0}$ di $|K|$ e il fatto che la maglia di $K^{(n)}$ tende a $0$.) Si pone allora
\[
h: (K')^0 \to L^0, \quad v \mapsto u \text{ tale che } f(\operatorname{st}(v)) \subseteq \operatorname{st}(u),
\]
e si estende $h$ a tutto $|K'| = |K|$ ponendo $h(x) = \sum_{v \in (K')^0} t_v(x)\, h(v)$. L'applicazione $h$ si chiama \textbf{approssimazione simpliciale} di $f$. Si dimostra che se $h, g$ sono entrambe approssimazioni simpliciali di $f$, allora $h$ e $g$ sono contigue e quindi $h_{*p} = g_{*p}$ per ogni $p$. Inoltre:
\begin{enumerate}
    \item[a)] l'identit\`a (continua) $\mathrm{Id}: |K'| \to |K|$ ha un'approssimazione simpliciale $\Phi$ tale che $\Phi_{*p}: H_p(K') \to H_p(K)$ \`e un \textbf{isomorfismo} per ogni $p$ (l'omologia non cambia per suddivisione baricentrica).
\end{enumerate}
Per definizione si pone allora
\[
f_{*p} = h_{*p} \circ \Phi_{*p}^{-1}: \ H_p(K) \xrightarrow{\ \Phi_{*p}^{-1}\ } H_p(K') \xrightarrow{\ h_{*p}\ } H_p(L).
\]
\begin{enumerate}
    \item[b)] L'identit\`a $\mathrm{Id}: |K| \to |K|$ induce l'identit\`a $\mathrm{Id}_{*p}: H_p(K) \to H_p(K)$ per ogni $p$, e date $f: |K| \to |L|$, $g: |L| \to |M|$ continue, $(g \circ f)_{*p} = g_{*p} \circ f_{*p}$ (funtorialit\`a).
\end{enumerate}

\begin{figure}[H]
    \centering
    \begin{tikzpicture}[scale=1.2]
        \begin{scope}
            \fill[blue!10] (0,0) -- (3,0) -- (1.3,2.2) -- cycle;
            \draw[thick] (0,0) -- (3,0) -- (1.3,2.2) -- cycle;
            \filldraw (0,0) circle (1.5pt) node[below left] {$v_0$};
            \filldraw (3,0) circle (1.5pt) node[below right] {$v_1$};
            \filldraw (1.3,2.2) circle (1.5pt) node[above] {$v_2$};
            \node at (1.5,-0.7) {$K$};
        \end{scope}
        \begin{scope}[shift={(5.2,0)}]
            \fill[blue!10] (0,0) -- (3,0) -- (1.3,2.2) -- cycle;
            \draw[thick] (0,0) -- (3,0) -- (1.3,2.2) -- cycle;
            \draw[thick] (0,0) -- (2.15,1.1); \draw[thick] (3,0) -- (0.65,1.1); \draw[thick] (1.3,2.2) -- (1.5,0);
            \foreach \p in {(0,0),(3,0),(1.3,2.2),(1.5,0),(0.65,1.1),(2.15,1.1),(1.3,0.73)} \filldraw \p circle (1.3pt);
            \node at (1.5,-0.7) {$K' = K^{(1)}$};
        \end{scope}
        \draw[->, >=stealth, thick, Purple] (3.3,1.0) -- (4.9,1.0) node[midway, above] {$\Phi \simeq \mathrm{Id}$};
    \end{tikzpicture}
    \caption{Il ciclo $\langle v_0,v_1\rangle + \langle v_1,v_2\rangle + \langle v_2,v_0\rangle = \partial_2 \langle v_0,v_1,v_2\rangle$ si manda in modo naturale nel ciclo analogo di $K'$ (suddiviso): $\Phi_*$ \`e un isomorfismo.}
\end{figure}

\textbf{IV passo: conclusione.} Sia $f: |K| \to |L|$ un omeomorfismo: $f$ e $f^{-1}$ sono continue e $f^{-1} \circ f = \mathrm{Id}_{|K|}$, $f \circ f^{-1} = \mathrm{Id}_{|L|}$. Per la funtorialit\`a e il punto b),
\[
(f^{-1} \circ f)_{*p} = (\mathrm{Id}_{|K|})_{*p}, \qquad (f \circ f^{-1})_{*p} = (\mathrm{Id}_{|L|})_{*p},
\]
cio\`e $f^{-1}_{*p} \circ f_{*p} = \mathrm{Id}_{H_p(K)}$ e $f_{*p} \circ f^{-1}_{*p} = \mathrm{Id}_{H_p(L)}$. Dunque $f_{*p}: H_p(K) \to H_p(L)$ \`e un \textbf{isomorfismo} (con inverso $(f^{-1})_{*p}$)\hfill$\square$

\begin{intuizione}
Il punto chiave \`e che l'omologia \`e definita a partire da una \emph{triangolazione}, ma applicazioni continue qualsiasi tra gli spazi $|K|$ e $|L|$ possono essere ``approssimate'' da applicazioni simpliciali (dopo aver suddiviso abbastanza), e le approssimazioni diverse inducono la stessa applicazione in omologia.
\end{intuizione}


% ======== NUOVO (pp. 175-176, 3/11) ========
\section{Spazi triangolabili}

\begin{definicion}[Spazio triangolabile]
Uno spazio topologico $X$ si dice \textbf{triangolabile} se esiste un complesso simpliciale $K$ tale che $X \cong |K|$ (nelle notazioni degli appunti: $X = |K|$, a meno di omeomorfismo). Un tale $K$ si chiama una \emph{triangolazione} di $X$.
\end{definicion}

\begin{observacion}
\begin{enumerate}
    \item Se $X$ \`e triangolabile, ha senso definire $H_p(X) = H_p(K)$ per ogni $p$, dove $X = |K|$: per il teorema di invarianza topologica il risultato non dipende dalla triangolazione scelta.
    \item Se $X$ \`e triangolabile, allora $X$ \`e \textbf{compatto} e \textbf{di Hausdorff} ($|K|$ \`e unione finita di simplessi chiusi compatti, ed \`e un sottospazio di $\mathbb{R}^n$). Dunque $\mathbb{R}^n$ o un aperto non possono essere triangolabili.
\end{enumerate}
\end{observacion}

\begin{observacion}[Il viceversa del teorema di invarianza \emph{non} vale]
Se $K, L$ sono complessi simpliciali, allora $|K| \cong |L| \Rightarrow H_p(K) \cong H_p(L)$ per ogni $p$, ma \textbf{non} vale il viceversa.
\end{observacion}

\begin{ejemplo}[Controesempio]
Consideriamo due complessi simpliciali $K$ e $L$ con $H_p(K) \cong H_p(L)$ per ogni $p$ ma $|K| \not\cong |L|$:
\begin{itemize}
    \item $K = \{ (v_0), (v_1), (v_0, v_1) \}$ (un segmento);
    \item $L = \{ (v_0), (v_1), (v_2), (v_3), (v_0, v_1), (v_0, v_2), (v_0, v_3) \}$ (una ``Y'': tre segmenti uscenti da $v_0$).
\end{itemize}
Entrambi hanno dimensione $1$, ma $|K| \not\cong |L|$: togliendo il punto $v_0$ da $|L|$ si ottengono tre componenti connesse, mentre togliendo un punto interno a un segmento si ottengono al pi\`u due.

\begin{figure}[H]
    \centering
    \begin{tikzpicture}[scale=1.2]
        \begin{scope}
            \draw[very thick, blue] (0,0) -- (2,0);
            \filldraw (0,0) circle (2pt) node[below] {$v_0$};
            \filldraw (2,0) circle (2pt) node[below] {$v_1$};
            \node at (1,-0.9) {$|K|$};
        \end{scope}
        \begin{scope}[shift={(5,0)}]
            \draw[very thick, red] (0,0) -- (-1.0,-0.8);
            \draw[very thick, red] (0,0) -- (1.0,-0.8);
            \draw[very thick, red] (0,0) -- (0,1.1);
            \filldraw (0,0) circle (2pt) node[right] {$v_0$};
            \filldraw (-1,-0.8) circle (2pt) node[below] {$v_1$};
            \filldraw (1,-0.8) circle (2pt) node[below] {$v_2$};
            \filldraw (0,1.1) circle (2pt) node[above] {$v_3$};
            \node at (0,-1.7) {$|L|$};
        \end{scope}
    \end{tikzpicture}
    \caption{Un segmento e una ``Y'' non sono omeomorfi, ma hanno la stessa omologia (sono entrambi contraibili).}
\end{figure}

Calcoliamo l'omologia.
\begin{itemize}
    \item $p = 0$: $H_0(K, \mathbb{Z}) = H_0(L, \mathbb{Z}) = \mathbb{Z}$, perch\'e $|K|$ e $|L|$ sono connessi.
    \item $p = 1$, $K$: $H_1(K, \mathbb{Z}) = Z_1(K, \mathbb{Z})$ (non ci sono $2$-simplessi, $B_1 = 0$). Se $c \in Z_1(K, \mathbb{Z})$, $c = \alpha \langle v_0, v_1 \rangle$ con $\alpha \in \mathbb{Z}$ e $0 = \partial_1 c = \alpha(\langle v_1 \rangle - \langle v_0 \rangle) = -\alpha \langle v_0 \rangle + \alpha \langle v_1 \rangle$, da cui $\alpha = 0$, $c = 0$. Quindi $H_1(K, \mathbb{Z}) = 0$.
    \item $p = 1$, $L$: se $c = \alpha \langle v_0, v_1 \rangle + \beta \langle v_0, v_2 \rangle + \gamma \langle v_0, v_3 \rangle \in Z_1(L, \mathbb{Z})$, allora
    \[
    0 = \partial_1 c = \alpha(\langle v_1 \rangle - \langle v_0 \rangle) + \beta(\langle v_2 \rangle - \langle v_0 \rangle) + \gamma(\langle v_3 \rangle - \langle v_0 \rangle) = -(\alpha + \beta + \gamma)\langle v_0 \rangle + \alpha \langle v_1 \rangle + \beta \langle v_2 \rangle + \gamma \langle v_3 \rangle,
    \]
    da cui $\alpha = \beta = \gamma = 0$ e $c = 0$. Quindi $H_1(L, \mathbb{Z}) = Z_1(L, \mathbb{Z}) = 0$.
    \item $p \geqslant 2$: $C_p = 0$, quindi $H_p = 0$ per entrambi.
\end{itemize}
Dunque $H_p(K, \mathbb{Z}) \cong H_p(L, \mathbb{Z})$ per ogni $p$, ma $|K| \not\cong |L|$: l'omologia \`e un invariante topologico \emph{non completo} (distingue meno degli omeomorfismi; \`e invece un invariante di tipo di omotopia).
\end{ejemplo}


% ======== NUOVO (pp. 177-178) ========
\section{Relazione tra $\Pi_1$ e $H_1$}

Sia $K$ un complesso simpliciale con $|K|$ \textbf{connesso} e fissiamo $v_0 \in K^0$.

\medskip
\textbf{Costruzione.} Si dimostra (con l'approssimazione simpliciale, cio\`e deformando il cappio fino a farlo correre lungo gli spigoli di $K$) che per ogni cappio $\sigma \in \Lambda_{v_0}(|K|)$ esiste una $1$-catena $c_\sigma \in C_1(K, \mathbb{Z})$ del tipo
\[
c_\sigma = \langle v_0, v_1 \rangle + \dots + \langle v_{k-1}, v_k \rangle + \langle v_k, v_0 \rangle
\]
(un \emph{cammino chiuso di spigoli}) tale che $\sigma \simeq c_\sigma$ rel $\partial I$ (qui $c_\sigma$ indica anche il cappio poligonale corrispondente).

\begin{figure}[H]
    \centering
    \begin{tikzpicture}[scale=1.3]
        \fill[blue!8] (0,0) -- (3,-0.5) -- (2.4,1.8) -- cycle;
        \draw[thick] (0,0) -- (3,-0.5) -- (2.4,1.8) -- cycle;
        \draw[thick, red] (0,0) to[out=30, in=200] (1.4,0.9) to[out=20, in=150] (2.2,0.7) to[out=-80, in=100] (2.3,-0.2) to[out=200, in=20] (0,0);
        \node[red] at (1.5,0.3) {$\sigma$};
        \draw[thick, ForestGreen, dashed] (0,0) -- (2.4,1.8) -- (3,-0.5) -- cycle;
        \node[ForestGreen] at (1.0,-0.8) {$c_\sigma$ (spezzata di spigoli)};
        \filldraw (0,0) circle (1.6pt) node[left] {$v_0$};
    \end{tikzpicture}
    \caption{Ogni cappio $\sigma$ \`e omotopo (rel $\partial I$) a un cammino chiuso lungo gli spigoli $c_\sigma$.}
\end{figure}

\begin{observacion}
$c_\sigma$ \`e un $1$-ciclo. Infatti
\[
\partial_1 c_\sigma = (\langle v_1 \rangle - \langle v_0 \rangle) + (\langle v_2 \rangle - \langle v_1 \rangle) + \dots + (\langle v_0 \rangle - \langle v_k \rangle) = 0 .
\]
\end{observacion}

Si dimostra inoltre che $\sigma \simeq \tau$ rel $\partial I \implies c_\sigma - c_\tau \in B_1(K, \mathbb{Z})$. Dunque ha senso definire
\[
\varphi: \Pi_1(|K|, v_0) \longrightarrow H_1(K, \mathbb{Z}), \qquad [\sigma]_{\simeq \text{ rel } \partial I} \longmapsto [c_\sigma]_{B_1} .
\]

\begin{prop}[Propriet\`a di $\varphi$]
\begin{enumerate}
    \item $\varphi$ \`e un omomorfismo di gruppi.
    \item $\varphi$ \`e suriettiva.
\end{enumerate}
\end{prop}
\begin{proof}
1) $c_{\sigma * \tau} = c_\sigma + c_\tau$ (il cammino chiuso di spigoli che percorre prima quello di $\sigma$ e poi quello di $\tau$), quindi
\[
\varphi([\sigma][\tau]) = \varphi([\sigma * \tau]) = [c_\sigma + c_\tau]_{B_1} = [c_\sigma]_{B_1} + [c_\tau]_{B_1} = \varphi([\sigma]) + \varphi([\tau]).
\]
2) Sia $[c] \in H_1(K, \mathbb{Z})$. Si dimostra che $[c]_{B_1}$ ha un rappresentante della forma $\langle v_0, v_1 \rangle + \dots + \langle v_k, v_0 \rangle$ (un ciclo di tipo ``cammino chiuso'', usando la connessione di $|K|$ per ricondursi a $v_0$). Allora il cappio poligonale $\sigma = \overline{v_0 v_1} \cup \dots \cup \overline{v_k v_0} \in \Lambda_{v_0}(|K|)$ soddisfa $\varphi([\sigma]) = [\langle v_0, v_1 \rangle + \dots + \langle v_k, v_0 \rangle] = [c]$\hfill$\square$
\end{proof}

Per il teorema di isomorfismo, $\varphi$ induce un isomorfismo di gruppi
\[
\widetilde{\varphi}: \Pi_1(|K|, v_0) / \operatorname{Ker} \varphi \xrightarrow{\ \cong\ } H_1(K, \mathbb{Z}) .
\]
Pi\`u precisamente, si dimostra che $\operatorname{Ker} \varphi$ \`e il derivato $\Pi_1'$ (il sottogruppo dei commutatori), cio\`e
\[
\boxed{ H_1(K, \mathbb{Z}) \cong \Pi_1 / \Pi_1' } ,
\]
\emph{$H_1$ \`e l'abelianizzato di $\Pi_1$} (teorema di Hurewicz in dimensione $1$).

\begin{coro}
Sia $K$ un complesso simpliciale con $|K|$ connesso. Allora $\Pi_1(|K|)$ \`e abeliano se e solo se $\Pi_1(|K|) \cong H_1(K, \mathbb{Z})$.
\end{coro}
\begin{proof}
$\Pi_1$ \`e abeliano $\iff$ $\Pi_1' = \{e\}$ $\iff$ $\Pi_1 / \Pi_1' = \Pi_1$\hfill$\square$
\end{proof}

\begin{ejemplo}[$S^1 \vee S^1$]
Sia $X = S^1 \vee S^1$ (il ``otto''), triangolabile: sia $K$ il complesso con vertici $v_0, v_1, v_2, v_3, v_4$ e spigoli $(v_0, v_1), (v_1, v_2), (v_2, v_0), (v_0, v_3), (v_3, v_4), (v_4, v_0)$ (due triangoli vuoti che condividono il vertice $v_0$), cos\`i $X = |K|$.

\begin{figure}[H]
    \centering
    \begin{tikzpicture}[scale=1.2]
        \draw[thick, blue] (0,0) -- (-1.3,0.9) -- (-1.3,-0.9) -- cycle;
        \draw[thick, blue] (0,0) -- (1.3,0.9) -- (1.3,-0.9) -- cycle;
        \filldraw (0,0) circle (2pt) node[above=2pt] {$v_0$};
        \filldraw (-1.3,0.9) circle (1.6pt) node[above] {$v_1$};
        \filldraw (-1.3,-0.9) circle (1.6pt) node[below] {$v_2$};
        \filldraw (1.3,0.9) circle (1.6pt) node[above] {$v_3$};
        \filldraw (1.3,-0.9) circle (1.6pt) node[below] {$v_4$};
    \end{tikzpicture}
    \caption{Una triangolazione di $S^1 \vee S^1$.}
\end{figure}

Allora $\Pi_1(X) = F_2$, che non \`e abeliano, dunque $\Pi_1(X) \not\cong H_1(K, \mathbb{Z})$. Infatti
\[
H_1(X, \mathbb{Z}) = H_1(K, \mathbb{Z}) = F_2 / (F_2)' \cong \mathbb{Z}^2
\]
(l'abelianizzato di $F_2 = \langle a, b \rangle$ \`e $\mathbb{Z}^2$). Verifica diretta: $K$ ha $l_0 = 5$ vertici, $l_1 = 6$ spigoli e nessun $2$-simplesso, quindi $H_1 = Z_1 = \operatorname{Ker}\partial_1$; poich\'e $|K|$ \`e connesso $\operatorname{rango}\operatorname{Im}\partial_1 = 5 - 1 = 4$ e $\operatorname{rango} Z_1 = 6 - 4 = 2$, e $Z_1$ \`e libero (sottogruppo di $\mathbb{Z}^6$): $Z_1 \cong \mathbb{Z}^2$, generato dai due triangoli $\langle v_0, v_1 \rangle + \langle v_1, v_2 \rangle + \langle v_2, v_0 \rangle$ e $\langle v_0, v_3 \rangle + \langle v_3, v_4 \rangle + \langle v_4, v_0 \rangle$.
\end{ejemplo}

\begin{intuizione}
$\Pi_1$ ricorda \emph{in che ordine} si percorrono i cappi (non commutativo: $ab \neq ba$ in $F_2$); $H_1$ ricorda solo \emph{quante volte} si percorre ciascuno (commutativo: $ab = ba$ in $\mathbb{Z}^2$). L'omologia \`e pi\`u debole ma molto pi\`u calcolabile.
\end{intuizione}


% ======== NUOVO (pp. 179-182) ========
\section{Calcolo dei gruppi di omologia: esempi}

\begin{ejercicio*}
Provare per calcolo diretto che $H_1(S^1 \vee S^1, \mathbb{Z}) \cong \mathbb{Z}^2$ (vedi la verifica nell'esempio della sezione precedente).
\end{ejercicio*}

\begin{ejercicio*}
Determinare i gruppi di omologia $H_p(X)$, dove $X$ \`e il seguente spazio triangolabile:
\end{ejercicio*}

\subsection{1) La circonferenza $X = S^1$}

$S^1 = |K|$ con $K = [v_0, v_1, v_2]^1$ (l'$1$-scheletro del triangolo, cio\`e il contorno senza l'interno):
\[
K = \{ (v_0), (v_1), (v_2), (v_0, v_1), (v_1, v_2), (v_0, v_2) \}, \qquad \dim K = 1 .
\]
Quindi dobbiamo calcolare $H_p(S^1) = H_p(K)$ per $p = 0, 1$.
\begin{itemize}
    \item $p = 0$: $H_0(S^1, \mathbb{Z}) = \mathbb{Z}$ (connesso).
    \item $p = 1$: $H_1(S^1, \mathbb{Z}) = \Pi_1(S^1)/\Pi_1(S^1)' = \mathbb{Z}$ (perch\'e $\Pi_1(S^1) = \mathbb{Z}$ \`e abeliano e quindi $\Pi_1' = 0$).
\end{itemize}
\textbf{Calcolo diretto.} Poich\'e $\dim K = 1$, $H_1(S^1, \mathbb{Z}) = H_1(K, \mathbb{Z}) = Z_1(K, \mathbb{Z})$. Sia $c \in C_1(K, \mathbb{Z})$: $c = \alpha \langle v_0, v_1 \rangle + \beta \langle v_1, v_2 \rangle + \gamma \langle v_2, v_0 \rangle$ con $\alpha, \beta, \gamma \in \mathbb{Z}$. Allora
\[
\partial_1 c = \alpha(\langle v_1 \rangle - \langle v_0 \rangle) + \beta(\langle v_2 \rangle - \langle v_1 \rangle) + \gamma(\langle v_0 \rangle - \langle v_2 \rangle) = (\gamma - \alpha)\langle v_0 \rangle + (\alpha - \beta)\langle v_1 \rangle + (\beta - \gamma)\langle v_2 \rangle .
\]
$c \in Z_1 \iff \gamma - \alpha = \alpha - \beta = \beta - \gamma = 0 \iff \alpha = \beta = \gamma$. Dunque
\[
Z_1(K, \mathbb{Z}) = \{ \alpha ( \langle v_0, v_1 \rangle + \langle v_1, v_2 \rangle + \langle v_2, v_0 \rangle ) \mid \alpha \in \mathbb{Z} \} \cong \mathbb{Z} = H_1(S^1, \mathbb{Z}).
\]
(Il generatore \`e il ciclo che percorre la circonferenza una volta.)

\textbf{Riassunto:} $H_0(S^1) = \mathbb{Z}$, $H_1(S^1) = \mathbb{Z}$, $H_p(S^1) = 0$ per $p \geqslant 2$.

\subsection{2) Il disco $X = B^2$}

$B^2 = |K|$ con $K = [v_0, v_1, v_2]$ il triangolo pieno:
\[
K = \{ (v_0), (v_1), (v_2), (v_0, v_1), (v_1, v_2), (v_2, v_0), (v_0, v_1, v_2) \}, \qquad \dim K = 2 .
\]
\begin{itemize}
    \item $p = 0$: $H_0(B^2, \mathbb{Z}) = \mathbb{Z}$ ($B^2$ \`e connesso).
    \item $p = 1$: gli $1$-cicli sono gli stessi di $S^1 = \partial B^2$: $Z_1(K, \mathbb{Z}) = \{ \alpha ( \langle v_0, v_1 \rangle + \langle v_1, v_2 \rangle + \langle v_2, v_0 \rangle ) \} = \{ \alpha\, \partial_2 \langle v_0, v_1, v_2 \rangle \mid \alpha \in \mathbb{Z} \} = B_1(K, \mathbb{Z})$. Ora il ciclo \`e un bordo (il disco ``riempie'' la circonferenza), quindi $H_1(K, \mathbb{Z}) = Z_1/B_1 = 0$. D'altra parte, coerentemente, $\Pi_1(B^2) = 0$ implica $H_1 = \Pi_1/\Pi_1' = 0$.
    \item $p = 2$: $H_2(B^2, \mathbb{Z}) = Z_2(K, \mathbb{Z})$ (non ci sono $3$-simplessi). $c = \alpha \langle v_0, v_1, v_2 \rangle$ e $\partial_2 c = \alpha( \langle v_1, v_2 \rangle - \langle v_0, v_2 \rangle + \langle v_0, v_1 \rangle )$ \`e nullo solo se $\alpha = 0$. Quindi $Z_2 = 0$ e $H_2(B^2, \mathbb{Z}) = 0$.
\end{itemize}
\textbf{Riassunto:} $H_0(B^2) = \mathbb{Z}$, $H_p(B^2) = 0$ per $p \geqslant 1$.

\subsection{3) La sfera $X = S^2$}

$S^2 = |K|$ con $K = [v_0, v_1, v_2, v_3]^2$ (il $2$-scheletro del tetraedro: $4$ vertici, $6$ spigoli, $4$ facce triangolari), $\dim K = 2$.
\begin{itemize}
    \item $p = 0$: $H_0(S^2, \mathbb{Z}) = \mathbb{Z}$ ($S^2$ \`e connessa).
    \item $p = 1$: $\Pi_1(S^2) = 0$, quindi $H_1(S^2, \mathbb{Z}) = \Pi_1/\Pi_1' = 0$ (esercizio: calcolo diretto).
    \item $p = 2 = \dim K$: $H_2(S^2, \mathbb{Z}) = Z_2(K, \mathbb{Z})$. Una $2$-catena \`e $c = \alpha \langle v_0, v_1, v_2 \rangle + \beta \langle v_0, v_1, v_3 \rangle + \gamma \langle v_0, v_2, v_3 \rangle + \delta \langle v_1, v_2, v_3 \rangle$. Poich\'e
    \begin{align*}
    \partial_2 \langle v_0, v_1, v_2 \rangle &= \langle v_1, v_2 \rangle - \langle v_0, v_2 \rangle + \langle v_0, v_1 \rangle, \\
    \partial_2 \langle v_0, v_1, v_3 \rangle &= \langle v_1, v_3 \rangle - \langle v_0, v_3 \rangle + \langle v_0, v_1 \rangle, \\
    \partial_2 \langle v_0, v_2, v_3 \rangle &= \langle v_2, v_3 \rangle - \langle v_0, v_3 \rangle + \langle v_0, v_2 \rangle, \\
    \partial_2 \langle v_1, v_2, v_3 \rangle &= \langle v_2, v_3 \rangle - \langle v_1, v_3 \rangle + \langle v_1, v_2 \rangle,
    \end{align*}
    si ha
    \[
    \partial_2 c = (\alpha + \delta)\langle v_1, v_2 \rangle + (\alpha + \beta)\langle v_0, v_1 \rangle + (\gamma - \alpha)\langle v_0, v_2 \rangle + (\beta - \delta)\langle v_1, v_3 \rangle - (\beta + \gamma)\langle v_0, v_3 \rangle + (\gamma + \delta)\langle v_2, v_3 \rangle .
    \]
    Imponendo $\partial_2 c = 0$: $\delta = -\alpha$, $\beta = -\alpha$, $\gamma = \alpha$ (le altre equazioni sono conseguenze). Dunque
    \[
    c = \alpha \big( \langle v_0, v_1, v_2 \rangle - \langle v_0, v_1, v_3 \rangle + \langle v_0, v_2, v_3 \rangle - \langle v_1, v_2, v_3 \rangle \big), \quad \alpha \in \mathbb{Z},
    \]
    e, non esistendo $3$-simplessi in $K$, $B_2 = 0$: $H_2(S^2, \mathbb{Z}) = Z_2 \cong \mathbb{Z}$.
\end{itemize}
\textbf{Riassunto:} $H_0(S^2) = \mathbb{Z}$, $H_1(S^2) = 0$, $H_2(S^2) = \mathbb{Z}$.

\begin{observacion}
\begin{enumerate}
    \item $B^2 \not\cong S^2$, perch\'e $H_2(S^2, \mathbb{Z}) = \mathbb{Z} \neq 0 = H_2(B^2, \mathbb{Z})$. (Qui l'omologia distingue due spazi che hanno lo stesso $\Pi_1 = 0$.)
    \item \textbf{La palla $B^3$.} $H_0(B^3) = \mathbb{Z}$ (connessa), $H_1(B^3) = 0$ (contraibile) e $H_3(B^3) = 0$ (da verificare). Inoltre $H_2(B^3) = 0$: $Z_2(B^3) = Z_2(S^2)$ (perch\'e $\partial B^3 = S^2$ e $B^3$ ha gli stessi $2$-simplessi) ma ogni generatore del tipo $c = \alpha(\langle v_0, v_1, v_2 \rangle - \langle v_0, v_1, v_3 \rangle + \langle v_0, v_2, v_3 \rangle - \langle v_1, v_2, v_3 \rangle)$ \`e un bordo: $c = -\alpha\, \partial_3 \langle v_0, v_1, v_2, v_3 \rangle$ (il solido ``riempie'' la sfera). Quindi $Z_2 = B_2$ e $H_2(B^3) = 0$.
    \item In generale si dimostra che
    \[
    H_0(S^n) = H_n(S^n) = \mathbb{Z}, \quad H_p(S^n) = 0 \ \ \forall p \neq 0, n; \qquad H_0(B^n) = \mathbb{Z}, \quad H_p(B^n) = 0 \ \ \forall p \neq 0.
    \]
    In particolare $S^n \not\cong B^n$ per ogni $n \geqslant 1$ (e $S^n \not\cong B^{n+1}$).
\end{enumerate}
\end{observacion}

\begin{figure}[H]
    \centering
    \begin{tikzpicture}[scale=1.1]
        \node at (0,2.3) {\textbf{omologia}};
        \node at (2.0,2.3) {$H_0$}; \node at (3.4,2.3) {$H_1$}; \node at (4.8,2.3) {$H_2$};
        \node[left] at (0.6,1.6) {$S^1$}; \node at (2.0,1.6) {$\mathbb{Z}$}; \node at (3.4,1.6) {$\mathbb{Z}$}; \node at (4.8,1.6) {$0$};
        \node[left] at (0.6,1.0) {$B^2$}; \node at (2.0,1.0) {$\mathbb{Z}$}; \node at (3.4,1.0) {$0$}; \node at (4.8,1.0) {$0$};
        \node[left] at (0.6,0.4) {$S^2$}; \node at (2.0,0.4) {$\mathbb{Z}$}; \node at (3.4,0.4) {$0$}; \node at (4.8,0.4) {$\mathbb{Z}$};
        \node[left] at (0.6,-0.2) {$B^3$}; \node at (2.0,-0.2) {$\mathbb{Z}$}; \node at (3.4,-0.2) {$0$}; \node at (4.8,-0.2) {$0$};
        \draw (-0.3,2.0) -- (5.3,2.0);
    \end{tikzpicture}
    \caption{Tabella riassuntiva dei gruppi di omologia di $S^1, B^2, S^2, B^3$.}
\end{figure}


% ======== NUOVO (pp. 183-186, 6/11) ========
\subsection{4) Il piano proiettivo $X = \mathbb{R}P^2$}

$\mathbb{R}P^2 = S^2/\{\pm\mathrm{Id}\}$ si rappresenta come un disco (o un quadrato) con i punti antipodali del bordo identificati. Triangolando il quadrato, si ottiene un complesso simpliciale $K$ con $\mathbb{R}P^2 = |K|$, che ha
\[
l_0 = 6 \ \text{($0$-simplessi)}, \qquad l_1 = 15 \ \text{($1$-simplessi)}, \qquad l_2 = 10 \ \text{($2$-simplessi)}.
\]

\begin{figure}[H]
    \centering
    \begin{tikzpicture}[scale=1.0]
        \draw[thick, fill=blue!6] (0,0) rectangle (3,3);
        \draw[thick, red, ->, >=stealth] (0,0) -- (1.5,0); \draw[thick, red] (1.5,0) -- (3,0);
        \draw[thick, red, ->, >=stealth] (3,3) -- (1.5,3); \draw[thick, red] (1.5,3) -- (0,3);
        \draw[thick, ForestGreen, ->, >=stealth] (0,3) -- (0,1.5); \draw[thick, ForestGreen] (0,1.5) -- (0,0);
        \draw[thick, ForestGreen, ->, >=stealth] (3,0) -- (3,1.5); \draw[thick, ForestGreen] (3,1.5) -- (3,3);
        \node[red, below] at (1.5,-0.1) {$a$}; \node[red, above] at (1.5,3.1) {$a$};
        \node[ForestGreen, left] at (-0.1,1.5) {$b$}; \node[ForestGreen, right] at (3.1,1.5) {$b$};
        \node at (1.5,1.5) {$\mathbb{R}P^2$};
    \end{tikzpicture}
    \caption{Il piano proiettivo come quadrato con i lati identificati (parola $abab$ nel senso: i lati opposti sono identificati con orientazione opposta, cio\`e con una torsione).}
\end{figure}

\begin{itemize}
    \item $H_0(\mathbb{R}P^2, \mathbb{Z}) = \mathbb{Z}$ (connesso).
    \item $H_1(\mathbb{R}P^2, \mathbb{Z}) = \Pi_1(\mathbb{R}P^2)/\Pi_1(\mathbb{R}P^2)' = \mathbb{Z}_2$ (perch\'e $\Pi_1(\mathbb{R}P^2) \cong \mathbb{Z}_2$ \`e abeliano). Osserviamo che $H_1(\mathbb{R}P^2, \mathbb{Z}) = 0 \oplus \mathbb{Z}_2$: la parte libera \`e nulla (rango $0$) e tutta la torsione \`e $\mathbb{Z}_2$.
    \item $H_2(\mathbb{R}P^2, \mathbb{Z}) = 0$. (Vedremo che $H_2(\mathbb{R}P^2, \mathbb{R}) = 0$; poich\'e $H_2(K, \mathbb{Z}) = Z_2(K, \mathbb{Z})$ \`e un sottogruppo del gruppo libero $C_2(K, \mathbb{Z})$, \`e libero, dunque se ha rango $0$ \`e nullo.)
\end{itemize}

\begin{intuizione}
$H_1(\mathbb{R}P^2) = \mathbb{Z}_2$ \`e il primo esempio di \emph{torsione}: la ``retta proiettiva'' $\mathbb{R}P^1 \subset \mathbb{R}P^2$ non \`e un bordo, ma $2 \cdot \mathbb{R}P^1$ \`e il bordo dell'intera $\mathbb{R}P^2$ (cio\`e del disco), e dunque rappresenta un elemento di ordine $2$.
\end{intuizione}

\section{La caratteristica di Eulero--Poincar\'e}

\begin{definicion}[Caratteristica di Eulero--Poincar\'e]
Sia $K$ un complesso simpliciale di dimensione $m$ e sia $b_p = \dim H_p(K, \mathbb{R}) = \operatorname{rango}\big(\beta_p(K, \mathbb{Z})\big)$ ($p = 0, \dots, m$) il $p$-esimo numero di Betti. Si chiama \textbf{caratteristica di Eulero--Poincar\'e} di $K$ il numero
\[
\chi(K) = \sum_{p=0}^m (-1)^p\, b_p(K) .
\]
\end{definicion}

\begin{observacion}
$\chi$ \`e un \textbf{invariante topologico}: se $X$ \`e uno spazio triangolabile, $X = |K|$, allora $\chi(X) = \chi(K)$ (non dipende dalla triangolazione, perch\'e i numeri di Betti sono invarianti per omeomorfismo).
\end{observacion}

\begin{lema}
Sia $K$ un complesso simpliciale e $l_p$ il numero dei $p$-simplessi di $K$. Allora
\[
b_p(K) = \dim Z_p(K, \mathbb{R}) + \dim Z_{p+1}(K, \mathbb{R}) - l_{p+1} .
\]
\end{lema}
\begin{proof}
$b_p(K) = \dim H_p(K, \mathbb{R}) = \dim \dfrac{Z_p(K,\mathbb{R})}{B_p(K, \mathbb{R})} = \dim Z_p(K, \mathbb{R}) - \dim B_p(K, \mathbb{R}) = \dim Z_p(K, \mathbb{R}) - \dim \operatorname{Im} \partial_{p+1}$.
Ora $\partial_{p+1}: C_{p+1}(K, \mathbb{R}) \to C_p(K, \mathbb{R})$ \`e lineare e $\dim C_{p+1}(K, \mathbb{R}) = l_{p+1}$; per il teorema nullit\`a $+$ rango,
\[
l_{p+1} = \dim \operatorname{Ker} \partial_{p+1} + \dim \operatorname{Im} \partial_{p+1} = \dim Z_{p+1}(K, \mathbb{R}) + \dim \operatorname{Im} \partial_{p+1} .
\]
Quindi $\dim \operatorname{Im} \partial_{p+1} = l_{p+1} - \dim Z_{p+1}$ e sostituendo si ottiene la tesi\hfill$\square$
\end{proof}

\begin{teo}[di Eulero--Poincar\'e]
Sia $K$ un complesso simpliciale, $m = \dim K$ e $l_p$ il numero di $p$-simplessi di $K$. Allora
\[
\chi(K) = \sum_{p=0}^m (-1)^p\, l_p .
\]
\end{teo}
\begin{proof}
Per il lemma (con la convenzione $Z_{m+1} = 0$, $l_{m+1} = 0$):
\[
\chi(K) = \sum_{p=0}^m (-1)^p b_p = \sum_{p=0}^m (-1)^p \dim Z_p + \sum_{p=0}^m (-1)^p \dim Z_{p+1} - \sum_{p=0}^m (-1)^p l_{p+1} .
\]
Con il cambio di indice $q = p+1$ nei due ultimi termini:
\[
\chi(K) = \sum_{p=0}^m (-1)^p \dim Z_p + \sum_{q=1}^m (-1)^{q-1} \dim Z_q - \sum_{q=1}^m (-1)^{q-1} l_q .
\]
Nella prima somma separiamo il termine $p = 0$: $\sum_{p=0}^m (-1)^p \dim Z_p = \dim Z_0 + \sum_{p=1}^m (-1)^p \dim Z_p$. La seconda somma, $\sum_{q=1}^m (-1)^{q-1} \dim Z_q = -\sum_{q=1}^m (-1)^q \dim Z_q$, cancella esattamente la parte $\sum_{p=1}^m (-1)^p \dim Z_p$. Rimane $\dim Z_0 = l_0$ (perch\'e $\partial_0 = 0$), e il terzo termine vale $-\sum_{q=1}^m (-1)^{q-1} l_q = \sum_{q=1}^m (-1)^q l_q$. Dunque
\[
\chi(K) = l_0 + \sum_{q=1}^m (-1)^q l_q = \sum_{p=0}^m (-1)^p l_p . \qquad \square
\]
\end{proof}

\begin{intuizione}
Per un poliedro ``classico'' ($l_0$ vertici, $l_1$ spigoli, $l_2$ facce) si ritrova la formula di Eulero $V - E + F$. Il teorema dice che questo numero \emph{non dipende dalla triangolazione}, perch\'e coincide con $\sum (-1)^p b_p$, che \`e un invariante topologico.
\end{intuizione}

\begin{ejemplo}
\begin{enumerate}
    \item \textbf{$\chi(S^1)$.} Sappiamo che $H_0(S^1, \mathbb{R}) = \mathbb{R}$ ($b_0 = 1$) e $H_1(S^1, \mathbb{R}) = \mathbb{R}$ ($b_1 = 1$), quindi $\chi(S^1) = b_0 - b_1 = 1 - 1 = 0$. Verifica: $S^1 = |K|$, $K = [v_0, v_1, v_2]^1$, $l_0 = 3$, $l_1 = 3$, quindi per il teorema di Eulero--Poincar\'e $\chi(S^1) = \chi(K) = 3 - 3 = 0$.
    \item \textbf{$\chi(B^2)$.} $K = [v_0, v_1, v_2]$ con $l_0 = 3$, $l_1 = 3$, $l_2 = 1$: $\chi(B^2) = 3 - 3 + 1 = 1$. Verifica: $H_0(B^2, \mathbb{R}) = \mathbb{R}$ ($b_0 = 1$), $H_1(B^2, \mathbb{R}) = H_2(B^2, \mathbb{R}) = 0$ ($b_1 = b_2 = 0$), quindi $\chi(B^2) = 1 - 0 + 0 = 1$.
    \item \textbf{$\chi(\mathbb{R}P^2)$.} Con la triangolazione precedente, $\chi(\mathbb{R}P^2) = 6 - 15 + 10 = 1$.
\end{enumerate}
\end{ejemplo}

\begin{teo}
Sia $p: \widetilde{X} \to X$ un'applicazione di rivestimento ad un numero finito $h$ di fogli. Se $X$ \`e triangolabile, allora $\widetilde{X}$ \`e triangolabile e
\[
\chi(\widetilde{X}) = h \cdot \chi(X) .
\]
\end{teo}


% ======== NUOVO (pp. 187-192, 6/11) ========
\begin{ejemplo}[Altri calcoli di $\chi$]
\begin{enumerate}
    \item[3)] \textbf{$\chi(S^2)$.} $S^2 = |K|$ con $K = [v_0, v_1, v_2, v_3]^2$: $l_0 = 4$, $l_1 = 6$, $l_2 = 4$, quindi $\chi(S^2) = 4 - 6 + 4 = 2$. Infatti $\chi(S^2) = b_0 - b_1 + b_2 = 1 - 0 + 1 = 2$.
    \item[4)] \textbf{$\chi(\mathbb{R}P^2)$ e $H_2(\mathbb{R}P^2, \mathbb{R})$.} Con la triangolazione di $\mathbb{R}P^2$ ($l_0 = 6$, $l_1 = 15$, $l_2 = 10$) si ha $\chi(\mathbb{R}P^2) = 6 - 15 + 10 = 1$. Allora $1 = \chi(\mathbb{R}P^2) = b_0 - b_1 + b_2 = 1 - 0 + b_2$ (qui $b_1 = 0$ perch\'e $H_1(\mathbb{R}P^2, \mathbb{Z}) = \mathbb{Z}_2$ ha rango $0$), da cui $b_2 = 0$, cio\`e $H_2(\mathbb{R}P^2, \mathbb{R}) = 0$.
    \item[5)] \textbf{$\chi(T^2)$ e $H_2(T^2, \mathbb{R})$.} Il toro \`e triangolabile: $T^2 = |K|$ con (ad esempio) $l_0 = 16$, $l_1 = 48$, $l_2 = 32$ (si triangola il quadrato con i lati identificati). Quindi $\chi(T^2) = 16 - 48 + 32 = 0$. Allora $0 = \chi(T^2) = b_0 - b_1 + b_2$ con $b_0 = 1$ ($T^2$ \`e connesso) e $H_1(T^2, \mathbb{Z}) = \Pi_1(T^2) = \mathbb{Z}^2$ (abeliano), dunque $b_1 = 2$. Ne segue $0 = 1 - 2 + b_2$, cio\`e $b_2 = 1$ e $H_2(T^2, \mathbb{R}) = \mathbb{R}$.
\end{enumerate}
\end{ejemplo}

\begin{figure}[H]
    \centering
    \begin{tikzpicture}[scale=1.0]
        \node at (0,3.3) {\textbf{Superficie}};
        \node at (3.2,3.3) {$\chi$}; \node at (5.0,3.3) {$b_0$}; \node at (6.4,3.3) {$b_1$}; \node at (7.8,3.3) {$b_2$};
        \draw (-1,3.0) -- (8.5,3.0);
        \node[left] at (0.6,2.5) {$S^2$}; \node at (3.2,2.5) {$2$}; \node at (5.0,2.5) {$1$}; \node at (6.4,2.5) {$0$}; \node at (7.8,2.5) {$1$};
        \node[left] at (0.6,1.9) {$T^2$}; \node at (3.2,1.9) {$0$}; \node at (5.0,1.9) {$1$}; \node at (6.4,1.9) {$2$}; \node at (7.8,1.9) {$1$};
        \node[left] at (0.6,1.3) {$\mathbb{R}P^2$}; \node at (3.2,1.3) {$1$}; \node at (5.0,1.3) {$1$}; \node at (6.4,1.3) {$0$}; \node at (7.8,1.3) {$0$};
        \node[left] at (0.6,0.7) {$B^2$}; \node at (3.2,0.7) {$1$}; \node at (5.0,0.7) {$1$}; \node at (6.4,0.7) {$0$}; \node at (7.8,0.7) {$0$};
        \node[left] at (0.6,0.1) {$S^1$}; \node at (3.2,0.1) {$0$}; \node at (5.0,0.1) {$1$}; \node at (6.4,0.1) {$1$}; \node at (7.8,0.1) {$0$};
    \end{tikzpicture}
    \caption{Tabella di caratteristiche di Eulero e numeri di Betti ($\chi = b_0 - b_1 + b_2$).}
\end{figure}

\begin{ejercicio*}
Calcolare $\chi(S^1 \vee S^1)$ e $\chi(S^1 \vee S^2)$.
\end{ejercicio*}
\textit{Soluzione.} Per $S^1 \vee S^1$ (due triangoli vuoti con un vertice in comune): $l_0 = 5$, $l_1 = 6$, $\chi = 5 - 6 = -1$ (conferma: $b_0 = 1$, $b_1 = 2$, $\chi = 1 - 2 = -1$). Per $S^1 \vee S^2$ si usa $b_0 = 1$, $b_1 = 1$, $b_2 = 1$: $\chi = 1 - 1 + 1 = 1$.

\begin{teo}
Per ogni $n \geqslant 1$:
\[
\chi(S^n) = \begin{cases} 0 & \text{se } n \text{ \`e dispari}, \\ 2 & \text{se } n \text{ \`e pari}. \end{cases}
\]
Inoltre $\chi(B^n) = 1$ per ogni $n$.
\end{teo}
\begin{proof}
$S^n$ \`e triangolabile: $S^n = |K|$ dove $K$ \`e l'$n$-scheletro di un $(n+1)$-simplesso (che ha $\binom{n+2}{p+1}$ simplessi di dimensione $p$). Per il teorema di Eulero--Poincar\'e,
\[
\begin{aligned}
\chi(S^n) = \chi(K) &= \sum_{p=0}^n (-1)^p \binom{n+2}{p+1} = - \sum_{q=1}^{n+1} (-1)^q \binom{n+2}{q} \\
&= - \Big( \underbrace{\sum_{q=0}^{n+2} (-1)^q \binom{n+2}{q}}_{=(1-1)^{n+2} = 0} - 1 - (-1)^{n+2} \Big) = 1 + (-1)^n ,
\end{aligned}
\]
che vale $0$ se $n$ \`e dispari e $2$ se $n$ \`e pari. Inoltre $B^n = |K'|$ con $K'$ un $n$-simplesso chiuso, che ha $\binom{n+1}{p+1}$ simplessi di dimensione $p$:
\[
\chi(B^n) = \sum_{p=0}^n (-1)^p \binom{n+1}{p+1} = \sum_{p=0}^{n-1} (-1)^p \binom{n+1}{p+1} + (-1)^n = \chi(S^{n-1}) + (-1)^n = 1 \quad \forall n .
\]
(se $n$ \`e dispari: $\chi(S^{n-1}) + (-1)^n = 2 - 1 = 1$; se $n$ \`e pari: $0 + 1 = 1$.)\hfill$\square$
\end{proof}

\begin{observacion}
Si pu\`o anche scrivere $B^n = \partial B^n \sqcup \mathring{B}^n = S^{n-1} \sqcup \{ (v_0, \dots, v_n) \}$ (il bordo pi\`u l'unico $n$-simplesso aperto, di segno $(-1)^n$): quindi $\chi(B^n) = \chi(S^{n-1}) + (-1)^n = 1$.
\end{observacion}

\begin{coro}
\[
\chi(\mathbb{R}P^n) = \begin{cases} 0 & \text{se } n \text{ \`e dispari}, \\ 1 & \text{se } n \text{ \`e pari}. \end{cases}
\]
\end{coro}
\begin{proof}
$S^n \to \mathbb{R}P^n$ \`e un rivestimento a $h = 2$ fogli, quindi $\chi(S^n) = 2\,\chi(\mathbb{R}P^n)$ per il teorema sulle caratteristiche dei rivestimenti, da cui $\chi(\mathbb{R}P^n) = \chi(S^n)/2 = 0$ ($n$ dispari) oppure $1$ ($n$ pari)\hfill$\square$
\end{proof}

\begin{observacion}[Idea della dimostrazione di $\chi(\widetilde{X}) = h\,\chi(X)$]
Se $K$ \`e una triangolazione di $X$, dopo un numero sufficiente di suddivisioni baricentriche ogni simplesso \`e contenuto in un aperto ben coperto (numero di Lebesgue). Poich\'e un simplesso \`e semplicemente connesso, si solleva in $h$ simplessi disgiunti di $\widetilde{X}$. Si ottiene una triangolazione di $\widetilde{X}$ con $\widetilde{l}_p = h\, l_p$ simplessi in ogni dimensione, e quindi $\chi(\widetilde{X}) = \sum (-1)^p h\, l_p = h\,\chi(X)$. (Questo spiega anche perch\'e $\chi(S^n) = 2\chi(\mathbb{R}P^n)$.)
\end{observacion}

\section{La formula di Eulero per i poliedri}

\begin{definicion}[Poliedro]
Un \textbf{poliedro} di $\mathbb{R}^3$ \`e un sistema finito di poligoni, giacenti su piani distinti, disposti in modo tale che in ogni spigolo si incontrino due facce, e che sia possibile passare da uno di tali poligoni a qualunque altro esclusivamente oltrepassando degli spigoli. Un poliedro $P$ si dice (metricamente) \textbf{regolare} se \`e formato da poligoni regolari congruenti tra loro; si dice \textbf{convesso} se lo \`e (come insieme di punti); si dice \textbf{semplice} se (come insieme di punti) \`e omeomorfo a $S^2$.
\end{definicion}

\begin{observacion}
Un poliedro convesso \`e semplice (omeomorfo a $S^2$ per proiezione radiale da un punto interno) e un poliedro semplice \`e in particolare un poliedro; i viceversa sono falsi (ad esempio un poliedro non convesso, o un toro poliedrico, che non \`e omeomorfo a $S^2$). I solidi platonici (poliedri regolari convessi) sono semplici. Inoltre ogni poligono \`e triangolabile, quindi ogni poliedro \`e triangolabile.
\end{observacion}

\begin{teo}[Formula di Eulero]
Sia $P$ un poliedro semplice e siano $V$, $S$, $F$ rispettivamente il numero di vertici, spigoli e facce di $P$. Allora
\[
V - S + F = 2 .
\]
\end{teo}
\begin{proof}
$P$ semplice $\Rightarrow$ $P \cong S^2$ $\Rightarrow$ $\chi(P) = \chi(S^2) = 2$. Senza perdere generalit\`a si pu\`o assumere $P$ convesso (per ogni poliedro semplice esiste un poliedro convesso con gli stessi $V$, $S$, $F$: la formula \`e puramente combinatoria). Indichiamo con $F_i$ ($i = 1, \dots, k$) il numero di facce di $P$ con $r_i$ lati; $F_1 + \dots + F_k = F$. Essendo $P$ convesso, ogni poligono di $P$ \`e convesso e, se ha $r_i$ lati, si pu\`o triangolare aggiungendo $(r_i - 3)$ diagonali uscenti da un suo vertice fissato.

\begin{figure}[H]
    \centering
    \begin{tikzpicture}[scale=1.0]
        \draw[thick] (0,0) -- (2,0) -- (2.4,1.4) -- (1,2.3) -- (-0.4,1.4) -- cycle;
        \draw[thick, red] (0,0) -- (2.4,1.4); \draw[thick, red] (0,0) -- (1,2.3);
        \node at (1,-0.6) {pentagono: $r_i = 5$, $r_i - 3 = 2$ diagonali};
    \end{tikzpicture}
    \caption{Triangolazione di una faccia convessa con $r_i - 3$ diagonali uscenti da un vertice.}
\end{figure}

Dunque $P = |K|$, dove $K$ \`e il complesso simpliciale ottenuto da $P$ aggiungendo le diagonali sui poligoni di $P$. Per il teorema di Eulero--Poincar\'e, $2 = \chi(P) = \chi(K) = l_0 - l_1 + l_2$ con $l_0 = V$, $l_1 = S + \sum_{i=1}^k (r_i - 3)F_i$, $l_2 = \sum_{i=1}^k (r_i - 2)F_i$. Perci\`o
\[
2 = V - S - \sum_i (r_i - 3)F_i + \sum_i (r_i - 2)F_i = V - S + \sum_{i=1}^k F_i = V - S + F. \qquad \square
\]
\end{proof}

\begin{teo}[Poliedri regolari]
Esistono esattamente $5$ classi distinte di poliedri regolari: tetraedro, cubo, ottaedro, dodecaedro e icosaedro.
\end{teo}
\begin{proof}[Dimostrazione (le condizioni necessarie)]
Sia $P$ regolare. Allora $P$ \`e semplice, quindi $V - S + F = 2$. Inoltre, essendo $P$ regolare, a ogni vertice converge lo stesso numero $n$ di spigoli, e ogni faccia ha lo stesso numero $r$ di lati. Contando le coppie (vertice, spigolo) e (faccia, spigolo): $nV = 2S$ e $rF = 2S$, cio\`e $V = 2S/n$ e $F = 2S/r$. Sostituendo nella formula di Eulero:
\[
\frac{2S}{n} - S + \frac{2S}{r} = 2 \implies \frac1n - \frac12 + \frac1r = \frac1S .
\]
Per il loro significato geometrico $n \geqslant 3$ e $r \geqslant 3$. Non pu\`o essere contemporaneamente $n \geqslant 4$ e $r \geqslant 4$: altrimenti $\frac1S = \frac1n - \frac12 + \frac1r \leqslant \frac14 - \frac12 + \frac14 = 0$, assurdo. Quindi $n = 3$ oppure $r = 3$.
\begin{itemize}
    \item $n = 3$: $\frac1S = \frac13 - \frac12 + \frac1r = \frac{6 - r}{6r} > 0$ implica $r < 6$, cio\`e $r \in \{3, 4, 5\}$.
    \item $r = 3$: analogamente $n \in \{3, 4, 5\}$.
\end{itemize}
Le coppie possibili $(n, r)$ sono dunque $(3,3), (3,4), (3,5), (4,3), (5,3)$, corrispondenti rispettivamente a tetraedro, cubo, dodecaedro, ottaedro e icosaedro (ciascuna esiste ed \`e unica a meno di similitudine)\hfill$\square$
\end{proof}

\begin{figure}[H]
    \centering
    \begin{tabular}{l|c|c|c|c|c}
    \textbf{Solido} & $n$ (spigoli per vertice) & $r$ (lati per faccia) & $V$ & $S$ & $F$ \\ \hline
    Tetraedro & 3 & 3 & 4 & 6 & 4 \\
    Cubo & 3 & 4 & 8 & 12 & 6 \\
    Ottaedro & 4 & 3 & 6 & 12 & 8 \\
    Dodecaedro & 3 & 5 & 20 & 30 & 12 \\
    Icosaedro & 5 & 3 & 12 & 30 & 20 \\
    \end{tabular}
    \caption{I cinque solidi platonici: in ogni caso $V - S + F = 2$.}
\end{figure}




\end{document}